\section{Introduction}
The analysis presented in this chapter is based on Ref.~\cite{Bhattacharyya:2024kxj}.
Building on the previous chapter—where we used the WQFT framework to analyze hyperbolic scattering beyond GR in a scalar--tensor setup and focused on observables such as impulse and waveform—this chapter turns to a different (and intrinsically spin-sensitive) deformation of Einstein gravity: \emph{parity-violating} dynamical Chern--Simons (dCS) gravity. Observationally, compact-object encounters remain a prime target for current and future gravitational-wave experiments. The network of ground-based detectors has established gravitational-wave astronomy in the high-frequency band \cite{LIGOScientific:2014oec,LIGOScientific:2016aoc,LIGOScientific:2016sjg,LIGOScientific:2016vlm,LIGOScientific:2017bnn,LIGOScientific:2019hgc}, and these tests will strengthen further with the next generation of ground- and space-based detectors \cite{hyp10}. Complementarily, pulsar timing arrays probe the nano-Hz window \cite{NANOGrav:2023gor,EPTA:2023fyk,NANOGrav:2023wsz,Verbiest:2024nid}, where low-frequency radiation from wide systems and burst-like signals from hyperbolic encounters can become accessible. In addition to astrophysical information, such measurements sharpen constraints on deviations from GR, particularly in regimes where strong-field dynamics and radiation are simultaneously relevant.

From the effective-theory viewpoint, GR is expected to receive corrections at higher energies/curvatures. In Weinberg’s EFT approach \cite{PhysRev.138.B988}, one systematically augments the Einstein--Hilbert action by higher-derivative operators. One may also extend the gravitational sector by introducing additional degrees of freedom (scalars, vectors, \emph{etc.}) \cite{Damour:1992we,Horbatsch:2015bua,Schon:2021pcv,Rainer:1996gw,DeFelice:2011bh}, or consider both mechanisms at once. Two widely studied higher-curvature extensions are Einstein--dilaton--Gauss--Bonnet (EdGB) and dynamical Chern--Simons (dCS) gravity \cite{Zwiebach:1985uq,Gross:1986mw,Jackiw:2003pm}. The EdGB interaction naturally appears in the low-energy effective action of heterotic string theory via a scalar non-minimally coupled to the Gauss--Bonnet invariant \cite{Zwiebach:1985uq,Gross:1986mw}. Dynamical Chern--Simons gravity instead introduces a parity-violating correction through a scalar coupled to a topological quadratic curvature invariant, yielding an effective CP-violating deformation of four-dimensional GR and featuring prominently in string-motivated discussions \cite{Alexander:2004us,Alexander:2004xd}. The crucial point for the present chapter is that, for \emph{non-spinning} black holes, the parity-violating dCS sector does not correct the conservative scattering dynamics at leading orders; hence \emph{spin is essential} to expose the dCS imprint in classical observables. This chapter therefore focuses on \emph{spinning} black-hole scattering in dCS gravity.

Our target observable is the \emph{spinning eikonal}, computed through \emph{3PM} order. This choice is motivated by two considerations. First, within the post-Minkowskian expansion (an expansion in Newton’s constant valid to all orders in velocity), the eikonal phase provides a compact encoding of conservative scattering data, directly tied to gauge-invariant quantities such as the scattering angle and its spin dependence (and hence relevant for EOB connections \cite{Buonanno:1998gg}). Second, already at 3PM the problem probes genuinely multiloop structures, making it a sharp testing ground for both formalism and technology.

High-precision results in GR—especially for waveform modeling and radiation reaction—have a vast literature \cite{Blanchet:2013haa,Schafer:2018kuf,Buonanno:1998gg,Blanchet:2004ek,Blanchet:2006gy,Blanchet:2023bwj,Blanchet:2023sbv,Blanchet:2023soy,Warburton:2024xnr,Levi:2018nxp,Wardell:2021fyy}, driven in part by template accuracy requirements \cite{Purrer:2019jcp}. In parallel, classical two-body observables have been efficiently extracted using QFT-inspired methods: direct amplitude-based approaches \cite{Brandhuber:2022qbk,Brandhuber:2023hhy,Kosower:2018adc,DeAngelis:2023lvf,Bjerrum-Bohr:2018xdl,Bjerrum-Bohr:2013bxa,Alessio:2024wmz,Brunello:2024ibk,Bern:2019crd,Bern:2019nnu,Bern:2020buy,Bern:2024adl,Bern:2023ity,Bern:2022kto,Bern:2021yeh,Bern:2021dqo,Bern:2020gjj,Brandhuber:2019qpg,AccettulliHuber:2019jqo,AccettulliHuber:2020oou,Barack:2023oqp,Bjerrum-Bohr:2023iey,Bjerrum-Bohr:2023jau,Bjerrum-Bohr:2022ows,Bjerrum-Bohr:2022blt,Bjerrum-Bohr:2021wwt,Bjerrum-Bohr:2021vuf,Bjerrum-Bohr:2020syg,Bjerrum-Bohr:2019kec,Cristofoli:2019neg,Bjerrum-Bohr:2016hpa,Chen:2024mmm,Alessio:2022kwv,Alessio:2023kgf,ashoke1,ashoke5,Gonzo:2024zxo,Aoude:2023vdk,Aoude:2023dui,Brandhuber:2023hhl,Adamo:2024oxy,Adamo:2023cfp,Adamo:2022ooq,Adamo:2021rfq,Georgoudis:2023eke,Georgoudis:2024pdz,Bini:2024rsy}\footnote{For a broad overview and further references, see \cite{DiVecchia:2023frv}.},
worldline/QFT hybrids such as WQFT \cite{Mogull:2020sak,Jakobsen:2021lvp,Jakobsen:2021zvh,Jakobsen:2023ndj,Jakobsen:2023hig,Jakobsen:2022psy,Jakobsen:2023oow,Wang:2022ntx,Klemm:2024wtd,Bhattacharyya:2024aeq,Driesse:2024xad,Adamo:2024oxy,DeAngelis:2023lvf,Cristofoli:2021vyo},
direct field-equation approaches \cite{1978ApJ...224...62K,Damour:1992we,Damour:2016gwp,Damour:2019lcq,Bini:2020rzn,Damour:2022ybd,Bini:2024ijq,DeVittori:2014psa,hyp16},
and EFT-based methods \cite{Porto:2007pw,Porto:2008jj,Porto:2012as,Levi:2018nxp,Cheung:2024jpo,Cheung:2023lnj,Cheung:2020gyp,Ivanov:2024sds,Bhattacharyya:2023kbh,Huang:2018pbu,Diedrichs:2023foj,Loebbert:2020aos,Mougiakakos:2021ckm,Riva:2021vnj,Mougiakakos:2022sic,Riva:2022fru,Bernard:2023eul,Porto:2007px,Porto:2017dgs,Dlapa:2024cje}.
In this chapter we adopt the WQFT viewpoint, leveraging the operator/S-matrix correspondence established by Mogull--Plefka--Steinhoff \cite{Mogull:2020sak,Jakobsen:2021lvp}, which has proven effective for computing classical scattering observables while bypassing certain bookkeeping subtleties that arise in more traditional amplitude extractions. The same philosophy has already enabled explorations beyond GR in scalar--tensor theories \cite{Bhattacharyya:2023kbh}; here we apply it to a higher-curvature, parity-violating deformation.

The genuinely new ingredient is \emph{spin}. In WQFT, spin is incorporated through additional anti-commuting worldline variables $\psi^{a}$, whose bilinears realize the spin tensor, schematically $S^{ab}\sim\bar\psi^{[a}\psi^{b]}$. Jakobsen \emph{et al.} showed that the spinning worldline theory underlying GR enjoys a hidden $\mathcal N=2$ structure \cite{Jakobsen:2021lvp,Jakobsen:2021zvh}, which constrains the couplings and clarifies the relationship to standard spinning worldline actions used in EFT treatments. For scattering, this formulation is particularly natural: classical conservative dynamics is efficiently packaged by the \emph{spinning eikonal}, and in the present dCS context it is precisely the spin sector that activates the parity-violating correction. A qualitative theme we emphasize is that the resulting dCS contribution can depend sensitively on spin orientation in the scattering configuration—an imprint of the CP-violating structure \cite{Yunes:2009hc}.

Technically, pushing PM calculations to higher orders confronts the multiloop integration bottleneck. Recent progress has revealed that high-PM integrals can probe rich analytic structures and connect to modern mathematical tools: intersection theory \cite{Mastrolia:2018uzb,Brunello:2023rpq,Brunello:2023fef}, Calabi--Yau geometries \cite{Frellesvig:2023bbf,Klemm:2024wtd}, modular graph forms \cite{Dorigoni:2022npe}, and function classes extending beyond multiple polylogarithms \cite{Weinzierl:2007cx}. In this chapter, we work at 3PM and develop a practical analytic pipeline based on integration-by-parts (IBP) reduction and differential equations for master integrals, which suffices to obtain closed-form control over the spinning eikonal in the dCS deformation.

The chapter is organised as follows. In Section~(\ref{ch3:sec2}) we formulate spinning WQFT in the presence of the effective parity-violating (dCS) interaction and derive the corresponding Feynman rules, including the vertices involving the extra spin degrees of freedom. Section~(\ref{ch3:sec3}) develops the technical tools used throughout, focusing on IBP reduction and differential-equation methods for the relevant multiloop master integrals. In Section~(\ref{ch3:sec4}) we apply these ingredients to compute the spinning eikonal through 3PM and connect it to the exponentiated WQFT partition function. Useful integral details are collected in the appendices.
\section{Chern-Simons gravity and worldline QFT}
\label{ch3:sec2}
In this section, we will review the dynamical Chern-Simons (dCS) model and derive all relevant Feynman rules. We also comment on the main differences between the higher curvature EFTs and Einstein GR.
\subsection*{Bulk gravitational action:}
The bulk action for dCS theory is given by,
\begin{align}
    \begin{split}
        S_{\textrm{dCS}}=\rmint d^4 x\sqrt{-g}\Big[-\frac{m_p^2}{2}R+\frac{1}{2}g^{\alpha\beta}\partial_{\alpha}\varphi\partial_{\beta}\varphi-\frac{l_{dCS}^2}{m_p^3}\varphi\, {}^{*}R R\Big]\,.
    \end{split}
\end{align}
%where, $\mathcal{G}(R^{(2)})$ is the Gauss-Bonnet term defined as,
%\begin{align}
  %  \begin{split}
    %    \mathcal{G}(R^{(2)}):=R_{\mu\nu\alpha\beta} R^{\mu\nu\alpha\beta}+R^2-4R_{\alpha\beta}R^{\alpha\beta}
  %  \end{split}
%\end{align}
Now we expand the bulk gravitational action around flat metric as,
\begin{align}
    \begin{split}
        g_{\mu\nu}=\eta_{\mu\nu}+\frac{h_{\mu\nu}}{m_p},\,\, \varphi\to 0+\varphi,
    \end{split}
\end{align}
Hence, the Chern-Simons term can be expanded in order of fluctuation $h$ as,
%\begin{align}
  %  \begin{split}
      %  \mathcal{G}[R^{(2)}]=\frac{1}{4}\Big(\partial_{\nu}\partial_{\beta}h_{\mu\alpha}+\partial_{\mu}\partial_{\alpha}h_{\nu\beta}-\partial_{\mu}\partial_{\beta}h_{\nu\alpha}-\partial_{\nu}\partial_{\alpha}h_{\mu\beta}\Big)^2+\Big(\Box h-\partial_{\alpha}\partial_{\beta}h^{\alpha\beta}\Big)^2
   % \end{split}
%\end{align}
\begin{align}
    *RR=\epsilon^{\chi\varepsilon\mu\nu}\Big(\partial_\mu\partial_\beta h_{\nu}^{\,\sigma}-\partial_\mu \partial^\sigma h_{\nu\beta}\Big)\Big(\partial_{\chi}\partial_{\sigma}h_{\delta}^{\,\beta}-\partial_{\chi}\partial^{\beta}h_{\delta\sigma}\Big)+\mathcal{O}(h^3).
\end{align}
Before going to the worldline action for the higher curvature EFTs, we first briefly review the case of General relativity.
\subsection*{SUSY worldline action in General relativity:}
As we will be considering spinning binaries in this paper, following \cite{Jakobsen:2021zvh}, we start with the $\mathcal{N}=2$ supersymmetric worldline action. It  has the following form:
\begin{align}
    \begin{split}
        S_{\mathcal{N}=2}=-\sum_{k=1}^2 
        \int d\tau\, \Big[\frac{1}{2e}g_{\mu\nu}\dot x_k^{\mu}\dot x_k^{\nu}+i\bar \psi_{k,a}\frac{\mathtt{D}\psi^a_k}{\mathtt{D}\tau}+\frac{e}{2}R_{abcd}\,\bar \psi^a_k \psi^b_k \bar \psi^c_k \psi^d_k+\frac{e}{2}m_k^2\Big]
    \end{split}\label{2.4a}
\end{align}
where,
\begin{align}
    \frac{\mathtt{D}\psi ^a_k}{\mathtt{D}\tau}= \dot\psi_k^a+\dot x^\mu\, \omega_{\mu}^{\,\,ab}\psi_{k,b}\,.
\end{align}
Moreover, in a convenient fashion, we can make the following gauge choice $e=1/m_k$ and make a rescaling of the fermions \cite{Jakobsen:2021zvh}:
\begin{align}
    \psi_k^a \to \sqrt{m_k}\,\psi_k^a,\,\,\bar\psi_k^a \to \sqrt{m_k}\,\bar\psi_k^a,
\end{align}
The action in \eqref{2.4a} is invariant under the following asymptotic $\mathcal{N}=2$ supersymmetry transformation.
\begin{align}
    \begin{split}
    &    \delta x^\mu=ie^{\mu}_a\,(\bar\epsilon\psi^a+\epsilon\bar\psi^ a),\,\,\,\delta \psi^a=-\epsilon e^{a}_\mu \dot x^\mu-\delta x^\mu\,{\omega_{\mu}}^{\,\,a}_{\,\,\,\,\,\,\,b}\,\psi^b\\ &
    \textrm{and,}\,\textcolor{black}{\delta}\bar\psi^a=-\bar\epsilon e^{a}_\mu \dot x^\mu-\delta x^\mu\,{\omega_{\mu}}^{\,\,a}_{\,\,\,\,\,\,\,b}\,\bar\psi^b.
    \end{split}
\end{align}
It also has global $U(1)$ symmetry.
\begin{align}
    \delta_{\hat\epsilon}\psi^a=i\hat\epsilon \psi^a,\,\,\delta_{\hat\epsilon}\bar\psi^a=i\hat\epsilon \bar\psi^a,\,\,\delta_{\hat\epsilon}x^\mu=0.
\end{align}
The Noether charges derived from these global transformations give different conservation laws along the worldline \cite{Jakobsen:2021zvh}.\par
In a perturbative framework, one can not distinguish between tetrad indices $(a,b,..)$ and curved indices $(\mu,\nu,..)$. In order to describe the scattering process, one can expand worldline DOF around undeflected trajectories as,
\begin{align}
    \begin{split}
        x_k^\mu\to b_k^\mu+v^\mu_k\tau+z_{k}^{\mu}(\tau),\,\,\, \psi^{a}_k(\tau)\to\Psi_k^{a}+{\psi}_k^a(\tau)\,.
    \end{split}
\end{align}
One can also introduce the constant spin tensor as,
\begin{align}
    \begin{split}
        \mathcal{S}_k^{ab}=-2i\bar \Psi_k^{[a}\Psi_k^{b]}\,.
    \end{split}
\end{align}
From the quadratic part of the action, one can find out the field and worldline propagators,
\begin{align}
    \begin{split}
     &    \begin{minipage}[h]{0.075\linewidth}
	\vspace{4pt}
	\scalebox{1.8}{\includegraphics[width=\linewidth]{graviton_p.png}}
\end{minipage}\hspace{0.45cm}=\frac{iP_{\mu\nu;\rho\sigma}}{k^2+i\varepsilon},\,\begin{minipage}[h]{0.075\linewidth}
	\vspace{4pt}
	\scalebox{1.6}{\includegraphics[width=\linewidth]{w_1.png}}
\end{minipage}\hspace{0.45cm}=-\frac{i \eta^{\mu\nu}}{2m_k}\left(\frac{1}{(\omega+i\varepsilon)^2}+\frac{1}{(\omega-i\varepsilon)^2}\right),\\ &
 \begin{minipage}[h]{0.075\linewidth}
	\vspace{4pt}
	\scalebox{1.8}{\includegraphics[width=\linewidth]{scalar_p.png}}
\end{minipage}\hspace{0.45cm}=\frac{i}{k^2+i\varepsilon},\,
        \begin{minipage}[h]{0.075\linewidth}
	\vspace{4pt}
	\scalebox{1.6}{\includegraphics[width=\linewidth]{w_2.png}}
\end{minipage}\hspace{0.45cm}=-\frac{i \eta^{\mu\nu}}{2m_k}\left(\frac{1}{\omega+i\varepsilon}+\frac{1}{\omega-i\varepsilon}\right).
    \end{split}
\end{align}
\textcolor{black}{Note that in this paper, we intend to compute the Eikonal phase, and hence we choose a time-symmetric/Feynman propagator. However, if one wants to compute the observables, then one needs to choose the retarded propagators, which gives the causally sensible classical solutions.}

One can eventually expand the tetrad and spin connection as,
{\small
\begin{align}
    \begin{split}
& e^a_{\,\mu}=\eta^{a\nu}\Big(\eta_{\mu\nu}+\frac{1}{2m_p}h_{\mu\nu}-\frac{1}{8m_p^2}h_{\mu\rho}h^{\rho}_{\,\nu}+\mathcal{O}(1/m_p^3)\Big)\,,\\ &
\omega_{\mu}^{\,\,\,ab}=-\frac{1}{m_p}\partial^{[a}h^{b]}_{\,\,\,\mu}-\frac{1}{2m_p^2}h^{\nu[a}\left(\partial^{b]}h_{\mu \nu}-\partial_\nu h^{b]}_{\,\,\,\,\mu}+\frac{1}{2}\partial_\mu h^{b]}_{\,\,\,\,\nu}\right)+\mathcal{O}(1/m_p^3).
    \end{split}
\end{align}
}
\subsection*{What happens for modified theories of gravity?}
In some EFT of gravity involving extra scalar degrees of freedom the scalar dipole moment must couple to the scalar field in the worldline Lagrangian. One way, mentioned in \cite{Loutrel:2018ydv}, is to demand the new canonical momenta $\mathcal{P}^\mu$ is now not only a function of graviton degrees of freedom but also a function of the extra degrees of freedom:
\begingroup
\small
\begin{align}
    \begin{split}
        \mathcal{P}_\mu=p_{\mu}+(\textrm{pure gravity contributions})+\frac{\partial \mathcal{P_{\mu}}}{\partial{\varphi}}\Bigg|_{\varphi=0}\,\varphi+\frac{\partial \mathcal{P_{\mu}}}{\partial{\nabla_{\alpha}\varphi}}\Bigg|_{\nabla_\alpha\varphi=0}\,\nabla_{\alpha}\varphi+\cdots\label{2.13k}
    \end{split}
\end{align}
\endgroup
In dCS theory, we have an extra bulk degree of freedom (in the form of a scalar field) along with graviton, and so we need to promote the masses of the black holes as a function of that scalar field $\varphi$, which  does not preserve the SUSY invariance for the following re-scaled fermions as,
\begingroup
\small
\begin{align}
    \psi_k^a \to \sqrt{m_k(\varphi)}\,\psi^a,\,\,\bar\psi_k^a \to \sqrt{m_k(\varphi)}\,\bar\psi^a,
\end{align}
\endgroup
 Extrapolating \cite{Loutrel:2018ydv}, the worldline action takes the form (up-to linear in spin),
\begin{align}
    \begin{split}
        \mathcal{S}_k\subset -\int d\tau \Big[\frac{m_k(\varphi)}{2}g_{\mu\nu}\dot x_k^\mu \dot x_k^\nu+i m_{k}(\varphi)\bar\psi_a^k\, D_{\tau}\psi^a_k-i\,\mathcal{C}_{dCS} \,\dot{x}^{\mu} \nabla_{\alpha}\varphi\,\epsilon_{\mu\,\,\,\rho\,\sigma}^{\,\,\alpha}\bar{\psi}_k^a\,\psi_k^{b}e^{\rho}_a e^{\sigma}_{b}\Big]\label{2.13d}
    \end{split}
\end{align}
where $\mathcal{C}_{dCS}=\frac{\partial \mathcal{P_{\mu}}}{\partial{\nabla_{\alpha}\varphi}}\Bigg|_{\nabla_\alpha\varphi=0}$ is an undetermined coefficient and refers to as scalar-dipole constant and is dependent on the coupling constant(s) of the theory. However, in our computation, we will ignore the finite size corrections. Hence, we can avoid the dipole term from the worldline action.  \par 
Now expanding the mass around $\varphi=0$, we get \footnote{In \cite{Wilson-Gerow:2025xhr}, the author asserts that the dynamical Chern-Simons (dCS) term contributes to the spinless scattering angle by treating the sensitivity parameters as functions of the theory's coupling. However, some contradictory claims regarding our computation have been made in \cite{Wilson-Gerow:2025xhr}, and we would like to provide a few clarifications to avoid any potential confusion. First and foremost, we do not compute the scattering angle, contrary to the assertion in \cite{Wilson-Gerow:2025xhr}. Moreover, deriving the scattering angle from the Eikonal phase is nontrivial at 3PM order and requires the careful considerations discussed in Section~\eqref{con}. Second, as we understand it, the author's approach primarily relies on the on-shell behaviour of the scalar field in the far region to establish a link between the sensitivity parameters and the couplings of the underlying theory. While it is true that all parameters in the model are implicitly (or explicitly) dependent on the theory's coupling, the precise functional relationship between these parameters and the theory itself remains ambiguous. For this reason, we adopt a specific prescription in our analysis, treating these parameters as independent. 
Our focus is on the diagrams where the explicit coupling of the dCS term is directly incorporated (and it can be easily seen that it contributes only for spinning case due its structure), and we defer the investigation of potential indirect effects emerging from the broader theoretical framework to future work. If one wishes to carry out the computation using an alternative prescription, one need only apply the transformation:  
   $ \{s_i, g_i\} \to \{s_i(l_{\textrm{dCS}}), g_i(l_{\textrm{dCS}})\}$
where the sensitivity parameters and couplings are explicitly treated as functions of the dCS coupling and get the appropriate results from the expressions provided in our paper. Also, in the \textit{spinless} part of the eikonal there are contributions form UT-2 polynomial as presented in \eqref{dCSne3waaa} which seems to be absent from the result presented in \cite{Wilson-Gerow:2025xhr}.  },
\begin{align}
    \begin{split}
        m_{k}(\varphi)=m_{k}(0)\Big(1+s_k \varphi +\cdots)
    \end{split}
\end{align}
Upon applying the SUSY transformation, we get, $\delta m_{k}(\varphi)=s_k \partial_{\mu}\varphi \delta x^\mu+\cdots$ and demanding that the sensitivity parameter $s_1$ is very small, we get $\delta m_k(\varphi)\approx \mathcal{O}(s_1)$ and the SUSY is preserved in an approximate sense. However, the breaking of SUSY is not a problem as the action in \eqref{2.13d} eventually recovers the standard spinning worldline action with dynamical mass \cite{Levi:2018nxp}. 
\par
\par
\subsubsection*{Primary differences between GR and modified EFTs:}
For GR, following \cite{Jakobsen:2021zvh}, one can argue that the worldline action is invariant under the SUSY transformation along the entire trajectory of the black holes as well as in the asymptotic regions, which eventually implies that $p_k^2,\ p_k\cdot \bar\psi_k,\,p_k\cdot \psi_k,\,\textrm{and,}\,\psi_k\cdot \bar\psi_k$ (asymptotic charges) are conserved between the initial and final asymptotic states. Using those constraints together with the extra condition $v_k\cdot \Psi_k=v_k\cdot \bar\Psi_k=0$ then implies that $p_{k,\mu} \cdot \mathcal{S}_k^{\mu\nu}$ is conserved throughout the scattering event, and without loss of generality we can choose it to vanish, which essentially recovers the SSC \cite{1967JMP.....8.1591D,Corinaldesi:1951pb,1964NCim...34..317D}:
\begin{align}
    \textrm{\bf SSC}: \, p_{k,\mu} \cdot \mathcal{S}_k^{\mu\nu}=0,\,\, \textrm{with,}\,\,
     p_k^\mu=m_k v_k^\mu+\mathcal{O}(\mathcal{S}^2).
    \label{SSC17}
\end{align}
\textcolor{black}{In our case, the worldline SUSY invariance is broken (or only approximate if we demand that the bulk scalar varies very slowly with respect to the proper time), and hence the conservation of supercharges is only approximate. As a consequence, the conservation of $p_k\cdot \mathcal{S}_k^{\mu\nu}(\tau)\,$ is also only approximate. Thus, the SSC throughout the dynamical process is only approximate, as shown in \cite{Loutrel:2018ydv}\footnote{If one chooses $m_k(\varphi)$ to be the renormalized mass such that it satisfies $D_{\tau} m_{k}(\varphi)=\mathcal{O}(S^4)$, then the SUSY invariance is recovered. Interested readers are referred to Appendix (B) of \cite{Loutrel:2018ydv}.}. However, the theory still enjoys asymptotic (background) SUSY invariance, since there is no gravity in the asymptotic states. Following \cite{Loutrel:2018ydv}, the SSC defined in \eqref{SSC17} holds for the asymptotic states, i.e. $v_{k,\mu} \cdot \mathcal{S}_k^{\mu\nu}(-\infty)=0$, which we will use in the subsequent computations, and the observables or eikonal depend only on the asymptotic parameters. Therefore, we can still use the WQFT formalism to compute the scattering observables. We will also work to linear order in the spin and in the small-coupling ($l_{dCS}$) limit so that we can replace the canonical momenta $p_k^\mu\to m\,v_k^\mu$. 
%In all of our subsequent computations we will restrict ourselves to $\mathcal{O}(\mathcal{S}^1)\,.$ 
}
\par
\subsection*{Main goal and computational procedure:}
In WQFT formalism $h_{\mu\nu}(x),\varphi(x),x^\mu(\tau),\psi^\mu(\tau)$ are promoted to the quantum operators. For this case, the partition function looks like,
\begin{align}
    \begin{split}
        Z_{\textrm{S-WQFT}}\equiv e^{i\chi}=\mathcal{N}&\times\rmint \boldsymbol{\mathcal{D}}[h_{\mu\nu},\varphi]\rmint \prod_{i=1}^n \boldsymbol{\mathcal{D}}[z_i,\psi_i,\bar{\psi}_i,\mathfrak{a}_i,\mathfrak{b}_i,\mathfrak{c}_i]\,,\\ &
        \times \exp\Big[i\Big(S_{EH}+S_{\textrm{dCS}}+\sum_i (S_{\textrm{sp}}^{i}+S^{i}_{\textrm{ghost}})\Big)\Big],
    \end{split}
\end{align}
where, $\chi$ is the eikonal phase and the extra ``ghost'' term comes from the metric dependent worldline measure,
\begin{align}
    \begin{split}
       \boldsymbol{ \mathcal{D}}[x]&=D[x]\prod_{0\le \sigma\le T}\sqrt{-g[x(\sigma)]}\,,\\ &
        =D[x]\underbrace{\rmint \boldsymbol{\mathcal{D}}[\mathfrak{a},\mathfrak{b},\mathfrak{c}]\exp \Big[-i\rmint d\tau\Big(\frac{1}{2}g_{\mu\nu}(\mathfrak{a}^\mu\mathfrak{a}^\mu+\mathfrak{b}^\nu\mathfrak{c}^\nu)\big)\Big]}_{\textrm{ghost term}}\,.
    \end{split}
\end{align}
In the classical computations, the ghost field does not appear and can be ignored. The relevant question, then, is how the WQFT partition function is related to the scattering amplitude. It has been shown that in the classical limit the eikonal phase is given by \cite{Amati:1987wq,Amati:1990xe},
\begin{align}
    e^{i\chi}=\frac{1}{4m_1 m_2}\int  \frac{d^4 q}{(2\pi)^4} \, \hat\delta(q\cdot v_1)\hat\delta(q\cdot v_2) e^{iq\cdot b}\,\langle \,\textrm{final}\,|\mathbb{S}|\,\textrm{initial}\,\rangle\xrightarrow[]{} Z_{\textrm{WQFT}}\label{2.18a}
\end{align}
So, the WQFT partition function is equivalent to computing the eikonal phase, which is nothing but the $2\to2$ scattering amplitude, Fourier transformed into impact parameter space.
Now, we have all the ingredients to compute the Feynman rules. Our primary goal in this paper is to compute the eikonal phase for the spinning binaries, especially to investigate the contribution due to the extra degrees of freedom (scalar field) present in the theory and see how the dCS modifies the phase. 
We will mainly concentrate on the spin part of the worldline action as the spinless part (from the scalar field) has already been done in \cite{Bhattacharyya:2024aeq}. For simplicity, we only focus on the terms where spins are linear. Now, to read off the interaction vertices from the supersymmetric worldline action, write the bulk fields in Fourier space.
\begin{align}
    \begin{split}
        \chi(x_i)=\sum_{n-0}^{\infty}\frac{i^n}{n!}\rmint_{k,\omega_1,\cdots \omega_n}e^{ik\cdot b_i}e^{i(k\cdot v_i+\sum_{j=1}^{n}\omega_j)\tau}\Big(\prod_{j=1}^n k\cdot z_i(-\omega_j)\Big)\chi(-k)
    \end{split}
\end{align}
and the worldline fields can be written in energy space as,
\begin{align}
    \begin{split}
        z_{i}^{\mu}(\tau)=\rmint_{\omega} e^{i\omega \tau}z_i^\mu{(-\omega)},\,\,\,\,\,\,\breve{\psi}_i^\mu(\tau)=\rmint_{\omega}e^{i\omega \tau}\breve{\psi}_i^\mu(-\omega)\,.
    \end{split}
\end{align}
The linear spin term in the worldline action takes the following form,
\begin{align}
    \begin{split}
        S_{sp}\sim& -i\rmint d\tau \Big(m_k(0)+\frac{m_k s_k}{m_p} \varphi+\frac{m_k g_k }{m_p^2}\varphi^2\Big)\Big(\bar\psi_{k,a}\dot \psi_k^a+\dot x_k^\mu \omega_{\mu}^{\,\,ab}\bar \psi_{k,a} \psi_{k,b}\Big)\\ &+(\textrm{dCS contribution to the worldline action})\,.\label{ch3:2.13}
    \end{split}
\end{align}
Note that in \eqref{ch3:2.13} at leading order ($\mathcal{O}(\varphi^0)$) we can add a total derivative term $-1/2\, d_{\tau}(\bar\psi^a\,\psi_a)$ which preserves the SUSY transformation. But when we go to the linear (or higher order in $\varphi$), this addition will be important so that we can have both $\varphi-\psi$ and $\varphi-\bar\psi$ vertices. Again, as mentioned earlier, the SUSY invariance is approximate, and we restrict ourselves to the linear order of $\varphi\,.$ 
Keeping terms that are linear in $\varphi$ we have the following,
\begin{align}
    \begin{split}
        S_{sp}\subset -i\frac{m_i s_i}{2 m_p}\int d\tau\,\varphi\, \,\,\Big(\bar\psi _a \dot \psi^a-\dot{\bar\psi}_a \psi^a\Big)\,. \label{2.19}
    \end{split}
\end{align}
\\We can now compute the Feynman rules in $\mathcal{O}(z^n,{\psi}^n,\chi^n)$. We will mainly focus on the diagrams in which we have an external $\varphi$ field.

%\textbullet $\,\,$ $\mathcal{O}(z^0,\Breve{\psi}^0,\varphi,h_{\mu\nu})$: 
%\begin{align}
    %\begin{split}
    %   -i \frac{m_1s_1}{m_p}\rmint d\tau \rmint_{k}e^{ik\cdot b}e^{i k\cdot v_i \,\tau} v^\mu \omega_{\mu}^{\,\,ab}(h)\bar\Psi_{a}\Psi_{b}&\to 2i\frac{m_1s_1}{m_p^2}\rmint d\tau \rmint_{k}e^{ik\cdot b}e^{i k\cdot v_i \,\tau} v^\mu \bar\Psi_{[a}\Psi_{b]}\partial^{[a}h^{b]}_{\,\,\mu}\,,\\ &
   %  \to 2i\frac{m_1s_1}{m_p^2}\rmint d\tau \rmint_{k}e^{ik\cdot b}e^{i k\cdot v_i \,\tau} v^\mu \bar\Psi_{[a}\Psi_{b]}\partial^{[a}\rmint_q e^{iq \cdot b}e^{iq\cdot v_i \tau}h^{b]}_{\mu}(-q) \varphi(-k)\,,\\ &
  %   \to -2 \frac{m_1s_1}{m_p^2}\rmint d\tau \rmint_{k}e^{i(k+q)\cdot b}e^{i (k+q)\cdot v_i \,\tau} v^\mu \mathcal{S}_{ab} \,iq^a\, h^{b}_{\,\,\mu}(-q)\varphi(-k)\,.
 %   \end{split}
%\end{align}
%Hence, the vertex factor looks like,
%\begin{align}
%\begin{split}
 %   \mathcal{V}[\mathcal{O}(z^0,\Breve{\psi}^0,\varphi,h_{\mu\nu})]:&
 %   \begin{minipage}[h]{0.08\linewidth}	\vspace{4pt}
%\scalebox{1.5}{\includegraphics[width=\linewidth]{phi_h_0.png}}
%\end{minipage}\\ &
%\hspace{0 cm}-4i\frac{m_is_i}{m_p^2}e^{i(k+q)\cdot b}\hat\delta(k\cdot v_i+q\cdot v_i)\Big(q\cdot \mathcal{S}\Big)^{(\mu}v^{\nu)}+\textrm{spinless part}\,.
%\end{split}
%\end{align}
%\vspace{-1.3 cm}
\subsection*{Feynman rules:}
 We list the additional vertices that we need in our computations. Other relevant vertices for the graviton-worldline coupling are derived in \cite{Jakobsen:2021zvh}, while the vertices involving the extra scalar and scalar-graviton couplings are derived in \cite{Bhattacharyya:2024aeq}.
\begingroup
\scriptsize
\begin{align}
    \begin{split}
      &  \mathcal{V}_1:   \begin{minipage}[h]{0.08\linewidth}	%\vspace{4pt}
\scalebox{1.5}{\includegraphics[width=\linewidth]{int_1.png}}
\end{minipage}\hspace{0.35cm}
\equiv 
\hspace{0 cm}i\frac{m_a s_a}{m_p} \,\bar\Psi^a_{\eta}\,\omega\,\hat\delta(\omega+k\cdot v_a) e^{ik\cdot b_a}\,,\\ &
\mathcal{V}_2:   \begin{minipage}[h]{0.08\linewidth}	%\vspace{4pt}
\scalebox{1.5}{\includegraphics[width=\linewidth]{v_1.png}}
\end{minipage}\hspace{0.35cm}
\equiv 
\hspace{0 cm}i\frac{m_a s_a}{2 m_p}\hat\delta(k\cdot v_a-\omega_1+\omega_2 )\,(\omega_2+\omega_1) e^{ik\cdot b_a}\,,\\ &
\mathcal{V}_3:   \begin{minipage}[h]{0.08\linewidth}	%\vspace{4pt}
\scalebox{1.3}{\includegraphics[width=\linewidth]{phi_h_0.png}}
\end{minipage}\hspace{0.35cm}
\equiv 
\hspace{0 cm}-i\frac{m_a s_a}{2m_p^2}e^{i(k_1+k_2)\cdot b}\,\hat\delta(k_1\cdot v_a+k_2\cdot v_a)\Big(v_a^\mu v_a^\nu+i (k_2\cdot \mathcal{S}_a)^{(\mu}v_a^{\nu)}\Big)\,,
\\ &
\hspace{0cm}\mathcal{V}_4:   \begin{minipage}[h]{0.08\linewidth}	%\vspace{4pt}
\scalebox{1.6}{\includegraphics[width=\linewidth]{h_z_2.png}}
\end{minipage}\hspace{0.35cm}
\equiv \, i\frac{m_a }{m_p} \,e^{ik\cdot b_a}\,\hat\delta(k\cdot v_a+\omega_1+\omega_2) \Bigg(\frac{1}{2} k_{\rho_1} k_{\rho_2} v_a^\mu v_a^\nu +\omega_1 k_{\rho_2}v_a^{(\mu}\delta^{\nu)}_{\rho_1}+\omega_2 k_{\rho_1}v_a^{(\mu}\delta^{\nu)}_{\rho_2}+\omega_1\omega_2 \delta^{(\mu}_{\rho_1}\delta^{\nu)}_{\rho_2}\\&
\hspace{0.8cm}+\left.\frac{i}{2}\Bigg(\left(\omega _1 \,k\cdot S^{(\mu }\delta ^{\nu )}_{\rho _1}k_{\rho _2}+\omega _2\,k \cdot S^{(\mu }\delta ^{\nu )}_{\rho _2}k_{\rho _1}\right)+\frac{1}{2}k\cdot S^{(\eta }\delta ^{\eta )}_{(\nu }v_{a,\mu )}k_{\rho _1}k_{\rho _2}\Bigg)\Bigg)\right]
    \end{split}
\end{align}
\endgroup
% \subsection*{\underline{Vertex for scalar field: \textcolor{red}{CITE our paper?}}}
%\begin{itemize}
  %  \textbullet $\,\,$ $\mathcal{O}{(z^0)}:- i\frac{m_a s_a}{2m_p} \,e^{ik\cdot b_a}\,\hat\delta(k\cdot v)\rightarrow  \begin{minipage}[h]{0.08\linewidth}
	%\vspace{4pt}
	%\scalebox{1.5}{\includegraphics[width=\linewidth]{feynmann_2_1.png}}
%\end{minipage}$\\
% \textbullet $\,\,$ $\mathcal{O}(z):\frac{m_a s_a}{2m_p}\,e^{ik\cdot b}\,\hat{\delta}(k\cdot v+\omega)(2\omega v_{\rho}+k_{\rho})\rightarrow \begin{minipage}[h]{0.08\linewidth}
%	\vspace{4pt}
	%\scalebox{1.5}{\includegraphics[width=\linewidth]{feynmann_3_1.png}}
%\end{minipage}$\\
 %\textbullet $\,\,$ $ \mathcal{O}(z): -\frac{m_a s_a}{2m_p^2}e^{i(k_1+k_2)\cdot b_a}\hat\delta(k_1\cdot v_1+k_2\cdot v_1+\omega)\{(k_{1\rho}+k_{2\rho})v^{\mu}v^{\nu}+2\omega v^{(\mu}\delta_{\rho}^{\nu)}\}\rightarrow
%\begin{minipage}[h]{0.08\linewidth}
%	\vspace{4pt}
%	\scalebox{1.8}%{\includegraphics[width=\linewidth]{phi_2_z_1.png}}
%\end{minipage}$\\
 %\textbullet $\,\,$ $ \mathcal{O}(z): -\frac{m_a g_a}{2m_p^2}e^{i(k_1+k_2)\cdot b_a}\hat\delta(k_1\cdot v_1+k_2\cdot v_1+\omega)\{(k_{1\rho}+k_{2\rho})v^{\mu}v^{\nu}+2\omega v_\rho\}\rightarrow
%\begin{minipage}[h]{0.08\linewidth}
	%\vspace{4pt}
	%\scalebox{1.8}{\includegraphics[width=\linewidth]{new_10.png}}
%\end{minipage}$\\
 %  \textbullet $\,\,$ $\mathcal{O}{(z^2)}:i\frac{m_a s_a}{2m_p} \,e^{ik\cdot b_a}\,\hat\delta(k\cdot v+\omega_1+\omega_2)  (\frac{1}{2} k_{\rho_1} k_{\rho_2} +\omega_1 k_{\rho_2}v_{\rho_1}+\omega_2 k_{\rho_1}v_{\rho_2}+\omega_1\omega_2 \eta_{\rho_1\rho_2})\rightarrow  \begin{minipage}[h]{0.08\linewidth}
	%\vspace{4pt}
	%\scalebox{1.6}{\includegraphics[width=\linewidth]{phi_z_2.png}}
%\end{minipage}$
%\end{itemize}
%\vspace{-1.6 cm}
%\newpage
%\subsection*{\underline{Vertex for gravitational field:\textcolor{red}{CITE ?}}}
%\begin{itemize}
   %  \textbullet $\,\,$ $\mathcal{O}{(z^0)}:- i\frac{m_a }{2 m_p} \,e^{ik\cdot b_a}\,\hat\delta(k\cdot v)v_a^\mu v_a^\nu \rightarrow  \begin{minipage}[h]{0.08\linewidth}
	%\vspace{4pt}
	%\scalebox{1.5}{\includegraphics[width=\linewidth]{h_z_0.png}}
%\end{minipage}$\\
 %\textbullet $\,\,$ $\mathcal{O}(z^1):\frac{m_a }{2 m_p}\,e^{ik\cdot b}\,\hat{\delta}(k\cdot v+\omega)(2\omega v^{(\mu}\delta^{\nu)}_\rho+v^\mu v^\nu k_{\rho})\rightarrow \begin{minipage}[h]{0.08\linewidth}
	%\vspace{4pt}
	%\scalebox{1.5}{\includegraphics[width=\linewidth]{h_z_1.png}}
%\end{minipage}$\\
% \textbullet $\,\,$ $\mathcal{O}{(z^2)}:i\frac{m_a }{m_p} \,e^{ik\cdot b_a}\,\hat\delta(k\cdot v+\omega_1+\omega_2) (\frac{1}{2} k_{\rho_1} k_{\rho_2} v^\mu v^\nu +\omega_1 k_{\rho_2}v^{(\mu}\delta^{\nu)}_{\rho_1}+\omega_2 k_{\rho_1}v^{(\mu}\delta^{\nu)}_{\rho_2}+\omega_1\omega_2 \delta^{(\mu}_{\rho_1}\delta^{\nu)}_{\rho_2})\rightarrow  \begin{minipage}[h]{0.08\linewidth}
	%\vspace{4pt}
	%\scalebox{1.5}{\includegraphics[width=\linewidth]{h_z_2.png}}
%\end{minipage}$
\section{Bootstrapping Post-Minkowskian (PM) physics from Post-Newtonian (PN) data: the idea, technology and example}
\label{ch3:sec3}
In this section, we briefly discuss the idea of bootstrapping of Post-Minkowskian (relativistic) observables from Post-Newtonian (non-relativistic) data \cite{Dlapa:2023hsl,Henn:2014qga},  which primarily based on differential equation technique for evaluating multiloop Feynman integrals \cite{Kotikov:1990kg,Gehrmann:1999as} and also illustrate the method with a specific two loop example relevant for our computations. In the upcoming sections, we widely encounter two loop integrals, which can be solved using IBP relations. In this section, we will present the results for the master integrals that we will use for our computations of the eikonal phase in the subsequent section.

\textbf{Identifying the integrands:} We encounter a general family of n-loop integrals that has the following form, which is essentially equivalent to computing $2\to 2$ scattering amplitude,
\begin{align}
           &\langle \,\textrm{final}\,|\mathbb{S}|\,\textrm{initial}\,\rangle\sim\mathbfcal{M}^{a_1\cdots a_n;\pm\cdots \pm}_{\alpha_1\cdots \alpha_n;\beta_1\cdots \beta_m}(|q|,\gamma)=\Bigg(\prod_{i=1}^n\int_{\ell_i}\frac{\hat\delta(\ell_i\cdot v_{a_i})}{(\pm \ell_i\cdot v_{\cancel{a}_i}+i\varepsilon)^{\alpha_i}}\Bigg)\frac{1}{D_1^{\beta _1 }D_2^{\beta_2}\cdots D_m^{\beta_m}}\,,
\end{align}
where $a_i$s are particle index label with $\cancel{1}=2,\,\cancel{2}=1$. $D_i$s are the graviton (scalar) propagators along with irreducible scalar products that form the numerators of the loops integrals and have the following form,
$$D_i=P_i^2,\,\,P_i=\tilde\Upsilon_{ij}\,l_j+\tilde\Upsilon_i\, q,\,\,\textrm{with},\,\,\tilde\Upsilon_{ij},\tilde\Upsilon_i\in(0,\pm 1)\,\,\forall\,\,\, 1\le i,j\le m$$
where $q$ is the external momentum and $v_{1,2}$ are the velocities of the two astrophysical objects. Furthermore, one would notice that each loop integral is supported by a Dirac-$\delta$ function, $\hat\delta(l_i\cdot v_{a_i})$. At this point, to utilize the multiloop collider physics techniques \cite{Weinzierl:2022eaz}, it is convenient to represent the $\delta$-functions by reverse-unitarity (sometimes called velocity-cut),
\begin{align}
    \begin{split}
       \frac{i}{(-1)^{s+1}} \hat\delta^{(s)}(D_i)=\frac{1}{(D_i+i\varepsilon)^{s+1}}-\frac{1}{(D_i-i\varepsilon)^{s+1}}
    \end{split}
\end{align}                 
which allows us to identify the $\delta$-functions as cut propagators and one can freely apply the Integration-by-parts (IBP) technique to evaluate the integrals.

\textbf{IBP reduction:}  After identifying the integrands, the next task is to find the IBP identities, first introduced in\cite{Chetyrkin:1981qh}, which help to identify the master integrals in a given family of Feynman integrals. IBP identities can be obtained by demanding that under dimensional regularisation, loop integrals of any total derivative identically vanish \cite{Weinzierl:2022eaz},
\begin{align}
    \int_{l_i}\frac{\partial}{\partial l_k^\mu}\Bigg(\frac{\eta^{\mu}}{(\pm l_i\cdot v_{\slashed{a}_i}+i\varepsilon)^{\alpha_i}\cdots D_1^{\beta_1}\,D_2^{\beta_2}\cdots \slashed{D}_1^{\gamma_1}\slashed{D}_2^{\gamma_2}\cdots}\Bigg)=0\,,\label{3.3a}
\end{align}
where, $l_k$s are loop momenta, $\eta^\mu$ is the linear combination of loop momenta and external momenta and $\slashed{D}$ are the `cut' propagators coming from the $\delta$-functions. The IBP reduction procedure generates a large number of homogeneous linear systems of equations of form,
\begin{align}
    \sum_k \gamma_k(|q|,d,\gamma)\mathbfcal{M}^{(\pm\pm)i}_{\alpha_1+\kappa_{k,1},\alpha_2+\kappa_{k,2},\cdots}(|q|,\gamma,d)=0\,.
\end{align}
However, all the equations are not independent but rather related by various symmetry relations. Ultimately, one needs to solve the equations in unique sectors. Fortunately, there are several publicly available packages for doing this task \textbf{LiteRed} \cite{Lee:2012cn,Lee:2013mka}, \textbf{Kira} \cite{Maierhofer:2017gsa}, \textbf{Reduze 2} \cite{vonManteuffel:2012np} and \textbf{FIRE6} \cite{Smirnov:2019qkx}. For our computations, we extensively use \textbf{LiteRed} .

\textbf{Master integrals and differential equations:}
One of the main tasks of the IBP procedure is to expand a general Feynman integral as a linear combination of master integrals. For example, let us take a look at the two-loop case. A general two-loop integral has the following form,
\begin{align}
    \begin{split}
        \mathbfcal{M}^{\pm,\pm,\textrm{2-loop}}_{\alpha_1,\alpha_2;\beta_1\cdots \beta_{5}}=\int_{l_{1,2}}\Bigg(\frac{\hat\delta^{(\gamma_1)}(\slashed D_1)\hat\delta^{(\gamma_2)}(\slashed D_2)}{(\pm D_1+i\varepsilon)^{\alpha_1}(\pm D_2+i\varepsilon)^{\alpha_2}}\Bigg)\frac{1}{D_3^{\beta_1}\cdots D_7^{\beta_5}}\label{3.4a}
    \end{split}
\end{align}
where, $\{D_i,\{\slashed D_k\}\}\in\Big(\ell_{1}\cdot v_1,\,\,\ell_{2}\cdot v_2,\,,\,\ell_{1}^2,\,\, \ell_{2}^2,\,\, (\ell_{1}+\ell_{2}-q)^2,\,\, (\ell_{1}-q)^2,\,\, (\ell_{2}-q)^2,\,\, \{\ell_{1}\cdot v_2,\,\, \ell_{2}\cdot v_1\}\Big) $. The set $\{D_i,\slashed D_k\}$ s are called basis function set and the number of independent basis functions (scalar-product) can be determined by noting the number of loop momenta and external momenta: $N_{\textrm{basis}}=\frac{L(L+1)}{2}+LE$, where $L$ and $E$ are the number of loop integrals and external momenta respectively. For the two-loop case, we have $2$ loop momenta and $3$ external momenta, hence the number of basis functions is, indeed, $N_{\textrm{basis}}^{2-\textrm{loop}}=9$. Once we have the basis function, we solve the IBP relation using \textbf{LiteRed}; after solving, one gets the unique master integral for this particular family of integrals. Now using built-in command \textbf{IBPReduce} \cite{Lee:2012cn,Lee:2013mka}, one can write the general two-loop integrals as,
\begin{align}
    \begin{split}
         \mathbfcal{M}^{\textrm{2-loop}}=\sum_{n=1}^{N_{m}} a_{n}(\gamma,d)\,\mathbfcal{M}^{\textrm{master}}_n(\gamma, d)\,.
    \end{split}
\end{align}
We face 16 independent master integrals (excluding the different signs of $i\varepsilon$-prescription, which can be derived further using appropriate field redefinition and invoking symmetry relations) after solving 17 unique sectors using \textbf{LiteRed}.

\textbf{Solving master integrals using differential equation:} In general, the master integrals do not factorize into the kinematic variable $\gamma$ and dimensional regularisation parameter $\epsilon$. In this case, it is difficult to find any closed-form expression. So, one needs to find the solutions order by order in $\epsilon$. For this purpose, the method of differential equations is very useful \cite{Remiddi:1997ny,Kotikov:1991pm, Henn:2013pwa}. We first define a column vector consisting of relevant master integrals and redefine the kinematic variable $\gamma$ to $x$ as $\gamma\to \frac{x^2+1}{2x}$,
\begin{align}
    \begin{split}
        \vec f(x,\epsilon)=\{\mathbfcal{M}_1^{\textrm{master}},\cdots, \mathbfcal{M}_{N_m}^{\textrm{master}} \}
    \end{split}
\end{align}
which satisfies the following differential equation,
\begin{align}
      \partial_x \vec f(x,\epsilon)=A(x,\epsilon)
     \vec f(x,\epsilon)\,.\label{3.7}
\end{align}
For all practical purposes, the differential equation is quite hard to solve. But it is convenient to find a suitable transformation \cite{Henn:2014qga},
\begin{align}
    \begin{split}
        \vec f(x,\epsilon)=\mathbb{T}(x,\epsilon)\vec{g}(x,\epsilon),
    \end{split}
\end{align}
such that the differential system can be brought to the $\epsilon$-form, which is easier to solve. Now, the system of differential equations has the form,
\begin{align}
    \begin{split}
        \partial_x \vec g(x,\epsilon)=\mathbb{T}^{-1}(A \mathbb{T}-\partial_x \mathbb{T})\vec g(x,\epsilon)\equiv \epsilon\, \mathbb{S}(x) \vec g(x,\epsilon)\,.\label{2.14a}
    \end{split}
\end{align}
The goal is to find the basis transformation matrix $\mathbb{T}(x,\epsilon)$ such that $\mathbb{T}(A \mathbb{T}-\partial_x \mathbb{T})=\epsilon \,\mathbb{S}(x)$. Note that $\epsilon$ is now completely factorized from the matrix $\mathbb{S}(x)\,.$, It was first conjectured by Henn \cite{Henn:2013pwa}. The $\mathbf {\epsilon}$-\textbf{factorization} can be systematically done using Lee's algorithm \cite{Lee:2014ioa} and is (semi-) automatized in the following packages \textbf{Fuchsia} \cite{Gituliar:2017vzm}, \textbf{epsilon} \cite{Prausa:2017ltv}, \textbf{Libra} \cite{Lee:2020zfb} and \textbf{CANONICA} \cite{Meyer:2017joq} and the matrix equation in $\bf{\epsilon}$-form has a general solution in terms of path-ordered exponential (Bootstrap equation), which can be solved systematically using \textbf{Libra} \cite{Lee:2020zfb},
\begin{align}
    \begin{split}
        \vec g(x,\epsilon)=\underbrace{\Bigg(\mathbfcal{P}e^{\epsilon\int_{x_0}^{x}S(x')\,dx'}\Bigg)}_{\mathbb{B}(x,x_0)}\vec g(x_0,\epsilon).\label{3.10a}
    \end{split}
\end{align}
Once we have \eqref{3.10a}, we need to use the boundary condition to fix $\vec g(x_0,\epsilon)\,.$ For relativistic two-body problem it is easy to find the solution of integration at `\textit{soft}' (near-static) limit which is at $\gamma\to 1^{+}\,(x\to 1^{-})$, which is nothing but the solution of the master integrals in Post-Newtonian (non-relativistic) regime: \textit{this is nothing but bootstrapping the complete relativistic integrals (PM-integrals) using the non-relativistic integrals (PN-integrals) \cite{Dlapa:2023hsl}.}

\par


\noindent
For illustration, we show the details of the computation in the \textbf{Sectors:\{0,0,0,0,1,1,1\}\&\{1,1,0,0,1,1,1\}}, where we encounter the following {\bf four} master integrals (the original basis is sometimes called the Laporta basis \cite{Laporta:2000dsw}). We take the following vector \footnote{For the sake of convenience, we will henceforth drop the subscript ``master'' from the $\mathbfcal{M}$'s, and the $|q|$ dependence can be recovered from dimensional analysis, since $|q|$ is the dimensionful scale.},
\begin{align}
\begin{split}
    \vec f(\gamma,\epsilon)=&\{\mathbfcal{M}_{0,0,0,0,1,1,1,\slashed 1,\slashed 1},\mathbfcal{M}_{0,0,0,0,2,1,1,\slashed 1,\slashed 1},\mathbfcal{M}_{0,0,0,0,1,2,1,\slashed 1,\slashed 1}, \mathbfcal{M}_{1,1,0,0,1,1,1,\slashed 1,\slashed 1}\}^{T},
    \end{split}
\end{align}
which satisfies, $\partial_x \vec f(x,\epsilon)=A(x,\epsilon)\vec f(x,\epsilon)$ with,
\begin{align}
      A(x,\epsilon)=  \left(
\begin{array}{cccc}
 \frac{3 \left(x^4+6 x^2+1\right) \epsilon -\left(x^2+1\right)^2}{x \left(x^4-1\right)} & -\frac{x^2-1}{2 \left(x^3+x\right)} & \frac{2 x (2 \epsilon +1)}{\epsilon -x^4 \epsilon } & 0 \\
 \frac{6 \left(x^2-1\right) \epsilon ^2}{x^3+x} & \frac{x^4 (\epsilon +1)+x^2 (6 \epsilon +2)+\epsilon +1}{x-x^5} & \frac{4 x (2 \epsilon +1)}{x^4-1} & 0 \\
 -\frac{6 \left(x^2-1\right) \epsilon ^2}{x^3+x} & \frac{\epsilon -x^2 \epsilon }{x^3+x} & \frac{\left(x^2-1\right) (2 \epsilon +1)}{x^3+x} & 0 \\
 0 & -\frac{4}{x^2-1} & 0 & \frac{2 \left(x^2+1\right)}{x-x^3} \\
\end{array}
\right)\,.\label{3.12}
\end{align}
The matrix equation can be brought to $\epsilon$-form (canonical form) using suitable basis transformation using the publicly available packages mentioned earlier. In the canonical basis/Uniform Transcendental (UT) basis, the differential equation takes the form, $\partial_x \vec g(x,\epsilon)=\epsilon \mathbb{S}(x)\vec g(x,\epsilon)$, where $\mathbb{S}(x)$ and the transformation matrix $\mathbb{T}(x,\epsilon)$ are given by,
\begin{align}
\begin{split}
   &\mathbb{T}(x,\epsilon)= \left(
\begin{array}{cccc}
 -\frac{7 x}{x^2-1} & \frac{x}{2-2 x^2} & \frac{3 x}{x^2-1} & 0 \\
 -\frac{6 x \epsilon }{x^2-1} & -\frac{5 x \epsilon }{x^2-1} & \frac{6 x \epsilon }{x^2-1} & 0 \\
 \frac{12 \left(x^2-1\right) \epsilon ^2}{2 x \epsilon +x} & \frac{2 \left(x^2-1\right) \epsilon ^2}{2 x \epsilon +x} & -\frac{12 x \epsilon ^2}{2 \epsilon +1} & 0 \\
 0 & 0 & 0 & \frac{2 x^2}{\left(x^2-1\right)^2} \\
\end{array}
\right),\quad\textrm{and}\\ &
\mathbb{S}(x)=
\left(
\begin{array}{cccc}
 \frac{9 \left(2 x^2+1\right)}{2 \left(x^3-x\right)} & \frac{2 x^2+3}{4 \left(x^3-x\right)} & -\frac{6 x}{x^2-1} & 0 \\
 -\frac{3 \left(2 x^2+5\right)}{x^3-x} & -\frac{6 x^2+5}{2 \left(x^3-x\right)} & \frac{12 x}{x^2-1} & 0 \\
 \frac{6}{x} & -\frac{1}{3 x} & -\frac{2}{x} & 0 \\
 \frac{12}{x} & \frac{10}{x} & -\frac{12}{x} & 0 \\
\end{array}
\right)=\sum_{i=1}^3\frac{\mathbb{S}_i}{x-x_i},\,(x_i=1,-1,0)
\end{split}
\end{align}
where the matrix residues take the form,
\begin{align}
\mathbb{S}_0=\left(\begin{array}{cccc}
         -  \frac{9}{2}  & -\frac{3}{4}& 0 &0  \\
            15 & \frac{5}{2} & 0 &0\\
            6&-\frac{1}{3}&-2&0\\
            12 & 10 & -12 &0 
     \end{array}
     \right)\,, \quad    
    \mathbb{S}_{\pm 1}=\left(\begin{array}{cccc}
           \frac{27}{4}  & \frac{5}{8}& -3 &0  \\
            -\frac{21}{2} & -\frac{11}{4} & 6 &0\\
            0&0&0&0\\
            0 & 0 & 0 &0 
     \end{array}
     \right)\,.  \label{2.17} 
\end{align}
In this fuchsified form, the equation looks like this,
\begin{align}
    \begin{split}
        \partial_x\vec g(x,\epsilon)=\epsilon \Bigg(\frac{S_0}{x}+\frac{S_1}{x+1}+\frac{S_{-1}}{x-1}\Bigg)\vec g(x,\epsilon)\label{3.15}
    \end{split}
\end{align}
%First we solve the \eqref{3.15} perturbatively as,
%\begin{align}
  %  \vec g(x,\epsilon)=\Bigg(1+\epsilon\int_{x_0}^{x}\mathbb{S}(x')dx'+\epsilon^2\int_{x_0}^x dx'\mathbb{S}(x')\int_{x_0}^{x'}dx'' \mathbb{S}(x'')+\mathcal{O}(\epsilon^3)\Bigg)\vec g(x_0,\epsilon)
%\end{align}
%Upto $\mathcal{O}(\epsilon)$ the solution is given by,
%\begin{align}
  %  \begin{split}
   %     \vec g(x,\epsilon)=\vec g(x_0,\epsilon)+\epsilon\Bigg(S_0\int_{x_0}^{x}d\log(x)+S_1\int_{x_0}^{x}d\log(x+1)+S_{-1}\int_{x_0}^xd\log(x-1)\Bigg)\vec g(x_0,\epsilon)\label{3.21m}
  %  \end{split}
%\end{align}
which has the solution in terms of a path ordered exponential as mentioned in \eqref{3.10a}. Finally, in  the Laporta basis, the formal solution takes the form,
\begin{align}
    \begin{split}
        \vec f(x,\epsilon)=\mathbb{T}(x,\epsilon)\,\Bigg(\mathbfcal{P}e^{\epsilon\int_{1^{-}}^x \,\mathbb{S}(x')\,dx'}\Bigg)\,\mathbb{T}^{-1}(1^{-},\epsilon) \,\vec f(1^{-},\epsilon)\,.\label{3.16}
    \end{split}
\end{align}
Taking the boundary limit is a bit tricky, and one needs to take the limit very carefully.First we expand the $\mathbb T^{-1}(x_0=1-v_{\infty})$ matrix  and  also expand the boundary vector $\vec f(x_0,\epsilon)$ around $v_\infty=0$ using the method of regions. For generic two-loop integrals, we will have both the contribution from potential and the radiation region \cite{Dlapa:2023hsl}. \textit{ However, at 3PM, the full conservative dynamics of the boundary integrals are captured by the contribution from the potential region only \cite{Dlapa:2023hsl}.}\\

\textbf{\textit{Potential region:}} In momentum space, the potential region is characterized by the following loop momenta scaling: $l_{i}\sim |\boldsymbol{q}|(v_{\infty},1)$ and, the velocity of black holes are parametrized as \cite{Dlapa:2023hsl}, 
\begin{align}
    \begin{split}
        v_2=(1,0,0,0),\,\, v_1=(1,v_{\infty},0,0).\label{2.23a}
    \end{split}
\end{align}
Using the above-mentioned scaling, the integral is reduced to,
\begin{align}
    \begin{split}
     \mathbfcal{M}\Big|_{v_{\infty}\to 0}^{\textrm{pot.}}\sim    v_{\infty}^{-\sum n_i}\Bigg(\prod_{i}\int_{\vec \ell_i}\frac{1}{(\pm\vec \ell_i\cdot \boldsymbol{n}+i\varepsilon)^{n_i}}\Bigg)\frac{1}{\boldsymbol {D_1}^{m_1}\cdots\boldsymbol{D_k}^{m_k}},\,\boldsymbol D\equiv \textrm{spatial part of the squared propagators.}
    \end{split}
\end{align}
For our case, we are interested in the following two loop integrals in potential regions.
\begin{align}
    \begin{split}
\mathbfcal{M}^{\textrm{2-loop},(\mp)}_{n_1,n_2;m_1,..,m_5}\Big|_{v_{\infty}\to 0}^{\textrm{pot.}}&\sim\int_{\ell_{1},\ell_{2}}\frac{\hat\delta(\ell_{1}^0)\hat\delta(\ell_{2}^0-v_{\infty}\ell_{2}^{(x)})}{(\ell_{1}^0-v_{\infty}\ell_{1}^{(x)}+i\varepsilon)^{n_1}(\ell_{2}^0+i\varepsilon)^{n_2}}\frac{1}{\boldsymbol{D_1}^{m_1}\cdots\boldsymbol{D_5}^{m_5} }+\mathcal{O}(v_{\infty}^{-1})\\ &
=v_{\infty}^{-n_1-n_2}\int_{\boldsymbol{\ell_1},\boldsymbol{\ell_2}}\frac{1}{(-\boldsymbol{\ell_1}\cdot \boldsymbol{n}+i\varepsilon)^{n_1}(\boldsymbol{\ell_2}\cdot \boldsymbol{n}+i\varepsilon)^{n_2}}\frac{1}{\boldsymbol{D_1}^{m_1}\cdots\boldsymbol{D_5}^{m_5} }+\mathcal{O}(v_{\infty}^{1-n_1-n_2})\,.
    \end{split}
\end{align}
The details of the computation of potential mode master integrals are given in Appendix~(\ref{ch3:app:B}).\par


\noindent
Now, the solution \eqref{3.16} to the differential equation to the differential equation can be presented in terms of multiple
polylogarithms (MPLs) \cite{Remiddi:1999ew,goncha},
\begin{align}
    G(a_1,a_2,\cdots a_n;z)=\int_{0}^z \frac{dt}{t-a_1}G(a_2,\cdots,a_n;t),\,\, G(\,;z)=1,\,\,\forall a_i,z\in \mathbb{C}\,.
\end{align}
If all $a_i's$ are zero, 
\begin{align}
    G(\vec 0_n;z)=\frac{1}{n!}\,\log^{n}(z)
\end{align}
and, for equal $a_i=a$ one can obtain,
\begin{align}
    G(\vec a_n,z)=\frac{1}{n! }\log^{n}\Big(1-\frac{z}{a}\Big),\,\,\forall a\ne 0\,.
\end{align}
Apart from ordinary logarithms, MPLs also contain classical polylogarithms, which are defined as,
\begin{align}
    \begin{split}
        G(\vec 0_{n-1},1;z)=- \textrm{{\bf Li}}_2(z)\,.
    \end{split}
\end{align}
Now, after collecting the ingredients discussed above and suitably regularizing the logarithmic divergences, we obtain the solutions for the Laporta master integrals needed in our subsequent computations, and we present the epsilon-form matrix $\mathbb{S}(x)$ in Appendix~\eqref{ch3:app:C} \footnote{{We match all the listed master integrals in \eqref{MasterM} with those of \cite{Jakobsen:2022fcj}. At leading order, up to which the integrals are given in \cite{Jakobsen:2022fcj}, the results agree perfectly with ours except for $\mathbfcal{M}_{0,0;1,1,2,1,1}$. However, we have cross-checked our result by explicitly matching it to the boundary value. S.G. would like to thank Gustav Jacobsen for confirming this point. }}.
%\usepackage{tikz,lipsum,lmodern}
%\usepackage[most]{tcolorbox}
%\newpage
%\newpage
%\newpage
%\newgeometry{top=2cm, bottom=3cm} % Change top and bottom margins
\begin{adjustwidth}{0cm}{0cm}
\begin{tcolorbox}[thesisresultbox, title=\textit{Master Integrals}, left=2.1mm, right=2.1mm]
\restorethesisbodyformat
\resizebox{0.955\linewidth}{!}{$\begin{aligned}
&
(-q^2)^{2\epsilon}\mathbfcal{M}_{0,0;0,0,1,1,1}=-\frac{x \log (x)}{32 \pi ^2 \left(x^2-1\right) \epsilon }+\frac{3x \left(\log^2(x)+\textbf{Li}_2(1-x^2)\right)}{32 \pi ^2 \left(x^2-1\right)}-\frac{ x \epsilon}{192 \pi ^2 \left(x^2-1\right)}  \Bigg[-7 \pi ^2 \log(x)\\ &\hspace{1.4 cm}+216 \boldsymbol{\mathscr{J}}(\{-1,-1,0\},x)-216 \boldsymbol{\mathscr{J}}(\{-1,0,0\},x)+216 \boldsymbol{\mathscr{J}}(\{-1,1,0\},x)-216 \boldsymbol{\mathscr{J}}(\{0,-1,0\},x)\\ &\hspace{-0.3 cm}+120 \boldsymbol{\mathscr{J}}(\{0,0,0\},x)-216 \boldsymbol{\mathscr{J}}(\{0,1,0\},x)+216 \boldsymbol{\mathscr{J}}(\{1,-1,0\},x)-216 \boldsymbol{\mathscr{J}}(\{1,0,0\},x)+216 \boldsymbol{\mathscr{J}}(\{1,1,0\},x)\Bigg]+\mathcal{O}(\epsilon^2)\,,\\ &
(-q^2)^{2\epsilon+1}\mathbfcal{M}_{0,0;0,0,1,2,1}=-\frac{x^2+1}{64 \pi ^2 x}+\frac{x^2+(1-x^2) \log (x)+1}{32 \pi ^2 x}\epsilon+\frac{\epsilon^2}{384 \pi ^2 x}\Big[120 \left(x^2-1\right) (\textbf{Li}_2(1-x)-\textbf{Li}_2(-x))\\ &\hspace{2.3 cm}-3 \left(8+\pi ^2\right) x^2+12 \log (x) \left(\left(9 x^2-1\right) \log (x)+\left(x^2-1\right) (2-10 \log (x+1))\right)+17 \pi ^2-24\Big]\,,\\ &
(-q^2)^{2\epsilon+1}\mathbfcal{M}_{0,0;0,0,2,1,1}=-\frac{x \log (x)}{16 \pi ^2 \left(x^2-1\right)}-\epsilon  \frac{x   \left(\log^2(x)+\textbf{Li}_2(1-x^2)\right)}{16 \pi ^2 \left(x^2-1\right)}+\mathcal{O}(\epsilon^2)\,,\\ &
     (-q^2)^{2\epsilon+1} \textcolor{black}{\mathbfcal{M}^{(++)}_{1,1;0,0,1,1,1}}=\frac{x^2}{8 \pi ^2 \left(x^2-1\right)^2 \epsilon ^2}-\frac{x^2 \left(\pi ^2-6 \log ^2(x)\right)}{48 \left(\pi ^2 \left(x^2-1\right)^2\right)}+\frac{x^2 \epsilon }{24 \pi ^2 \left(x^2-1\right)^2}\Big[ (-12 (\boldsymbol{\mathscr{J}}(\{0,-1,0\},x)\\ &\hspace{0.8cm}-\boldsymbol{\mathscr{J}}(\{0,0,0\},x)+\boldsymbol{\mathscr{J}}(\{0,1,0\},x))-32 \zeta (3))\Big]+O\left(\epsilon ^2\right)=  (-q^2)^{2\epsilon+1}\frac{1}{2}\mathbfcal{M}^{(+-)}_{1,1;0,0,1,1,1}\,,
 \\     &(-q^2)^{1+2\epsilon}\mathbfcal{M}^{(+\pm)}_{0,0;1,1,0,1,1}=\frac{1}{(4\pi)^{3-2\epsilon}}\frac{\Gamma^4(1/2-\epsilon)\Gamma^2(1/2+\epsilon)}{\Gamma^2(1-2\epsilon)},\\ &
(-q^2)^{2+2\epsilon}\mathbfcal{M}^{(+\pm)}_{0,0;1,1,1,1,1}=\frac{x \log (x)}{16 \pi ^2 (x^2-1) \epsilon }+\frac{\Big(\log^2(x)+\textbf{Li}_2(1-x^2)\Big)}{16\pi^2(x^2-1)}-\frac{\epsilon }{96 \pi ^2 \left(x^2-1\right)} \Bigg(24 x (3 \boldsymbol{\mathscr{J}}(\{-1,-1,0\},x)\\&\hspace{2.4 cm} -3 \boldsymbol{\mathscr{J}}(\{-1,0,0\},x)+3 \boldsymbol{\mathscr{J}}(\{-1,1,0\},x)-3 \boldsymbol{\mathscr{J}}(\{0,-1,0\},x)+\boldsymbol{\mathscr{J}}(\{0,0,0\},x)-3 \boldsymbol{\mathscr{J}}(\{0,1,0\},x)\\ & \hspace{2.4 cm}+3 \boldsymbol{\mathscr{J}}(\{1,-1,0\},x)-3 \boldsymbol{\mathscr{J}}(\{1,0,0\},x)+3 \boldsymbol{\mathscr{J}}(\{1,1,0\},x))-5 \pi^2 x \log (x)\Bigg)+\mathcal{O}(\epsilon^2)\,,\\ &
(-q^2)^{3+2\epsilon}\mathbfcal{M}^{(+\pm)}_{0,0;1,1,2,1,1}=\frac{x \left(x^4-4 x^2 \log (x)-1\right)}{16 \pi ^2 \left(x^2-1\right)^3 \epsilon }+\frac{x}{16 \pi ^2 \left(x^2-1\right)^3}\Bigg[5 \left(x^4-1\right)-2 \left(x^4+4 x^2+1\right) \log (x)\\ & \hspace{0.2 cm}+4x^2 \Big(\log^2 x+\textbf{Li}_2(1-x^2)\Big)\Bigg]+\epsilon\frac{x}{96 \pi ^2 \left(x^2-1\right)^3}\Big[288 x^2\Bigg( \boldsymbol{\mathscr{J}}(\{-1,-1,0\},x)- \boldsymbol{\mathscr{J}}(\{-1,0,0\},x)+ \boldsymbol{\mathscr{J}}(\{-1,1,0\},x)\\ &\hspace{0.2 cm}- \boldsymbol{\mathscr{J}}(\{0,-1,0\},x)+ \frac{1}{3} \boldsymbol{\mathscr{J}}(\{0,0,0\},x)- \boldsymbol{\mathscr{J}}(\{0,1,0\},x)+ \boldsymbol{\mathscr{J}}(\{1,-1,0\},x)- \boldsymbol{\mathscr{J}}(\{1,0,0\},x)+ \boldsymbol{\mathscr{J}}(\{1,1,0\},x)\Bigg)\\ &\hspace{0.2 cm}+24 \left(x^4+12 x^2+1\right) \textbf{Li}_2(1-x)-24 \left(x^4+12 x^2+1\right) \textbf{Li}_2(-x)+54 \left(x^4-1\right)+36 x^4 \log ^2(x)-24 x^4 \log (x) \log (x+1)\\ &\hspace{0.2 cm}-96 x^4 \log (x)+\pi ^2 \left(3 x^2 \left(x^2-8\right)-7\right)+144 x^2 \log ^2(x)-288 x^2 \log (x) \log (x+1)-20 \pi ^2 x^2 \log (x)\\ &\hspace{1.2 cm}+24 x^2 \log (x)-12 \log ^2(x)-24 \log (x) \log (x+1)-96 \log (x)\Big]+\mathcal{O}(\epsilon^2)\,,\\ &
 (-q^2)^{2\epsilon} \textcolor{black}{\mathbfcal{M}^{(+\pm)}_{0,0;0,1,1,0,1}}=0,\,\,
 (-q^2)^{2\epsilon} \textcolor{black}{\mathbfcal{M}^{(+\pm)}_{0,1;0,0,1,1,1}}=-\frac{i\,x}{32\pi \epsilon (x^2-1)}+\mathcal{O}(\epsilon^0)
\end{aligned}$}\label{MasterM}
\end{tcolorbox}
\restorethesisbodyformat
\end{adjustwidth}
\vspace{0.15cm}
%\restoregeometry
The iterative integral in \eqref{MasterM} is defined as,
\begin{align}
    \boldsymbol{\mathscr{J}}(\{i,j,k\},x)\equiv \int_{1^-}^x \frac{dt}{t-i}\int_{1^-}^{t}\frac{dt'}{t'-j}\int_{1^-}^{t'}\frac{dt''}{t''-k} \boldsymbol{\mathscr{J}}(\{\},t'');~~~~~~\boldsymbol{\mathscr{J}}(\{\},t'')=1,\label{3.43k}
\end{align}
and, the relevant 3-point iterative UT integrals are defined in Appendix~\eqref{ch3:app:D}\footnote{Note that the integrals defined in \eqref{3.43k} are not exactly Goncharov MPLs (which is defined for iterative integrals with lower limit `$0$') because here the lower limit is `$1$'. However one can write \eqref{3.43k} as a difference of Goncharov MPLs.}. 
%\vspace{0.5cm}
In \eqref{MasterM} we extensively use the following identity,
\begin{align}
    \textbf{Li}_2(-x)-\textbf{Li}_2(1-x)=-\frac{\pi^2}{12}-\log(x)\log(1+x)-\frac{1}{2}\textbf{Li}_2(1-x^2),
\end{align}
which can be derived from the fundamental identities of dilogarithm, namely \textit{Reflection formula and Abel's duplication formula}  \cite{goncha}. From now on, we make the $(++)$ sign of the last master integral of \eqref{MasterM} implicit for all computations. All the derivations are done by implicitly making the $i\varepsilon$ prescription positive for all linear propagators using proper field redefinition.\par
Now, we provide a bit of detail regarding the computation of master integrals. Let's focus on the following $\mathbfcal{M}_{0,1;0,0,1,1,1}$ and $\mathbfcal{M}_{0,0;0,1,0,1,1}$ that we do not include in the differential system mentioned in \eqref{3.12}. Fortunately, the first integral can be found separately using the differential equation method as it closes itself, and the second one identically vanishes because of the tadpole nature of $\ell_{1}$ loop integral,
\begin{align}
    \begin{split}
     &   \frac{d}{d\gamma}\mathbfcal{M}_{0, 1; 0, 0, 1, 1, 1, \slashed 1, \slashed 1}(\gamma,\epsilon)=-\frac{\gamma}{\gamma^2-1}\mathbfcal{M}_{0, 1; 0, 0, 1, 1, 1, \slashed 1, \slashed 1}(\gamma,\epsilon)+\frac{1}{\gamma^2-1}{\mathbfcal{M}_{0, 0; 0, 0, 1, 1, 1, \slashed 2, \slashed 1}}\,.
     \label{3.29}
    \end{split}
\end{align}
It appears that the master $\mathbfcal{M}_{0, 1; 0, 0, 1, 1, 1, \slashed 1, \slashed 1}(\gamma,\epsilon)$ does not close, but we now sketch that the master $\mathbfcal{M}_{0, 0; 0, 0, 1, 1, 1, \slashed 2, \slashed 1}$, that appears in the second term in \eqref{3.29} does not contribute and can be proved by noting that the integral has a flip symmetry $\ell_i\to -\ell_i,\,q\to -q$.
\begin{align}
\begin{split}
    \mathbfcal{M}_{0, 0; 0, 0, 1, 1, 1, \slashed 2, \slashed 1}(|q|,\gamma)&\propto \int_{\ell_{1},\ell_{2}}\frac{\hat\delta(\ell_{2}\cdot v_1)}{\ell_{1}^2 \ell_{2}^2 (\ell_{1}+\ell_{2}-q)^2}\Bigg(\frac{1}{(\ell_{1}\cdot v_2+i\varepsilon)^2}-\frac{1}{(\ell_{1}\cdot v_2-i\varepsilon)^2}\Bigg)\,,\\ &\xrightarrow[l_i\to -\ell_{1}]{q\to -q}\int_{\ell_{1},\ell_{2}}\frac{\hat\delta(\ell_{2}\cdot v_1)}{\ell_{1}^2 \ell_{2}^2 (\ell_{1}+\ell_{2}-q)^2}\Bigg(\frac{1}{(\ell_{1}\cdot v_2-i\varepsilon)^2}-\frac{1}{(\ell_{1}\cdot v_2+i\varepsilon)^2}\Bigg)\,,\\ &
    =-  \mathbfcal{M}_{0, 0; 0, 0, 1, 1, 1, \slashed 2, \slashed 1}\implies   \mathbfcal{M}_{0, 0; 0, 0, 1, 1, 1, \slashed 2, \slashed 1}\to 0\,.
    \end{split}
\end{align}
Therefore, the solution can be written straightforwardly and is given by,
\begin{align}
    \begin{split}
     &   \mathbfcal{M}_{0, 1; 0, 0, 1, 1, 1}(\gamma,\epsilon)=\frac{C_1(\epsilon,\gamma_0)}{\sqrt{\gamma^2-1}},
     \\ &
   \mathbfcal{M}_{0,0;0,1,0,1,1}(\gamma,\epsilon)  =0. \label{MasterC}
    \end{split}
\end{align}
To fix the boundary data, we evaluate the integrals in the following limit,
\begin{align}
    \begin{split}
   \lim_{\gamma\to 1^{++}}     \sqrt{\gamma^2-1} \,\mathbfcal{M}_{0, 1; 0, 0, 1, 1, 1, \slashed 1, \slashed 1}^{(++)}(\gamma,\epsilon)&=\lim_{\gamma\to 1^{+}}\sqrt{\gamma^2-1}\int_{\ell_{1},\ell_{2}}\frac{{\hat\delta(\ell_{1}^{(0)})\hat\delta(\gamma \ell_{2}^{0}-\gamma\beta \ell_{2}^{(x)})}}{(\ell_{2}^{(0)}+i\varepsilon)(\ell_{1}+\ell_{2}-q)^2(\ell_{1}-q)^2(\ell_{2}-q)^2}\,,\\ &
   =-\int_{\boldsymbol{\ell_1},\boldsymbol{\ell_2}}\frac{1}{(\boldsymbol{\ell_2}\cdot \boldsymbol{n}+ i \delta)\,\boldsymbol{\ell_1}^2\,\boldsymbol{\ell_2}^2 \,(\boldsymbol{\ell_1}+\boldsymbol{\ell_2}-\boldsymbol{q})^2}\,,\\ &
   %=\frac{2^{2 \epsilon -3} \pi ^{\epsilon -\frac{3}{2}} \Gamma \left(\frac{1}{2}-\epsilon \right)^2 \Gamma \left(\epsilon +\frac{1}{2}\right)}{\Gamma (1-2 \epsilon )}\int_{\boldsymbol{\ell_2}}\frac{1}{(\boldsymbol{\ell_2}\cdot \boldsymbol{n}\pm i0^{+})\,\boldsymbol{\ell_2}^2(\boldsymbol{\ell_2}-\boldsymbol{q})^{2(\epsilon+1/2)}}\\ &
   =\frac{i 4^{2 \epsilon -3} \pi ^{2 \epsilon -\frac{5}{2}} \Gamma \left(\frac{1}{2}-2 \epsilon \right) \Gamma (-\epsilon ) \Gamma \left(2 \epsilon +\frac{1}{2}\right)
   \Gamma^2(\frac{1}{2}-\epsilon)}{\Gamma \left(\frac{1}{2}-3 \epsilon \right) \Gamma(1-2\epsilon)}\frac{1}{(-q^2)^{1/2+2\epsilon}}\,.\label{3.37h}
    \end{split}
\end{align}
The appearance of negative sign in \eqref{3.37h} can be justified by noticing that $\mathbfcal{M}_{0,1;0,0,1,1,1}^{(+)}(+i\varepsilon)=-\mathbfcal{M}_{0,1;0,0,1,1,1}^{(-)}(-i\varepsilon)$, which can be achieved by momentum flip $\ell_i\to -\ell_i,\,q\to -q$. 
Now, we can extract the constant $C_1$ mentioned in \eqref{MasterC} and is given by,
\begin{align}
    \begin{split}
        C_1(\gamma_0\to 1^{+},\epsilon)=\frac{i\times 4^{2 \epsilon -3} \pi ^{2 \epsilon -\frac{5}{2}} \Gamma \left(\frac{1}{2}-2 \epsilon \right) \Gamma (-\epsilon ) \Gamma \left(2 \epsilon +\frac{1}{2}\right)
   \Gamma^2(\frac{1}{2}-\epsilon)}{\Gamma \left(\frac{1}{2}-3 \epsilon \right) \Gamma(1-2\epsilon)}\,.
    \end{split}
\end{align}
Therefore, the master integral is given by,
\begin{align}
    \begin{split}
       (-q^2)^{1/2+2\epsilon} \mathbfcal{M}_{0, 1; 0, 0, 1, 1, 1}^{(++)}(|q|,\gamma,\epsilon)&=\frac{i\times 4^{2 \epsilon -3} \pi ^{2 \epsilon -\frac{5}{2}}\Gamma \left(\frac{1}{2}-2 \epsilon \right) \Gamma (-\epsilon ) \Gamma \left(2 \epsilon +\frac{1}{2}\right)
   \Gamma^2(\frac{1}{2}-\epsilon)}{\sqrt{\gamma^2-1}\Gamma \left(\frac{1}{2}-3 \epsilon \right) \Gamma(1-2\epsilon)}\,,\\ &
       =-\frac{i}{64\pi \epsilon\sqrt{\gamma^2-1}}+\mathcal{O}(\epsilon^0)=  (-q^2)^{1/2+2\epsilon} \mathbfcal{M}_{0, 1; 0, 0, 1, 1, 1}^{(+-)}\,.
    \end{split}
\end{align}
\par 
%For completeness we list down full set of master integrals,
%\begin{align}
   %%&\hspace{0cm}(-q^2)^{1+2\epsilon}\mathbfcal{M}^{(+\pm)}_{0,0;1,1,0,1,1}=\frac{1}{(4\pi)^{3-2\epsilon}}\frac{\Gamma^4(1/2-\epsilon)\Gamma^2(1/2+\epsilon)}{\Gamma^2(1-2\epsilon)},\\ &
%\hspace{0cm}(-q^2)^{2+2\epsilon}\mathbfcal{M}^{(+\pm)}_{0,0;1,1,1,1,1}=\frac{x \log (x)}{16 \pi ^2 (x^2-1) \epsilon }+\frac{1}{16\pi^2(x^2-1)}\Big(\log^2(x)+\textbf{Li}_2(1-x^2)\Big)-\frac{\epsilon }{96 \pi ^2 \left(x^2-1\right)} \Bigg(24 x (3 \boldsymbol{\mathscr{J}}(\{-1,-1,0\},x)\\&\hspace{2.4 cm} -3 \boldsymbol{\mathscr{J}}(\{-1,0,0\},x)+3 \boldsymbol{\mathscr{J}}(\{-1,1,0\},x)-3 \boldsymbol{\mathscr{J}}(\{0,-1,0\},x)+\boldsymbol{\mathscr{J}}(\{0,0,0\},x)-3 \boldsymbol{\mathscr{J}}(\{0,1,0\},x)\\ & \hspace{2.4 cm}+3 \boldsymbol{\mathscr{J}}(\{1,-1,0\},x)-3 \boldsymbol{\mathscr{J}}(\{1,0,0\},x)+3 \boldsymbol{\mathscr{J}}(\{1,1,0\},x))-5 \pi^2 x \log (x)\Bigg)+\mathcal{O}(\epsilon^2)\\ &
%\hspace{0cm}(-q^2)^{3+2\epsilon}\mathbfcal{M}^{(+\pm)}_{0,0;1,1,2,1,1}=\frac{x \left(x^4-4 x^2 \log (x)-1\right)}{16 \pi ^2 \left(x^2-1\right)^3 \epsilon }+\frac{x}{16 \pi ^2 \left(x^2-1\right)^3}\Bigg[5 \left(x^4-1\right)-2 \left(x^4+4 x^2+1\right) \log (x)\\ & \hspace{2.4 cm}+4x^2 \Big(\log^2 x+\textbf{Li}_2(1-x^2)\Big)\Bigg]+\epsilon\frac{x}{96 \pi ^2 \left(x^2-1\right)^3}\Bigg[288 x^2\Bigg( \boldsymbol{\mathscr{J}}(\{-1,-1,0\},x)- \boldsymbol{\mathscr{J}}(\{-1,0,0\},x)\\ &\hspace{2.4 cm}+ \boldsymbol{\mathscr{J}}(\{-1,1,0\},x)- \boldsymbol{\mathscr{J}}(\{0,-1,0\},x)+ \frac{1}{3} \boldsymbol{\mathscr{J}}(\{0,0,0\},x)- \boldsymbol{\mathscr{J}}(\{0,1,0\},x)+ \boldsymbol{\mathscr{J}}(\{1,-1,0\},x)\\ &\hspace{2.4 cm}- \boldsymbol{\mathscr{J}}(\{1,0,0\},x)+ \boldsymbol{\mathscr{J}}(\{1,1,0\},x)\Bigg)+24 \left(x^4+12 x^2+1\right) \textbf{Li}_2(1-x)-24 \left(x^4+12 x^2+1\right) \textbf{Li}_2(-x)\\ &\hspace{2.4 cm}+54 \left(x^4-1\right)+36 x^4 \log ^2(x)-24 x^4 \log (x) \log (x+1)-96 x^4 \log (x)+\pi ^2 \left(3 x^2 \left(x^2-8\right)-7\right)\\ &\hspace{2.4 cm}+144 x^2 \log ^2(x)-288 x^2 \log (x) \log (x+1)-20 \pi ^2 x^2 \log (x)+24 x^2 \log (x)-12 \log ^2(x)\\ & \hspace{2.4 cm}-24 \log (x) \log (x+1)-96 \log (x)+\mathcal{O}(\epsilon^2)\label{3.42i}
 %   \end{split}
%\end{align}
At two loops we encounter another family of integrals defined in \eqref{3.4a} with basis functions,
 $\{D_i,\slashed D_k\}\in\Big(\ell_{1}\cdot v_1,\,\,\ell_{2}\cdot v_1,\,,\,\ell_{1}^2,\,\, \ell_{2}^2,\,\, (\ell_{1}+\ell_{2}-q)^2,\,\, (\ell_{1}-q)^2,\,\, (\ell_{2}-q)^2,\,\, \ell_{1}\cdot v_2,\,\, \ell_{2}\cdot v_2\Big) $, where the last two entries are the cut propagators. They are just `potential region' counterparts of the $\mathbfcal{M}$-type master integrals discussed earlier. After carrying out the  IBP reduction, we found 8 master integrals, and they are completely factorized in $\epsilon$ and $\gamma$, and we do not need differential equation techniques to solve them. Of these, we will need only two for our subsequent computations, and we list the results for those two below:
 \begin{align}
     \begin{split}
& (-q^2)^{2\epsilon}\mathbfcal{L}_{0,0;0,0,1,1,1}\equiv \frac{1}{(4\pi)^{3-2\epsilon}}\frac{\Gamma^3(1/2-\epsilon)\Gamma(2\epsilon)}{\Gamma(3/2-3\epsilon)}\,,\\ &
(-q^2)^{2\epsilon+1}\mathbfcal{L}^{(++)}_{1,1;0,0,1,1,1}=(-q^2)^{2\epsilon+1}2\mathbfcal{L}^{(+-)}_{1,1;0,0,1,1,1}\equiv \frac{1}{(4\pi)^{2-2\epsilon}}\frac{\Gamma^3(-\epsilon)\Gamma(1+2\epsilon)}{3(\gamma^2-1)\Gamma(-3\epsilon)}\,. \label{MasterL}
\end{split}
 \end{align}
At one loop, we encounter the following family of integral (we always look for the case when $\gamma=1$),
\begin{align}
    \begin{split}
        \mathbfcal{G}^{\pm,\textrm{1-loop}}_{\alpha_1;\beta_1 \beta_{2}}=\int_{l_{1,2}}\frac{\hat\delta^{(\gamma_1)}(\slashed D_1)}{(\pm D_1+i\varepsilon)^{\alpha_1}}\frac{1}{D_2^{\beta_1}D_3^{\beta_2}}\label{3.41a}
    \end{split}
\end{align}
with $\{D_i\,\lceil\slashed D_k\rceil\}\in \Big(\ell_{1}\cdot v_1,\,\ell_{1}^2,\,(\ell_{1}-q)^2,\,\lceil \ell_{1}\cdot v_2\rceil\Big)$. After IBP reduction we found two master integrals: ($\mathbfcal{G}_{0;1,1},\,\mathbfcal{G}_{1;1,1}$). We will require only one of them, and it is given by,
\begin{align}
    \begin{split}
        &(-q^2)^{\epsilon+1/2}\mathbfcal{G}_{0;1,1}=\frac{2^{2 \epsilon -3} \pi ^{\epsilon -\frac{3}{2}} \Gamma \left(\frac{1}{2}-\epsilon \right)^2 \Gamma \left(\epsilon +\frac{1}{2}\right)}{\Gamma (1-2 \epsilon )}\,. \label{3.4aa}
    \end{split}
\end{align}
{\color{black}
\textbf{\textit{Comment on the \(i\varepsilon\) prescription for the bulk propagators:}}

In the computation of the eikonal phase, and in its relation to classical observables such as the impulse and the waveform, the choice of \(i\varepsilon\) prescription plays an important role. The \(i\varepsilon\) is not merely a regulator; it specifies how the poles of the propagators are bypassed and therefore determines the boundary conditions of the corresponding classical solution. In the conservative sector, which is naturally associated with the elastic eikonal phase, one uses the standard in-out prescription with Feynman propagators,
\begin{equation}
    \frac{1}{P^2+i\varepsilon}\, .
\end{equation}
The conservative contribution is then obtained from the real part of the in-out effective action or, equivalently, from the real part of the corresponding scattering amplitude. This prescription is appropriate for time-symmetric observables such as the conservative scattering angle and the conservative part of the eikonal phase.
For causal observables, such as the full impulse and the gravitational waveform, the situation is more subtle. These observables depend on the physical radiation emitted during the scattering process and therefore require retarded boundary conditions in their fully causal formulation. In the worldline QFT description this is implemented through the Schwinger--Keldysh, or in-in, formalism. The corresponding graviton propagators are retarded or advanced propagators, whose momentum-space prescription may be written schematically as
\begin{equation}
    P^2
    \;\longrightarrow\;
    (P^0\pm i\varepsilon)^2-\mathbf{P}^2 ,
\end{equation}
where the sign selects the retarded or advanced Green's function, depending on the Fourier-transform convention.

However, for the conservative part of the impulse at two-loop order, it is sufficient to use the Feynman prescription for the bulk propagators. In this case, diagrams in which a single internal propagator goes on shell are associated with radiative, dissipative dynamics, rather than with the conservative impulse. Thus the conservative two-loop impulse may still be extracted from the Feynman-prescribed in-out amplitude, after taking the appropriate real, time-symmetric contribution.

Starting at three loops and beyond, the distinction becomes more delicate. Multiple internal propagators can go on shell simultaneously, and such multi-cut contributions may also enter the conservative dynamics. In that regime, the relation between the Feynman prescription, the retarded prescription, and the separation into conservative and radiative sectors must be handled with greater care. Equivalently, the dissipative part of the impulse can be obtained either from the full in-in result with retarded propagators after subtracting the conservative part, or by explicitly accounting for the difference between Feynman and retarded Green's functions~\cite{Jakobsen:2022psy}.

Therefore, throughout this thesis, whenever the eikonal phase is discussed we use the Feynman prescription appropriate to the conservative elastic sector. For the conservative impulse at two loops, the same bulk Feynman prescription is sufficient. For causal observables, such as the full impulse and waveform, one must impose the retarded prescription required by causality, while at three loops and higher the possible conservative contribution of multiple on-shell internal modes requires a more careful treatment of the \(i\varepsilon\) prescription.
}
\section{Computation of eikonal phase}
\label{ch3:sec4}
In this section, we provide the detailed computation of the eikonal phase up to third Post-Minkowskian (PM) order. We mainly focus on the scalar-field and dCS corrections to the phase, since the contribution of the Einstein-Hilbert term has already been investigated in \cite{Jakobsen:2022fcj, Jakobsen:2021zvh}. We restrict our analysis to linear order in spin, leaving the study beyond that order for future work.

%\textcolor{red}{Check the overall factor of $\frac{1}{\sqrt{\gamma^2-1}}$ for every integral due to delta function!}\\
\textbf{Spinning eikonal at 2PM:}
In this subsection, we compute the eikonal phase at 2PM order by considering only the scalar-graviton interaction terms by noting that spin only appears when there are scalar-graviton vertices or scalar-fermion vertices in the worldline. The spinless diagrams involving scalar and scalar-graviton have been done in \cite{Bhattacharyya:2024aeq}.\\

\textbullet $\,\,$ We start with the diagram of the following topology where we have a scalar-superfield vertex,
\begin{align}
    \begin{split}
         \begin{minipage}[h]{0.08\linewidth}
	\vspace{4pt}
	\scalebox{1.8}{\includegraphics[width=\linewidth]{2PM_1.png}}
 \end{minipage}&
    \end{split}
\end{align}
The corresponding amplitude is given by
\begin{align}
    \begin{split}
      i \chi\Big|_{\mathcal{S}^1}= -i\frac{s_1 s_2 m_1 m_2^2}{8 m _p^4} \int e^{-i(k_1+k_2)\cdot b}\frac{\hat\delta(\omega-k_1\cdot v_1)\hat\delta(\omega+k_2\cdot v_1)\hat\delta(k_1\cdot v_2)\hat\delta(k_2\cdot \textcolor{black}{v_2})}{\omega\,k_1^2\, k_2^2}\mathcal{N}
    \end{split}
\end{align}
where the numerator takes the form:
\begin{align}
    \begin{split}
        \mathcal{N}=\omega \,\bar\Psi_{1\eta}\,\Psi_1^{\sigma}\eta^{\eta\delta}k_{2[\delta}\delta_{\sigma]}^{(\mu}v_1^{\nu)}P_{\mu\nu;\alpha\beta}v_2^\alpha v_2^\beta\,.
    \end{split}
\end{align}
Now, using the spin-supplementary condition (\ref{SSC17}) and integrating over the worldline energy we get,
\begin{align}
    \begin{split}
      i  \chi\Big|_{\mathcal{S}^1}&= \frac{\gamma s_1 s_2 m_1 m_2^2}{32 m _p^4} v_{2\eta}\mathcal{S}_1^{\eta\delta}\int e^{-i(k_1+k_2)\cdot b}\frac{\hat\delta(k_1\cdot v_1+k_2\cdot v_1)\hat\delta(k_1\cdot v_2)\hat\delta(k_2\cdot v_2)\,k_{2\delta}}{k_1^2 k_2^2}\,,\\ &
        =\frac{\gamma s_1 s_2 m_1 m_2^2}{32 m _p^4}v_{2\eta}\mathcal{S}_1^{\eta\delta}\int_{q} e^{-iq\cdot b}\hat\delta(q\cdot v_1)\hat\delta(q\cdot v_2)\int_{k_1}\frac{\hat\delta(k_1\cdot v_2)\,(q-k_1)_\delta}{k_1^2(k_1-q)^2}\,.
    \end{split}
\end{align}
Now, using Passarino-Veltman (PV) reduction and inserting the vertex factors, we get,


\begin{align}
    \begin{split} 
      i  \chi\Big|_{\mathcal{S}^1}&= \frac{\gamma s_1 s_2 m_1 m_2^2}{512 m _p^4}v_{2\eta}\mathcal{S}_1^{\eta i}\int e^{-iq\cdot b}\frac{\hat\delta(q\cdot v_1)\hat\delta(q\cdot v_2)\, q_i}{(-q^2)^{1/2}}\,,\\ &
        =-i\frac{ s_1 s_2 m_1 m_2^2}{1024 \pi m _p^4|\boldsymbol{b}|^2}\Big(\frac{1+x^2}{1-x^2}\Big)\Big(\hat b\cdot \mathcal{S}_1\cdot v_2\Big)\,.
        \end{split}
\end{align}
Here we have replaced $\gamma$ by $\frac{x^2+1}{2 x}.$ All the subsequent expressions for the eikonal phase will be expressed in terms of $x\,.$ For the $q$ integral, we have used \eqref{int} of Appendix~(\ref{ch3:app:A}). In all the subsequent computations, whenever we will encounter this kind of vector integral, we will use the result mentioned in \eqref{int}.

\textbullet $\,\,$ Another possible diagram in this order can be obtained by replacing the graviton line with the scalar line in the previous diagram.
\begin{align}
    \begin{split}
         \begin{minipage}[h]{0.08\linewidth}
	\vspace{4pt}
	\scalebox{1.8}{\includegraphics[width=\linewidth]{2PM_1_1.png}}
 \end{minipage}&
    \end{split}
\end{align}
The corresponding phase takes the following form,
\begin{align}
    \begin{split}
        i\chi\sim\int e^{-i(k_1+k_2)\cdot b}\frac{\hat\delta(\omega-k_1\cdot v_1)\hat\delta(\omega+k_2\cdot v_1)\hat\delta(k_1\cdot v_2)\hat\delta(k_2\cdot v_2)}{\omega\,k_1^2\,k_2^2}\mathcal{N}
    \end{split}
\end{align}
where the numerator takes the form,
\begin{align}
    \begin{split}
        \mathcal{N}=\omega^2\, \bar\Psi_{\eta}\Psi_{\delta}\,\eta^{\eta\delta}\,.
    \end{split}
\end{align}
Now, using the fact that another diagram with opposite spinor flow will eventually add up to this diagram such that we will eventually have terms that are antisymmetric in of $(\eta,\delta)$ and therefore, this diagram will not contribute to the eikonal phase.\textit{This structure eventually ensures that only the scalar line along with worldline propagators does not contribute to the spinning eikonal.}
\par
\textbullet $\,\,$ Now we show the detailed computation of the following 2PM diagram with $\mathbb{U}$-type topology. Note that we only focus on the spinning part of the diagram as the spinless part has been done in \cite{Bhattacharyya:2024aeq}.
\begin{align}
    \begin{split}
         \begin{minipage}[h]{0.08\linewidth}
	\vspace{4pt}
	\scalebox{1.7}{\includegraphics[width=\linewidth]{2PM_2.png}}
 \end{minipage}&
    \end{split}
\end{align}
The contribution to the eikonal phase from this diagram takes the following form,
\begin{align}
    \begin{split}
       i \chi\Big|_{\mathcal{S}^1} =-i\frac{s_1 s_2 m_1 m_2^2}{16m_p^4} \int e^{-i(k_1+k_2)\cdot b}\frac{\hat\delta(k_1\cdot v_1+k_2\cdot v_1)\hat\delta(k_1\cdot v_2)\hat\delta(k_2\cdot v_2)}{\omega^2 \,k_1^2\, k_2^2}\mathcal{N}\,,
    \end{split}
\end{align}
where the numerator $\mathcal{N}$ has the following form,
\begin{align}
    \begin{split}
        \mathcal{N}=&i\gamma (k_2\cdot \mathcal{S}_1\cdot v_2)\Big[k_1\cdot k_2-2 (k_1\cdot v_1)^2-2\, k_1\cdot v_1\,k_2\cdot v_1\Big]-4i\gamma \,(k_2\cdot \mathcal S_2\cdot v_1)(k_1\cdot v_1)^2+i\gamma (k_2\cdot \mathcal{S}_2\cdot k_1)(k_1\cdot v_1)\\ &
        -\frac{i}{2}(k_2\cdot \mathcal{S}_1\cdot k_1)(k_1\cdot v_1)\,.
    \end{split}
\end{align}
Now, redefining $k_1+k_2\to q$ and ignoring the tadpoles, we will get \footnote{We will always throw away the tadpoles before performing the IBP reduction. From now on, we will not mention it explicitly. It should be understood.},
\begin{align}
    \begin{split}
      i  \chi\Big|_{\mathcal{S}^1}\sim \int e^{-iq\cdot b}\hat\delta(q\cdot v_1)\hat\delta(q\cdot v_2)&\int \frac{\hat\delta(k_1\cdot v_2)}{(k_1\cdot v_1)^2\,k_1^2\,(k_1-q)^2}\Bigg[\frac{\gamma}{2}q^2\Big(q\cdot \mathcal{S}_1\cdot v_2-k_1\cdot \mathcal{S}_1\cdot v_2\Big)-4\gamma\Big(q\cdot \mathcal{S}_2\cdot v_1-k_1\cdot \mathcal S_2\cdot v_1\Big)\\ &
    \times (k_1\cdot v_1)^2+\gamma(q\cdot \mathcal S_2\cdot k_1)(k_1\cdot v_1) -\frac{1}{2}(q\cdot \mathcal{S}_1\cdot k_1)(k_1\cdot v_1)   \Bigg]\,.\label{4.9}
    \end{split}
\end{align}
To proceed further, we need to evaluate the following tensor integrals:
\begin{align}
    \begin{split}
        \theta_1^{\eta}=\int_{k_1}\frac{\hat\delta(k_1\cdot v_2)\,k_1^\eta}{(k_1\cdot v_1)^2\,k_1^2\,(k_1-q)^2}&=\frac{v_1^\eta-\gamma v_2^\eta}{1-\gamma^2}\int_{k_1}\frac{\hat\delta(k_1\cdot v_2)}{(k_1\cdot v_1)\,k_1^2\,(k_1-q)^2}+\frac{q^\eta}{2}\int_{k_1}\frac{\hat\delta(k_1\cdot v_2)}{(k_1\cdot v_1)^2\,k_1^2\,(k_1-q)^2}\,,\\ &
        =\frac{v_1^\eta-\gamma v_2^\eta}{1-\gamma^2} \,\mathbfcal{G}_{1;1,1}+\frac{q^\eta}{2}\frac{8 \epsilon } {(\gamma^2 -1) q^2}\mathbfcal{G}_{0;1,1}\,,\\ &
        \xrightarrow[]{\epsilon\to 0}\frac{v_1^\eta-\gamma v_2^\eta}{1-\gamma^2} \,\mathbfcal{G}_{1;1,1},\\ &
        \hspace{0cm}
        \theta_2^\eta =\int_{k_1}\frac{\hat\delta(k_1\cdot v_2)\,k_1^\eta}{(k_1\cdot v_1)\,k_1^2\,(k_1-q)^2}=\frac{v_1^\eta-\gamma v_2^\eta}{1-\gamma^2}\mathbfcal{G}_{0;1,1}+\frac{q^\eta}{2} \mathbfcal{G}_{1;1,1},\,\,\\&  \hspace{0cm}\theta_3^\eta=\int_{k_1}\frac{\hat\delta(k_1\cdot v_2)\,k_1^\eta}{k_1^2\,(k_1-q)^2}=\frac{q^\eta}{2}\mathbfcal{G}_{0;1,1}\,.
    \end{split}\label{ch3:4.10}
\end{align}
 In \eqref{ch3:4.10}, we have used the IBP reduction using a suitable choice of basis functions and write the full integral in terms of two master integrals: ($\mathbfcal{G}_{1;1,1},\,\mathbfcal{G}_{0;1,1}$). These are defined in \eqref{3.41a}. 
%Therefore, the eikonal phase is given by,
%\begin{align}
 %   \begin{split}
  %   i   \chi\sim \int e^{-iq\cdot b}\hat\delta(q\cdot v_1)\hat\delta(q\cdot v_2)\int_{k_1}\frac{\hat\delta(k_1\cdot v_2)}{(k_1\cdot v_1)^2\,k_1^2\,(k_1-q)^2}\Bigg[&-\frac{\gamma}{2}q^2\,\mathcal{S}_{1\eta\delta}v_2^\delta\, k_1^\eta+4\gamma \mathcal{S}_{2\eta\delta} v_2^\delta \,k_1^\eta (k_1\cdot v_1)^2+\gamma q^\delta \mathcal S_{2\delta\eta} k_1^\eta\\ &
   %     \times (k_1\cdot v_1)-\frac{1}{2} q^\delta \mathcal S_{1\delta\eta} k_1^\eta (k_1\cdot v_1)\Bigg]\,.\label{4.11}
    %\end{split}
%\end{align}
Now, using \eqref{ch3:4.10} as well as the anti-symmetric nature of the spin tensor and the spin supplementary condition {\eqref{SSC17}}, we get,

\begin{align}
    \begin{split}
    i    \chi\Big|_{\mathcal{S}^1}=&\frac{s_1 s_2 m_1 m_2^2}{16m_p^4} \int e^{-iq\cdot b}\hat\delta(q\cdot v_1)\hat\delta(q\cdot v_2)\Bigg(\frac{\gamma}{1-\gamma^2}\,(q\cdot \mathcal{S}_2\cdot v_1)\,+\frac{\gamma}{2(1-\gamma^2)}(q\cdot \mathcal{S}_1\cdot v_2)\Bigg)\mathbfcal{G}_{0;1,1}\,,\\&
    =i\frac{ s_1 s_2 m_1 m_2^2\,x^2(1+x^2)}{64\, \pi\, |\boldsymbol{b}|^2 m_p^4\left(1-x^2\right)^3} \Big(\hat b\cdot \mathcal{S}_2\cdot v_1 +\frac{1}{2}\hat{b}\cdot \mathcal{S}_1\cdot v_2\Big)\,.
    \end{split}
\end{align}
As we can see, although the intermediate step is tedious, the result is remarkably simple and depends on only one master integral $\mathbfcal{G}_{0;1,1}$ defined in \eqref{3.4aa} \footnote{Note that we have only shown the regularization independent finite part after doing an expansion in small $\epsilon\,.$ We will follow this same strategy for all the diagrams.}. We will use a similar methodology to evaluate the rest of the diagrams (at 3PM also).

\textbullet $\,\,$ Another scalar-graviton worldline interaction vertex having a $\mathbb{V}$-type topology will contribute to the 2PM eikonal phase. Again, for this case, the spinless part has been done in \cite{Bhattacharyya:2024aeq}\,.
\begin{align}
    \begin{split}
         \begin{minipage}[h]{0.08\linewidth}
	\vspace{4pt}
	\scalebox{1.7}{\includegraphics[width=\linewidth]{2PM_3.png}}
 \end{minipage}&
    \end{split}
\end{align}
The eikonal phase contribution is given by,
\begin{align}
    \begin{split}
        i\,\chi\Big|_{\mathcal{S}^1} & =-i\frac{s_1 s_2 m_1 m_2^2}{8m_p^4} \int e^{-i(k_1+k_2)\cdot b}\frac{\hat\delta(k_1\cdot v_1+k_2\cdot v_1)\hat\delta(k_1\cdot v_2)\hat\delta(k_2\cdot v_2)}{k_1^2 k_2^2}\\ &
        \hspace{0.8cm}\times\Big(v_1^\mu v_1^\nu-i(k_2\cdot \mathcal{S}_1)^{(\mu}v_1^{\nu)}\Big)P_{\mu\nu;\alpha\beta}\Big(v_2^\alpha v_2^\beta+i (k_2\cdot \mathcal{S}_2)^{(\alpha}v_2^{\beta)}\Big)\,,\\ &
        =\frac{ \gamma s_1 s_2 m_1 m_2^2}{8m_p^4} \int e^{-iq \cdot b}\hat\delta(q\cdot v_1)\hat\delta(q\cdot v_2)\int \frac{\hat\delta(k_1\cdot v_2)}{k_1^2 (k_1-q)^2}(-k_1\cdot 
    \mathcal{S}_2\cdot v_1+k_1\cdot \mathcal{S}_1\cdot v_2+q\cdot \mathcal{S}_2\cdot v_1-q\cdot \mathcal{S}_1\cdot v_2),\\ &
    =\frac{ \gamma s_1 s_2 m_1 m_2^2}{16m_p^4} \int e^{-iq\cdot b}\hat\delta(q\cdot v_1)\hat\delta(q\cdot v_2)\Big((q\cdot \mathcal S_2\cdot v_1)-(q\cdot \mathcal S_1\cdot v_2)\Big)\,\mathbfcal{G}_{0;1,1}\,,\\&
    =i\frac{ s_1 s_2 m_1 m_2^2\,(1+x^2)}{256\,\pi\,  |b|^2 m_p^4\left(1-x^2\right)} \Big(\hat b\cdot \mathcal{S}_2\cdot v_1 -\hat{b}\cdot \mathcal{S}_1\cdot v_2\Big)\,.
    \end{split}
\end{align}
Hence the final form of the eikonal phase is,

\begin{align}
    \begin{split}
i\chi\Big|_{\mathcal{S}^{1}}=  i\frac{ s_1 s_2 m_1 m_2^2\,(1+x^2)}{256\,\pi\,  |\boldsymbol{b}|^2 m_p^4\left(1-x^2\right)} \Big(\hat b\cdot \mathcal{S}_2\cdot v_1 -\hat{b}\cdot \mathcal{S}_1\cdot v_2\Big)\,.
    \end{split}
\end{align}
At this point, we would like to remind the readers that one also needs to add the contributions from 2PM diagrams obtained by interchanging the worldlines one and two in all the diagrams mentioned above. We will not show them explicitly as the contributions from those diagrams will be similar to those we have already computed and can be obtained easily by exchanging labels one and two.

Finally adding all the $\mathcal{O}(\mathcal{S}^1)$ at 2PM contributions we have,

\begin{tcolorbox}[thesisresultbox, title=\textit{Spinning Eikonal at 2PM}]
\restorethesisbodyformat
\begin{align*}   \chi\Big|_{\mathcal{S}^1}^{\textrm{2PM}}=&\frac{m_1m_2^2}{m_p^4}\Bigg[\Rho_{01}\Big(\hat b\cdot \mathcal{S}_{1}\cdot v_2\Big)+\Rho_{02}\Big(\hat b\cdot \mathcal{S}_2\cdot v_1 \Big)\Bigg]+(1\leftrightarrow 2)\,,
\end{align*}
\end{tcolorbox}
\restorethesisbodyformat
where, $$\Rho_{01}=\frac{s_1\,s_2\,\left(x^2+1\right) \left(5 x^4-18 x^2+5\right)}{1024 \pi  |\boldsymbol{b}|^2  \left(x^2-1\right)^3}\,,\quad \Rho_{02}=-\frac{s_1\,s_2\,\left(x^2+1\right)^3}{256 \pi  |\boldsymbol{b}|^2  \left(x^2-1\right)^3}\,.$$
%\newpage
\subsection*{Spinning eikonal at 3PM}
Now, we will compute the contributions to the eikonal phase at 3PM order. We will focus on the correction to the eikonal phase due to the extra scalar degree of freedom.


\textbullet $\,\,$ First, we will compute the Feynman diagrams coming purely from the scalar part. We start with the first comb-type diagram. 
\begin{center}
\includegraphics[width=0.38\linewidth]{3PM_1.png}
\end{center}
The contribution to the eikonal phase looks like,
\begin{align}
    \begin{split}
       i\, \chi\Big|_{\mathcal{S}^0}&=-i \frac{s_1^3 s_2^3 m_1 m_2^3}{64 m_p^6} \rmint_{k_i,\omega_i}\hat\delta(k_1\cdot v_2)\hat\delta(k_2\cdot v_2)\hat\delta(k_3\cdot v_2)\hat\delta(-k_1\cdot v_1+\omega_1)\hat\delta(-k_2\cdot v_1+\omega_2-\omega_1)\hat\delta(\omega_2+k_3\cdot v_1)\\ &
       \hspace{2 cm} \times\frac{ \mathcal{N}}{\prod_i k_i^2}\left(\sum_{(\pm)} \frac{1}{(\pm\omega_1+i\varepsilon)^2(\pm\omega_2+i\varepsilon)^2}\right)\,e^{-i(k_1+k_2+k_3)\cdot b}
    \end{split}
\end{align}
Now relabeling the momenta $k_1\to \ell_1,\,k_2\to q-\ell_1-\ell_2\,,k_3\to\ell_2$ we can write the numerator as,\footnote{We have performed all the tensor contractions used in this paper using \textbf{FeynCalc} \cite{Mertig:1990an, Shtabovenko:2020gxv,Shtabovenko:2023idz}\,.},
\begin{align}
    \begin{split}
        \mathcal{N}=&4(\ell_1\cdot v_1)(\ell_2\cdot v_1)\left(\ell_1\cdot v_1+\ell_2\cdot v_1\right)^2
    \end{split}
\end{align}
%The integral in \eqref{2.3a} can be written in terms of the family of integrals,
%\begin{align}
 %   \begin{split}
        %\mathcal{M}^{(\pm,\pm)}_{n_1,n_2\cdots n_7}(q)=\int_{k_1,l}\frac{1}{D_1^{n_1}D_2^{n_2}\cdots D_7^{n_7}D_8^{(\pm)n_8}D_9^{(\pm)n_9}}
   % \end{split}
%\end{align}
%with,
%\begin{align}
  %  \begin{split}
  %    &  D_1=k_1\cdot v_1,\,D_2=l\cdot v_1,\,D_3=k_1^2,\,D_4=l^2,\,D_5=(l-k_1)^2,\,D_{6}=(k_1-q)^2,\,D_7=(l-q)^2,\, \slashed D_8^{(\pm)}=k_1\cdot v_2\pm i \sigma_1,\\ &\,\slashed{D}_9^{(\pm)}=l\cdot v_2\pm i \sigma_2.
  %  \end{split}
%\end{align}
%After doing IBP using \textbf{LiteRed} \cite{Lee:2013mka} to reduce the twoloop integral into linear combination of master integrals as,
%\begin{align}
 %   \begin{split}
      %  \chi&\sim \rmint_{q,k_i}e^{-i q\cdot b}\hat\delta(q\cdot v_1)\hat\delta(q\cdot v_2)\sum a(n_1,n_2,\cdots,n_7)\,\mathcal{G}_{n_1,n_2\cdots n_7,1,1}(q)\\ &
      %  \to \rmint_{q,k_i}e^{-i q\cdot b}\hat\delta(q\cdot v_1)\hat\delta(q\cdot v_2)\sum' \Tilde{a}(n_1,..,n_7)\mathcal{M}_{n_1,..,n_7,1,1}^{\textrm{m.i.}}(q)
    %\end{split}
%\end{align}
%where, $\sum'$ denotes the sum over relevant master integrals.
\noindent
Now, taking care of the $\pm i\varepsilon$ prescription and then doing IBP reduction using \textbf{LiteRed} \cite{Lee:2012cn,Lee:2013mka}, one can find the eikonal phase in terms of two master integrals and is given by, %[$\color{red}{\gamma=\frac{x^2+1}{2x}}$],
\begin{align}
    \begin{split}
      &  i\chi\Big|_{\mathcal{S}^0}= -i \frac{s_1^3 s_2^3 m_1 m_2^3}{64 m_p^6}  \int_{q}e^{-iq\cdot b}\hat\delta(q\cdot v_1)\hat\delta(q\cdot v_2)\Bigg[  \boldsymbol{\tilde h_1}\mathbfcal{L}_{0,0;0,0,1,1,1}(q) -\boldsymbol{\boldsymbol{\tilde h_2}}\,q^2\mathbfcal{L}^{(+-)}_{1,1;0,0,1,1,1}(q)\Bigg]\,.
    \end{split}
\end{align}
where, $\boldsymbol{\tilde h_1}$ and $\boldsymbol{\tilde h_2}$ are defined in Appendix~(\ref{ch3:app:E})\,.
%\begin{align}
%   \boldsymbol{\tilde h_1}=4,\,\,  \boldsymbol{\tilde h_2}=0
%\end{align}
Then using \eqref{MasterL}, we finally get,

%\hfsetfillcolor{gray!10}
%\hfsetbordercolor{black!150}
\begin{align}
    \begin{split} \label{eq1}
\chi\Big|_{\mathcal{S}^0}=\frac{ m_1 m_2^3}{m_p^6}\Bigg(\frac{s_1^3 s_2^3 x}{256 \pi ^3 |\boldsymbol{b}|^2 \left(x^2-1\right)}\Bigg) ,
    \end{split}
\end{align}
where $\gamma_E$ is the Euler-Mascheroni constant. For the $q$ integral, we have used \eqref{int1} of Appendix~(\ref{ch3:app:A}). In all the subsequent computations, whenever we will encounter this kind of scalar integral, we will use the result mentioned in \eqref{int1}.
The second diagram gives the following contribution,
\begin{align}
    \begin{split}
      i\chi\Big|_{\mathcal{S}^0}=-i \frac{s_1^3 s_2^3 m_1 m_2^3}{256 m_p^6} \int e^{-iq\cdot b}\hat\delta(q\cdot v_1)\hat\delta(q\cdot v_2)\int \sum_{(\pm)}\frac{\hat\delta(\ell_1\cdot v_2)\hat\delta(\ell_2\cdot v_2)\,\mathcal{N}}{(\pm\ell_1\cdot v_1+i\varepsilon)^2\,(\pm\ell_2\cdot v_1+i\varepsilon)^2\,\ell_1^2\,\ell_2^2\,(\ell_1+
        \ell_2-q)^2}
    \end{split}
\end{align}
where the numerator has the following form,
\begin{align}
    \begin{split}
        \mathcal{N}=&\frac{1}{8}\Bigg[(\ell_1-q)^2(\ell_2-q)^2+4q^2 (\ell_1\cdot v_1)(\ell_2\cdot v_1)\Bigg]\,.
    \end{split}
\end{align}
Therefore, the eikonal phase is given by,
\begin{align}
    \begin{split} 
    i\chi\Big|_{\mathcal{S}^0}=-i \frac{s_1^3 s_2^3 m_1 m_2^3}{64 m_p^6}\int e^{-iq\cdot b}\hat\delta(q\cdot v_1)\hat\delta(q\cdot v_2)\Bigg[\boldsymbol{\tilde h}_3\,\mathbfcal{L}_{0,0;0,0,1,1,1}+\frac{1}{4}\boldsymbol{\tilde h}_4\,q^2\,\mathbfcal{L}^{(++)}_{1,1;0,0,1,1,1}\Bigg]
    \end{split}
\end{align}
where, $ \boldsymbol{\tilde h_3},  \boldsymbol{\tilde h_4}$ are defined in Appendix~(\ref{ch3:app:E}).
Using \eqref{MasterL}, we finally get,

\begin{align}
    \begin{split} \label{eq2}
 \chi\Big|_{\mathcal{S}^0}=\frac{m_1 m_2^3}{m_p^6}\Bigg(\frac{ s_1^3 s_2^3 x^3 \left(x^4+1\right) (6 \log (|\boldsymbol{b}|)+5 \gamma_E +\log \left(\frac{\pi }{16}\right))}{512 \pi ^3 |\boldsymbol{b}|^2 \left(x^2-1\right)^5}\Bigg)\,.
    \end{split}
\end{align}
%Then we get, 
%\begin{align}
%\begin{split}
%\textcolor{red}{ i\chi\Big|_{\mathcal{S}^0} \sim \Big(\frac{2 x}{1-x^2}\Big)\frac{1}{12 \pi ^3 |\boldsymbol{b}|^2 x \left(x^2-1\right)^4}}\Bigg[& 36 \left(3 x^4+2 x^2+3\right) x^3 \log (|\boldsymbol{b}|)+x \big(3 x^9-3 x^8-3 x^7+ \\& x^6 (54 \gamma_E +6+32 \log (2)+70 \log (\pi ))+2 x^4 (18 \gamma_E +9+16 \log (2)\\& +26 \log (\pi ))+x^2 (54 \gamma_E +6+32 \log (2)+70 \log (\pi ))-3 x-3\big)+3\Bigg]\,.
%\end{split}
%\end{align}
Then, the total contribution to the eikonal phase coming from these two diagrams after summing (\ref{eq1}) and (\ref{eq2}) is the following, 
\begin{align}
    \begin{split} \label{eqt1}
 \chi\Big|_{\mathcal{S}^0}=\frac{m_1 m_2^3}{m_p^6}\Delta_{\mathfrak{01}}^{\textrm{poly}}\,,
    \end{split}
\end{align}
where, $$\Delta_{\mathfrak{01}}^{\textrm{poly}}=\Bigg(\frac{s_1^3 s_2^3 x \left(\left(x^6+x^2\right) \log \left(\frac{\pi  |\boldsymbol{b}|^6}{16}\right)+2 x^8+(5 \gamma_E -8) x^6+12 x^4+(5 \gamma_E -8) x^2+2\right)}{512 \pi ^3 |\boldsymbol{b}|^2 \left(x^2-1\right)^5}\Bigg)\,.$$
%\hfsetfillcolor{gray!10}
%\hfsetbordercolor{black!150}
%\begin{align}
%    \begin{split}
%     \tikzmarkin[disable rounded corners=false]{eq}(3.2,-0.55)(-0.5,1) 
%\chi\Big|_{\mathcal{S}^0} =\frac{m_1m_2^3}{m_p^6}\Bigg(\frac{s_1^3s_2^3}{2048 \pi ^3 |\boldsymbol{b}|^2 x^3 \left(x^2-1\right)^5}\Bigg)&\Bigg[6 \left(x^4+1\right) \left(8 x^4-3 \left(x^2-1\right)^4\right) x^2 \log \left(\pi  b^2\right)\\&+6 x^{16}-3 (19+6 \gamma_E ) x^{14}+6 (35+12 \gamma_E ) x^{12}\\&-3 (141+26 \gamma_E ) x^{10}+16 (34+9 \gamma_E ) x^8-3 (141+26 \gamma_E ) x^6\\&+6 (35+12 \gamma_E ) x^4-3 (19+6 \gamma_E ) x^2+6\Bigg]\,.
%\tikzmarkend{eq}
%\end{split}
%\end{align}


\textbullet $\,\,$ Apart from the previous one, we have the following two comb-type diagrams where one of the scalar propagators is replaced by the graviton propagator. 
\begin{center}
\includegraphics[width=0.38\linewidth]{newo_1.png}
\end{center}
%From the first diagram we have the following contribution to the eikonal phase,
%\begin{align}
%    \begin{split}
%        \chi=-i\frac{m_1m_2^3 s_1^2 s_2^2}{32 m_p^6} \int e^{-iq\cdot b}\hat\delta(q\cdot v_1)\hat\delta(q\cdot v_2)\int \frac{\hat\delta(\ell_1\cdot v_2)\hat\delta(\ell_2\cdot v_2) \,\,(\mathcal{N}|_{\mathcal{S}^0}+\mathcal{N}|_{\mathcal{S}^1})}{(\ell_1\cdot v_1+i\varepsilon)^2(\ell_2\cdot v_1-i\varepsilon)^2\,\ell_1^2\,\ell_2^2\,(\ell_1+\ell_2-q)^2}
%    \end{split}
%\end{align}
%where,
%\begin{align}
%    \begin{split}
%&\mathcal{N}|_{\mathcal{S}^0}=-\frac{1}{8 (d-2)}\Bigg[(q-\ell_1)^2 ((1-\gamma ^2 (d-2)) (q-\ell_2)^2+4 \gamma ^2 (d-2)(\ell_1\cdot v_1)^2+4 (2-\gamma ^2 (d-2)) (\ell_1\cdot v_1) (\ell_2\cdot v_1))
%\\&+4 (\ell_1\cdot v_1) (\ell_2\cdot v_1) (2 (1-\gamma ^2 (d-2)) (q-\ell_2)^2+4 \gamma ^2 (d-2)(\ell_1\cdot v_1)^2+8 \gamma ^2 (d-2) (\ell_1\cdot v_1) (\ell_2\cdot v_1)\\&+q^2 (\gamma ^2 (d-2)-1))\Bigg]\,,\\&
% \mathcal{N}|_{\mathcal{S}^1}=\frac{i}{8}\,\ell_1\cdot \mathcal{S}_1\cdot v_2\,\Bigg[4 \gamma  (\ell_1\cdot v_1) (\ell_2\cdot v_1) (q^2-2 (q-\ell_2){}^2)-\gamma  (q-\ell_1){}^2 ((q-\ell_2){}^2+4 (\ell_1\cdot v_1) (\ell_2\cdot v_1))\Bigg]\\&-\frac{i}{4 (d-2)}\ell_1\cdot\mathcal{S}_1\cdot q\,\Bigg[{ (\ell_1\cdot v_1) ((q-\ell_1){}^2+4 (\ell_1\cdot v_1) (\ell_2\cdot v_1))}\Bigg]\\&+\frac{i}{4(d-2)}\ell_1\cdot \mathcal{S}_1\cdot \ell_2\,\Bigg[ (\ell_1\cdot v_1) ((q-\ell_1){}^2+8 (\ell_1\cdot v_1) (\ell_2\cdot v_1))\Bigg]\\\ &
 %       -\frac{i\gamma}{8}\ell_1\cdot \mathcal{S}_2\cdot v_1\Bigg[  ((q-\ell_1){}^2 (-(q-\ell_2){}^2+4 (\ell_1\cdot v_1){}^2-4 (\ell_2\cdot v_1) (\ell_1\cdot v_1))+\\& 4 (\ell_1\cdot v_1) (\ell_2\cdot v_1) (-2 (q-\ell_2){}^2+4 (\ell_1\cdot v_1){}^2+8 (\ell_1\cdot v_1) (\ell_2\cdot v_1)+q^2))\Bigg]\\&-\frac{i\gamma}{4}\ell_1\cdot \mathcal{S}_2\cdot q\,\Bigg[  (\ell_1\cdot v_1) ((q-\ell_1){}^2+4 (\ell_1\cdot v_1) (\ell_2\cdot v_1))\Bigg]
%     +   \frac{i\gamma}{4}\ell_1\cdot \mathcal{S}_2\cdot \ell_2\,\Bigg[  (\ell_1\cdot v_1) ((q-\ell_1){}^2+8 (\ell_1\cdot v_1) (\ell_2\cdot v_1))\Bigg]\,.
%    \end{split}
%\end{align}
We proceed as before. The non-spinning part of the eikonal phase is given by, %\textcolor{black}{AB: Need to update the coefficients!},
\begin{align}
    \begin{split} \label{eq3}
      i\chi\Big|_{\mathcal{S}^0}=-i\frac{m_1m_2^3 s_1^2 s_2^2}{32 m_p^6}\int e^{-iq\cdot b}\hat\delta(q\cdot v_1)\hat\delta(q\cdot v_2)\,\Bigg[ \boldsymbol{\tilde h}_5 \,\mathbfcal{L}_{0,0;0,0,1,1,1}-\frac{1}{2}\boldsymbol{\tilde h}_6 \,q^2\,\mathbfcal{L}^{(+-)}_{1,1;0,0,1,1,1}\Bigg]
    \end{split}
\end{align}
$ \boldsymbol{\tilde h_5},  \boldsymbol{\tilde h_6}$ are defined in Appendix~(\ref{ch3:app:E}).
Similarly, the spinning part is given by, 
\begin{align}
    \begin{split}  \label{eq4}
      i\chi\Big|_{\mathcal{S}^1}=-i\frac{m_1m_2^3 s_1^2 s_2^2}{32 m_p^6}\int e^{-iq\cdot b}\hat\delta(q\cdot v_1)\hat\delta(q\cdot v_2)\,\Bigg[ &i\,\Bigg(\boldsymbol{\tilde h}_7 \,\mathbfcal{L}_{0,0;0,0,1,1,1}-\boldsymbol{\tilde h}_8 \,q^2\,\mathbfcal{L}^{(+-)}_{1,1;0,0,1,1,1}\Bigg)(q\cdot \mathcal{S}_1\cdot v_2)+\\&i\,\Bigg(\boldsymbol{\tilde h}_9\,\mathbfcal{L}_{0,0;0,0,1,1,1}-\frac{1}{2}\boldsymbol{\tilde h}_{10}\,q^2\,\mathbfcal{L}^{(+-)}_{1,1;0,0,1,1,1}\Bigg)(q\cdot \mathcal{S}_2\cdot v_1)\Bigg]\,.
    \end{split}
\end{align}

%\begin{align}
 %   \begin{split}
% &   \hspace{0cm} \
%    \end{split}
%\end{align}
%For the computation of spinning Eikonal we encounter the following vector integrals,
%\begin{align}
%    \begin{split}
 %   \boldsymbol{\mathscr{M}}_1^{\eta}&=\int \frac{\hat\delta(\ell_1\cdot v_2)\hat\delta(\ell_2\cdot v_2) \,\,\ell_1^\eta}{(\ell_1\cdot v_1+i\varepsilon)^2(\ell_2\cdot v_1-i\varepsilon)^2\,\ell_1^2\,\ell_2^2\,(\ell_1+\ell_2-q)^2}\Bigg[\Bigg]\\ & =\boldsymbol{\mathfrak{m}}_1\,q^\eta+\boldsymbol{\mathfrak{m}}_2(\gamma v_2^\mu-v_1^\mu),
%    \end{split}
%\end{align}
%where,
%\begin{align}
%   \begin{split}
 %   &\boldsymbol{\mathfrak{m}}_1=\frac{\gamma  \left(32-24 \gamma ^2\right) \epsilon +\gamma  \left(6 \gamma ^2-7\right)-48 \gamma  \epsilon ^3-4 \gamma  \epsilon ^2}{3 \left(\gamma ^2-1\right)^2}\mathbfcal{L}_{0,0;0,0,1,1,1}\\&+\frac{2 \gamma  \left(-9 \gamma ^2+9 \left(\gamma ^2-1\right)+7\right) \epsilon +2 \gamma  \left(3 \gamma ^2-3 \left(\gamma ^2-1\right)-2\right)-4 \gamma  \epsilon ^2}{3 \left(\gamma ^2-1\right)^2}q^2\mathbfcal{L}_{1,1;0,0,1,1,1},\,\,\\&\boldsymbol{\mathfrak{m}}_2=-\frac{-2 \gamma  \left(1-2 \gamma ^2\right)-24 \gamma  \epsilon ^2-8 \gamma  \epsilon}{\left(\gamma ^2-1\right) (6 \epsilon -1)}q^2\,\mathbfcal{L}^{(+-)}_{0,1;0,0,1,1,1}
 %   \end{split}
%\end{align}
%\begin{align}
 %   \begin{split}
        %\boldsymbol{\mathscr{M}}_2^{\eta}&=\int \frac{\hat\delta(\ell_1\cdot v_2)\hat\delta(\ell_2\cdot v_2) \,\,\ell_1^\eta\,(\ell_1\cdot v_1)}{(\ell_1\cdot v_1+i\varepsilon)^2(\ell_2\cdot v_1-i\varepsilon)^2\,\ell_1^2\,\ell_2^2\,(\ell_1+\ell_2-q)^2}\Bigg[\Bigg]\\ &
%=\boldsymbol{\mathfrak{m}_3}q^{\eta}+\boldsymbol{\mathfrak{m}_4}(\gamma v_2^\eta-v_1^\eta)
 %   \end{split}
%\end{align}
%where,
%\begin{align}
 %   \begin{split}
 %      \boldsymbol{\mathfrak{m}_3}= -\frac{\left(6-8 \gamma ^2\right) \epsilon +2 \gamma ^2-24 \epsilon ^2-1}{2 (\gamma -1) (\gamma +1) (\epsilon -1) (6 \epsilon -1)}\mathbfcal{L}^{(+-)}_{0,1;0,0,1,1,1},\,\,  \boldsymbol{\mathfrak{m}_4}=\frac{-2 \gamma ^2+6 \epsilon +1}{\gamma^2-1}\mathbfcal{L}_{0,0;0,0,1,1,1}
   % \end{split}
%.\end{align}
%\begin{align}
%    \begin{split}
%\boldsymbol{\mathscr{M}}_{4}^\eta=&\int \frac{\hat\delta(\ell_1\cdot v_2)\hat\delta(\ell_2\cdot v_2) \,\,\ell_1^\eta}{(\ell_1\cdot v_1+i\varepsilon)^2(\ell_2\cdot v_1-i\varepsilon)^2\,\ell_1^2\,\ell_2^2\,(\ell_1+\ell_2-q)^2}\Bigg[\Bigg]\\ &
%=\boldsymbol{\mathfrak{m}_5}q^\eta+\boldsymbol{\mathfrak{m}_6}(\gamma v_2^\eta-v_1^\eta)
 %   \end{split}
%\end{align}
%where,
%\begin{align}
 %   \begin{split}
%\boldsymbol{\mathfrak{m}_5}=&\frac{2 \left(24 \gamma ^2+20\right) \epsilon ^3+2 \left(7-6 \gamma ^2\right)+2 \left(24 \gamma ^4-96 \gamma ^2+84\right) \epsilon ^2+2 \left(-20 \gamma ^4+76 \gamma ^2-65\right) \epsilon -96 \epsilon ^4}{3 \left(\gamma ^2-1\right)^2}\,\mathbfcal{L}_{0,0;0,0,1,1,1}\\ &-\frac{4 \left(-9 \gamma ^2+9 \left(\gamma ^2-1\right)+7\right) \epsilon +4 \left(3 \gamma ^2-3 \left(\gamma ^2-1\right)-2\right)-8 \epsilon ^2}{3 \left(\gamma ^2-1\right)^2}q^2\mathbfcal{L}^{(+-)}_{1,1;0,0,1,1,1}\\ &
 %    \hspace{0cm}   %\boldsymbol{\mathfrak{m}_6}=-\frac{\left(
%\begin{array}{c}
 %2 \gamma ^4-4 \left(\gamma ^2-5\right) \epsilon ^2-7 \gamma ^2+\left(-4 \gamma %^4+12 \gamma ^2+2\right) \epsilon -24 \epsilon ^3+3 \\
%\end{array}
%\right)}{(\gamma -1) (\gamma +1) (\epsilon -1) (6 \epsilon -1)}q^2 %\mathbfcal{L}^{(+-)}_{0,1,0,0,1,1,1}
%    \end{split}
%\end{align}
%\begin{align}
%\begin{split}
        %\boldsymbol{\mathscr{M}}_{5}^\eta=&\int \frac{\hat\delta(\ell_1\cdot v_2)\hat\delta(\ell_2\cdot v_2) \,\,\ell_1^\eta}{(\ell_1\cdot v_1+i\varepsilon)^2(\ell_2\cdot v_1-i\varepsilon)^2\,\ell_1^2\,\ell_2^2\,(\ell_1+\ell_2-q)^2}(\ell_1\cdot v_1)\Bigg[\Bigg]\\ &
%=\boldsymbol{\mathfrak{m}_7}q^\eta+\boldsymbol{\mathfrak{m}_8}(\gamma v_2^\eta-v_1^\eta)
 %   \end{split}
%\end{align}
%where,
%\begin{align}
  %  \begin{split}
 %    \boldsymbol{\mathfrak{m}_7}=  -\frac{\left(6-8 \gamma ^2\right) \epsilon +2 \gamma ^2-24 \epsilon ^2-1}{2 (\gamma -1) (\gamma +1) (\epsilon -1) (6 \epsilon -1)}\mathbfcal{L}^{(+-)}_{0,1;0,0,1,1,1},\,\, %\boldsymbol{\mathfrak{m}_8}=\frac{-2 \gamma ^2+6 \epsilon +1}{(\gamma ^2-1)}
  %   \mathbfcal{L}_{0,0;0,0,1,1,1}
 %   \end{split}
%\end{align}
%Apart from the above mentioned vector integrals we will have the following tensor integral,
%\begin{align}
 %   \begin{split}
%\boldsymbol{\mathscr{M}}_3^{\eta\delta}&=\int\frac{\hat\delta(\ell_1\cdot v_2)\hat\delta(\ell_2\cdot v_2) \,\,\ell_1^\eta\ell_2^\delta}{(\ell_1\cdot v_1+i\varepsilon)^2(\ell_2\cdot v_1-i\varepsilon)^2\,\ell_1^2\,\ell_2^2\,(\ell_1+\ell_2-q)^2}(\ell_1\cdot v_1)\Bigg[\Bigg]\\ &
%        =\boldsymbol{\mathfrak{m}}_9 q^{[\eta}v_1^{\delta]}+\boldsymbol{\mathfrak{m}}_{10} q^{[\eta}v_2^{\delta]}
%    \end{split}
%\end{align}
%where,
%\begin{align}
 %   \begin{split}
  %      &\boldsymbol{\mathfrak{m}}_9=-\frac{24 \left(\gamma ^2-1\right)-12 \epsilon +2}{6 \left(1-\gamma ^2\right)^2}\mathbfcal{L}_{0,0;0,0,1,1,1}+\frac{\epsilon q^2 }{6 \left(1-\gamma ^2\right) (1-3 \epsilon)}\mathbfcal{L}^{(+-)}_{1,1;0,0,1,1,1}\\ &
        %\boldsymbol{\mathfrak{m}}_{10}=\frac{2 \gamma  \left(12 \gamma ^2-11\right)-12 \gamma  \epsilon}{6 \left(\gamma ^2-1\right)^2}\mathbfcal{L}_{0,0;0,0,1,1,1}-\frac{\gamma  \epsilon q^2 }{6\left(\gamma ^2-1\right) (3 \epsilon -1)}\mathbfcal{L}^{(+-)}_{1,1;0,0,1,1,1}
  %  \end{split}
%\end{align}



%\begin{align}
   % \begin{split}
%\boldsymbol{\mathscr{M}}_6^{\eta\delta}&=\int\frac{\hat\delta(\ell_1\cdot v_2)\hat\delta(\ell_2\cdot v_2) \,\,\ell_1^\eta\ell_2^\delta}{(\ell_1\cdot v_1+i\varepsilon)^2(\ell_2\cdot v_1-i\varepsilon)^2\,\ell_1^2\,\ell_2^2\,(\ell_1+\ell_2-q)^2}(\ell_1\cdot v_1)\Bigg[\Bigg]\\ &
        %=\boldsymbol{\mathfrak{p}}_{11}q^{[\eta}v_1^{\delta]}+\boldsymbol{\mathfrak{p}}_{12} q^{[\eta}v_2^{\delta]}
   % \end{split}
%\end{align}
%where,
%\begin{align}
  %  \begin{split}
  %      &\boldsymbol{\mathfrak{p}}_{11}=-\frac{24 \left(\gamma ^2-1\right)-12 \epsilon +2}{6 \left(1-\gamma ^2\right)^2}\mathbfcal{L}_{0,0;0,0,1,1,1}+\frac{\epsilon q^2 }{6 \left(1-\gamma ^2\right) (1-3 \epsilon)}\mathbfcal{L}^{(+-)}_{1,1;0,0,1,1,1}\\ &
        %\boldsymbol{\mathfrak{p}}_{12}=\frac{2 \gamma  \left(12 \gamma ^2-11\right)-12 \gamma  \epsilon}{6 \left(\gamma ^2-1\right)^2}\mathbfcal{L}_{0,0;0,0,1,1,1}-\frac{\gamma  \epsilon q^2 }{6\left(\gamma ^2-1\right) (3 \epsilon -1)}\mathbfcal{L}^{(+-)}_{1,1;0,0,1,1,1}
  %  \end{split}
%\end{align}
Now, we compute the contribution to the eikonal phase from the second diagram.
%\begin{align}
 %   \begin{split} 
 %       \chi=-i\frac{m_1m_2^3 s_1^2 s_2^2}{32 m_p^6} \int e^{-iq\cdot b}\hat\delta(q\cdot v_1)\hat\delta(q\cdot v_2)\int \frac{\hat\delta(\ell_1\cdot v_2)\hat\delta(\ell_2\cdot v_2)(\mathcal{N}|_{\mathcal{S}^0}+\mathcal{N}|_{\mathcal{S}^1})}{(\ell_1\cdot v_1+i\varepsilon)^2\,(\ell_2\cdot v_1+i\varepsilon)^2\,\ell_1^2\,\ell_2^2\,(\ell_1+
 %       \ell_2-q)^2}
 %   \end{split}
%\end{align}
%where,
%\begin{align}
%    \begin{split} 
%     &   \mathcal{N}|_{\mathcal S^0}= -\frac{1}{8 (d-2)}\Bigg[(q-\ell_1)^2 \Bigg((1-%\gamma ^2 (d-2)) (q-\ell_2)^2+4 \gamma ^2 (d-2) ( \ell_1 \cdot v_1)^2+4 \gamma ^2 (d-2) \ell_1 \cdot v_1 \ell_2\cdot v_1\Bigg)\\&+4 (\ell_1 \cdot v_1) (\ell_2\cdot v_1) \Bigg(4 \gamma ^2 (d-2) ( \ell_1 \cdot v_1)^2+q^2 (1-\gamma ^2 (d-2))\Bigg)\Bigg]\,,\\&
%       \mathcal{N}|_{\mathcal{S}^1}=-\frac{i\gamma}{8}\,\ell_1\cdot \mathcal{S}_1\cdot v_2\,  \Bigg[\left(q-\ell_1\right)^2 \left(\left( q-\ell_2\right)^2-4 \left(\ell_1\cdot v_1\right) \left(\ell_2\cdot v_1\right)\right)+4 q^2 \left(\ell_1\cdot v_1\right) \left(\ell_2\cdot v_1\right)\Bigg]\\&-\frac{i}{4 (d-2)}\ell_1\cdot\mathcal{S}_1\cdot q\,\Bigg[{ (\ell_1\cdot v_1) ((q-\ell_1){}^2+4 (\ell_1\cdot v_1) (\ell_2\cdot v_1))}\Bigg]\\&+\frac{i}{4(d-2)}\ell_1\cdot \mathcal{S}_1\cdot \ell_2\,\Bigg[ (\ell_1\cdot v_1) ((q-\ell_1)^2\Bigg]\\ &
%        -\frac{i\gamma}{8}\ell_1\cdot \mathcal{S}_2\cdot v_1\Bigg[  ((q-\ell_1){}^2 ((q-\ell_2){}^2-4 (\ell_1\cdot v_1){}^2-4 (\ell_2\cdot v_1) (\ell_1\cdot v_1))+\\& 4 (\ell_1\cdot v_1) (\ell_2\cdot v_1) (-4 (\ell_1\cdot v_1){}^2+q^2))\Bigg]\\&-\frac{i\gamma}{4}\ell_1\cdot \mathcal{S}_2\cdot q\,\Bigg[  (\ell_1\cdot v_1) ((q-\ell_1){}^2+4 (\ell_1\cdot v_1) (\ell_2\cdot v_1))\Bigg]
%     +   \frac{i\gamma}{4}\ell_1\cdot \mathcal{S}_2\cdot \ell_2\,\Bigg[  (\ell_1\cdot v_1) (q-\ell_1)^2\Bigg]\,.
%    \end{split}
%\end{align}
The spinless part of the eikonal phase is given by,
\begin{align}
    \begin{split}\label{eq5}
        i\chi\Big|_{\mathcal{S}^0}=-i\frac{m_1m_2^3 s_1^2 s_2^2}{32 m_p^6}\int e^{-iq\cdot b}\hat\delta(q\cdot v_1)\hat\delta(q\cdot v_2)\,\Bigg[\boldsymbol{\tilde h}_{11}\,\mathbfcal{L}_{0,0;0,0,1,1,1}+\frac{1}{4}\boldsymbol{\tilde h}_{12}q^2\,\mathbfcal{L}^{(++)}_{1,1;0,0,1,1,1}\,\Bigg]\,.
    \end{split}
\end{align}
Similarly, the spinning part is given by, 
\begin{align}
    \begin{split} \label{eq6}
      i\chi\Big|_{\mathcal{S}^1}=-i\frac{m_1m_2^3 s_1^2 s_2^2}{32 m_p^6}\int e^{-iq\cdot b}\hat\delta(q\cdot v_1)\hat\delta(q\cdot v_2)\,\Bigg[ &i\,\Bigg(\boldsymbol{\tilde h}_{13} \,\mathbfcal{L}_{0,0;0,0,1,1,1}+ \frac{1}{4}\boldsymbol{\tilde h}_{14} \,q^2\,\mathbfcal{L}^{(++)}_{1,1;0,0,1,1,1}\Bigg)(q\cdot \mathcal{S}_1\cdot v_2)+\\ &i\,\Bigg(\boldsymbol{\tilde h}_{15}\,\mathbfcal{L}_{0,0;0,0,1,1,1}+ \frac{1}{4}\boldsymbol{\tilde h}_{16}\,q^2\,\mathbfcal{L}^{(++)}_{1,1;0,0,1,1,1}\Bigg)(q\cdot \mathcal{S}_2\cdot v_1)\Bigg]\,.
    \end{split}
\end{align}
Again all the $ \boldsymbol{\tilde h_i}'s\,, i=7,\cdots 16\,.$ are defined in Appendix~(\ref{ch3:app:E}). Note that, at $\mathcal{O}(\mathcal{S})$, we encounter tensor integrals. We first have to perform a Veltman-Passarino Reduction (PV) to write it in terms of a set of scalar integrals. Then, we do the IBP reduction for each of those scalar integrals. We will use this same strategy throughout this paper.  \par
Finally, the total contribution at $\mathcal{O}(\mathcal{S}^0)$ to the eikonal phase coming from these two diagrams after summing (\ref{eq3}) and (\ref{eq5}) is the following, 
\begin{align}
    \begin{split}  \label{eqt2}
   \chi\Big|_{\mathcal{S}^0} =\frac{m_1m_2^3}{m_p^6}\Delta_{\mathfrak{02}}^{\textrm{poly}}\,,
 \end{split}
\end{align}
\enlargethispage{2\baselineskip}
where, 
\begin{thesiscompactdisplay}
\begin{align}
    \begin{split} 
\Delta_{\mathfrak{02}}^{\textrm{poly}}= &\Bigg(\frac{s_1^2s_2^2}{12288 \pi ^3 |\boldsymbol{b}|^2 x \left(x^2-1\right)^7}\Bigg)\Bigg[6 \left(-x^{12}+10 x^{10}+x^8+44 x^6+x^4+10 x^2-1\right) x^2 \log \left(\pi  |\boldsymbol{b}|^6\right)\\&\hspace{0.8cm}-9 x^{16}+x^{14} (-36 \gamma_E +56+\log (4096))+4 x^{12} (82 \gamma_E  -25-46 \log (2))\\&\hspace{0.8cm}-4 x^{10} (7 \gamma_E -6+35 \log (2))+2 x^8 (696 \gamma_E  -35-456 \log (2))\\&\hspace{0.8cm}-4 x^6 (7 \gamma_E  -6+35 \log (2))+4 x^4 (82 \gamma_E  -25-46 \log (2))\\&\hspace{0.8cm}+x^2 (-36 \gamma_E  +56+\log (4096))-9\Bigg]
     \end{split}
\end{align}
\end{thesiscompactdisplay}
and the total contribution at $\mathcal{O}(\mathcal{S}^1)$ to the eikonal phase coming from these two diagrams after summing (\ref{eq4}) and (\ref{eq6}) is the following, 
\begin{align}
    \begin{split}   \label{eqt3}
     \chi\Big|_{\mathcal{S}^1}=\frac{m_1m_2^3}{m_p^6}\Bigg[\Delta_{\mathfrak{03}}^{\textrm{poly}}(\hat{b}\cdot \mathcal{S}_1\cdot v_2)+\Delta_{\mathfrak{04}}^{\textrm{poly}}(\hat{b}\cdot \mathcal{S}_2\cdot v_1)\Bigg]\,,
     \end{split}
\end{align}
\enlargethispage{2\baselineskip}
where, 
\begin{thesiscompactdisplay}
\begin{align}
    \begin{split}
    \Delta_{\mathfrak{03}}^{\textrm{poly}}= \Bigg(\frac{s_1^2s_2^2\,x} {3072 \pi ^3 |\boldsymbol{b}|^3 \left(x^2-1\right)^5 }\Bigg)&\Bigg[-6 \left(x^8+6 x^7-2 x^6-26 x^5-6 x^4-26 x^3-2 x^2+6 x+1\right) \log (|\boldsymbol{b}|)\\&\hspace{0 cm}+5 x^8+x^8 \log (64)-x^8 \log (\pi )+x^7 \log (4096)-6 x^7 \log (\pi )-6 x^6\\ &+2 x^6 \log \left(\frac{\pi }{64}\right)-81 x^5+10 x^5 \log (4)-26 x^5 \log \left(\frac{16}{\pi }\right)+2 x^4\\ &-6 x^4 \log (64)+6 x^4 \log (\pi )-81 x^3+10 x^3 \log (4)-26 x^3 \log \left(\frac{16}{\pi }\right)-6 x^2\\ &+2 x^2 \log \left(\frac{\pi }{64}\right)-4 \gamma_E  \left(x^8+9 x^7-2 x^6-35 x^5-6 x^4-35 x^3-2 x^2+9 x+1\right)\\ &+x \log \left(\frac{4096}{\pi ^6}\right)+5+\log (64)-\log (\pi )\Bigg]\,,
        \end{split}
\end{align}
\begin{align}
    \begin{split}    \Delta_{\mathfrak{04}}^{\textrm{poly}}=\Bigg(\frac{s_1^2s_2^2\,x^2(1+x^2)} {768  \pi ^3 |\boldsymbol{b}|^3 \left(x^2-1\right)^7 }\Bigg)&\Bigg[12 \left(x^8+x^6-8 x^4+x^2+1\right) \log (|\boldsymbol{b}|)+13 x^8+x^8 \log \left(\frac{\pi ^2}{64}\right)-159 x^6\\ &-12 x^6 \log (2)+2 x^6 \log (\pi )+392 x^4+68 x^4 \log (2)-16 x^4 \log (\pi )\\ &-159 x^2-12 x^2 \log (2)+2 x^2 \log (\pi )+\gamma_{E}  \left(11 x^8+8 x^6-78 x^4+8 x^2+11\right)\\ &+13+\log \left(\frac{\pi ^2}{64}\right)\Bigg]\,.
    \end{split}
\end{align}
\end{thesiscompactdisplay}
%Numerator resposible for the spinning eikonal has the following form,
%\begin{align}
 %   \begin{split}
%       
 %   \end{split}
%\end{align}
%We need the following vector integrals for computing the spinning Eikonal phase.
%\begin{align}
%    \begin{split}
%    &    \boldsymbol{\mathscr{M}}_{7}^{\eta}=\int \frac{\hat\delta(\ell_1\cdot v_2)\hat\delta(\ell_2\cdot v_2)\ell_1^\eta}{(\ell_1\cdot v_1+i\varepsilon)^2\,(\ell_2\cdot v_1+i\varepsilon)^2\,\ell_1^2\,\ell_2^2\,(\ell_1+
%        \ell_2-q)^2}[]\\ &
%=\boldsymbol{\mathfrak{m}}_{11} q^\eta +\boldsymbol{\mathfrak{m}}_{12} (\gamma v_2^\eta-v_1^\eta)
%    \end{split}
%\end{align}
%where,
%\begin{align}
%    \begin{split}
%   &     \boldsymbol{\mathfrak{m}}_{11}=\frac{-4 \gamma  \left(6 \gamma ^2-13\right) \epsilon ^2+\gamma  \left(14 \gamma ^2-25\right) \epsilon +\gamma  \left(3-2 \gamma ^2\right)-144 \gamma  \epsilon ^4+36 \gamma  \epsilon ^3}{3 \left(\gamma ^2-1\right)^2  \epsilon  (3 \epsilon -1)}\,\mathbfcal{L}_{0,0;0,0,1,1,1}\\&+\frac{2 \gamma  \left(\gamma ^2 (4 \epsilon -2)+2 \epsilon ^2-2 \epsilon +1\right)}{3 \left(\gamma ^2-1\right)^2  \epsilon  (3 \epsilon -1)}q^2\,\mathbfcal{L}^{(++)}_{1,1;0,0,1,1,1}\\ &
 %   \boldsymbol{\mathfrak{m}}_{12}=\frac{2 \gamma  \left(-\gamma ^2+9 \left(\gamma ^2-1\right)-1\right) \epsilon +2 \gamma  \left(\gamma ^2-3 \left(\gamma ^2-1\right)\right)-4 \gamma  \epsilon ^2}{\left(\gamma ^2-1\right) (6 \epsilon -1)}q^2\,\mathbfcal{L}^{(++)}_{0,1;0,0,1,1,1}
%    \end{split}
%\end{align}
%\begin{align}
  %  \begin{split}
 %        \boldsymbol{\mathscr{M}}_{8}^{\eta}&=\int \frac{\hat\delta(\ell_1\cdot v_2)\hat\delta(\ell_2\cdot v_2)\ell_1^\eta}{(\ell_1\cdot v_1+i\varepsilon)^2\,(\ell_2\cdot v_1+i\varepsilon)^2\,\ell_1^2\,\ell_2^2\,(\ell_1+
%        \ell_2-q)^2}\left[\right]\\ &
%=\boldsymbol{\mathfrak{m}}_{13} q^\eta +\boldsymbol{\mathfrak{m}}_{14} (\gamma v_2^\eta-v_1^\eta)\equiv  \boldsymbol{\mathscr{M}}_{11}^{\eta}
 %   \end{split}
%\end{align}
%where,
%\begin{align}
%    \begin{split}
 %       \boldsymbol{\mathfrak{m}}_{13}=\frac{\gamma ^2 (2-8 \epsilon )+6 \epsilon  (1-4 \epsilon )-1}{2(\gamma ^2-1)  (6 \epsilon -1)}\mathbfcal{L}^{(++)}_{0,1;0,0,1,1,1},\,\, \boldsymbol{\mathfrak{m}}_{14}=\frac{-2 \gamma ^2+6 \epsilon +1}{(\gamma ^2-1)}\,\mathbfcal{L}_{0,0;0,0,1,1,1}.
%    \end{split}
%\end{align}
%\begin{align}
  %  \begin{split}
 %        \boldsymbol{\mathscr{M}}_{10}^{\eta}&=\int \frac{\hat\delta(\ell_1\cdot v_2)\hat\delta(\ell_2\cdot v_2)\ell_1^\eta}{(\ell_1\cdot v_1+i\varepsilon)^2\,(\ell_2\cdot v_1+i\varepsilon)^2\,\ell_1^2\,\ell_2^2\,(\ell_1+
 %       \ell_2-q)^2}\Big[\Big]\\ &
  %      =\boldsymbol{\mathfrak{m}}_{15} q^\eta +\boldsymbol{\mathfrak{m}}_{16} (\gamma v_2^\eta-v_1^\eta)
%    \end{split}
%\end{align}
%where,
%\begin{align}
 %   \begin{split}
 %   &   \boldsymbol{\mathfrak{m}}_{15}=\frac{1}{3 \left(\gamma ^2-1\right)^2 %(\epsilon -1)(3\epsilon -1) \epsilon }\Big[-2 \left(24 \gamma ^2+20\right) \epsilon ^3-2 \left(3-12 \gamma ^2\right) \epsilon -2 \left(2 \gamma ^2-1\right)-8 \left(\gamma ^4-4 \gamma ^2+6\right) \epsilon ^2+96 \epsilon ^4\Big]\\&\mathbfcal{L}_{0,0;0,0,1,1,1}
 %   -\frac{4 (\gamma -1) (\gamma +1) \left(7-5 \gamma ^2\right) \epsilon +4 (\gamma -1) (\gamma +1) \left(\gamma ^2-2\right)+8 (\gamma -1) (\gamma +1) \epsilon ^2}{3 \left(\gamma ^2-1\right)^2 (\epsilon -1) \epsilon  (3 \epsilon -1)}q^2\,\mathbfcal{L}^{(++)}_{1,1;0,0,1,1,1}\\ &
%\boldsymbol{\mathfrak{m}}_{16}=\frac{1}{(\gamma^2 -1)  (\epsilon -1) (6 \epsilon -1)}\Big[\left(
%\begin{array}{c}
% -2 \gamma ^4+\left(20 \gamma ^2-36\right) \epsilon ^2+7 \gamma ^2+\left(4 \gamma ^4-28 \gamma ^2+14\right) \epsilon +24 \epsilon ^3-3 \\
%\end{array}
%\right)\Big]q^2\,\mathbfcal{L}^{(++)}_{0,1;0,0,1,1,1}
%    \end{split}
%\end{align}
%And, we have the following tensor integral,
%\begin{align}
%    \begin{split}
%\boldsymbol{\mathscr{M}}_9^{\eta\delta}&=\int\frac{\hat\delta(\ell_1\cdot v_2)\hat\delta(\ell_2\cdot v_2) \,\,\ell_1^\eta\ell_2^\delta}{(\ell_1\cdot v_1+i\varepsilon)^2(\ell_2\cdot v_1-i\varepsilon)^2\,\ell_1^2\,\ell_2^2\,(\ell_1+\ell_2-q)^2}(\ell_1\cdot v_1)\Bigg[\Bigg]\\ &
%=\boldsymbol{\mathfrak{o}}_{11}q^{[\eta}v_1^{\delta]}+\boldsymbol{\mathfrak{o}}_{12} %q^{[\eta}v_2^{\delta]}\equiv  \boldsymbol{\mathscr{M}}_{12}^{\eta\delta}
%    \end{split}
%\end{align}
%where,
%\begin{align}
%    \begin{split}
   %     &\boldsymbol{\mathfrak{o}}_{11}=-\frac{18 \epsilon ^2-9 \epsilon +1}{3 %\left(\gamma ^2-1\right)^2 (3 \epsilon -1)}\mathbfcal{L}_{0,0;0,0,1,1,1}+\frac{\epsilon q^2 }{3 \left(\gamma ^2-1\right)^2 (3 \epsilon -1)}\mathbfcal{L}^{(+-)}_{1,1;0,0,1,1,1}\\ &
%\boldsymbol{\mathfrak{o}}_{12}=\frac{2 \gamma(1-  6  \epsilon)}{6 \left(\gamma ^2-1\right)^2}\mathbfcal{L}_{0,0;0,0,1,1,1}-\frac{\gamma  \epsilon q^2 }{6\left(\gamma ^2-1\right) (3 \epsilon -1)}\mathbfcal{L}^{(+-)}_{1,1;0,0,1,1,1}
%    \end{split}
%\end{align}
%\textcolor{red}{Final $q$ integral.}\\
%where the two master integral tooks the following form,
%\begin{align}
%    \begin{split}
 %       &\mathcal{M}_{0,0,0,1,1,1,0}(q)\equiv \int_{k_1,l}\frac{\hat\delta(k_1\cdot v_2)\hat\delta(l\cdot v_2)}{l^2(l-k_1)^2(k_1-q)^2}\\ &
  %      \mathcal{M}_{1,1,0,1,1,1,0}(q)\equiv\int_{k_1,l}\frac{\hat\delta(k_1\cdot v_2)\hat\delta(l\cdot v_2)}{l^2(l-k_1)^2(k_1-q)^2(l\cdot v_1)(k_1\cdot v_1)}
   % \end{split}
%\end{align}
%\newpage
\par
\textbullet $\,\,$ We have another two comb-type diagrams with two graviton line 
%\textcolor{black}{AB: Need to update the coefficients from the spin part !}
\begin{center}
\includegraphics[width=0.38\linewidth]{comb.png}
\end{center}
%\textcolor{black}{for the first diagram do the following change of variable,
%$k_1 \to \ell_1,\,k_3\to \ell_2,\,k_2\to q-\ell_1-\ell_2$.}\\

%Hence the contribution to the eikonal phase Numerator is given by,

%\begin{align}
%    \begin{split}&
%\textcolor{red}{\mathcal{N}|_{\mathcal{S}^0}}=\frac{1}{32 (\epsilon -1)^2}\Bigg[16 (\ell_1\cdot v_1)^2 \Bigg(2 \left(2 \gamma ^2 (\epsilon -1) (4 \epsilon -3)-1\right) (\ell_2\cdot v_1)^2+\gamma ^2 ((3-2 \epsilon ) \epsilon -1) q^2\Bigg)\\&+4 \left(\ell_1\cdot v_1\right) \left(\ell_2\cdot v_1\right) \left(q^2 \left(1-4 \gamma ^4 (\epsilon -1)^2\right)-8 \left(3 \gamma ^2 (\epsilon -1)+2\right) \left(\left(\ell_2\cdot v_1\right)\right)^2\right)\\&+64 \gamma ^2 (\epsilon -1) (2 \epsilon -1) \left(\ell_2\cdot v_1\right) (\ell_1\cdot v_1)^3+q^2 \left(2 \gamma ^2 (\epsilon -1)+1\right)\Bigg( (q^2 \left(2 \gamma ^2 (\epsilon -1)+1\right)-\\&8 \left(3 \gamma ^2 (\epsilon -1)+2\right) \left(\ell_2\cdot v_1\right)^2)\Bigg)\Bigg]
%  \end{split}
%\end{align}
For the first diagram, we have, 

\begin{align}
    \begin{split} \label{eq8}
        i\chi\Big|_{\mathcal{S}^0} =i\frac{m_1 m_2^3 s_1 s_2}{64 m_p^6}\int e^{-iq\cdot b}\hat\delta(q\cdot v_1)\hat\delta(q\cdot v_2)\Bigg( \boldsymbol{\tilde h}_{17}\mathbfcal{L}_{0,0;0,0,1,1,1}-\frac{1}{2}\boldsymbol{\tilde h}_{18}q^2\,\mathbfcal{L}^{(+-)}_{1,1;0,0,1,1,1}\Bigg)\,,
    \end{split}
\end{align}
and 
\begin{align}
    \begin{split} \label{eq9}
        i\chi\Big|_{\mathcal{S}^1} =i\frac{m_1 m_2^3 s_1 s_2}{64 m_p^6}\int e^{-iq\cdot b}\hat\delta(q\cdot v_1)\hat\delta(q\cdot v_2)&\Bigg[i\Bigg( \boldsymbol{\tilde h}_{19}\mathbfcal{L}_{0,0;0,0,1,1,1}-\frac{1}{2}\boldsymbol{\tilde h}_{20}q^2\,\mathbfcal{L}^{(+-)}_{1,1;0,0,1,1,1}\Bigg)(q\cdot \mathcal{S}_1\cdot v_2)\\&i\,\Bigg( \boldsymbol{\tilde h}_{21}\mathbfcal{L}_{0,0;0,0,1,1,1}-\frac{1}{2}\boldsymbol{\tilde h}_{22}q^2\,\mathbfcal{L}^{(+-)}_{1,1;0,0,1,1,1}\Bigg)(q\cdot \mathcal{S}_2\cdot v_1)\Bigg]\,,
    \end{split}
\end{align}
%where,






%\newpage
%\begin{align}
%   \begin{split}&\textcolor{red}{\mathcal{N}|_{\mathcal{S}^1}}=
%\frac{\ell_1\cdot\mathcal{S}_1\cdot \ell_2}{4(d-2)^2}\Bigg[i \Bigg((\ell_1\cdot v_1) (4 (\gamma ^2 (d^2-12 d+20)+4) (\ell_2\cdot v_1)^2+q^2 (-16 \gamma ^2-\gamma ^2 d^2+10 \gamma ^2 d+d-8))\\&+4 \gamma ^2 (d^2-10 d+16) (\ell_2\cdot v_1) (\ell_1\cdot v_1)^2+q^2 (-14 \gamma ^2-\gamma ^2 d^2+9 \gamma ^2 d+d-7) (\ell_2\cdot v_1)\Bigg)\Bigg]\\&+\textcolor{red}{\frac{\ell_1\cdot\mathcal{S}_1\cdot q}{4(d-2)^2}}\Bigg[-i (\ell_1\cdot v_1) \Bigg(4 (20 \gamma ^2+\gamma ^2 d^2-(12 \gamma ^2+1) d+11) (\ell_2\cdot v_1)^2+q^2 (-16 \gamma ^2-\gamma ^2 d^2+10 \gamma ^2 d+d-8)\\&+4 (d-8) (\gamma ^2 (d-2)-1) (\ell_1\cdot v_1) (\ell_2\cdot v_1)\Bigg)\Bigg]\\&\hspace{0cm}+\textcolor{red}{\frac{\ell_1\cdot\mathcal{S}_2\cdot q}{2(d-2)}}\Bigg[i \gamma  (\ell_1\cdot v_1) \Bigg(2 (5 \gamma ^2 (d-2)-6) (\ell_2\cdot v_1)^2+4 (\gamma ^2 (d-2)-1) (\ell_1\cdot v_1) (\ell_2\cdot v_1)+q^2 (1-\gamma ^2 (d-2))\Bigg)\Bigg]\\&\hspace{0cm}+
% \textcolor{red}{\frac{\ell_1 \cdot \mathcal{S}_1 \cdot v_2}{2(d-2)}}\Bigg[i \gamma  (-(\ell_1\cdot v_1)^2 (4 (\ell_2\cdot v_1)^2+(d-6) q^2)+4 (d-6) (\ell_2\cdot v_1)(\ell_1\cdot v_1)^3+q^2 (\gamma ^2 (d-2)-2) (\ell_2\cdot v_1)^2+\\&8 (\ell_1\cdot v_1)(\ell_2\cdot v_1)^3)\Bigg]+ \textcolor{red}{\ell_1 \cdot \mathcal{S}_2 \cdot v_1}\Bigg[-\frac{i \gamma  (\ell_2\cdot v_1)^2 (8 (\ell_1\cdot v_1)^2+12 (\ell_1\cdot v_1) (\ell_2\cdot v_1)-q^2)}{2 (d-2)}\Bigg]\\&+ \textcolor{red}{\ell_2 \cdot \mathcal{S}_1 \cdot q }\Bigg[\frac{i (d-7) (\ell_2\cdot v_1) (4 (\ell_1\cdot v_1)^2+q^2 (1-\gamma ^2 (d-2))+4 (\ell_1\cdot v_1) (\ell_2\cdot v_1))}{4 (d-2)^2}\Bigg]\\&+\textcolor{red}{\ell_2 \cdot \mathcal{S}_2 \cdot q}\Bigg[\frac{i \gamma  (\ell_2\cdot v_1) (-4 (d-2) (\ell_1\cdot v_1)^2+q^2 (\gamma ^2 (d-2)-1)-4 (\ell_1\cdot v_1) (\ell_2\cdot v_1))}{4 (d-2)}\Bigg]\\&+\textcolor{red}{\frac{\ell_2 \cdot \mathcal{S}_1 \cdot v_2}{8(d-2)}}\Bigg[-i \gamma  \Bigg(-4 (\ell_1\cdot v_1) (8(\ell_2\cdot v_1)^3-\gamma ^2 (d-2) q^2 (\ell_2\cdot v_1))-4 (\ell_1\cdot v_1)^2 (4 (2 d-7) (\ell_2\cdot v_1)^2+\\&(9-2 d) q^2)-16 (2 d-9) (\ell_2\cdot v_1)(\ell_1\cdot v_1)^3+q^2 (\gamma ^2 (d-2)-1) (8 (\ell_2\cdot v_1)^2-q^2)\Bigg)\Bigg]\\&+ \textcolor{red}{\frac{\ell_2 \cdot \mathcal{S}_2 \cdot v_1}{8(d-2)}} \Bigg[i \gamma  \Bigg((\ell_1\cdot v_1) (4 \gamma ^2 (d-2) q^2 (\ell_2\cdot v_1)-48(\ell_2\cdot v_1)^3)-4 (\ell_1\cdot v_1)^2 (8 (d-2) (\ell_2\cdot v_1)^2-\\&(d-3) q^2)-16 (d-3) (\ell_2\cdot v_1)(\ell_1\cdot v_1)^3+q^2 (\gamma ^2 (d-2)-1) (12 (\ell_2\cdot v_1)^2-q^2)\Bigg)\Bigg]\\&+\textcolor{red}{\frac{q \cdot \mathcal{S}_1 \cdot v_2}{8(d-2)}}\Bigg[i \gamma\Bigg(-4 (\ell_1\cdot v_1) (8(\ell_2\cdot v_1)^3-\gamma ^2 (d-2) q^2 (\ell_2\cdot v_1))-4 (\ell_1\cdot v_1)^2 (4 (2 d-7) (\ell_2\cdot v_1)^2\\&+(9-2 d) q^2)-16 (2 d-9) (\ell_2\cdot v_1)(\ell_1\cdot v_1)^3+q^2 (\gamma ^2 (d-2)-1) (8 (\ell_2\cdot v_1)^2-q^2)\Bigg)\Bigg]\nonumber\\&
%\end{split}
 %\end{align}
 
 
% \begin{align}\begin{split}&+ \textcolor{red}{\frac{q \cdot \mathcal{S}_2 \cdot v_1}{8(d-2)}}\Bigg[-i \gamma  \Bigg((\ell_1\cdot v_1) (4 \gamma ^2 (d-2) q^2 (\ell_2\cdot v_1)-48(\ell_2\cdot v_1)^3)-4 (\ell_1\cdot v_1)^2 (8 (d-2) (\ell_2\cdot v_1)^2-\\&(d-3) q^2)-16 (d-3) (\ell_2\cdot v_1)(\ell_1\cdot v_1)^3+q^2 (\gamma ^2 (d-2)-1) (12 (\ell_2\cdot v_1)^2-q^2)\Bigg)\Bigg]\\&- \frac{\ell_1\cdot \mathcal{S}_2\cdot \ell_2}{4 (d-2)}\Bigg[i \gamma  \left(2 \left(\gamma ^2 (d-2)-1\right) \left(\ell_1 \cdot v_1\right) \left(10 \left( \ell_2 \cdot v_1\right)^2-q^2\right)\\&+4 \left(2 \gamma ^2 (d-2)+d-3\right) \left(\ell_2 \cdot v_1\right)  \left(\ell_1 \cdot v_1\right)^2+q^2 \left(1-\gamma ^2 (d-2)\right) \left(\ell_2 \cdot v_1\right)\right)\Bigg] 
% \end{split}
%\end{align}
%The spinning eikonal contribution from the first diagram looks like,



%\begin{align}
 %   \begin{split}
 %   &    \boldsymbol{\mathscr{M}}_{1}^{\eta\gamma}=\int \frac{\hat\delta(\ell_1\cdot %v_2)\hat\delta(\ell_2\cdot v_2)\ell_1^\eta\ell_2^\gamma}{(\ell_1\cdot v_1+i\varepsilon)^2\,(\ell_2\cdot v_1+i\varepsilon)^2\,\ell_1^2\,\ell_2^2\,(\ell_1+
 %       \ell_2-q)^2}[]\\ &
%=\boldsymbol{\mathfrak{n}}_{0} (q^{[\eta} v_1^{\gamma]})+\boldsymbol{\mathfrak{n}}_{1}(q^{[\eta} v_2^{\gamma]}) 
%    \end{split}
%\end{align}
%Where ,
%\begin{align}
%    \begin{split}
%   \boldsymbol{\mathfrak{n}}_{0}&=\frac{i}{6 \left(\gamma ^2-1\right)^2 \epsilon}\Bigg[-2 \gamma ^2-24 \left(4 \gamma ^4-\gamma ^2\right) \epsilon ^3-36 \left(4 \gamma ^4-6 \gamma ^2+1\right) \epsilon ^2+2 \left(120 \gamma ^4-143 \gamma ^2+24\right) \epsilon +1\Bigg]\mathbfcal{L}_{0,0;0,0,1,1,1}\\&+\frac{i}{6 \left(\gamma ^2-1\right)^2 \epsilon}\Bigg[2 \gamma ^2 \epsilon -2 \gamma ^2+1\Bigg]q^2\mathbfcal{L}_{1,1;0,0,1,1,1}
%    \end{split}
%\end{align}

%\begin{align}
%    \begin{split}
%   \boldsymbol{\mathfrak{n}}_{1}&=\frac{i}{6 \left(\gamma ^2-1\right)^2 \epsilon}\Bigg[\gamma  \left(2 \gamma ^2-1\right)+24 \left(4 \gamma ^5-\gamma ^3\right) \epsilon ^3+36 \left(4 \gamma ^5-6 \gamma ^3+\gamma \right) \epsilon ^2-2 \left(120 \gamma ^5-143 \gamma ^3+24 \gamma \right) \epsilon\Bigg]\mathbfcal{L}_{0,0;0,0,1,1,1}\\&+\frac{i}{6 \left(\gamma ^2-1\right)^2 \epsilon}\Bigg[\gamma  \left(2 \gamma ^2-1\right)-2 \gamma ^3 \epsilon\Bigg]q^2\mathbfcal{L}_{1,1;0,0,1,1,1}
%    \end{split}
%\end{align}

%\begin{align}
%    \begin{split}
%    &    \boldsymbol{\mathscr{M}}_{12}^{\eta\gamma}=\int \frac{\hat\delta(\ell_1\cdot v_2)\hat\delta(\ell_2\cdot v_2)\ell_1^\eta\ell_2^\gamma}{(\ell_1\cdot v_1+i\varepsilon)^2\,(\ell_2\cdot v_1+i\varepsilon)^2\,\ell_1^2\,\ell_2^2\,(\ell_1+
 %       \ell_2-q)^2}[]\\ &
%=\boldsymbol{\mathfrak{n}}_{12} (q^{[\eta} %v_1^{\gamma]})+\boldsymbol{\mathfrak{n}}_{13}(q^{[\eta} v_2^{\gamma]}) 
%    \end{split}
%\end{align}
%Where ,
%\begin{align}
%    \begin{split}
%   \boldsymbol{\mathfrak{n}}_{13}&=\frac{i}{6 \left(\gamma ^2-1\right)^2 \epsilon}\Bigg[-72 \gamma ^3 \epsilon ^3+\gamma  \left(1-2 \gamma ^2\right)-12 \left(14 \gamma ^5-18 \gamma ^3+\gamma \right) \epsilon ^2+2 \left(84 \gamma ^5-107 \gamma ^3+24 \gamma \right) \epsilon\Bigg]\mathbfcal{L}_{0,0;0,0,1,1,1}\\&+\frac{i}{6 \left(\gamma ^2-1\right)^2 \epsilon}\Bigg[2 \gamma ^3 \epsilon +\left(1-2 \gamma ^2\right) \gamma\Bigg]q^2\mathbfcal{L}_{1,1;0,0,1,1,1}
 %   \end{split}
%\end{align}

%\begin{align}
%    \begin{split}
%   \boldsymbol{\mathfrak{n}}_{12}&=\frac{i}{6 \left(\gamma ^2-1\right)^2 \epsilon}\Bigg[72 \gamma ^2 \epsilon ^3+2 \gamma ^2+12 \left(14 \gamma ^4-18 \gamma ^2+1\right) \epsilon ^2-2 \left(84 \gamma ^4-107 \gamma ^2+24\right) \epsilon -1\Bigg]\mathbfcal{L}_{0,0;0,0,1,1,1}\\&+\frac{i}{6 \left(\gamma ^2-1\right)^2 \epsilon}\Bigg[-2 \gamma ^2 \epsilon +2 \gamma ^2-1\Bigg]q^2\mathbfcal{L}_{1,1;0,0,1,1,1}
%    \end{split}
%\end{align}

%\begin{align}
%    \begin{split}
%    &    \boldsymbol{\mathscr{M}}_{2}^{\eta}=\int \frac{\hat\delta(\ell_1\cdot v_2)\hat\delta(\ell_2\cdot v_2)\ell_1^\eta}{(\ell_1\cdot v_1+i\varepsilon)^2\,(\ell_2\cdot v_1+i\varepsilon)^2\,\ell_1^2\,\ell_2^2\,(\ell_1+
%        \ell_2-q)^2}[]\\ &
%=\boldsymbol{\mathfrak{n}}_{2} (v_1^\eta-\gamma v_2^\eta)
%    \end{split}
%\end{align}
%Where ,
%\begin{align}
 %   \begin{split}
%   \boldsymbol{\mathfrak{n}}_{2}&=\frac{1}{\left(\gamma ^2-1\right) (2 \epsilon +1)}\Bigg[144 i \gamma ^2 \epsilon ^4-8 i \left(2 \gamma ^4-20 \gamma ^2-9\right) \epsilon ^3-4 i \left(14 \gamma ^4+61 \gamma ^2-38\right) \epsilon ^2+\\&2 i \left(20 \gamma ^4-44 \gamma ^2+21\right) \epsilon +4 i \left(8 \gamma ^4-11 \gamma ^2+4\right)\Bigg]\mathbfcal{L}_{0,0;0,0,1,1,1}
 %   \end{split}
%\end{align}


%\begin{align}
%    \begin{split}
%    &    \boldsymbol{\mathscr{M}}_{3}^{\eta}=\int \frac{\hat\delta(\ell_1\cdot v_2)\hat\delta(\ell_2\cdot v_2)\ell_1^\eta}{(\ell_1\cdot v_1+i\varepsilon)^2\,(\ell_2\cdot v_1+i\varepsilon)^2\,\ell_1^2\,\ell_2^2\,(\ell_1+
%        \ell_2-q)^2}[]\\ &
%=\boldsymbol{\mathfrak{n}}_{3} (v_1^\eta-\gamma v_2^\eta)
%    \end{split}
%\end{align}

%Where ,
%\begin{align}
%    \begin{split}
%   \boldsymbol{\mathfrak{n}}_{3}&=\frac{1}{\left(\gamma ^2-1\right) (2 \epsilon +1)}\Bigg[72 i \gamma ^2 \epsilon ^3+i \left(3-4 \gamma ^2\right)^2-4 i \left(8 \gamma ^4+10 \gamma ^2-9\right) \epsilon ^2+2 i \left(8 \gamma ^4-19 \gamma ^2+10\right) \epsilon\Bigg]\mathbfcal{L}_{0,0;0,0,1,1,1}
%    \end{split}
%\end{align}




%\begin{align}
%    \begin{split}
%    &    \boldsymbol{\mathscr{M}}_{4}^{\eta}=\int \frac{\hat\delta(\ell_1\cdot v_2)\hat\delta(\ell_2\cdot v_2)\ell_1^\eta}{(\ell_1\cdot v_1+i\varepsilon)^2\,(\ell_2\cdot v_1+i\varepsilon)^2\,\ell_1^2\,\ell_2^2\,(\ell_1+
%        \ell_2-q)^2}[]\\ &
%=\boldsymbol{\mathfrak{n}}_{4} q^\eta 
%    \end{split}
%\end{align}
%Where ,
%\begin{align}
%    \begin{split}
%   \boldsymbol{\mathfrak{n}}_{4}&=\frac{1}{6 (\gamma -1) (\gamma +1) (\epsilon -1) (2 \epsilon +1)}\Bigg[8 i \left(14 \gamma -41 \gamma ^3\right) \epsilon ^3-12 i \gamma ^3+144 i \gamma  \left(\gamma ^2-1\right) \epsilon ^4+36 i \gamma  \left(6 \gamma ^2-1\right) \epsilon ^2-\\&16 i \gamma  \left(2 \gamma ^2-5\right) \epsilon\Bigg]\mathbfcal{L}_{0,0;0,0,1,1,1}
 %   \end{split}
%\end{align}
%\begin{align}
%    \begin{split}
 %   &    \boldsymbol{\mathscr{M}}_{5}^{\eta}=\int \frac{\hat\delta(\ell_1\cdot v_2)\hat\delta(\ell_2\cdot v_2)\ell_1^\eta}{(\ell_1\cdot v_1+i\varepsilon)^2\,(\ell_2\cdot v_1+i\varepsilon)^2\,\ell_1^2\,\ell_2^2\,(\ell_1+
%        \ell_2-q)^2}[]\\ &
%=\boldsymbol{\mathfrak{n}}_{5} q^\eta 
%    \end{split}
%\end{align}
%Where ,
%\begin{align}
%    \begin{split}
%   \boldsymbol{\mathfrak{n}}_{5}&=\frac{1}{6 (\gamma -1) (\gamma +1) (\epsilon -1) (2 \epsilon +1)}\Bigg[28 i \gamma  \left(2 \gamma ^2-5\right) \epsilon ^2-4 i \gamma  \left(10 \gamma ^2-13\right) \epsilon -34 i \gamma  \left(\gamma ^2-1\right)+72 i \gamma  \epsilon ^3\Bigg]\mathbfcal{L}_{0,0;0,0,1,1,1}
 %   \end{split}
%\end{align}

%\begin{align}
%    \begin{split}
 %   &    \boldsymbol{\mathscr{M}}_{6}^{\eta}=\int \frac{\hat\delta(\ell_1\cdot v_2)\hat\delta(\ell_2\cdot v_2)\ell_2^\eta}{(\ell_1\cdot v_1+i\varepsilon)^2\,(\ell_2\cdot v_1+i\varepsilon)^2\,\ell_1^2\,\ell_2^2\,(\ell_1+
 %       \ell_2-q)^2}[]\\ &
%=\boldsymbol{\mathfrak{n}}_{6} (v_2^\eta-\gamma v_1^\eta)
%    \end{split}
%\end{align}
%Where ,
%\begin{align}
%    \begin{split}
 %   &    
%\boldsymbol{\mathfrak{n}}_{6}=\frac{i}{\left(\gamma ^2-1\right) (2 \epsilon +1)}\Bigg[-144 \gamma ^2 \epsilon ^4-72 \left(\gamma ^2+1\right) \epsilon ^3+4 \left(53 \gamma ^2-25\right) \epsilon ^2+\left(18-14 \gamma ^2\right) \epsilon -12 \gamma ^2+9\Bigg]\mathbfcal{L}_{0,0;0,0,1,1,1}
%\end{split}
%\end{align}
%\begin{align}
 %   \begin{split}
 %   &    \boldsymbol{\mathscr{M}}_{7}^{\eta}=\int \frac{\hat\delta(\ell_1\cdot %v_2)\hat\delta(\ell_2\cdot v_2)\ell_2^\eta}{(\ell_1\cdot v_1+i\varepsilon)^2\,(\ell_2\cdot v_1+i\varepsilon)^2\,\ell_1^2\,\ell_2^2\,(\ell_1+
%        \ell_2-q)^2}[]\\ &
%=\boldsymbol{\mathfrak{n}}_{7} (v_2^\eta-\gamma v_1^\eta)
%    \end{split}
%\end{align}
%Where ,
%\begin{align}
%    \begin{split}
%    &    
%\boldsymbol{\mathfrak{n}}_{7}=\frac{i\gamma}{\left(\gamma ^2-1\right) (2 \epsilon +1)}\Bigg[-72 \gamma ^3 \epsilon ^3+\left(88 \gamma ^3-52 \gamma \right) \epsilon ^2+\left(4 \gamma -2 \gamma ^3\right) \epsilon -8 \gamma ^3+7 \gamma\Bigg]\mathbfcal{L}_{0,0;0,0,1,1,1}
%\end{split}
%\end{align}





%\begin{align}
 %   \begin{split}
 %   &    \boldsymbol{\mathscr{M}}_{8}^{\eta}=\int \frac{\hat\delta(\ell_1\cdot v_2)\hat\delta(\ell_2\cdot v_2)\ell_2^\eta}{(\ell_1\cdot v_1+i\varepsilon)^2\,(\ell_2\cdot v_1+i\varepsilon)^2\,\ell_1^2\,\ell_2^2\,(\ell_1+
 %       \ell_2-q)^2}[]\\ &
%=\boldsymbol{\mathfrak{n}}_{8} q^\eta 
%    \end{split}
%\end{align}
%Where ,
%\begin{align}
 %   \begin{split}
%   &    
%\boldsymbol{\mathfrak{n}}_{8}=\frac{1}{3 \left(\gamma ^2-1\right)^2 (\epsilon +1)}\Bigg[
%864 i \gamma  \left(\gamma ^2-3\right) \epsilon ^6+144 i \gamma  \left(4 \gamma ^4-18 \gamma ^2+11\right) \epsilon ^5-4 i \gamma  \left(37 \gamma ^4+409 \gamma ^2-620\right) \epsilon ^4-\\&4 i \gamma  \left(167 \gamma ^4-353 \gamma ^2+183\right) \epsilon ^3+i \gamma  \left(51 \gamma ^4+5 \gamma ^2-74\right) \epsilon ^2-4 i \gamma  \left(\gamma ^4+2 \gamma ^2-3\right) \epsilon +i \gamma ^3 \left(\gamma ^2-1\right)\Bigg]\mathbfcal{L}_{0,0;0,0,1,1,1}\\&
%+\frac{1}{3 \left(\gamma ^2-1\right)^2 (\epsilon +1)}\Bigg[12 i \gamma  \left(\gamma ^2-3\right) \epsilon ^3-6 i \gamma  \left(7 \gamma ^2-8\right) \epsilon -2 i \gamma  \left(\gamma ^4+14 \gamma ^2-15\right) \epsilon ^2+2 i \gamma  \left(\gamma ^4-7 \gamma ^2+6\right)\Bigg]q^2\mathbfcal{L}_{1,1;0,0,1,1,1}
%\end{split}
%\end{align}

%\begin{align}
%    \begin{split}
%    &    \boldsymbol{\mathscr{M}}_{9}^{\eta}=\int \frac{\hat\delta(\ell_1\cdot v_2)\hat\delta(\ell_2\cdot v_2)\ell_2^\eta}{(\ell_1\cdot v_1+i\varepsilon)^2\,(\ell_2\cdot v_1+i\varepsilon)^2\,\ell_1^2\,\ell_2^2\,(\ell_1+
%        \ell_2-q)^2}[]\\ &
%=\boldsymbol{\mathfrak{n}}_{9} q^\eta
%    \end{split}
%\end{align}
%Where ,
%\begin{align}
 %   \begin{split}
%    &    
%\boldsymbol{\mathfrak{n}}_{9}=\frac{1}{3 \left(\gamma ^2-1\right)^2 (\epsilon -1) \epsilon  (\epsilon +1) (2 \epsilon +1)}\Bigg[3456 i \gamma ^3 \epsilon ^7-144 i \gamma  \left(12 \gamma ^4+19 \gamma ^2-16\right) \epsilon ^6+24 i \gamma  \left(72 \gamma ^4-212 \gamma ^2+19\right) \epsilon ^5\\&+4 i \gamma  \left(486 \gamma ^4+419 \gamma ^2-512\right) \epsilon ^4-2 i \gamma  \left(910 \gamma ^4-1320 \gamma ^2+371\right) \epsilon ^3-2 i \gamma  \left(92 \gamma ^4-76 \gamma ^2+5\right) \epsilon ^2\\&+16 i \gamma  \left(8 \gamma ^4-15 \gamma ^2+7\right) \epsilon +4 i \gamma ^3 \left(\gamma ^2-1\right)\Bigg]\mathbfcal{L}_{0,0;0,0,1,1,1}\\&
%+\frac{1}{\left(\gamma ^2-1\right)^2 (\epsilon +1)}\Bigg[48 i \gamma ^3 \epsilon ^3-12 i \gamma  \left(3 \gamma ^2-2\right) \epsilon +2 i \gamma  \left(-8 \gamma ^4+8 \gamma ^2+12\right) \epsilon ^2+2 i \gamma  \left(8 \gamma ^4-14 \gamma ^2+3\right)\Bigg]q^2\mathbfcal{L}_{1,1;0,0,1,1,1}
%\end{split}
%\end{align}



%\begin{align}
%    \begin{split}
%    &    \boldsymbol{\mathscr{M}}_{10}=\int \frac{\hat\delta(\ell_1\cdot v_2)\hat\delta(\ell_2\cdot v_2)}{(\ell_1\cdot v_1+i\varepsilon)^2\,(\ell_2\cdot v_1+i\varepsilon)^2\,\ell_1^2\,\ell_2^2\,(\ell_1+
%   \ell_2-q)^2}[]\\ &
%   =\boldsymbol{\mathfrak{n}}_{10}
%    \end{split}
%\end{align}
%Where ,
%\begin{align}
%    \begin{split}
%    &    
%\boldsymbol{\mathfrak{n}}_{10}=\frac{1}{\left(\gamma ^2-1\right)^2 (\epsilon +1)}\Bigg[1728 i \gamma ^3 \epsilon ^5-4 i \gamma  \left(\gamma ^2-1\right) \left(6 \gamma ^2-5\right)-288 i \gamma  \left(2 \gamma ^4+\gamma ^2-5\right) \epsilon ^4+64 i \gamma  \left(\gamma ^4-41 \gamma ^2+28\right) \epsilon ^3\\&+8 i \gamma  \left(84 \gamma ^4-145 \gamma ^2+59\right) \epsilon ^2+4 i \gamma  \left(2 \gamma ^4+13 \gamma ^2-10\right) \epsilon\Bigg]\mathbfcal{L}_{0,0;0,0,1,1,1}\\&
%+\frac{1}{\left(\gamma ^2-1\right)^2 (\epsilon +1)}\Bigg[24 i \gamma ^3 \epsilon ^3-4 i \gamma  \left(\gamma ^2+1\right) \left(2 \gamma ^2-3\right) \epsilon ^2-2 i \gamma  \left(8 \gamma ^2-5\right) \epsilon +2 i \left(4 \gamma ^5-6 \gamma ^3+\gamma \right)\Bigg]q^2\mathbfcal{L}_{1,1;0,0,1,1,1}
%\end{split}
%\end{align}


%\begin{align}
%    \begin{split}
 %   &    \boldsymbol{\mathscr{M}}_{11}=\int \frac{\hat\delta(\ell_1\cdot %v_2)\hat\delta(\ell_2\cdot v_2)}{(\ell_1\cdot v_1+i\varepsilon)^2\,(\ell_2\cdot v_1+i\varepsilon)^2\,\ell_1^2\,\ell_2^2\,(\ell_1+
%   \ell_2-q)^2}[]\\ &
%   =\boldsymbol{\mathfrak{n}}_{11}
 %   \end{split}
%\end{align}
%Where ,
%\begin{align}
 %   \begin{split}
 %   &    
%\boldsymbol{\mathfrak{n}}_{11}=\frac{1}{\left(\gamma ^2-1\right)^2 (\epsilon +1)}\Bigg[-1728 i \gamma ^3 \epsilon ^5+288 i \gamma  \left(\gamma ^2+1\right) \left(3 \gamma ^2-4\right) \epsilon ^4+8 i \gamma  \left(\gamma ^2-1\right) \left(7 \gamma ^2-5\right)-\\&48 i \gamma  \left(2 \gamma ^4-47 \gamma ^2+29\right) \epsilon ^3-8 i \gamma  \left(121 \gamma ^4-165 \gamma ^2+42\right) \epsilon ^2+4 i \gamma  \left(12 \gamma ^4-38 \gamma ^2+21\right) \epsilon\Bigg]\mathbfcal{L}_{0,0;0,0,1,1,1}\\&
%+\frac{1}{\left(\gamma ^2-1\right)^2 (\epsilon +1)}\Bigg[-24 i \gamma ^3 \epsilon ^3+4 i \gamma  \left(\gamma ^2+1\right) \left(2 \gamma ^2-3\right) \epsilon ^2+2 i \gamma  \left(8 \gamma ^2-5\right) \epsilon -2 i \left(4 \gamma ^5-6 \gamma ^3+\gamma \right)\Bigg]q^2\mathbfcal{L}_{1,1;0,0,1,1,1}
%\end{split}
%\end{align}





















We now compute the eikonal phase for the second diagram.
%\begin{align}
%    \begin{split}
%        \chi=-i\frac{m_1 m_2^3 s_1 s_2}{64 m_p^6}\int e^{-iq\cdot b}\hat\delta(q\cdot v_1)\hat\delta(q\cdot v_2)\int \frac{\hat\delta(\ell_1\cdot v_2)\hat\delta(\ell_2\cdot v_2)\mathcal{N}}{(\ell_1\cdot v_1+i\varepsilon)^2\,(\ell_2\cdot v_1+i\varepsilon)^2\,\ell_1^2\,\ell_2^2\,(\ell_1+\ell_2-q)^2},
%    \end{split}
%\end{align}
%where the spinless part of the numerator takes the following form,
%\begin{align}
%    \begin{split}
%\mathcal{N}|_{\mathcal{S}^0}=
%       & 
%    \end{split}
%\end{align}
The spinless eikonal phase is given by,
\begin{align}
    \begin{split} \label{eq10}
        \chi\Big|_{\mathcal{S}^0} =-i\frac{m_1 m_2^3 s_1 s_2}{64 m_p^6} \int e^{-iq\cdot b}\hat\delta(q\cdot v_1)\,\hat\delta(q\cdot v_2)\,\Bigg[\boldsymbol{\tilde h}_{23} \,\mathbfcal{L}_{0,0;0,0,1,1,1}+\frac{1}{4}\,\boldsymbol{\tilde h}_{24}\,q^2\,\mathbfcal{L}^{(++)}_{1,1;0,0,1,1,1}\Bigg]
    \end{split}
\end{align}
The spinning eikonal phase is given by,
\begin{align}
    \begin{split} \label{eq11}
        i\chi\Big|_{\mathcal{S}^1} =-i\frac{m_1 m_2^3 s_1 s_2}{64 m_p^6}\int e^{-iq\cdot b}\hat\delta(q\cdot v_1)\hat\delta(q\cdot v_2)&\Bigg[i\Bigg( \boldsymbol{\tilde h}_{25}\mathbfcal{L}_{0,0;0,0,1,1,1}+\frac{1}{4}\boldsymbol{\tilde h}_{26}q^2\,\mathbfcal{L}^{(++)}_{1,1;0,0,1,1,1}\Bigg)(q\cdot \mathcal{S}_1\cdot v_2)\\&i\,\Bigg( \boldsymbol{\tilde h}_{27}\mathbfcal{L}_{0,0;0,0,1,1,1}+\frac{1}{4}\boldsymbol{\tilde h}_{28}q^2\,\mathbfcal{L}^{(++)}_{1,1;0,0,1,1,1}\Bigg)(q\cdot \mathcal{S}_2\cdot v_1)\Bigg]\,,
    \end{split}
\end{align}
Again, all the coefficients $\boldsymbol{\tilde h}'s$ are defined in Appendix ~(\ref{ch3:app:E}).
Finally, the total contribution at $\mathcal{O}(\mathcal{S}^0)$ to the eikonal phase coming from these two diagrams after summing (\ref{eq8}) and (\ref{eq10}) is the following, 
\begin{align}
    \begin{split}  \label{eqt4}
     \chi\Big|_{\mathcal{S}^0}=\frac{m_1m_2^3}{m_p^6}\Delta_{\mathfrak{05}}^{\textrm{poly}}
     \end{split}
\end{align}
where,
\begin{align}
    \begin{split}
\Delta_{\mathfrak{05}}^{\textrm{poly}}=\frac{s_1s_2}{24576 \pi ^3 |\boldsymbol{b}|^2 x \left(x^2-1\right)^5}\Bigg[&24 \left(x^8+x^4\right) \log \left(\pi  |\boldsymbol{b}|^6\right)+x^2 \Big(-5 x^6 (37+20 \log (2))-4 x^4 (\log (4)-47)\\ &-5 x^2 (37+20 \log (2))+x^8 \left(x^2 (\log (16)-43)+134+\log (16)\right)\\ &+134+\log (16)\Big)+2 \gamma_{E}  \left(x^{12}+x^{10}+59 x^8-2 x^6+59 x^4+x^2+1\right)\\ &-43+\log (16)\Bigg]\,.
        \end{split}
\end{align}
The total contribution at $\mathcal{O}(\mathcal{S}^1)$ to the eikonal phase coming from these two diagrams after summing (\ref{eq9}) and (\ref{eq11}) is the following,
\begin{align}
    \begin{split}  \label{eqt5}
     \chi\Big|_{\mathcal{S}^1}=\frac{m_1m_2^3}{m_p^6}\Bigg[\Delta_{\mathfrak{06}}^{\textrm{poly}}(\hat{b}\cdot \mathcal{S}_1\cdot v_2)+\Delta_{\mathfrak{07}}^{\textrm{poly}}(\hat{b}\cdot \mathcal{S}_2\cdot v_1)\Bigg]\,,
     \end{split}
\end{align}
where, 
\begin{align}
    \begin{split}   &\Delta_{\mathfrak{06}}^{\textrm{poly}}=\frac{s_1s_2}{147456 \pi ^3 |\boldsymbol{b}|^3  \left(x^2-1\right)^7}\Bigg[3\alpha_0-3\alpha_1(x^2+x^{12})+3\alpha_2x^{14}-3\alpha_3(x^6+x^{8})+3\alpha_4(x^4+x^{10})\\&\hspace{0.8cm}-13824 x^{11}+41472 x^9-55296 x^7+41472 x^5-13824 x^3\Bigg]\,,\\
& \Delta_{\mathfrak{07}}^{\textrm{poly}}=-\frac{s_1 s_2}{24576 \pi ^3 |\boldsymbol{b}|^3 \left(x^2-1\right)^7 }\Bigg[\hat{\alpha}_0+8\hat{\alpha}_1(x^5+x^{9})+\hat{\alpha}_0x^{14}+\alpha_2(x^2+x^{12})+\hat{\alpha}_3(x^4+x^{10})\\&\hspace{0.8cm}+\hat{\alpha}_4(x^6+x^{8})+8\hat{\alpha}_5(x^3+x^{11})-360 x^{13}-360 x\\&\hspace{0.8cm}-16 x^7 (18 \log (|\boldsymbol{b}|)+14 \gamma_E -161-14 \log (2)+3 \log (\pi ))\Bigg]
\end{split}
\end{align}
\enlargethispage{2\baselineskip}
with, 
\begin{thesiscompactdisplay}
\begin{align}
    \begin{split} 
   & \alpha_0= 8 \log \left(\frac{\pi  |\boldsymbol{b}|^6}{16}\right)+72 \log ^2(b)+24 \left(5 \gamma_E -3+\log \left(\frac{\pi }{16}\right)\right) \log (|\boldsymbol{b}|)+50 \gamma_E ^2+\pi ^2-19 \gamma_E +89+\\&\hspace{1cm} 2 \log ^2(\pi )+48 \log (2)+16 \log (2) \log (4)+\log (4)+20 \gamma_E  \log \left(\frac{\pi }{16}\right)-16 \log (2) \log (\pi )-12 \log (\pi )\,,\\&
   \alpha_1=24 \log (|\boldsymbol{b}|) (3 \log (b)+5 \gamma_E +24)+4 \log \left(\frac{\pi }{16}\right) (6 \log (b)+5 \gamma_E )+\pi ^2+\gamma_E  (473+50 \gamma )+573+32 \log ^2(2)\\&\hspace{1cm}-2 (199+8 \log (\pi )) \log (2)+2 \log (\pi ) (48+\log (\pi ))\,,\\
    &\alpha_2=24 \left(\log \left(\frac{\pi  |\boldsymbol{b}|^3}{16}\right)+5 \gamma_E -1\right) \log (|\boldsymbol{b}|)+50 \gamma_E ^2+\pi ^2+89+\log (4) \left(9+8 \log \left(\frac{4}{\pi }\right)\right)+2 (\log (\pi )-2) \log (\pi )\\&\hspace{1cm}+\gamma_E  (-19-80 \log (2)+20 \log (\pi ))\,,\nonumber
     \end{split}
\end{align}
 \begin{align}
    \begin{split} 
& \alpha_3=-504 \log \left(\frac{\pi  |\boldsymbol{b}|^6}{16}\right)+24 \log (|\boldsymbol{b}|) (3 \log (b)+5 \gamma_E -3)+4 \log \left(\frac{\pi }{16}\right) (6 \log (|\boldsymbol{b}|)+5 \gamma_E)+\gamma_E (50 \gamma_E -2583)+\pi ^2\\&\hspace{1cm}+2245+(21+16 \log (2)) \log (4)+2 \log (\pi ) (-6-8 \log (2)+\log (\pi ))\,,\\&
\alpha_4=-20 \log \left(\frac{\pi  |\boldsymbol{b}|^6}{16}\right)+36 \left(6 \log (|\boldsymbol{b}|) (3 \log (|\boldsymbol{b}|)+5 \gamma_E -3)+\log \left(\frac{\pi }{16}\right) (6 \log (|\boldsymbol{b}|)+5 \gamma_E )\right)+450 \gamma_E ^2+9 \pi ^2-651 \gamma_E \\&\hspace{1cm}+2697+(205+144 \log (2)) \log (4)+18 \log (\pi ) (-6-8 \log (2)+\log (\pi ))\,,\\&\hat{\alpha}_0=11 \log \left(\frac{\pi  |\boldsymbol{b}|^6}{16}\right)+65 \gamma_E +54+10 \log (4)\,,\\&\hat{\alpha}_1=72 \log ^2(|\boldsymbol{b}|)+12 (10 \gamma_E-7-8 \log (2)+2 \log (\pi )) \log (|\boldsymbol{b}|)+50 \gamma_E ^2+\pi ^2-271+4 \log (2) (13+\log (256))+\\&\hspace{1cm}2 \left(\log \left(\frac{\pi }{256}\right)-7\right) \log (\pi )-4 \gamma_E  (18+20 \log (2)-5 \log (\pi))\,,\\&\hat{\alpha}_2=32 \log \left(\frac{\pi }{16}\right) (6 \log (|\boldsymbol{b}|)+5 \gamma_E)+6 \log (|\boldsymbol{b}|) (96 \log (|\boldsymbol{b}|)+160 \gamma_E -85)+400 \gamma_E ^2+8 \pi ^2-435 \gamma_E +684+\\&\hspace{1cm}16 \left(\log ^2(16)+\log ^2(\pi )+\log (2) (20-8 \log (\pi ))\right)-85 \log (\pi )\,,\\&\hat{\alpha}_3=-275 \log \left(\frac{\pi  |\boldsymbol{b}|^6}{16}\right)-64 \left(6 \log (|\boldsymbol{b}|) (3 \log (|\boldsymbol{b}|)+5 \gamma_E -3)+\log \left(\frac{\pi }{16}\right) (6 \log (|\boldsymbol{b}|)+5 \gamma_E )\right)-\\&\hspace{1cm}\gamma_E  (429+800 \gamma_E )-4 \left(4 \pi ^2+398+\log (2) (199+128 \log (2))+8 \log (\pi ) (-6-8 \log (2)+\log (\pi ))\right)\,,\\&\hat{\alpha}_4=-131 \log \left(\frac{\pi |\boldsymbol{b}|^6}{16}\right)-160 \left(6 \log (|\boldsymbol{b}|) (3 \log (|\boldsymbol{b}|)+5 \gamma_E -3)+\log \left(\frac{\pi }{16}\right) (6 \log (|\boldsymbol{b}|)+5 \gamma_E )\right)+\\&\hspace{1cm}\gamma_E  (1759-2000 \gamma_E)-40 \pi ^2+1814-4 \log (2) (473+320 \log (2))-80 \log (\pi ) (-6-8 \log (2)+\log (\pi ))\,,\\&
\hat{\alpha}_5=\log \left(\frac{\pi |\boldsymbol{b}|^6}{4}\right)+6 \gamma_E +163\,.
    \end{split}
\end{align}
\end{thesiscompactdisplay}
%\begin{align}   
%\begin{align}
 %  \begin{split}
 %     \hspace{0cm}\mathcal{N}|_{\mathcal{S}^1}&=\frac{i(\ell_1\cdot %\mathcal{S}_1\cdot q)}{4(d-2)^2}
%\Bigg[\left(1-\gamma^2(d-2)\right)(\ell_1-q)^2(\ell_1\cdot v_1)+4(\ell_1\cdot v_1)(\ell_2\cdot v_1)\left(\gamma^2(d-2)(\ell_2\cdot v_1)+\ell_1\cdot v_1\right)\Bigg]\\ &
%-\frac{i\gamma\,(\ell_1\cdot \mathcal{S}_1\cdot v_2)}{8(d-2)}\Bigg[\cdots\Bigg]+i\frac{\ell_1\cdot \mathcal{S}_1\cdot \ell_2}{4(d-2)^2}\Bigg[\cdots\Bigg]+i\frac{\ell_2\cdot \mathcal{S}_1\cdot q}{4(d-2)^2}\Bigg[\cdots\Bigg]-\frac{i\gamma \,(\ell_2\cdot \mathcal{S}_1\cdot v_2)}{8(d-2)}\Bigg[\cdots\Bigg]+\frac{i(\ell_2\cdot \mathcal{S}_1\cdot \ell_1)}{4(d-2)^2}\Bigg[\cdots\Bigg]\\ &
%+i\frac{\gamma (\ell_1\mathcal{S}_2\cdot q)}{4(d-2)}\Bigg[\cdots\Bigg]+i\frac{\gamma (\ell_1\cdot\mathcal{S}_2\cdot v_1)}{8(d-2)}\Bigg[\Bigg]-\frac{i\gamma (\ell_1\cdot \mathcal{S}_2\cdot \ell_2)}{4(d-2)}\Bigg[\cdots\Bigg]+\frac{i\gamma (\ell_2\cdot \mathcal{S}_2\cdot q)}{4(d-2)}\Bigg[\cdots\Bigg]+\frac{i\gamma (\ell_2\cdot \mathcal{S}_2\cdot v_1)}{8(d-2)}\Bigg[\cdots\Bigg]\\ &
%-\frac{i\gamma (\ell_2\cdot \mathcal{S}_2\cdot \ell_1)}{4(d-2)}\Bigg[\cdots\Bigg]
%    \end{split}
%\end{align}
%The integrals we need to evaluate are the following \footnote{We only focus on the parts of the integrals that gives non vanishing contributions to the eikonal phase.},
%\begin{align}
  %  \begin{split}
  %      &\mathbfcal{W}_1^\eta=\int \frac{\hat\delta(\ell_1\cdot v_2)\hat\delta(\ell_2\cdot v_2)\,\ell_1^\eta}{(\ell_1\cdot v_1+i\varepsilon)^2\,(\ell_2\cdot v_1+i\varepsilon)^2\,\ell_1^2\,\ell_2^2\,(\ell_1+\ell_2-q)^2}\Bigg[\left(1-\gamma^2(d-2)\right)(\ell_1-q)^2(\ell_1\cdot v_1)+4(\ell_1\cdot v_1)(\ell_2\cdot v_1)\\ &\hspace{0.8cm}\left(\gamma^2(d-2)(\ell_2\cdot v_1)+\ell_1\cdot v_1\right)\Bigg]
   %     \\ &
  %      \hspace{0.7 cm}=\boldsymbol{\mathfrak{u}_1 %}\,\mathbfcal{L}_{0,0;0,0,1,1,1}\,\frac{\gamma v_2^\eta-v_1^\eta}{\gamma^2-1}\\ &
 %     \mathbfcal{W}_2^\eta  =\int \frac{\hat\delta(\ell_1\cdot %v_2)\hat\delta(\ell_2\cdot v_2)\,\ell_1^\eta}{(\ell_1\cdot v_1+i\varepsilon)^2\,(\ell_2\cdot v_1+i\varepsilon)^2\,\ell_1^2\,\ell_2^2\,(\ell_1+\ell_2-q)^2}[]\\ &
 %    \hspace{0.8 cm} \sim \Big(\boldsymbol{\mathfrak{u}_2 }\mathbfcal{L}_{0,0;0,0,1,1,1}+\boldsymbol{\mathfrak{u}_3 }\,q^2\,\mathbfcal{L}^{(++)}_{1,1;0,0,1,1,1}\Big)\,q^\eta\\ &
 %        \mathbfcal{W}_3^{\eta\delta}  =\int \frac{\hat\delta(\ell_1\cdot v_2)\hat\delta(\ell_2\cdot v_2)\,\ell_1^\eta\,\ell_2^\delta}{(\ell_1\cdot v_1+i\varepsilon)^2\,(\ell_2\cdot v_1+i\varepsilon)^2\,\ell_1^2\,\ell_2^2\,(\ell_1+\ell_2-q)^2}[]
 %        \\ &\hspace{0.8 cm}=\frac{2}{\gamma^2-1}\left(\boldsymbol{\mathfrak{u}_4\,\mathbfcal{L}_{0,0;0,0,1,1,1} }+\boldsymbol{\mathfrak{u}_5\,q^2\,\mathbfcal{L}^{(++)}_{1,1;0,0,1,1,1} }\right) %(\gamma q^{[\eta}v_2^{\delta]}-q^{[\eta}v_1^{\delta]})\\ &
%\mathbfcal{W}_4^{\eta}  =\int \frac{\hat\delta(\ell_1\cdot v_2)\hat\delta(\ell_2\cdot v_2)\,\ell_2^\eta}{(\ell_1\cdot v_1+i\varepsilon)^2\,(\ell_2\cdot v_1+i\varepsilon)^2\,\ell_1^2\,\ell_2^2\,(\ell_1+\ell_2-q)^2}[]\\ &
%           \hspace{0.8 cm}=\boldsymbol{\mathfrak{u}_6\,\mathbfcal{L}_{0,0;0,0,1,1,1} }\,\frac{\gamma v_2^\eta-v_1^\eta}{\gamma^2-1}\\ &
 %          \mathbfcal{W}_5^{\eta}  =\int \frac{\hat\delta(\ell_1\cdot v_2)\hat\delta(\ell_2\cdot v_2)\,\ell_2^\eta}{(\ell_1\cdot v_1+i\varepsilon)^2\,(\ell_2\cdot v_1+i\varepsilon)^2\,\ell_1^2\,\ell_2^2\,(\ell_1+\ell_2-q)^2}[]\\ &
 %          \hspace{0.8 cm}\sim \left(\boldsymbol{\mathfrak{u}_7\,\mathbfcal{L}_{0,0;0,0,1,1,1} }+q^2\,\boldsymbol{\mathfrak{u}_8\,\mathbfcal{L}^{(++)}_{1,1;0,0,1,1,1}}\right)\,q^\eta\\ &
 %          \mathbfcal{W}_6^{\eta\delta}  =\int \frac{\hat\delta(\ell_1\cdot v_2)\hat\delta(\ell_2\cdot v_2)\,\ell_2^\eta\,\ell_1^\delta}{(\ell_1\cdot v_1+i\varepsilon)^2\,(\ell_2\cdot v_1+i\varepsilon)^2\,\ell_1^2\,\ell_2^2\,(\ell_1+\ell_2-q)^2}[]\\ &
%           \hspace{0.8 cm}=\frac{2}{\gamma^2-1}\left(\boldsymbol{\mathfrak{u}_9\,\mathbfcal{L}_{0,0;0,0,1,1,1} }+\boldsymbol{\mathfrak{u}_{10}\,q^2\,\mathbfcal{L}^{(++)}_{1,1;0,0,1,1,1} }\right) (\gamma q^{[\eta}v_2^{\delta]}-q^{[\eta}v_1^{\delta]})\\ &
 %           \mathbfcal{W}_7^{\eta}  =\int \frac{\hat\delta(\ell_1\cdot v_2)\hat\delta(\ell_2\cdot v_2)\,\ell_1^\eta}{(\ell_1\cdot %v_1+i\varepsilon)^2\,(\ell_2\cdot v_1+i\varepsilon)^2\,\ell_1^2\,\ell_2^2\,(\ell_1+\ell_2-q)^2}[]\\ &
 %           \hspace{0.8 cm}=\boldsymbol{\mathfrak{u}_{11}\,\mathbfcal{L}_{0,0;0,0,1,1,1}}\,\frac{\gamma v_2^\eta-v_1^\eta}{\gamma^2-1}\\ &
  %           \mathbfcal{W}_8^{\eta}  =\int \frac{\hat\delta(\ell_1\cdot v_2)\hat\delta(\ell_2\cdot v_2)\,\ell_1^\eta}{(\ell_1\cdot v_1+i\varepsilon)^2\,(\ell_2\cdot v_1+i\varepsilon)^2\,\ell_1^2\,\ell_2^2\,(\ell_1+\ell_2-q)^2}[]\\ &
 %            \hspace{0.8 cm}=\left(\boldsymbol{\mathfrak{u}_{12}}\,\mathbfcal{L}_{0,0;0,0,1,1,1}+q%^2\,\boldsymbol{\mathfrak{u}_{13}}\,\mathbfcal{L}^{(++)}_{1,1;0,0,1,1,1}\right)\,q^\eta\\ &
  %            \mathbfcal{W}_9^{\eta}  =\int \frac{\hat\delta(\ell_1\cdot %v_2)\hat\delta(\ell_2\cdot v_2)\,\ell_1^\eta\,\ell_2^\delta}{(\ell_1\cdot v_1+i\varepsilon)^2\,(\ell_2\cdot v_1+i\varepsilon)^2\,\ell_1^2\,\ell_2^2\,(\ell_1+\ell_2-q)^2}[]\\ &
 %          \hspace{0.8 cm}   =\frac{2}{\gamma^2-%1}\left(\boldsymbol{\mathfrak{u}_{14}\,\mathbfcal{L}_{0,0;0,0,1,1,1} }+\boldsymbol{\mathfrak{u}_{15}\,q^2\,\mathbfcal{L}^{(++)}_{1,1;0,0,1,1,1} }\right) (\gamma q^{[\eta}v_2^{\delta]}-q^{[\eta}v_1^{\delta]})
%           \end{split}
%\end{align}
%\begin{align}
%    \begin{split}
 %     &    \mathbfcal{W}_{10}^{\eta}= \int \frac{\hat\delta(\ell_1\cdot v_2)\hat\delta(\ell_2\cdot v_2)\,\ell_2^\eta}{(\ell_1\cdot v_1+i\varepsilon)^2\,(\ell_2\cdot v_1+i\varepsilon)^2\,\ell_1^2\,\ell_2^2\,(\ell_1+\ell_2-q)^2}[]\\ &
%\hspace{0.8 cm}=  \boldsymbol{\mathfrak{u}_{16}}\,\frac{\gamma v_2^\eta-v_1^\eta}{\gamma^2-1}\,\mathbfcal{L}_{0,0;0,0,1,1,1}\\ &
%\mathbfcal{W}_{11}^{\eta}= \int \frac{\hat\delta(\ell_1\cdot v_2)\hat\delta(\ell_2\cdot v_2)\,\ell_2^\eta}{(\ell_1\cdot v_1+i\varepsilon)^2\,(\ell_2\cdot v_1+i\varepsilon)^2\,\ell_1^2\,\ell_2^2\,(\ell_1+\ell_2-q)^2}[]\\ &
%\hspace{0.8 cm} = \left(\boldsymbol{\mathfrak{u}_{17}}\,\mathbfcal{L}_{0,0;0,0,1,1,1}+q^2\,\boldsymbol{\mathfrak{u}_{18}}\,\mathbfcal{L}^{(++)}_{1,1;0,0,1,1,1}\right)\,q^\eta\\ &
%\mathbfcal{W}_{12}^{\eta\delta}= \int \frac{\hat\delta(\ell_1\cdot v_2)\hat\delta(\ell_2\cdot v_2)\,\ell_2^\eta\,\ell_1^\delta}{(\ell_1\cdot v_1+i\varepsilon)^2\,(\ell_2\cdot v_1+i\varepsilon)^2\,\ell_1^2\,\ell_2^2\,(\ell_1+\ell_2-q)^2}[]\\ &
%\hspace{0.8 cm}=\frac{2}{\gamma^2-1}\left(\boldsymbol{\mathfrak{u}_{19}\,\mathbfcal{L}_{0,0;0,0,1,1,1} }+\boldsymbol{\mathfrak{u}_{20}\,q^2\,\mathbfcal{L}^{(++)}_{1,1;0,0,1,1,1} }\right) (\gamma q^{[\eta}v_2^{\delta]}-q^{[\eta}v_1^{\delta]})
%    \end{split}
%\end{align}
%\textbullet $\,\,$ Now we go on to compute the zig-zag diagram coming purely from the scalar part. 
%We will face the following kind of integral family.
%\begin{align}
% \begin{split}
        %\mathcal{M}^{(\pm,\pm)}_{n_1,n_2\cdots n_7}(q)=\int_{\ell_{1},\ell_{2}}\frac{1}{D_1^{n_1}D_2^{n_2}\cdots D_7^{n_7}\slashed D_8^{(\pm)n_8}\slashed D_9^{(\pm)n_9}}
   % \end{split}
%\end{align}
%where,
%$D_1=\ell_{1}\cdot v_1,\,D_2=\ell_{2}\cdot v_2,\,D_3=\ell_{1}^2,\,D_4=\ell_{2}^2,\,D_5=(q-\ell_{1}-\ell_{2})^2\,D_6=(q-\ell_{1})^2,\,D_7=(q-\ell_{2})^2,\,\slashed D_8^{\pm}=\ell_{1}\cdot v_2\pm i\sigma_1\,\slashed D_9=\ell_{2}\cdot v_1\pm i\sigma_2$.
%One can notice that, due to the asymmetry in the cut propagator the kinematic variable $\gamma$ and the regularisation parameter $\epsilon$ do not factorize. So one need to compute the integrals perturbatively in $\epsilon$. The method of differential equation provides useful way to compute the integrlas and extract out the kinematic dependence of the integrals. This
%method provides us an interesting way to bootstrap relativistic data using boundary
%information computed in the static limit of $\gamma\to 1$. For our purpose we readily need four master integrtals and the set is denoted by a vector as follows,
%\begin{align}
  %  \vec f(\gamma,\epsilon)=\{\mathbfcal{M}_{0,0,0,0,1,1,1,\slashed 1,\slashed 1},\mathbfcal{M}_{0,0,0,0,2,1,1,\slashed 1,\slashed 1},\mathbfcal{M}_{0,0,0,0,1,2,1,\slashed 1,\slashed 1},\mathbfcal{M}_{1,1,0,0,1,1,1,\slashed 1,\slashed 1}\}^{T}
%\end{align}
%By IBP identities the vector $\vec f$ satisfies an ordinary differential equation. It is always useful to rationalize the kinematic variable $\gamma$ as, $\gamma\to \frac{x^2+1}{2x}$.
%\begin{align}
%    \begin{split}
  %   &   \partial_x \vec f(x,\epsilon)=A(x,\epsilon)
  %   \vec f(x,\epsilon)
  %  \end{split}
%\end{align}
%It is quite hard to solve the differential equation in its current structure. But using suitable basis trasformation as,
%\begin{align}
  %  \begin{split}
   %     \vec %f(x,\epsilon)=\mathbb{T}(x,\epsilon)\vec{g}(x,\epsilon),
   % \end{split}
%\end{align}
%the differential system can be brought to the $\epsilon$-form which is more easier to solve. Now the system of differential equation has the form,
%\begin{align}
  %  \begin{split}
   %     \partial_x \vec g(x,\epsilon)=\mathbb{T}^{-1}(A \mathbb{T}-\partial_x \mathbb{T})\vec g(x,\epsilon)\equiv \epsilon\, \mathbb{S}(x) \vec g(x,\epsilon)\label{2.14a}
   % \end{split}
%\end{align}
%Our goal is to find the basis transformation matrix $\mathbb{T}^{-1}(x,\epsilon)$ such that $\mathbb{T}(A \mathbb{T}-\partial_x \mathbb{T})=\epsilon \,\mathbb{S}(x)$. Note that $\epsilon$ completely factorized from the transformed matrix. The matrix equation has general solution (Bootstrap equation),
%\begin{align}
  %  \begin{split}
   %     \vec g(x,\epsilon)=\underbrace{\Bigg(\mathcal{P}e^{\epsilon\int_{x_0}^{x}S(x')dx'}\Bigg)}_{\mathbb{B}(x,x_0)}\vec g(x_0,\epsilon).\labe%l{2.15a}
  %  \end{split}
%\end{align}
%and can be solved iterated way in orders of $\epsilon$. We find $\mathbb{T}(x,\epsilon)$ and $\mathbb{S}(x)$ using the publicly aviliable package \textbf{Fuchsia} \cite{Gituliar:2017vzm} and are given by,
%\begin{align}
%\begin{split}
 %  &\mathbb{T}(x,\epsilon)= \left(
%\begin{array}{cccc}
% -\frac{7 x}{x^2-1} & \frac{x}{2-2 x^2} & \frac{3 x}{x^2-1} & 0 \\
% -\frac{6 x \epsilon }{x^2-1} & -\frac{5 x \epsilon }{x^2-1} & \frac{6 x \epsilon }{x^2-1} & 0 \\
% \frac{12 \left(x^2-1\right) \epsilon ^2}{2 x \epsilon +x} & \frac{2 \left(x^2-1\right) \epsilon ^2}{2 x \epsilon +x} & -\frac{12 x \epsilon ^2}{2 \epsilon +1} & 0 \\
% 0 & 0 & 0 & \frac{2 x^2}{\left(x^2-1\right)^2} \\
%\end{array}
%\right),\textrm{and},\\ &
%\mathbb{S}(x)=
%\left(
%\begin{array}{cccc}
 %\frac{9 \left(2 x^2+1\right)}{2 \left(x^3-x\right)} & \frac{2 x^2+3}{4 \left(x^3-x\right)} & -\frac{6 x}{x^2-1} & 0 \\
 %-\frac{3 \left(2 x^2+5\right)}{x^3-x} & -\frac{6 x^2+5}{2 \left(x^3-x\right)} & \frac{12 x}{x^2-1} & 0 \\
 %\frac{6}{x} & -\frac{1}{3 x} & -\frac{2}{x} & 0 \\
% \frac{12}{x} & \frac{10}{x} & -\frac{12}{x} & 0 \\
%\end{array}
%\right):=\sum_{i=1}^3\frac{\mathbb{S}_i}{x-x_i},\,(x_i=1,-1,0)
%\end{split}
%\end{align}
%\begin{align}
%    \begin{split}
 %       \mathcal{M}\sim \prod_{i,j} \Bigg(\int_{l_{i,j}}\frac{\hat\delta(l_{i}\cdot v_1)\hat\delta(l_{j}\cdot v_2)}{
 %    (\pm l_j\cdot v_1+i\sigma)^{n_j}(\pm l_i\cdot v_2+i\sigma)^{\bar n_i} }\Bigg)\frac{1}{D_{1}^{m_1}\cdots D_{k}^{m_k}}\label{2.22a}
%    \end{split}
%\end{align}
%\par
%\textbf{\textit{Radiation region:}}
%Here we directly focus on the two loop integrals. The radiation regions are characterized by the propagators that can go on-shell. For 3PM/2-loop case we do have the following integral,
%\begin{align}
   % \begin{split}
     %   \mathcal{M}\sim \int_{\ell_{1},\ell_{2}}\frac{\hat\delta(\ell_{1}\cdot v_2)\hat\delta(\ell_{2}\cdot v_1)}{(\ell_{1}\cdot v_1+i\sigma)(\ell_{2}\cdot v_2+i\sigma)\ell_{1}^2 \ell_{2}^2 (\ell_{1}+\ell_{2}-q)^2(q-\ell_{1})^2(q-\ell_{2})^2}
  %  \end{split}
%\end{align}
%Due to the presence of the delta functions $(\ell_{1}+\ell_{2}-q)^2$ propgator can go on shell (in principle the 3 propagators except $\ell_{1}^2,\ell_{2}^2$ can go on shell). Now re-labelling $k=\ell_{1}+\ell_{2}-q$ and $\ell_{1}=l$ we left with,
%\begin{align}
  %  \begin{split}
    %    \mathcal{M}\sim %\int_{k,l}\frac{\hat\delta(l\cdot v_2)\hat\delta[(k-l+q)\cdot v_1]}{(l\cdot v_1+i\sigma)[(k-l+q)\cdot v_2+i\sigma]l^2 (k-l+q)^2 k^2(q-l)^2(\ell_{1}-k)^2}
   % \end{split}
%\end{align}
%Now going to the rest frame of the black hole 2 , using the parametrization mentioned in \eqref{2.23a}  and also applying the scalaing for the radiation region $k\sim v_{\infty}(1,1)|q|,\,l \sim (v_{\infty},1)|q|$, we will get,
%%\begin{align}
   % \begin{split}
        %\mathcal{M}^{(+)}_{n_1,n_2;m_1,\cdots m_5}\Big|_{v_{\infty}}^{\textrm{rad.}}\sim\int_{\boldsymbol l}\frac{v_{\infty}^{d-n_1-n_2-2m_3}}{(-\boldsymbol{l}\cdot \boldsymbol{n}+i\sigma)^{n_1+n_2}(-\boldsymbol{l}^2)^{m_1+m_5}[-(\boldsymbol{l}-\boldsymbol{q})^2]^{m_2+m_4}}\int_{\boldsymbol k}\frac{1}{[(-\boldsymbol{l}\cdot \boldsymbol{n}+i\sigma)^2-\boldsymbol{k}^2]^{m_3}}+\mathcal{O}()\label{2.28}
  %  \end{split}
%\end{align}
%The $\boldsymbol{k}$ integral can be evaluated straightforwardly using the retarded propagator and for $m_3=1$. The higher $m_3$ integrals subsequently computed using IBP reduction.
%\begin{align}
  %  \begin{split}
    %    \int_{\boldsymbol k}\frac{1}{[(-\boldsymbol{l}\cdot\boldsymbol{n}+i\sigma)^2-\boldsymbol{k}^2]}=\frac{\Gamma[1-\frac{d}{2}]}{(-\boldsymbol{l}\cdot\boldsymbol{n}+i\sigma)^{2-d}}e^{-i\pi d/2}
   % \end{split}
%\end{align}
%Now we are in a position to compute the boundary integrals. First we start with computing the potential region integrals.\\ \\
%\textbf{\textit{Computations of the boundary integrals:}}\\
%\\
%From the radiation region we have the following contribution to the boundary master integral,
%\begin{align}
   % \begin{split}
        %\mathcal{M}_{0,0;0,0,1,1,1}\Big|_{v_\infty\to 0}^{\textrm{rad.}}&=\int_{\boldsymbol{l}}\frac{v_{\infty}^{1-2\epsilon}}{(\boldsymbol{l}^2)(\boldsymbol{l-q})^2}\int_{\boldsymbol k}\frac{1}{[(-\boldsymbol{l}\cdot\boldsymbol{n}+i\sigma)^2-\boldsymbol{k}^2]}\\ &
        %=v_{\infty}^{1-2\epsilon}\int_{\boldsymbol{l}}\frac{1}{(-\boldsymbol{l\cdot n}+i\sigma)^{-1+2\epsilon}\boldsymbol{l^2}\boldsymbol{(l-q)^2}}
  %  \end{split}
%\end{align}
%Therefore the solution of the canonical master integral $\vec g(x,\epsilon)$ in different orders of $\epsilon$ is given by,
%\begin{align}
   % \begin{split}
    %    g^{(1)}(\gamma,\epsilon)&=-\frac{1}{1024 \pi ^2 \epsilon ^2}-\frac{6 \textrm{arccosh} (\gamma )+1}{1024 \pi ^2 \epsilon }+\frac{1}{6144 \pi ^2}\Bigg[24\, \textrm{arctanh} \left(\gamma -\sqrt{\gamma ^2-1}\right)\\ &-168 \log \left(-\sqrt{\gamma ^2-1}+\gamma +1\right) \cosh ^{-1}(\gamma )+168 \log \left(\sqrt{\gamma ^2-1}-\gamma +1\right) \textrm{arccosh}(\gamma )\\ &+36 \left(1-\textrm{arccosh}(\gamma )\right) \textrm{arccosh}(\gamma )+84 \,\textbf{Li}_2\left(\left(\gamma -\sqrt{\gamma ^2-1}\right)^2\right)+7 \pi ^2+72\Bigg]+\mathcal{O}(\epsilon)\\ &
    %   \hspace{0cm} %g^{(2)}(\gamma,\epsilon)=-\frac{1}{512 \pi ^2 \epsilon ^2}+\frac{7-10 \log (x)}{512 \pi ^2 \epsilon }+\frac{1}{3072\pi^2}\Big[60 \log \left(\gamma -\sqrt{\gamma ^2-1}\right)\\ &+12 \log \left(\gamma -\sqrt{\gamma ^2-1}\right) \left(-2 \log \left(-\sqrt{\gamma ^2-1}+\gamma +1\right)+2 \log \left(\sqrt{\gamma ^2-1}-\gamma +1\right)+5 \log \left(\gamma -\sqrt{\gamma ^2-1}\right)\right)\\ &-168 \tanh ^{-1}\left(\gamma -\sqrt{\gamma ^2-1}\right)-12 \textbf{Li}_2\left(\left(\gamma -\sqrt{\gamma ^2-1}\right)^2\right)+7 \pi ^2-120\Big]
   % \end{split}
%\end{align}
%\begin{align}
    %\begin{split}
      %  \vec f(x,\epsilon)= \left(\begin{array}{cc}
          %   &  -\frac{\textrm{arccosh}(\gamma)}{64 \pi ^2 \sqrt{\gamma ^2-1} \epsilon }+\mathcal{O}(\epsilon^0)\\
           %  & 
           % \frac{\textrm{arccosh}(\gamma)}{32 \pi ^2 \sqrt{\gamma ^2-1}}+\mathcal{O}(\epsilon^1)\\ &
         %   -\frac{2 \sqrt{\gamma ^2-1}\, \epsilon  \,\textrm{arccosh}(\gamma)-2 \gamma  \epsilon +\gamma }{32 \pi ^2}+\mathcal{O}(\epsilon^2)\\ &
        %  -\frac{1}{16 \pi^2 (\gamma^2-1)} + \frac{ \left(64\, \zeta (3) \,\epsilon +\pi ^2\right)}{96 \pi ^2 \left(\gamma ^2-1\right) }+\mathcal{O}(\epsilon^2)
      %  \end{array} 
      %   \right)
   % \end{split}
%\end{align}
%\\
\textbullet $\,\,$ Next we compute the following two diagrams. 
\begin{center}
\includegraphics[width=0.38\linewidth]{zig-zag.png}
\end{center}
The amplitude for the first diagram is given by,
\begin{align}
    \begin{split}
        i\chi&=  i\,\frac{s_1^3 s_2^3 m_1^2 m_2^2} {128 m_p^6}\int_q e^{-iq\cdot b}\hat\delta(q\cdot v_1)\hat\delta(q\cdot v_2)\int_{\ell_{1},\ell_{2}}\left(\sum_{(\pm)}\frac{1}{(\pm\omega_1+i\varepsilon)^2(\pm\omega_2+i\varepsilon)^2}\right)\frac{\hat\delta(\ell_{1}\cdot v_2)\hat\delta(\ell_{2}\cdot v_1)}{\ell_{1}^2 \ell_{2}^2(q-\ell_{1}-\ell_{2})^2}\\&\hspace{0.8cm}\times\Big((q-\ell_{1}-\ell_{2})\cdot \ell_{1}\Big)\Big( (q-\ell_{1}-\ell_{2})\cdot \ell_{2}\Big)\Bigg|_{\omega_1\to \ell_1\cdot v_1,\omega_2\to\ell_2\cdot v_2}\,.\\ &
    \end{split}
\end{align}
After IBP reduction, we are left with the following integral,
\begin{align}
    \begin{split}
       i\, \chi\Big|_{\mathcal{S}^0} &=
        %i\,\frac{s_1^3 s_2^3 m_1^2 m_2^2} {512 m_p^6} \int_q e^{-iq\cdot b}\hat\delta(q\cdot v_1)\hat\delta(q\cdot v_2)\sum' \Tilde{a}(n_1,..,n_7)\mathbfcal{M}_{n_1,..,n_7,1,1}^{\textrm{m.i.}}(q)\\ &
         i\,\frac{s_1^3 s_2^3 m_1^2 m_2^2} {128 m_p^6}\int_q e^{-iq\cdot b}\hat\delta(q\cdot v_1)\hat\delta(q\cdot v_2)\Big[\boldsymbol{\tilde h}_{29}\mathbfcal{M}_{0,0;0,0,1,1,1}+\boldsymbol{\tilde h}_{30}\,q^2\,\mathbfcal{M}_{0,0;0,0,2,1,1}  +\boldsymbol{\tilde h}_{31}\,q^2\,\mathbfcal{M}_{0,0;0,0,1,2,1}\\ &\hspace{0.8cm}-\frac{1}{2}\boldsymbol{\tilde h}_{32} q^2\,\mathbfcal{M}^{(++)}_{1,1;0,0,1,1,1} \Big]\,.
    \end{split}
\end{align}
%where,
%\begin{align}
    %\begin{split}
    %    & \tilde a[111]=\frac{4}{ \left(\gamma ^2-1\right) }  \epsilon  (2 \epsilon +1),\,
     %   \tilde a[111]=-\frac{2 (6 \epsilon +5)}{3 \left(\gamma ^2-1\right)}q^2\\ &
    %    \tilde a[121]=\frac{2 \left(12 \epsilon ^2+8 \epsilon +1\right)}{3 \left(\gamma ^2-1\right)^2 \epsilon }q^2,\,\tilde a[11111]=-\frac{2 \gamma  (2 \epsilon +3)}{3 \left(\gamma ^2-1\right)} q^2
  %  \end{split}
%\end{align}
%Then performing the remaining integral and keeping only $\mathcal{O}(\epsilon^0)$ term we get,\\
%\hfsetfillcolor{gray!10}
%\hfsetbordercolor{black!150}
%\begin{align}
%    \begin{split}
%     \tikzmarkin[disable rounded corners=false]{e5}(4,-1)(-0.1,0.7)
  %      i\,\chi\Big|_{\mathcal{S}^0} = & i\,\frac{s_1^3 s_2^3 m_1^2 m_2^2} {128 %m_p^6}\int_q e^{-iq\cdot b}\hat\delta(q\cdot v_1)\hat\delta(q\cdot v_2)\Big[\frac{4}{ \left(\gamma ^2-1\right) }  \epsilon  (2 \epsilon +1)\mathbfcal{M}_{0,0;0,0,1,1,1}-\frac{2 (6 \epsilon +5)}{3 \left(\gamma ^2-1\right)}q^2\,\mathbfcal{M}_{0,0;0,0,2,1,1}\\ &
  %   \hspace{2 cm}   +\frac{2 \left(12 \epsilon ^2
  %      +8 \epsilon +1\right)}{3 \left(\gamma ^2-1\right)^2 \epsilon }q^2\,\mathbfcal{M}_{0,0;0,0,1,2,1}\hspace{0 cm}-\frac{2 \gamma  (2 \epsilon +3)}{3 \left(\gamma ^2-1\right)} q^2\,\mathbfcal{M}_{1,1;0,0,1,1,1} \Big],\\&
 %       =i\,\frac{s_1^3 s_2^3 m_1^2 m_2^2} {128 m_p^6}\Big(\frac{2 x}{1-%x^2}\Big)\Bigg(\frac{x^3 \left(1+x^2\right) \left(6 \log \left(|\boldsymbol{b}|\right)+5 \gamma_{E} +1+\log \left(\frac{\pi }{16}\right)\right)}{\pi ^3 |\boldsymbol{b}|^2 \left(1-x^2\right)^4}\Bigg)\,.
 %       \tikzmarkend{e5}
%    \end{split}
%\end{align}
The contribution to the phase from the second diagram has the following form,
\begin{align}
\begin{split}
i \chi\Big|_{\mathcal{S}^0}
=&\,i\frac{s_1^3 s_2^3 m_1^2 m_2^2} {128 m_p^6} \int_{k_i}e^{-i(k_1+k_2+k_3)\cdot b}\hat\delta(k_1\cdot v_2)\hat\delta(k_2\cdot v_1)\hat\delta(k_1\cdot v_1+k_3\cdot v_1)\hat\delta(k_2\cdot v_2+k_3\cdot v_2)\\
&\times\frac{\mathcal{N}}{(k_1\cdot v_1+i\varepsilon)^2(k_3\cdot v_2+i\varepsilon)^2\,k_1^2 \,k_2^2\,k_3^2}\,.
\end{split}
\end{align}
Now, doing the following variable transformation, $k_1\to \ell_1-q,\,k_2\to (\ell_2-q)\,,k_3\to (q-\ell_1-\ell_2)$, we will get,
\begin{align}
    \begin{split}
   \hspace{0cm}   i\,  \chi \Big|_{\mathcal{S}^0}=i\frac{s_1^3 s_2^3 m_1^2 m_2^2} {128 m_p^6}  \int e^{i q\cdot b}\hat\delta(q\cdot v_1)\hat\delta(q\cdot v_2)\int_{\ell_1,\ell_2}\sum_{\mp}\frac{\hat\delta(\ell_1\cdot v_2)\hat\delta(\ell_2\cdot v_1)\mathcal{N}}{(\pm\ell_1\cdot v_1+i\varepsilon)^2(\mp\ell_2\cdot v_2+i\varepsilon)^2\,(\ell_1-q)^2(\ell_2-q)^2(\ell_1+\ell_2-q)^2},
    \end{split}
\end{align}
where, the numerator $\mathcal{N}$ takes the form (ignoring the tadpoles),
\begin{align}
    \begin{split}
        \mathcal{N}=\frac{1}{8}\ell_1^2\,\ell_2^2-\frac{1}{16}(\ell_1^4+\ell_2^4)+\frac{q^4}{16}\,.
    \end{split}
\end{align}
Therefore, the contribution to the eikonal phase from the second diagram is given by,
\begin{align}
    \begin{split}
       i\, \chi\Big|_{\mathcal{S}^0}&=i\frac{s_1^3 s_2^3 m_1^2 m_2^2} {128 m_p^6}  \int e^{iq\cdot b}\hat\delta(q\cdot v_1)\hat\delta(q\cdot v_2)\Bigg[\boldsymbol{\tilde h}_{33}\,\mathbfcal{M}_{0,0;0,0,1,1,1}+\boldsymbol{\tilde h}_{34}\,q^2 \mathbfcal{M}_{0,0;0,0,2,1,1}+\boldsymbol{\tilde h}_{35}\,q^2 \mathbfcal{M}_{0,0;0,0,1,2,1}\\ &\hspace{0.8cm}+\frac{1}{4}\boldsymbol{\tilde h}_{36}\,q^2 \mathbfcal{M}^{(+-)}_{1,1;0,0,1,1,1}\Bigg]\,.
    \end{split}
\end{align}
We finally get the total contribution to the eikonal phase from these two diagrams at $\mathcal{O}(\mathcal{S}^0)$, 
\begin{align}
\begin{split} \label{eqt6}
 \chi\Big|_{\mathcal{S}^0} =\frac{m_1^2m_2^2}{m_p^6}\Bigg(\Delta_{\mathfrak{08}}^{\textrm{poly}}+\Delta_{\mathfrak{01}}^{\textrm{log}}\log (x)\Bigg)
\end{split}
\end{align}
where,
\begin{align}
\begin{split}
&\Delta_{\mathfrak{08}}^{\textrm{poly}}= -\frac{ s_1^3 s_2^3 x^4 \left(x^2+1\right) (258 \log (|\boldsymbol{b}|)+215 \gamma_E -12-172 \log (2)+43 \log (\pi ))}{6144 \pi ^3 |\boldsymbol{b}|^2 \left(x^2-1\right)^5 }\,,\\&
\Delta_{\mathfrak{01}}^{\textrm{log}}=-\frac{ s_1^3 s_2^3 x^4}{512 \pi ^3 |\boldsymbol{b}|^2 \left(x^2-1\right)^4 }\,.
\end{split}
\end{align}
We have used the results for the master integrals mentioned in \eqref{MasterM}. 
\par
%\textcolor{red}{Final $q$ integral.}\\
\textbullet $\,\,$ Now, we move on to calculating  Feynman topology with two scalar lines and one graviton line, which appear at 3PM. In these diagrams, the spin of the black holes appears because of the presence of graviton. %\textcolor{red}{AB: Need to put this digram before!}% \textcolor{black}{AB: $\log[x]$ behaviour for spinless part due to the $\boldsymbol{\tilde h}_{43}$ coefficient.}
\begin{center}
\includegraphics[width=0.38\linewidth]{spin_scalar_graviton_1.png}
\end{center}
The corresponding phase is given by,
\begin{align}
    \begin{split}
        i\chi&=i\,\frac{s_1^2 s_2^2 m_1^2 m_2^2} {128 m_p^6} \int_{k_i,\omega_i}e^{-i(k_1+k_2+k_3)\cdot b}\hat\delta(k_1\cdot v_2)\hat\delta(k_3\cdot v_1)\hat\delta(\omega_1-k_1\cdot v_1)\hat\delta(\omega_1+k_2\cdot v_1)\\ &
\hspace{2 cm}\times \hat\delta(\omega_2+k_2\cdot v_2)\hat\delta(\omega_2-k_3\cdot v_2)\,\frac{\mathcal{N}}{k_1^2 \,k_2^2\, k_3^2}\left(\sum_{(\pm)}\frac{1}{(\pm\omega_1+i\varepsilon)^2(\pm\omega_2+i\varepsilon)^2}\right)
    \end{split}
\end{align}
%where the numerator takes the form,
%\begin{align}
   % \begin{split}
        %\mathcal{N}=&\Big(-4\omega_1^2-2\omega_1 k_2\cdot v_1+2\omega_1 k_1\cdot v_1+k_1\cdot k_2\Big)\Big(-4\omega_2^2v_2^\mu v_2^\nu+2\omega_2 v_2^\mu v_2^\nu k_3\cdot v_2\\ &+2i \omega_2 k_3\cdot v_2(k_3\cdot \mathcal{S}_2)^{(\mu}v_2^{\nu)}-2i\omega_2^2(k_3\cdot \mathcal{S}_2)^{(\mu}v_2^{\nu)}+\omega_2(k_3\cdot v_2)(k_3\cdot \mathcal{S}_2)^{\mu}(\mathcal{S}_2\cdot k_3)^{\nu}\\ &
  %    -2\omega_2 v_2^{(\mu}k_3^{\nu)}+v_2^\mu v_2^\nu k_3^2+ik_3^2(k_3\cdot \mathcal{S}_2)^{(\mu}v_2^{\nu)} -i\omega_2(k_3\cdot \mathcal{S}_2)^{(\mu}k_3^{\nu)}+\frac{1}{2}k_3^2 (k_3\cdot \mathcal{S}_2)^{\mu}(\mathcal{S}_2\cdot k_3)^{\nu} \Big)\\ &
   %   \times \frac{1}{2}(\eta_{\mu\alpha}\eta_{\nu\beta}+\eta_{\mu\beta}\eta_{\nu\alpha}-\eta_{\mu\nu}\eta_{\alpha\beta})\Big(v_1^\alpha v_1^\beta-i(k_3\cdot \mathcal{S}_1)^{(\alpha}v_1^{\beta)}+\frac{1}{2}(k_3\cdot \mathcal{S}_1)^{\alpha}(k_3\cdot \mathcal{S}_1)^{\beta}\Big)\,.
  %  \end{split}
%\end{align}
Now redefining $k_1\to \ell_{1}, k_3\to \ell_{2}, q=\ell_{1}+k_2+\ell_{2}$, we can write the numerator up to linear order in spin in the following way ($\mathcal{N}= \mathcal{N}|_{\mathcal{S}^0}+ \mathcal{N}|_{\mathcal{S}^1}$).
\begin{align}
    \begin{split}
    &   \mathcal{N}\Big|_{\mathcal{S}^0}= -\frac{\gamma^2(d-2)-2}{d-2}(\ell_2-q)^2(\ell_2\cdot v_2)^2\,,\\ &
     \mathcal{N}\Big|_{\mathcal{S}^1}=i\gamma \Big(\ell_2\cdot \mathcal{S}_1\cdot v_2\Big)(\ell_2\cdot v_2)^2(\ell_2-q)^2\,.
     %\\ &
    % \mathcal{O}(\mathcal{S}^2):(\ell_{2}\cdot v_2)^2\Bigg[(\ell_{2}\cdot \mathcal{S}_1)^2-(\ell_{2}\cdot \mathcal{S}_1\cdot v_2)^2+\frac{1}{2}(\ell_{2}\cdot \mathcal{S}_2)^2-(\ell_{2}\cdot \mathcal{S}_2\cdot v_1)^2\Big]+\textrm{tadpole terms}
    \end{split}
\end{align}
The spinless part of the eikonal phase is given by,
\begin{align}
    \begin{split}
        i\chi\Big|_{\mathcal{S}^0}&
        =-i\,\frac{s_1^2 s_2^2 m_1^2 m_2^2} {128 m_p^6} \frac{\gamma^2(d-2)-2}{d-2}\int e^{-iq\cdot b}\hat\delta(q\cdot v_1)\hat\delta(q\cdot v_2)\Bigg(\boldsymbol{\tilde h}_{37}\mathbfcal{M}_{0,0;0,0,1,1,1}+\boldsymbol{\tilde h}_{38}\,q^2\mathbfcal{M}_{0,0;0,0,2,1,1}\\ &\hspace{0.8cm}+\boldsymbol{\tilde h}_{39}\,q^2\mathbfcal{M}_{0,0;0,0,1,2,1}\Bigg)\,.
    \end{split}
\end{align}
%=-i\,\frac{s_1^2 s_2^2 m_1^2 m_2^2} {128 m_p^6} \frac{\gamma^2(d-2)-2}{d-2}\int e^{-iq\cdot b}\hat\delta(q\cdot v_1)\hat\delta(q\cdot v_2)\int_{\ell_{1},\ell_{2}}\frac{\hat\delta(\ell_{1}\cdot v_2)\hat\delta(\ell_{2}\cdot v_1)\,(q-\ell_{2})^2}{(\ell_{1}\cdot v_1)^2 \ell_{1}^2 \ell_{2}^2 (\ell_{1}+\ell_{2}-q)^2}\,,\\ &
%Then using \eqref{MasterM} we get, \\
%\hfsetfillcolor{gray!10}
%\hfsetbordercolor{black!150}
%\begin{align}
 %   \begin{split}
 %        \tikzmarkin[disable rounded corners=false]{e46}(0.2,-1)%(-0.1,1)i\chi\Big|_{\mathcal{S}^0}&=\,.
  %       \tikzmarkend{e46}
 %        \end{split}
%\end{align}
At $\mathcal{O}(\mathcal{S})$ we get, 
%\begin{align}
 %   \begin{split}
  %     \textcolor{black}{ i\chi}\Big|_{\mathcal{S}^1}&=i\,\frac{s_1^2 s_2^2 m_1^2 m_2^2} {128 m_p^6} \int_{q}e^{-i q\cdot b}\hat\delta(q\cdot v_1)\hat\delta(q\cdot v_2)\int_{\ell_{1},\ell_{2}}\frac{\hat\delta(\ell_{1}\cdot v_2)\hat\delta(\ell_{2}\cdot v_1)\,\mathcal{N}\Big|_{\mathcal{S}^{1}}}{(\ell_{1}\cdot v_1)^2\,(\ell_{2}\cdot v_2)^2\, \ell_{1}^2\,\ell_{2}^2\,(\ell_{1}+\ell_{2}-q)^2}\,,\\ &
   %     =i\,\frac{s_1^2 s_2^2 m_1^2 m_2^2} {256 m_p^6} \int_{q}e^{-i q\cdot b}\hat\delta(q\cdot v_1)\hat\delta(q\cdot v_2)\int_{\ell_{1},\ell_{2}}\frac{\hat\delta(\ell_{1}\cdot v_2)\hat\delta(\ell_{2}\cdot v_1)(\ell_{2}-q)^2}{(\ell_{1}\cdot v_1)^2\,(\ell_{2}\cdot v_2)^2\, \ell_{1}^2\,\ell_{2}^2\,(\ell_{1}+\ell_{2}-q)^2}\Big[2i\gamma \,\mathcal{S}_{1\mu\nu}v_2^\nu \,\ell_{2}^{\mu}(\ell_{2}\cdot v_2)^2\Big]\,.
    %    %-i\gamma \underbrace{\ell_{2}^2(\cdots)}_{\textrm{tadpole}}\Big]
    %\end{split}\label{3.44a}
%\end{align}
 %One may notice that the term at $\mathcal{O}(\mathcal{S})$ can not be written in terms of the integral basis. So, first we need to do a PV reduction before doing the IBP reduction. Hence, we first evaluate the following vector integral with the help of PV reduction. We will use this same strategy for other other integrals from from the spinning part. We will henceforth simply quote the final result.
%\begin{align}
 %   \begin{split}
  %      \mathcal{J}^{\mu}\equiv \int_{\ell_{1},\ell_{2}}\frac{\hat\delta(\ell_{1}\cdot v_2)\hat\delta(\ell_{2}\cdot v_1)\,(\ell_{2}-q)^2\,\ell_{2}^\mu}{(\ell_{1}\cdot v_1)^2\, \ell_{1}^2\,\ell_{2}^2\,(\ell_{1}+\ell_{2}-q)^2}:=a_1 \frac{q^\mu}{q^2} +a_2(v_2^\mu-\gamma v_1^\mu)
   % \end{split}
%\end{align}
%where,
%\begin{align}
%\begin{split}
 %   a_1&=\frac{1}{2}\int\frac{\hat\delta(\ell_{1}\cdot v_2)\hat\delta(\ell_{2}\cdot v_1)(\ell_{2}-q)^2}{(\ell_{1}\cdot v_1)^2\, \ell_{1}^2\,\ell_{2}^2\,(\ell_{1}+\ell_{2}-q)^2}\Big[q^2+\ell_{2}^2-(q-\ell_{2})^2
%\Big]\,,\\ &
%=\boldsymbol{\tilde h}_{40}\,\mathbfcal{M}_{0,0;0,0,1,1,1}q^2+\boldsymbol{\tilde h}_{41}\,\mathbfcal{M}_{0,0;0,0,2,1,1}q^4
%+\boldsymbol{\tilde h}_{42}\,\mathbfcal{M}_{0,0;0,0,1,2,1}q^4
%\end{split}
%\end{align}
%and 
%\begin{align}
 %   \begin{split}
  %      a_2&=\frac{1}{1-\gamma^2}\int \frac{\hat\delta(\ell_{1}\cdot v_2)\hat\delta(\ell_{2}\cdot v_1)(\ell_{2}\cdot v_2)(\ell_{2}-q)^2}{(\ell_{1}\cdot v_1)^2\,\ell_{1}^2 \ell_{2}^2 (\ell_{1}+\ell_{2}-q)^2}
   %     =\Bigg(-\frac{1}{(\gamma^2-1)}\frac{\gamma   \left(4 \epsilon ^2-1\right)}{4 \left(6 \epsilon ^2-7 \epsilon +1\right)}\Bigg)\,\mathbfcal{M}_{0, 1; 0, 0, 1, 1, 1}q^2 \,.
    %\end{split}
%\end{align}
%Finally, putting all these into \eqref{3.44a} we get,\\
%\hfsetfillcolor{gray!10}
%\hfsetbordercolor{black!150}
\begin{align}
    \begin{split}
      %  \tikzmarkin[disable rounded corners=false]{e6}(15,-1)(-0.1,1)&
      i\chi\Big|_{\mathcal{S}^1}  = i\,\frac{s_1^2 s_2^2 m_1^2 m_2^2} {128 m_p^6} \int_{q}e^{-i q\cdot b}\hat\delta(q\cdot v_1)\,\hat\delta(q\cdot v_2)
      &\Bigg[i\gamma\Bigg(\boldsymbol{\tilde h}_{40}\mathbfcal{M}_{0,0;0,0,1,1,1}+\boldsymbol{\tilde h}_{41} q^2\,\mathbfcal{M}_{0,0;0,0,2,1,1}\\&+\boldsymbol{\tilde h}_{42}q^2\,\mathbfcal{M}_{0,0;0,0,    1,2,1}\Bigg)\Bigg]\Big(q\cdot \mathcal{S}_1\cdot v_2\Big)\,,\\&
%=i\,\frac{s_1^2 s_2^2 m_1^2 m_2^2} {1024 \pi^3\,|\boldsymbol{b}|^3\,m_p^6 }\,\Bigg(\frac{x \left(1+x^2\right)^2 }{ \left(1-x^2\right)^3  }\Bigg)\,\Bigg(\hat{b}\cdot \mathcal{S}_1\cdot v_2\Bigg)\,.\\&
       % =\Bigg(\frac{-i x (x^2+1)\log(x) }{16\pi^3(1-x^2)^2|\boldsymbol{b}|^4}-\frac{i x \left(x^2+1\right)^2}{16 \pi ^3 |\boldsymbol{b}|^2 \left(1-x^2\right)^3}+\frac{i x \left(x^2+1\right) \log (x)}{16 \pi ^3 |\boldsymbol{b}|^4 \left(1-x^2\right)^2} \Bigg)\,(b\cdot \mathcal{S}_{1}\cdot v_2)  \\&
       % = -i \frac{ x (x^2+1)^2}{16 \pi^3 |\boldsymbol{b}|^4 (1-x^2)^3}\,(b\cdot \mathcal{S}_{1}\cdot v_2)\,.
     %  \tikzmarkend{e6}
    \end{split}
\end{align}
%Now, move on to the computation of the $\mathcal{S}^2$ order contribution.
%\begin{align}
 %   \begin{split}
     %   \chi\Big|_{\mathcal{S}^2}&=\int_{q}e^{-i q\cdot %b}\hat\delta(q\cdot v_1)\hat\delta(q\cdot v_2)\int_{\ell_{1},\ell_{2}}\frac{\hat\delta(\ell_{1}\cdot v_2)\hat\delta(\ell_{2}\cdot v_1)\,\mathcal{N}(q,l_i,v_i)\Big|_{\mathcal{S}^{2}}}{(\ell_{1}\cdot v_1)^2\,(\ell_{2}\cdot v_2)^2\, \ell_{1}^2\,\ell_{2}^2\,(\ell_{1}+\ell_{2}-q)^2}\\ &
   %     =\int_{q}e^{-i q\cdot b}\hat\delta(q\cdot v_1)\hat\delta(q\cdot v_2)\int_{\ell_{1},\ell_{2}}\frac{\hat\delta(\ell_{1}\cdot v_2)\hat\delta(\ell_{2}\cdot v_1)}{(\ell_{1}\cdot v_1)^2\,\ell_{1}^2\,\ell_{2}^2\,(\ell_{1}+\ell_{2}-q)^2}\Bigg[(\ell_{2}\cdot \mathcal{S}_1)\cdot(\ell_{2}\cdot \mathcal{S}_1)-(\ell_{2}\cdot \mathcal{S}_1\cdot v_2)^2\\ &\hspace{0.8cm}+\frac{1}{2}(\ell_{2}\cdot \mathcal{S}_2)\cdot(\ell_{2}\cdot \mathcal{S}_2)-(\ell_{2}\cdot \mathcal{S}_2\cdot v_1)^2\Big]
 %   \end{split}
%\end{align}
%It is evident that we need to evaluate the following tensor integral,
%\begin{align}
%    \begin{split}
%    \mathcal{J}^{\mu\nu}&\equiv \int_{\ell_{1},\ell_{2}}\frac{\hat\delta(\ell_{1}\cdot v_2)\hat\delta(\ell_{2}\cdot v_1)\,\ell_{2}^\mu \ell_{2}^\nu}{(\ell_{1}\cdot v_1)^2\, \ell_{1}^2\,\ell_{2}^2\,(\ell_{1}+\ell_{2}-q)^2}\\ &
%    =\frac{\left(\gamma ^2-1\right) J_2-q^2\, J_1}{q^2\,\left(\gamma ^2-1\right) (2 \epsilon -1)}\eta^{\mu\nu}+\frac{J_1\,q^2\, \left(\gamma ^2 (2 \epsilon -1)-1\right)+\left(\gamma ^2-1\right) J_2}{q^2\left(\gamma ^2-1\right)^2 (2 \epsilon -1)}v_1^\mu v_1^\nu+\frac{\left(\gamma ^2-1\right) J_2+2\,q^2 J_1 (\epsilon -1)}{q^2\,\left(\gamma ^2-1\right)^2 (2 \epsilon -1)}v_2^\mu v_2^\nu\\ &
%  \hspace{0.2 cm}-\frac{2 \gamma  \left(\left(\gamma ^2-1\right) J_2+2\,q^2 J_1 (\epsilon -1)\right)}{q^2\left(\gamma ^2-1\right)^2 (2 \epsilon -1)}v_1^{(\mu}v_2^{\nu)}+  \frac{2 \gamma  J_3}{\left(\gamma ^2-1\right) q^2}q^{(\mu}v_1^{\nu)}-\frac{2 J_3}{\left(\gamma ^2-1\right) q^2}q^{(\mu}v_2^{\nu)}\\&+\frac{2 \left(\gamma ^2-1\right) J_2 (\epsilon -1)+J_1\,q^2}{\left(\gamma ^2-1\right) q^4 (2 \epsilon -1)}q^\mu q^\nu
%    \end{split}
%\end{align}
%where,
%\begin{align}
%    \begin{split}
%        J_1:&=
      %  \int_{\ell_{1},\ell_{2}}\frac{\hat\delta(\ell_{1}\cdot v_2)\hat\delta(\ell_{2}\cdot v_1)(\ell_{2}\cdot v_2)^2}{(\ell_{1}\cdot v_1)^2\ell_{1}^2 \ell_{2}^2(\ell_{1}+\ell_{2}-q)^2}\\ &
 %       \frac{\epsilon  \left(-\left(\gamma ^2 (\epsilon -1)\right)+5 \epsilon +1\right)}{4 \epsilon ^2-6 \epsilon +2}\mathbfcal{M}_{0,0;0,0,1,1,1,\slashed 1,\slashed 1}-\frac{q^2 \left(2 \epsilon ^2+\epsilon +1\right)}{4 (\epsilon -1) \epsilon  (2 \epsilon -1)}\mathbfcal{M}_{0,0;0,0,1,2,1,\slashed 1,\slashed 1}+\frac{\left(\gamma ^2-1\right) q^2}{8 \epsilon -4}\mathbfcal{M}_{0,0;0,0,2,1,1,\slashed 1,\slashed 1}\\ &
%        \hspace{0cm}
%        J_{2}:
        %\int_{\ell_{1},\ell_{2}}\frac{\hat\delta(\ell_{1}\cdot v_2)\hat\delta(\ell_{2}\cdot v_1)\,(q\cdot \ell_{2})^2}{(\ell_{1}\cdot v_1)^2\,\ell_{1}^2\,\ell_{2}^2\,(\ell_{1}+\ell_{2}-q)^2}\\ &
       % =\frac{1}{4}\int_{\ell_{1},\ell_{2}}\frac{\hat\delta(\ell_{1}\cdot v_2)\hat\delta(\ell_{2}\cdot v_1)}{(\ell_{1}\cdot v_1)^2\,\ell_{1}^2\,\ell_{2}^2\,(\ell_{1}+\ell_{2}-q)^2}\Big[q^4+\ell_{2}^4+(q-\ell_{2})^4+2q^2 \ell_{2}^2-2q^2 (q-\ell_{2})^2-2\ell_{2}^2(q-\ell_{2})^2\Big] \\ &
%        =\frac{(3-5 \gamma ^2)\epsilon+(21 \gamma ^2-17)\epsilon^2+(40-24 \gamma ^2)\epsilon^3+(8 \gamma ^2-8)\epsilon^4}{4  (\gamma^2 -1) (\epsilon -1) (2 \epsilon -1) (3 \epsilon -1)}q^2 \mathbfcal{M}_{0,0;0,0,1,1,1,\slashed 1,\slashed 1},\\ & 
%        \hspace{0.35cm}+
%\frac{ 2 \left(6 \epsilon ^2+\epsilon -1\right)+\left(-20 \epsilon ^2+14 \epsilon -3\right)}{8 \left(6 \epsilon ^2-5 \epsilon +1\right)}q^4\mathbfcal{M}_{0,0;0,0,2,1,1,\slashed 1,\slashed 1},\\ &\hspace{0.35cm}+ \frac{16 \epsilon ^4-40 \epsilon ^3+18 \epsilon ^2-7 \epsilon +1}{8 (\gamma^2 -1)  (\epsilon -1) \epsilon  (2 \epsilon -1) (3 \epsilon -1)}q^4\mathbfcal{M}_{0,0;0,0,1,2,1,\slashed 1,\slashed 1}\\ &
%  \textrm{and,}\\ &
%\hspace{0cm}
%J_3:
%\frac{1}{2}\int_{\ell_{1},\ell_{2}}\frac{\hat\delta(\ell_{1}\cdot v_2)\hat\delta(\ell_{2}\cdot v_1)}{(\ell_{1}\cdot v_1)^2\ell_{1}^2\ell_{2}^2(\ell_{1}+\ell_{2}-q)^2}(\ell_{2}\cdot v_2)[q^2+\ell_{2}^2-(q-\ell_{2})^2]\\ &
%=-\frac{\gamma q^2  \left( \left(-6 \epsilon ^2+\epsilon +1\right)+2  \epsilon  (\epsilon +1)\right)}{4 (2 \epsilon -1) (3 \epsilon +1)}\mathbfcal{M}_{0,1;0,0,1,1,1,\slashed 1,\slashed 1}
%    \end{split}
%\end{align}
%Then we get, 
%\begin{align}
 %   \begin{split}
 %       \chi\Big|_{\mathcal{S}^2} =\int_{q}e^{-i q\cdot b}\hat\delta(q\cdot v_1)\hat\delta(q\cdot v_2)&\Big[a_1\, \mathcal{S}_1^2+a_3 (v_2\cdot \mathcal{S}_1)\cdot (v_2\cdot \mathcal{S}_1)-a_6 (q\cdot \mathcal{S}_1)\cdot (v_2\cdot \mathcal{S}_1)+a_7 (q\cdot \mathcal{S}_1)\cdot (q\cdot \mathcal{S}_1)\\&-a_1 (v_2\cdot \mathcal{S}_1)\cdot (v_2\cdot \mathcal{S}_1)-a_7 (q\cdot \mathcal{S}_1\cdot v_2)^2+\frac{a_1}{2}\mathcal{S}_2^2+\frac{a_2}{2} (v_1\cdot \mathcal{S}_2)\cdot (v_1\cdot \mathcal{S}_2)\\&+\frac{a_5}{2}(q\cdot \mathcal{S}_2)\cdot (v_1\cdot \mathcal{S}_2)+\frac{a_7}{2}(q\cdot \mathcal{S}_2)\cdot (q\cdot \mathcal{S}_2)-a_1 (v_1\cdot \mathcal{S}_2)\cdot (v_1\cdot \mathcal{S}_2)\\&-a_7 (q\cdot \mathcal{S}_2\cdot v_1)^2\Big]\,,
%    \end{split}
%\end{align}
%where, 
%\begin{align}
%    \begin{split}
%   & a_1=\frac{\left(\gamma ^2-1\right) J_2-q^2\, J_1}{q^2\,\left(\gamma ^2-1\right) (2 \epsilon -1)}, \quad a_2=\frac{J_1\,q^2\, \left(\gamma ^2 (2 \epsilon -1)-1\right)+\left(\gamma ^2-1\right) J_2}{q^2\left(\gamma ^2-1\right)^2 (2 \epsilon -1)}\,,\quad a_3=\frac{\left(\gamma ^2-1\right) J_2+2\,q^2 J_1 (\epsilon -1)}{q^2\,\left(\gamma ^2-1\right)^2 (2 \epsilon -1)}\,,\\&
%    a_5= \frac{2 \gamma  J_3}{\left(\gamma ^2-1\right) q^2},\quad a_6=\frac{2 J_3}{\left(\gamma ^2-1\right) q^2},\quad a_7=\frac{2 \left(\gamma ^2-1\right) J_2 (\epsilon -1)+J_1\,q^2}{\left(\gamma ^2-1\right) q^4 (2 \epsilon -1)}\,.
%      \end{split}
%\end{align}
%Then neglecting the tadpole terms we get, 
%\begin{align}
%    \begin{split}
   %     \chi\Big|_{\mathcal{S}^2} =\int_{q}e^{-i q\cdot %b}\hat\delta(q\cdot v_1)\hat\delta(q\cdot v_2)&\Big[-a_6 (q\cdot \mathcal{S}_1)\cdot (v_2\cdot \mathcal{S}_1)+a_7 (q\cdot \mathcal{S}_1)\cdot (q\cdot \mathcal{S}_1)-a_7 (q\cdot \mathcal{S}_1\cdot v_2)^2\\&+\frac{a_5}{2}(q\cdot \mathcal{S}_2)\cdot (v_1\cdot \mathcal{S}_2)+\frac{a_7}{2}(q\cdot \mathcal{S}_2)\cdot (q\cdot \mathcal{S}_2)-a_7 (q\cdot \mathcal{S}_2\cdot v_1)^2\Big]\,,
 %   \end{split}
%\end{align}
%\newpage
%\noindent
%\textbullet\textcolor{red}{{\bf Bhai diagram of 4.32:}}
The contribution from the second diagram to the eikonal phase is given by,
%\begin{align}
  %  \begin{split}
 %      i\, \chi\Big|_{\mathcal{S}^0}&=-i \frac{s_1^2 s_2^2 m_1^2 m_2^2}{64 m_p^6} \rmint_{k_i,\omega_i}\hat\delta(k_1\cdot v_2)\hat\delta(k_2\cdot v_1)\hat\delta(k_3\cdot v_2-\omega_2)\hat\delta(-k_1\cdot v_1+\omega_1)\hat\delta(k_2\cdot v_2+\omega_2)\hat\delta(\omega_1+k_3\cdot v_1)\\ &
%       \hspace{2 cm} \times\frac{ \mathcal{N}|_{\mathcal{S}_0}}{\prod_i k_i^2\prod_j \omega_j^2}\,e^{-i(k_1+k_2+k_3)\cdot b}
 %   \end{split}
%\end{align}
%Where
%\begin{align}
%\begin{split}\textcolor{red}{\mathcal{N}|_{\mathcal{S}_0}}&=\gamma ^2 \left(2\ell_2^2+q^2\right) \left(\ell_1\cdot v_1\right)^2+\gamma ^2 \left(2\ell_1^2+q^2\right) \left(\ell_2\cdot v_2\right)^2+4 \gamma ^2 \left(\ell_1\cdot v_1\right)^2 \left(\ell_2\cdot v_2\right)^2+\frac{1}{4} \gamma ^2 \left(2\ell_1^2+q^2\right) \left(2\ell_2^2+q^2\right)-\\&\gamma  \left(\ell_1^2+\ell_2^2-q^2\right) \left(\ell_1\cdot v_1\right) \left(\ell_2\cdot v_2\right)+\frac{\left(2\ell_1^2+q^2\right) \left(2\ell_2^2+q^2\right)}{8 (\epsilon -1)}
%\end{split}
%\end{align}
%The eikonal phase is given by,
\begin{align}
    \begin{split}
i\chi\Big|_{\mathcal{S}^0}=-\frac{i m_1^2 m_2^2 s_1^2 s_2^2}{64 m_p^6}\,\int e^{-i q \cdot b}\, \hat\delta(q\cdot v_1)\hat\delta(q\cdot v_2)&\Big(\boldsymbol{\tilde h}_{43}\,\,\mathbfcal{M}_{0,0;0,0,1,1,1}+\boldsymbol{\tilde h}_{44}\,\,q^2\mathbfcal{M}_{0,0;0,0,2,1,1}\\&+\boldsymbol{\tilde h}_{45}\,\,q^2\mathbfcal{M}_{0,0;0,0,1,2,1}+\frac{1}{4}\boldsymbol{\tilde h}_{46}\,\,q^2\mathbfcal{M}^{(+-)}_{1,1;0,0,1,1,1}\Big)\,,
   \end{split}
\end{align}
%where, $h_i$'s are defined in \eqref{E.1}.
Similarly,
\begin{align}
    \begin{split}
       i\, \chi\Big|_{\mathcal{S}^1}=-i \frac{s_1^2 s_2^2 m_1^2 m_2^2}{64 m_p^6} \int_{q}e^{i q\cdot b}\hat\delta(q\cdot v_1)\,\hat\delta(q\cdot v_2)&\Bigg[i\,\Bigg(\boldsymbol{\tilde h}_{47}\,\,\mathbfcal{M}_{0,0;0,0,1,1,1}+\boldsymbol{\tilde h}_{48}\,\,q^2\mathbfcal{M}_{0,0;0,0,2,1,1}\\&+\boldsymbol{\tilde h}_{49}\,\,q^2\mathbfcal{M}_{0,0;0,0,1,2,1}+\frac{1}{4}\boldsymbol{\tilde h}_{50}\,\,q^2\mathbfcal{M}^{(+-)}_{1,1;0,0,1,1,1}\Bigg)\big(q\cdot\mathcal{S}_1\cdot v_2\big)+\\&i\,\Bigg(\boldsymbol{\tilde h}_{51}\,\,\mathbfcal{M}_{0,0;0,0,1,1,1}+\boldsymbol{\tilde h}_{52}\,\,q^2\mathbfcal{M}_{0,0;0,0,2,1,1}\\&+\boldsymbol{\tilde h}_{53}\,\,q^2\mathbfcal{M}_{0,0;0,0,1,2,1}+\frac{1}{4}\boldsymbol{\tilde h}_{54}\,\,q^2\mathbfcal{M}^{(+-)}_{1,1;0,0,1,1,1}\Bigg)\big(q\cdot\mathcal{S}_2\cdot v_1\big)\Bigg]
    \end{split}
\end{align}
We finally get the total contribution to the eikonal phase from these two diagrams at $\mathcal{O}(\mathcal{S}^0)$, 
\begin{align}
\begin{split} \label{eqt7}
 \chi\Big|_{\mathcal{S}^0} =\frac{m_1^2m_2^2}{m_p^6}\Bigg(\Delta_{\mathfrak{09}}^{\textrm{poly}}+\Delta_{\mathfrak{02}}^{\textrm{log}}\log (x)\Bigg)
\end{split}
\end{align}
\enlargethispage{2\baselineskip}
where, 
\begin{thesiscompactdisplay}
\begin{align}
\begin{split}
&\Delta_{\mathfrak{09}}^{\textrm{poly}}=-\frac{s_1^2 s_2^2}{4096 \pi ^3 |\boldsymbol{b}|^2 \left(x^2-1\right)^5}\Bigg(\left(x^2+1\right) \big(-12 \left(11 x^4-6 x^2+11\right) x^2 \log (|\boldsymbol{b}|)+5 x^8-22 x^6 \left(5 \gamma_E +2+\log \left(\frac{\pi }{16}\right)\right)\\&\hspace{0.8cm}+6 x^4 (10 \gamma_E +9-8 \log (2)+2 \log (\pi ))-22 x^2 \left(5 \gamma_E +2+\log \left(\frac{\pi }{16}\right)\right)+5\big)\Bigg)\,,\\&
\Delta_{\mathfrak{02}}^{\textrm{log}}=-\frac{s_1^2 s_2^2 \left(x^4+6 x^2+1\right)}{512 \pi ^3 |\boldsymbol{b}|^2 \left(x^2-1\right)^2}\,.
\end{split}
\end{align}
\end{thesiscompactdisplay}
Also the total contribution at $\mathcal{O}(\mathcal{S}^1)$ is, 
\enlargethispage{2\baselineskip}
\begin{thesiscompactdisplay}[\scriptsize]
\begin{align}
\begin{split}
\chi\Big|_{\mathcal{S}^1} =\frac{m_1^2m_2^2}{m_p^6}&\Bigg[\Big(\Delta_{\mathfrak{10}}^{\textrm{poly}}+\Delta_{\mathfrak{03}}^{\textrm{log}}\log (x)+\Delta_{\mathfrak{01}}^{\textrm{Polylog}}\left(\textbf{Li}_2(1-x^2)+\log^2(x)\right)+\boldsymbol{\mathfrak{U}}_{\mathfrak{1}}^{\textrm{LI}_3}\Big\{54\, \textbf{Li}_3\left(x^2\right)-216\, \textbf{Li}_3\left(\frac{x+1}{1-x}\right)
\\ &+432\, \textbf{Li}_3(1-x)+216\, \textbf{Li}_3\left(\frac{x+1}{x-1}\right)+432\, \textbf{Li}_3(x+1)+216 \Big(-\textbf{Li}_2\left(\frac{x+1}{1-x}\right)-\textbf{Li}_2(1-x)\\ &+\textbf{Li}_2\left(\frac{x+1}{x-1}\right)+\textbf{Li}_2(x+1)\Big) \log \left(\frac{2}{x+1}-1\right)-54 \left(\pi ^2-2 \log ^2(x)\right) \log \left(1-x^2\right)-20 \log ^3(x)\\ &+108 i \pi  \log ^2(1-x)+108 i \pi  \log ^2(x+1)-11 \pi ^2 \log (x)-432 \zeta (3)\Big\}\Big)\Big(\hat{b}\cdot \mathcal{S}_1\cdot v_2\Big)+\\&\Big(\Delta_{\mathfrak{11}}^{\textrm{poly}}+\Delta_{\mathfrak{04}}^{\textrm{log}}\log (x)+\Delta_{\mathfrak{02}}^{\textrm{Polylog}}\left(\textbf{Li}_2(1-x^2)+\log^2(x)\right)+\boldsymbol{\mathfrak{U}}_{\mathfrak{2}}^{\textrm{LI}_3}\Big\{54\, \textbf{Li}_3\left(x^2\right)-216\, \textbf{Li}_3\left(\frac{x+1}{1-x}\right)
\\ &+432\, \textbf{Li}_3(1-x)+216\, \textbf{Li}_3\left(\frac{x+1}{x-1}\right)+432\, \textbf{Li}_3(x+1)+216 \Big(-\textbf{Li}_2\left(\frac{x+1}{1-x}\right)-\textbf{Li}_2(1-x)\\ &+\textbf{Li}_2\left(\frac{x+1}{x-1}\right)+\textbf{Li}_2(x+1)\Big) \log \left(\frac{2}{x+1}-1\right)-54 \left(\pi ^2-2 \log ^2(x)\right) \log \left(1-x^2\right)-20 \log ^3(x)\\ &+108 i \pi  \log ^2(1-x)+108 i \pi  \log ^2(x+1)-11 \pi ^2 \log (x)-432 \zeta (3)\Big\}\Big)\Big)\Big(\hat{b}\cdot \mathcal{S}_2\cdot v_1\Big)\Bigg]
\end{split}
\end{align}
\end{thesiscompactdisplay}
%\newpage
\enlargethispage{2\baselineskip}
where,
\begin{thesiscompactdisplay}[\scriptsize]
\begin{align}
\begin{split}
\Delta_{\mathfrak{10}}^{\textrm{poly}}=\frac{xs_1^2 s_2^2}{12288 \pi ^3 |\boldsymbol{b}|^3 \left(x^2-1\right)^5 }\Bigg[&-12 \left(12 x^6+77 x^5+72 x^4+154 x^3+72 x^2+77 x+12\right) x \log(|\boldsymbol{b}|)\\
&-13 x^8-120 \gamma_E  x^7+60 x^7+24 x^7 \log \left(\frac{16}{\pi }\right)-770 \gamma_E  x^6+398 x^6+154 x^6 \log \left(\frac{16}{\pi }\right)\\
&-720 \gamma_E  x^5+372 x^5+144 x^5 \log \left(\frac{16}{\pi }\right)-1540 \gamma_E  x^4+630 x^4+308 x^4 \log \left(\frac{16}{\pi }\right)\\
&-720 \gamma_E  x^3+372 x^3+144 x^3 \log \left(\frac{16}{\pi }\right)-770 \gamma_E  x^2+398 x^2+154 x^2 \log \left(\frac{16}{\pi }\right)\\
&-120 \gamma_E  x+60 x+24 x \log \left(\frac{16}{\pi }\right)-13\Bigg]\,,\\ &
\hspace{-4.5 cm}\Delta_{\mathfrak{3}}^{\textrm{log}}=\frac{1}{18432 \pi ^3 |\boldsymbol{b}|^3 \left(x^2-1\right)^6 \left(x^2+1\right) }\Bigg[-48(1+x^{12})+(-24 x^8-24 x^4)\beta_1+2(x^5+x^7)\beta_2+3(x^3+x^9)\beta_3\\
& +(x+x^{11})\beta_4-24x^6\beta_5+6(x^2+x^{10})\beta_6\Bigg]\,,\nonumber
\end{split}
\end{align}
\begin{align}
\begin{split}
&\Delta_{\mathfrak{01}}^{\textrm{Polylog}}=\frac{x^2s_1^2 s_2^2}{6144 \pi ^3 |\boldsymbol{b}|^3 \left(x^2-1\right)^6 \left(x^2+1\right) }\Bigg[72 \left(x^{10}+24 x^9+3 x^8-10 x^6-48 x^5-10 x^4+3 x^2+24 x+1\right) \log(|\boldsymbol{b}|)\\ &\hspace{0.8cm}+9 x^{10}-12 x^{10} \log \left(\frac{16}{\pi }\right)-1236 x^9-288 x^9 \log \left(\frac{16}{\pi }\right)+9 x^8\\ &\hspace{0.8cm}-36 x^8 \log \left(\frac{16}{\pi }\right)-392 x^7+222 x^6+120 x^6 \log \left(\frac{16}{\pi }\right)+3224 x^5\\ &\hspace{0.8cm}+576 x^5 \log \left(\frac{16}{\pi }\right)+222 x^4+120 x^4 \log \left(\frac{16}{\pi }\right)-392 x^3+9 x^2\\ &\hspace{0.8cm}-36 x^2 \log \left(\frac{16}{\pi }\right)+60 \gamma_E  (x^{10}+24 x^9+3 x^8-10 x^6-48 x^5-10 x^4\\ &\hspace{0.8cm}+3 x^2+24 x+1)-1236 x-288 x \log \left(\frac{16}{\pi }\right)+9-12 \log \left(\frac{16}{\pi }\right)\Bigg]\,,\\
&\boldsymbol{\mathfrak{U}_1}^{\textrm{LI}_3}=\frac{ \pi ^2 s_1^2 s_2^2 x}{9216 \pi ^5 |\boldsymbol{b}|^3 \left(x^2-1\right)^6}\Bigg[x^8+24 x^7+2 x^6-24 x^5-12 x^4-24 x^3+2 x^2+24 x+1\Bigg],
\\ &
\Delta_{\mathfrak{04}}^{\textrm{log}}=\frac{1}{36864 \pi ^3 |\boldsymbol{b}|^3 (x^2-1)^6 \left(x^2+1\right) }\Bigg[96(1+x^{12})+48(x^3+x^{9})\delta_1+3(x^2+x^{10})\delta_2+48(x^5+x^7)\delta_3\\ &\hspace{0.8cm}+4(x^4+x^8)\delta_4+x^6\delta_5\Bigg]\,,\\&
\end{split}
\end{align}
\begin{align}
\Delta_{\mathfrak{11}}^{\textrm{poly}}
&= \frac{x^2 s_1^2 s_2^2}{6144 \pi ^3 |\boldsymbol{b}|^3 \left(x^2-1\right)^5 }\Bigg[
12 (42 x^4-36 x^3+89 x^2-36 x+42) x \log(|\boldsymbol{b}|)-12 x^6+420 \gamma_E  x^5-323 x^5
\notag\\
&\quad -84 x^5 \log \left(\frac{16}{\pi }\right)-360 \gamma_E  x^4+108 x^4+72 x^4 \log \left(\frac{16}{\pi }\right)+890 \gamma_E  x^3-558 x^3
\notag\\
&\quad -178 x^3 \log \left(\frac{16}{\pi }\right)-360 \gamma_E  x^2+108 x^2+72 x^2 \log \left(\frac{16}{\pi }\right)+420 \gamma_E  x-323 x
\notag\\
&\quad -84 x \log \left(\frac{16}{\pi }\right)-12\Bigg]\,, \notag\\[2pt]
\Delta_{\mathfrak{02}}^{\textrm{Polylog}}
&= \frac{x^3 s_1^2 s_2^2}{768 \pi ^3 |\boldsymbol{b}|^3 (x^2-1)^6 \left(x^2+1\right) }\Bigg[
18 \left(12 x^8-7 x^6-12 x^5-38 x^4-12 x^3-7 x^2+12\right) \log (|\boldsymbol{b}|)-150 x^8
\notag\\
&\quad -36 x^8 \log \left(\frac{16}{\pi }\right)-15 x^7+30 x^6+21 x^6 \log \left(\frac{16}{\pi }\right)+39 x^5+36 x^5 \log \left(\frac{16}{\pi }\right)
\notag\\
&\quad +554 x^4+114 x^4 \log \left(\frac{16}{\pi }\right)+39 x^3+36 x^3 \log \left(\frac{16}{\pi }\right)+30 x^2+21 x^2 \log \left(\frac{16}{\pi }\right)
\notag\\
&\quad +15 \gamma_E  \left(12 x^8-7 x^6-12 x^5-38 x^4-12 x^3-7 x^2+12\right)-15 x-150
\notag\\
&\quad -36 \log \left(\frac{16}{\pi }\right)\Bigg]\,, \notag\\[2pt]
\boldsymbol{\mathfrak{U}}_{\mathfrak{2}}^{\textrm{LI}_3}
&= \frac{s_1^2 s_2^2 x^2}{4608 \pi ^3 |\boldsymbol{b}|^3 \left(x^2-1\right)^6}\Bigg[
12 x^6-19 x^4-12 x^3-19 x^2+12\Bigg]
\end{align}
\end{thesiscompactdisplay}
with,
\begin{align}
    \begin{split}
       & \beta_1=90 \log (|\boldsymbol{b}|)+75 \gamma_E -26-15 \log \left(\frac{16}{\pi }\right),\\ &
       \beta_2=-1080 \log ^2(|\boldsymbol{b}|)-1800 \gamma_E  \log (|\boldsymbol{b}|)+666 \log (|\boldsymbol{b}|)+3 \log \left(\frac{\pi }{16}\right) \left(-120 \log (|\boldsymbol{b}|)+37+10 \log \left(\frac{16}{\pi }\right)\right)\\ &\hspace{0.35cm}-750 \gamma_E ^2-25 \pi ^2+555 \gamma_E +216+300 \gamma_E  \log \left(\frac{16}{\pi }\right),\\ &
       \beta_3=216 \log ^2(|\boldsymbol{b}|)+18 (20 \gamma_E +1-16 \log (2)+4 \log (\pi )) \log (|\boldsymbol{b}|)+5 \pi ^2+176+15 \gamma_E  (10 \gamma_E +1-16 \log (2)\\ &\hspace{0.35cm}+4 \log (\pi ))+3 \log \left(\frac{16}{\pi }\right) \left(\log \left(\frac{256}{\pi ^2}\right)-1\right)\,,\\ &
       \beta_4=216 \log ^2(|\boldsymbol{b}|)+18 (20 \gamma_E +3-16 \log (2)+4 \log (\pi )) \log (|\boldsymbol{b}|)+5 \pi ^2+144+15 \gamma_E  (10 \gamma_E +3-16 \log (2)\\ &\hspace{0.35cm}+4 \log (\pi ))+3 \log \left(\frac{16}{\pi }\right) \left(\log \left(\frac{256}{\pi ^2}\right)-3\right)\,,\\ &
       \beta_5=432 \log ^2(|\boldsymbol{b}|)+6 (120 \gamma_E -137-96 \log (2)+24 \log (\pi )) \log (|\boldsymbol{b}|)+10 \pi ^2+151+12 \log ^2\left(\frac{\pi }{16}\right)+548 \log (2)\\ &\hspace{0.35cm}-137 \log (\pi )+5 \gamma_E  (60 \gamma_E -137-96 \log (2)+24 \log (\pi ))\,,\\ &
       \beta_6=2 \Big(432 \log ^2(|\boldsymbol{b}|)+6 (120 \gamma_E -103-96 \log (2)+24 \log (\pi )) \log (|\boldsymbol{b}|)+10 \pi ^2+12 \log ^2\left(\frac{\pi }{16}\right)\\ &\hspace{0.35cm}+103 (1+\log (16)-\log (\pi ))+5 \gamma_E  (60 \gamma_E -103-96 \log (2)+24 \log (\pi ))\Big)
    \end{split}
\end{align}
and,
\begin{align}
    \begin{split}
        &\delta_1=30 \log(|\boldsymbol{b}|)+25 \gamma_E +6-5 \log \left(\frac{16}{\pi }\right),\\ &
        \delta_2=-192 \log(|\boldsymbol{b}|) (18 \log(|\boldsymbol{b}|)+30 \gamma_E -25)+96 \log \left(\frac{16}{\pi }\right) \left(12 \log(|\boldsymbol{b}|)+10 \gamma_E +\log \left(\frac{\pi }{16}\right)\right)\\ &\hspace{0.7 cm}-800 \gamma_E  (3 \gamma_E -5)-80 \pi ^2-1199-800 \log \left(\frac{16}{\pi }\right)\,,\\ &
        \delta_3=216 \log ^2(|\boldsymbol{b}|)+6 (60 \gamma_E -13-48 \log (2)+12 \log (\pi )) \log(|\boldsymbol{b}|)+5 \left(\gamma_E  (30 \gamma_E -13)+\pi ^2-14\right)\\ &
        \hspace{0.7 cm}+\log \left(\frac{\pi }{16}\right) (60 \gamma_E -13-24 \log (2)+6 \log (\pi ))\,,\\ &
        \delta_4=84 \log \left(\frac{\pi }{16}\right) (6 \log(|\boldsymbol{b}|)+5 \gamma_E )+72 \log(|\boldsymbol{b}|) (21 \log(|\boldsymbol{b}|)+35 \gamma_E -10)+5 \left(30 \gamma_E  (7 \gamma_E -4)+7 \pi ^2-54\right)\\ &\hspace{0.7 cm}+6 \log \left(\frac{16}{\pi }\right) (20+28 \log (2)-7 \log (\pi ))\,,\\ &
        \delta_5=96 \left(19 \log \left(\frac{\pi }{16}\right) (6 \log(|\boldsymbol{b}|)+5 \gamma_E )+6 \log(|\boldsymbol{b}|) (57 \log(|\boldsymbol{b}|)+95 \gamma_E -91)\right)+760 \pi ^2+240 \gamma_E  (95 \gamma_E -182)\\ &\hspace{0.7 cm}+9642+14592 \log ^2(2)+7296 \log (8)+3264 \log \left(\frac{16}{\pi }\right)+912 \left(\log \left(\frac{\pi }{256}\right)-6\right) \log (\pi )\,.
    \end{split}
\end{align}
%Where
%\begin{align}
%\begin{split}
%\textcolor{red}{\mathcal{N}|_{\mathcal{S}^1}}=
 %  &\frac{-i}{4}(\ell_1+\ell_2-q)\cdot\mathcal{S}_1\cdot v_2\,\, \Bigg[ \Bigg(2\ell_1^2\left(\gamma  \left(4 (\ell_2\cdot v_2)^2+2\ell_2^2+q^2\right)-\left(\ell_1\cdot v_1\right) \left(\ell_2\cdot v_2\right)\right)\Bigg)\\&+\gamma  q^2 \left(4(\ell_2\cdot v_2)^2+2\ell_2^2+q^2\right)+2 \left(q^2-\ell_2^2\right) \left(\ell_1\cdot v_1\right) \left(\ell_2\cdot v_2\right)\Bigg]\\&-\frac{i}{4}(q-\ell_1-\ell_2)\cdot \mathcal{S}_2\cdot v_1\Bigg[2\ell_1^2(\gamma  \left(4(\ell_2\cdot v_2)^2+2\ell_2^2+q^2\right)\\&-\left(\ell_1\cdot v_1\right) \left(\ell_2\cdot v_2\right))+\gamma  q^2 \left(4(\ell_2\cdot v_2)^2+2\ell_2^2+q^2\right)+2 \left(q^2-\ell_2^2\right) \left(\ell_1\cdot v_1\right) \left(\ell_2\cdot v_2\right)\Bigg]\\&-\frac{i\gamma}{2}\ell_2\cdot \mathcal{S}_1\cdot (q-\ell_1 )\Bigg[\left(2\ell_1^2+q^2\right) \left(\ell_2\cdot v_2\right)\Bigg]+\frac{i}{4(\epsilon-1)}\ell_2\cdot\mathcal{S}_2\cdot(q-\ell_1)\Bigg[\left(2\ell_1^2+q^2\right) \left(\ell_2\cdot v_2\right)\Bigg]\\&+\frac{-i}{4(\epsilon-1)}q\cdot \mathcal{S}_1\cdot (\ell_1+\ell_2)\Bigg[\left(2\ell_2^2+q^2\right) \left(\ell_1\cdot v_1\right)+2 \gamma  (\epsilon -1) \left(2\ell_1^2+q^2\right) \left(\ell_2\cdot v_2\right)\Bigg]\\&+\frac{i}{4}q\cdot \mathcal{S}_2\cdot (\ell_1+\ell_2)\Bigg[ 2 \gamma  \left(2 \ell_2^2+q^2\right) \left( \ell_1 \cdot v_1\right)+\frac{\left(2 \ell_1^2+q^2\right) \left(  \ell_2 \cdot v_2\right)}{\epsilon -1}\Bigg]\\&+\ell_1\cdot\mathcal{S}_2\cdot(q-\ell_2)\Bigg[\frac{i\gamma}{2}(\ell_1\cdot v_1)(2\ell_2^2+q^2)\Bigg]+\frac{-i}{4(\epsilon-1)}\ell_1\cdot\mathcal{S}_1\cdot(q-\ell_2)\Bigg[\left(2\ell_2^2+q^2\right) \left(\ell_1\cdot v_1\right)\Bigg]
 %  \end{split}
%\end{align}
%In this Feynman topology we have the following vector integral,
%\begin{align}
 %   \begin{split}
 %      \textcolor{black}{\boldsymbol{\mathscr O}_{(\ell_1+\ell_2-q)\cdot S_1\cdot v_2}^\gamma}&=\frac{-i}{4}\int\frac{\hat\delta(\ell_1\cdot v_2)\hat\delta(\ell_2\cdot v_1)(\ell_1+\ell_2-q)^{\gamma}}{(\ell_1\cdot v_1+i\varepsilon)^2\,(\ell_2\cdot v_2-i\varepsilon)^2\,(\ell_1+\ell_2-q)^2\,(\ell_1-q)^2\,(\ell_2-q)^2}\Bigg[\Bigg]\\ &
%=\frac{i}{4}(\chi_2 v_2^\gamma+\chi_1 v_1^\gamma +\chi_3 q^\gamma)
 %   \end{split}
%\end{align}
%\begin{align}
%\begin{split}
% &
% \chi_1=\frac{-2 \gamma ^2 (\gamma +7)-4 \left(6 \gamma ^3+4 \gamma ^2-6 \gamma -1\right) \epsilon +400 (\gamma -2) \gamma  \epsilon ^3-8 (\gamma  (5 (\gamma -4) \gamma +12)-1) \epsilon ^2}{(\gamma -1) (\gamma +1)^2 (2 \epsilon -1) (2 \epsilon +1) (6 \epsilon -1)}\,q^2\mathbfcal{M}_{0,1;0,0,1,1,1}\\ &
 % \chi_2=\frac{1}{(\gamma -1) (\gamma +1)^2 (2 \epsilon -1) (2 \epsilon +1) (6 \epsilon -1)}\Bigg[2 \left(-2 (\gamma +1) \left(2 \gamma ^3-11 \gamma -1\right)-2 (\gamma -2) \gamma \right) \epsilon -\\&8 \left(\gamma  \left(5 \gamma  \left(\gamma ^2+\gamma -7\right)+2\right)-1\right) \epsilon ^2+400 (\gamma -2) \gamma  \epsilon ^3+2 \gamma  (\gamma  (3 \gamma  (\gamma +1)-8)-2)\Bigg]\,q^2\mathbfcal{M}_{0,1;0,0,1,1,1}\\ &
 %\chi_3=-\frac{1}{3 \left(\gamma ^2-1\right)^2 \epsilon ^2 (\gamma  \epsilon  (6 \epsilon -5)+\gamma )}\Bigg[-400 (\gamma -1) \gamma ^2 (\gamma +1) \epsilon ^4-4 (\gamma -1) (\gamma +1) \left(-37 \gamma ^2-2\right) \epsilon -24 (\gamma -1) \gamma ^2 \\&(\gamma +1)-4 (\gamma -1) (\gamma +1) \left(10 \gamma ^4-88 \gamma ^2\right) \epsilon ^3-4 (\gamma -1) (\gamma +1) \left(-4 \gamma ^4+77 \gamma ^2+7\right) \epsilon ^2\Bigg]\,\mathbfcal{M}_{0,0;0,0,1,1,1}\\&+\frac{1}{3 \left(\gamma ^2-1\right)^2 \epsilon ^2 (\gamma  \epsilon  (6 \epsilon -5)+\gamma )}\Bigg[-1200 \gamma ^2 \epsilon ^5-2 \left(-3 \gamma ^2-2\right) \epsilon -6 \gamma ^2-2 \left(60 \gamma ^4-428 \gamma ^2\right)\\& \epsilon ^4-2 \left(-102 \gamma ^4+138 \gamma ^2+14\right) \epsilon ^3-2 \left(24 \gamma ^4-37 \gamma ^2+3\right) \epsilon ^2\Bigg]\,q^2\mathbfcal{M}^{}_{0,0;0,0,1,2,1}\\&+\frac{1}{3 \left(\gamma ^2-1\right)^2 \epsilon ^2 (\gamma  \epsilon  (6 \epsilon -5)+\gamma )}\Bigg[600 \left(\gamma ^4-\gamma ^2\right) \epsilon ^5-2 \left(-27 \gamma ^4+25 \gamma ^2+2\right) \epsilon ^2-6 \left(\gamma ^4-\gamma ^2\right) \epsilon \\&-2 \left(-30 \gamma ^6+194 \gamma ^4-164 \gamma ^2\right) \epsilon ^4-2 \left(12 \gamma ^6+23 \gamma ^4-28 \gamma ^2-7\right) \epsilon ^3\Bigg]\,q^2\mathbfcal{M}^{}_{0,0;0,0,2,1,1}\\&-\frac{200 \left(\gamma ^5-\gamma ^3\right) \epsilon ^5-2 \left(-10 \gamma ^5+8 \gamma ^3+2 \gamma \right) \epsilon ^4-2 \left(60 \gamma ^5-59 \gamma ^3-\gamma \right) \epsilon ^3+30 \left(\gamma ^5-\gamma ^3\right) \epsilon ^2}{3 \left(\gamma ^2-1\right)^2 \epsilon ^2 (\gamma  \epsilon  (6 \epsilon -5)+\gamma )} \,q^2\mathbfcal{M}^{(+-)}_{1,1;0,0,1,1,1}
 %\end{split}
%\end{align}
%\begin{align}
 %   \begin{split}
  %     \textcolor{black}{\boldsymbol{\mathscr O}_{(\ell_1+\ell_2-q)\cdot S_2\cdot v_1}^\gamma}&=\frac{-i}{4}\int\frac{\hat\delta(\ell_1\cdot v_2)\hat\delta(\ell_2\cdot v_1)(\ell_1+\ell_2-q)^{\gamma}}{(\ell_1\cdot v_1+i\varepsilon)^2\,(\ell_2\cdot v_2-i\varepsilon)^2\,(\ell_1+\ell_2-q)^2\,(\ell_1-q)^2\,(\ell_2-q)^2}\Bigg[\Bigg]\\ &
%=\frac{-i}{4}(\chi_5 v_2^\gamma+\chi_4 v_1^\gamma +\chi_6 q^\gamma)
%    \end{split}
%\end{align}
%\begin{align}
%\begin{split}
% &
% \chi_4=\frac{1}{\left(\gamma ^2-1\right)^2 (2 \epsilon -1) (2 \epsilon +1) (6 \epsilon -1)}\Bigg[-2 (\gamma -1) \left(2 (\gamma +1) \left(2 \gamma ^3-7 \gamma -1\right)+2 (\gamma -2) \gamma \right) \epsilon\\& -8 (\gamma -1) \left(\gamma  \left(5 \gamma  \left(\gamma ^2+\gamma -7\right)+2\right)-1\right) \epsilon ^2+6 (\gamma -1) \gamma  \left(\gamma  \left(\gamma ^2+\gamma -4\right)-2\right)+400 (\gamma -2) (\gamma -1) \gamma  \epsilon ^3\Bigg]\,q^2\mathbfcal{M}_{0,1;0,0,1,1,1}\\ &
%  \chi_5=\frac{1}{\left(\gamma ^2-1\right)^2 (2 \epsilon -1) (2 \epsilon +1) (6 \epsilon -1)}\Bigg[2 \left(2 \left(\gamma  \left(-2 \gamma ^2+\gamma +4\right)+1\right)-2 (\gamma -2) \gamma \right) \epsilon +6 (\gamma -1) \gamma ^2\\&+400 (\gamma -2) \gamma  \epsilon ^3+8 (\gamma  (-5 (\gamma -4) \gamma -12)+1) \epsilon ^2\Bigg]\,q^2\mathbfcal{M}_{0,1;0,0,1,1,1}\\ &
% \chi_6=-\frac{1}{3 \left(\gamma ^2-1\right)^2 \epsilon ^2 (\gamma  \epsilon  (6 \epsilon -5)+\gamma )}\Bigg[-400 (\gamma -1) \gamma ^2 (\gamma +1) \epsilon ^4-4 (\gamma -1) (\gamma +1) \left(-37 \gamma ^2-2\right) \epsilon\\& -24 (\gamma -1) \gamma ^2 (\gamma +1)-4 (\gamma -1) (\gamma +1) \left(10 \gamma ^4-88 \gamma ^2\right) \epsilon ^3-4 (\gamma -1) (\gamma +1) \left(-4 \gamma ^4+77 \gamma ^2+7\right) \epsilon ^2\Bigg]\,\mathbfcal{M}_{0,0;0,0,1,1,1}\\&+\frac{1}{3 \left(\gamma ^2-1\right)^2 \epsilon ^2 (\gamma  \epsilon  (6 \epsilon -5)+\gamma )}\Bigg[-1200 \gamma ^2 \epsilon ^5-2 \left(-3 \gamma ^2-2\right) \epsilon -6 \gamma ^2-2 \left(60 \gamma ^4-428 \gamma ^2\right) \epsilon ^4\\&-2 \left(-102 \gamma ^4+138 \gamma ^2+14\right) \epsilon ^3-2 \left(24 \gamma ^4-37 \gamma ^2+3\right) \epsilon ^2\Bigg]\,q^2\mathbfcal{M}^{}_{0,0;0,0,1,2,1}\\&+\frac{1}{3 \left(\gamma ^2-1\right)^2 \epsilon ^2 (\gamma  \epsilon  (6 \epsilon -5)+\gamma )}\Bigg[600 \left(\gamma ^4-\gamma ^2\right) \epsilon ^5-2 \left(-27 \gamma ^4+25 \gamma ^2+2\right) \epsilon ^2-6 \left(\gamma ^4-\gamma ^2\right) \epsilon \\&-2 \left(-30 \gamma ^6+194 \gamma ^4-164 \gamma ^2\right) \epsilon ^4-2 \left(12 \gamma ^6+23 \gamma ^4-28 \gamma ^2-7\right) \epsilon ^3\Bigg]\,q^2\mathbfcal{M}^{}_{0,0;0,0,2,1,1}\\&-\frac{200 \left(\gamma ^5-\gamma ^3\right) \epsilon ^5-2 \left(-10 \gamma ^5+8 \gamma ^3+2 \gamma \right) \epsilon ^4-2 \left(60 \gamma ^5-59 \gamma ^3-\gamma \right) \epsilon ^3+30 \left(\gamma ^5-\gamma ^3\right) \epsilon ^2}{3 \left(\gamma ^2-1\right)^2 \epsilon ^2 (\gamma  \epsilon  (6 \epsilon -5)+\gamma )} \,q^2\mathbfcal{M}^{(+-)}_{1,1;0,0,1,1,1}
 %\end{split}
%\end{align}
%\begin{align}
   % \begin{split}
   %    \textcolor{black}{\boldsymbol{\mathscr O}_{\ell_2\cdot S_1\cdot q}^\gamma}&=\frac{-i\gamma}{2}\int\frac{\hat\delta(\ell_1\cdot v_2)\hat\delta(\ell_2\cdot v_1)(\ell_2)^{\gamma}}{(\ell_1\cdot v_1+i\varepsilon)^2\,(\ell_2\cdot v_2-i\varepsilon)^2\,(\ell_1+\ell_2-q)^2\,(\ell_1-q)^2\,(\ell_2-q)^2}\Bigg[\Bigg]\\ &
%=\frac{-i\gamma}{2}(\chi_7 \,(-\gamma v_1^\gamma +v_2^\gamma) +\chi_8 q^\gamma)
  %  \end{split}
%\end{align}
%where,
%\begin{align}
%\begin{split}
% &
% \chi_7=\frac{  (10 \epsilon -1)}{2 \left(\gamma ^2-1\right) \epsilon }\,q^2\mathbfcal{M}_{0,0;0,0,1,2,1}-\frac{5}{2}q^2\mathbfcal{M}_{0,0;0,0,2,1,1}+5\epsilon \mathbfcal{M}_{0,0;0,0,1,1,1}\\ &
% \chi_8=-\frac{ \left(-160 (\gamma -2) \epsilon ^3+(16-60 \gamma ) \epsilon ^2+3 \gamma -8 \epsilon \right)}{(\gamma -1) (\gamma +1) (2 \epsilon +1) (6 \epsilon -1)} \,\mathbfcal{M}^{(+-)}_{0,1;0,0,1,1,1}
 %\end{split}
%\end{align}
%\begin{align}
 %   \begin{split}
    %    \boldsymbol{\mathscr O}_{\text{l2}.\text{S2}.q}^\gamma&=\frac{i}{4(\epsilon-1)}\int\frac{\hat\delta(\ell_1\cdot v_2)\hat\delta(\ell_2\cdot v_1)(\ell_2)^{\gamma}}{(\ell_1\cdot v_1+i\varepsilon)^2\,(\ell_2\cdot v_2-i\varepsilon)^2\,(\ell_1+\ell_2-q)^2\,(\ell_1-q)^2\,(\ell_2-q)^2}\Bigg[\Bigg]\\ &
%=\frac{1}{4(\epsilon-1)}\left(\chi_9 (v_2^\gamma -\gamma v_1^\gamma) +\chi_{10} q^\gamma\right)
 %   \end{split}
%\end{align}
%,(starting from $\tilde{h}_{21}$)
%\begin{align}
%    \begin{split}
%   &  
%   \chi_9=\frac{10 i \epsilon -i }{2 \left(\gamma ^2-1\right) \epsilon }\,q^2\mathbfcal{M}_{0,0;0,0,1,2,1}-\frac{5 i \epsilon  \left(\gamma ^2 -1
%\right)}{2}q^2\mathbfcal{M}_{0,0;0,0,2,1,1}+\frac{5 i \epsilon}{8 \left(\gamma ^2-1\right) (\epsilon -1) \epsilon}\mathbfcal{M}_{0,0;0,0,1,1,1}\\ &
  % \chi_{10}=\frac{-160 i (\gamma -2) \epsilon ^3+i (16-60 \gamma ) \epsilon ^2+3 i \gamma -8 i \epsilon}{4 (\gamma -1) (\gamma +1) (\epsilon -1) (2 \epsilon +1) (6 \epsilon -1)}\mathbfcal{M}_{0,1;0,0,1,1,1}
  %  \end{split}
%\end{align}
%\begin{align}
  %  \begin{split}
  %   \hspace{0cm} \boldsymbol{  \mathscr{O}}_{\ell_1\cdot S_2\cdot (q-\ell_2)}^{\gamma}&=\frac{i\gamma}{2}\int\frac{\hat\delta(\ell_1\cdot v_2)\hat\delta(\ell_2\cdot v_1)\ell_1^{\gamma}}{(\ell_1\cdot v_1+i\varepsilon)^2\,(\ell_2\cdot v_2-i\varepsilon)^2\,(\ell_1+\ell_2-q)^2\,(\ell_1-q)^2\,(\ell_2-q)^2}\Bigg[\Bigg]\\&
    %  =\left(\chi_{12} q^\gamma+\chi_{11} (v_1^\gamma-\gamma v_2^\gamma)\right)
   % \end{split}
%\end{align}
%where,
%begin{align}
 %   \begin{split}
 %  &  
 %  \chi_{11}=\frac{i \gamma  q^2 (10 \epsilon -1)}{4 (\gamma ^2-1)\epsilon}\,q^2\mathbfcal{M}_{0,0;0,0,1,2,1}-\frac{5 i \gamma}{4 \epsilon}q^2\mathbfcal{M}_{0,0;0,0,2,1,1}+\frac{5 i\gamma \epsilon}{4 \epsilon}\mathbfcal{M}_{0,0;0,0,1,1,1}\\ &
  % \chi_{12}=\frac{-80 i\gamma  (\gamma -2) \epsilon ^3-4 i\gamma  (6 \gamma +1) \epsilon ^2+2 i\gamma  (\gamma -1) \epsilon +i\gamma }{(\gamma -1) (\gamma +1) (2 \epsilon +1) (6 \epsilon -1)}\mathbfcal{M}_{0,1;0,0,1,1,1}
 %   \end{split}
%\end{align}
%begin{align}
 %   \begin{split}
 %    \hspace{0cm} \boldsymbol{  \mathscr{O}}_{\ell_1\cdot S_1\cdot (q-\ell_2)}^{\gamma}&=\frac{-i}{4(\epsilon-1)}\int\frac{\hat\delta(\ell_1\cdot v_2)\hat\delta(\ell_2\cdot v_1)\ell_1^{\gamma}}{(\ell_1\cdot v_1+i\varepsilon)^2\,(\ell_2\cdot v_2-i\varepsilon)^2\,(\ell_1+\ell_2-q)^2\,(\ell_1-q)^2\,(\ell_2-q)^2}\Bigg[\Bigg]\\&
 %     =\left(\chi_{14}q^\gamma+\chi_{13}(v_2^\gamma-\gamma v_1^\gamma)\right)
%    \end{split}
%\end{align}
%where,
%\begin{align}
 %   \begin{split}
 %  &  
%   \chi_{13}=\frac{i (1-10 \epsilon ) }{8 \left(\gamma ^2-1\right) (\epsilon -1) \epsilon}\,q^2\mathbfcal{M}_{0,0;0,0,1,2,1}-\frac{5 i \left(\gamma ^2-1\right) \epsilon}{8 \left(\gamma ^2-1\right) (\epsilon -1) \epsilon}q^2\mathbfcal{M}_{0,0;0,0,2,1,1}+\frac{-5 i \epsilon}{8 \left(\gamma ^2-1\right) (\epsilon -1) \epsilon}\mathbfcal{M}_{0,0;0,0,1,1,1}\\&
 %  \chi_{14}=\frac{80 i (\gamma -2) \epsilon ^3+4 i (6 \gamma +1) \epsilon ^2-2 i (\gamma -1) \epsilon -i \gamma}{2 (\gamma -1) (\gamma +1) (\epsilon -1) (2 \epsilon +1) (6 \epsilon -1)}\mathbfcal{M}_{0,1;0,0,1,1,1}
 %   \end{split}
%\end{align}
%\begin{align}
 %   \begin{split}
 %    \hspace{0cm} \boldsymbol{  \mathscr{O}}_{q
  %   \cdot S_1\cdot(\ell_1+\ell_2)}^{\gamma}&=\frac{-i}{4(\epsilon-1)}\int\frac{\hat\delta(\ell_1\cdot v_2)\hat\delta(\ell_2\cdot v_1)(\ell_1+\ell_2)^{\gamma}}{(\ell_1\cdot v_1+i\varepsilon)^2\,(\ell_2\cdot v_2-i\varepsilon)^2\,(\ell_1+\ell_2-q)^2\,(\ell_1-q)^2\,(\ell_2-q)^2}\Bigg[\Bigg]\\&
   %   =\left(\chi_{17} q^\gamma+\chi_{15} v_1^\gamma+\chi_{16} v_2^\gamma\right)
  %  \end{split}
%\end{align}
%\begin{align}
%    \begin{split}
%   &  
%   \chi_{15} =\frac{10 i \gamma ^2 \epsilon ^2+i \left(-10 \gamma ^2-1\right) \epsilon -6 i}{24 \left(\gamma ^2-1\right)^2 (\epsilon -1) \epsilon ^2}\,q^2\mathbfcal{M}_{0,0;0,0,1,1,1}\\&-\frac{60 i \gamma ^2 q^2 \epsilon ^3-22 i \epsilon ^2 \left(3 \gamma ^2 q^2+q^2\right)+i \epsilon  \left(6 \gamma ^2 q^2-5 q^2\right)-6 i q^2}{24 \left(\gamma ^2-1\right)^2 (\epsilon -1) \epsilon ^2}q^2\mathbfcal{M}_{0,0;0,0,1,2,1}\\&+\frac{-6 i \epsilon  \left(\gamma ^2 q^2-q^2\right)-30 i \epsilon ^3 \left(\gamma ^4 q^2-\gamma ^2 q^2\right)+i \epsilon ^2 \left(30 \gamma ^4 q^2-11 \gamma ^2 q^2-19 q^2\right)}{24 \left(\gamma ^2-1\right)^2 (\epsilon -1) \epsilon ^2}\mathbfcal{M}_{0,0;0,0,2,1,1}
%\\&+\frac{30 i \epsilon ^2 \left(\gamma ^3 q^2-\gamma  q^2\right)-20 i \epsilon ^3 \left(\gamma ^3 q^2-\gamma  q^2\right)}{24 \left(\gamma ^2-1\right)^2 (\epsilon -1) \epsilon ^2}\mathbfcal{M}_{1, 1; 0, 0, 1, 1, 1}\\ &
%\chi_{16} =\frac{10 i \left(\gamma ^4-\gamma ^2\right) \epsilon ^2-11 i \left(\gamma ^4-\gamma ^2\right) \epsilon +i \left(-4 \gamma ^4+2 \gamma ^2+2\right)}{24 \left(\gamma ^2-1\right)^2 (\epsilon -1) \epsilon ^2}\,q^2\mathbfcal{M}_{0,0;0,0,1,1,1}\\&-\frac{60 i \gamma ^2 q^2 \epsilon ^3-88 i \gamma ^2 q^2 \epsilon ^2+i \epsilon  \left(5 \gamma ^2 q^2-4 q^2\right)-2 i \left(2 \gamma ^2 q^2+q^2\right)}{24 \left(\gamma ^2-1\right)^2 (\epsilon -1) \epsilon ^2}q^2\mathbfcal{M}_{0,0;0,0,1,2,1}\\&+\frac{-30 i \epsilon ^3 \left(\gamma ^4 q^2-\gamma ^2 q^2\right)+49 i \epsilon ^2 \left(\gamma ^4 q^2-\gamma ^2 q^2\right)+i \epsilon  \left(-4 \gamma ^4 q^2+2 \gamma ^2 q^2+2 q^2\right)}{24 \left(\gamma ^2-1\right)^2 (\epsilon -1) \epsilon ^2}\mathbfcal{M}_{0,0;0,0,2,1,1}
%\\&+\frac{i \epsilon ^2 \left(20 \gamma ^5 q^2-10 \gamma ^3 q^2-10 \gamma  q^2\right)-20 i \epsilon ^3 \left(\gamma ^5 q^2-\gamma ^3 q^2\right)}{24 \left(\gamma ^2-1\right)^2 (\epsilon -1) \epsilon ^2}\mathbfcal{M}_{1, 1; 0, 0, 1, 1, 1}\\&
 %   \end{split}
%\end{align}
%\begin{align}
 %   \begin{split}
 %    \hspace{0cm} \boldsymbol{  \mathscr{O}}_{q
 %    \cdot S_2\cdot(\ell_1+\ell_2)}^{\gamma}&=\frac{i}{4}\int\frac{\hat\delta(\ell_1\cdot v_2)\hat\delta(\ell_2\cdot v_1)(\ell_1+\ell_2)^{\gamma}}{(\ell_1\cdot v_1+i\varepsilon)^2\,(\ell_2\cdot v_2-i\varepsilon)^2\,(\ell_1+\ell_2-q)^2\,(\ell_1-q)^2\,(\ell_2-q)^2}\Bigg[\Bigg]\\&
  %    =\left(\chi_{20}  q^\gamma+\chi_{18} v_1^\gamma+\chi_{19} v_2^\gamma\right)
  %  \end{split}
%\end{align}
%where,
%\begin{align}
 %   \begin{split}
 %  &  
 % \chi_{18}=\frac{4 i \gamma  (\gamma +3) \epsilon ^2+i (\gamma  (-10 \gamma -3)+2) \epsilon -6 i}{24 \left(\gamma ^2-1\right)^2 (\epsilon -1) \epsilon ^2}\,\mathbfcal{M}_{0,0;0,0,1,1,1}\\&+\frac{24 i \epsilon ^3 \left(\gamma ^2 q^2+3 \gamma  q^2\right)+i \epsilon ^2 \left(-72 \gamma ^2 q^2-6 \gamma  q^2-4 q^2\right)+i \epsilon  \left(6 \gamma ^2 q^2-3 \gamma  q^2-2 q^2\right)-6 i q^2}{24 \left(\gamma ^2-1\right)^2 (\epsilon -1) \epsilon ^2}\mathbfcal{M}_{0,0;0,0,1,2,1}\\&+\frac{-6 i \epsilon  \left(\gamma ^2 q^2-q^2\right)-12 i \epsilon ^3 \left(\gamma ^4 q^2+3 \gamma ^3 q^2-\gamma ^2 q^2-3 \gamma  q^2\right)+i \epsilon ^2 \left(30 \gamma ^4 q^2+9 \gamma ^3 q^2-20 \gamma ^2 q^2-9 \gamma  q^2-10 q^2\right)}{24 \left(\gamma ^2-1\right)^2 (\epsilon -1) \epsilon ^2}\\&\mathbfcal{M}_{0,0;0,0,2,1,1}
%+\frac{30 i \epsilon ^2 \left(\gamma ^3 q^2-\gamma  q^2\right)-20 i \epsilon ^3 \left(\gamma ^3 q^2-\gamma  q^2\right)}{24 \left(\gamma ^2-1\right)^2 (\epsilon -1) \epsilon ^2}\mathbfcal{M}_{1, 1; 0, 0, 1, 1, 1}\\ &
%\chi_{19}=\frac{-16 i (\gamma -1) (\gamma +1) \gamma ^2 \epsilon ^2+17 i (\gamma -1) (\gamma +1) \gamma ^2 \epsilon -i (\gamma -1) (\gamma +1) \left(-4 \gamma ^2-2\right)}{24 \gamma  \left(\gamma ^2-1\right)^2 (\epsilon -1) \epsilon ^2}\,\mathbfcal{M}_{0,0;0,0,1,1,1}\\&+\frac{-96 i \gamma ^2 q^2 \epsilon ^3+118 i \gamma ^2 q^2 \epsilon ^2+i \epsilon  \left(\gamma ^2 q^2+4 q^2\right)+2 i \left(2 \gamma ^2 q^2+q^2\right)}{24 \gamma  \left(\gamma ^2-1\right)^2 (\epsilon -1) \epsilon ^2}\mathbfcal{M}_{0,0;0,0,1,2,1}\\&+\frac{48 i \epsilon ^3 \left(\gamma ^4 q^2-\gamma ^2 q^2\right)-67 i \epsilon ^2 \left(\gamma ^4 q^2-\gamma ^2 q^2\right)+i \epsilon  \left(4 \gamma ^4 q^2-2 \gamma ^2 q^2-2 q^2\right)}{24 \gamma  \left(\gamma ^2-1\right)^2 (\epsilon -1) \epsilon ^2}\mathbfcal{M}_{0,0;0,0,2,1,1}
%\\&+\frac{20 i \epsilon ^3 \left(\gamma ^5 q^2-\gamma ^3 q^2\right)+i \epsilon ^2 \left(-20 \gamma ^5 q^2+10 \gamma ^3 q^2+10 \gamma  q^2\right)}{24 \gamma  \left(\gamma ^2-1\right)^2 (\epsilon -1) \epsilon ^2}\mathbfcal{M}_{1, 1; 0, 0, 1, 1, 1}\\ &
 %  \textcolor{red}{\chi_{20}=\frac{8 i \left(59 \gamma ^2-210 \gamma +80\right) \epsilon ^3+4 i \left(48 \gamma ^2-3 \gamma -10\right) \epsilon ^2-2 i \left(9 \gamma ^2-10 \gamma +2\right) \epsilon -640 i (\gamma -2) \gamma  \epsilon ^4+3 i \gamma  (1-2 \gamma )}{2 (\gamma -1) (\gamma +1) (\epsilon -1) (2 \epsilon +1) (6 \epsilon -1)}\\&\mathbfcal{M}_{0,1;0,0,1,1,1}}
  %  \end{split}
%\end{align}
%\begin{align}
 %   \begin{split}
 %    \hspace{0cm} \boldsymbol{  \mathscr{O}}_{3/4,\ell_2\cdot S_1\cdot \ell_1 / \ell_2\cdot S_2 \cdot\ell_1}^{\eta\delta}&=\int\frac{\hat\delta(\ell_1\cdot v_2)\hat\delta(\ell_2\cdot v_1)\ell_1^{\eta}\ell_2^{\delta}}{(\ell_1\cdot v_1+i\varepsilon)^2\,(\ell_2\cdot v_2-i\varepsilon)^2\,(\ell_1+\ell_2-q)^2\,(\ell_1-q)^2\,(\ell_2-q)^2}\Bigg[\Bigg]\\&
 %     =\left(\Delta_1 q^{[\eta} v_1^{\delta]}+\Delta_2 q^{[\eta} v_2^{\delta]}\right)
 %   \end{split}
%\end{align}
%where,
%\begin{align}
%\begin{split}
%&\Delta_1=\tilde h_{61}\,\mathbfcal{M}_{0,0;0,0,1,1,1}+\tilde h_{62}\,q^2\,\mathbfcal{M}_{0,0;0,0,2,1,1}+\tilde h_{63}\,q^2\,\mathbfcal{M}_{0,0;0,0,1,2,1}-\tilde h_{64}\,q^2\,\mathbfcal{M}^{(+-)}_{1,1;0,0,1,1,1}\\ &
   %     \Delta_2=\tilde h_{65}\,\mathbfcal{M}_{0,0;0,0,1,1,1}+\tilde h_{66}\,
%q^2\,\mathbfcal{M}_{0,0;0,0,2,1,1}+\tilde h_{67}\,q^2\,\mathbfcal{M}_{0,0;0,0,1,2,1}-\tilde h_{68}\,q^2\,\mathbfcal{M}^{(+-)}_{1,1;0,0,1,1,1}
  %  \end{split}
%\end{align}
%\begin{align}
 %   \begin{split}
 %    \hspace{0cm} \boldsymbol{  \mathscr{O}}_{8/9,\ell_1\cdot S_1\cdot \ell_2 / \ell_1\cdot S_2 \cdot\ell_2}^{\eta\delta}&=\int\frac{\hat\delta(\ell_1\cdot v_2)\hat\delta(\ell_2\cdot v_1)\ell_1^{\eta}\ell_2^\delta}{(\ell_1\cdot v_1+i\varepsilon)^2\,(\ell_2\cdot v_2-i\varepsilon)^2\,(\ell_1+\ell_2-q)^2\,(\ell_1-q)^2\,(\ell_2-q)^2}\Bigg[\Bigg]\\&
  %    =\left(\Delta_1 q^{[\eta} v_1^{\delta]}+\Delta_2 q^{[\eta} v_2^{\delta]}\right)
  %  \end{split}
%\end{align}
%where,
%\begin{align}
%\begin{split}
%&\Delta_1=\tilde h_{51}\,\mathbfcal{M}_{0,0;0,0,1,1,1}+\tilde h_{52}q^2\mathbfcal{M}_{0,0;0,0,2,1,1}+\tilde h_{53}\,q^2\,\mathbfcal{M}_{0,0;0,0,1,2,1}-\tilde h_{54}\,q^2\mathbfcal{M}^{(+-)}_{1,1;0,0,1,1,1}\\
  %     & \Delta_2=\tilde h_{55}\,\mathbfcal{M}_{0,0;0,0,1,1,1}+\tilde h_{56}\,q^2\mathbfcal{M}_{0,0;0,0,2,1,1}+
   %   \tilde h_{57}\,q^2\,\mathbfcal{M}_{0,0;0,0,1,2,1}-\tilde h_{58}\,q^2\mathbfcal{M}^{(+-)}_{1,1;0,0,1,1,1}
 %   \end{split}
%\end{align}
%Finally we get, 
%\begin{align}
 %   \begin{split}
 %      i\, \chi\Big|_{\mathcal{S}^1}=-i \frac{s_1^3 s_2^3 m_1 m_2^3}{64 m_p^6}&\Bigg[\frac{2x\,\Bigg(\hat{b}\cdot \mathcal{S}_1\cdot v_2\Bigg)}{192 \pi ^3 b^3 \left(x^2-1\right)^5 \left(1-x^4\right)}\Bigg(\left(x^4-1\right) \big(57 x^8-72 x^7-12 x^6-8 x^5 \left(5 \gamma -3+\log \left(\frac{\pi }{16}\right)\right)\\&+2 x^4 (40 \gamma -117-32 \log (2)+8 \log (\pi ))-8 x^3 \left(5 \gamma -3+\log \left(\frac{\pi }{16}\right)\right)-12 x^2-72 x+57\big)\\&-48 (x-1)^3 x^3 \left(x^3+x^2+x+1\right) \log (|\boldsymbol{b}|)+8 \big(6 x^{12}-6 x^{11}+24 x^{10}-17 x^9-2 x^8+11 x^7\\&-56 x^6+11 x^5-2 x^4-17 x^3+24 x^2-6 x+6\big)\log (x)\Bigg)+\\&\frac{2x\,\Bigg(\hat{b}\cdot \mathcal{S}_2\cdot v_1\Bigg)}{64 \pi ^3 b^3 \left(x^2-1\right)^5 \left(1-x^4\right)}\Bigg((x-1) x (x+1) \left(x^2+1\right) \big(2 x (x (3 x+2)+3) (x-1)^2 \log \left(\pi  b^6\right)\\&+x \big(15 x^5+5 x^4+41 x^3+34 x^2+41 x+10 \gamma_E  (x-1)^2 (x (3 x+2)+3)\\&-8 (x-1)^2 (x (3 x+2)+3) \log (2)+5\big)+15\big)+\big(2 x^2 \big(x^8+6 x^7-12 x^6+90 x^5\\&-90 x^4+90 x^3-12 x^2+6 x+1\big)\big)\log(x)\Bigg)\Bigg]\,.
 %   \end{split}
%\end{align}
\textbullet $\,\,$ We now compute the 3PM diagrams with a graviton propagator in the middle. The topology has the following form,
\begin{center}
\includegraphics[width=0.38\linewidth]{neww_3.png}
\end{center}
%The contribution to the eikonal phase is given by,
%\begin{align}
    %i\chi=-\frac{i m_1^2 m_2^2 s_1^2 s_2^2}{64 m_p^6}\int_{q}e^{-iq\cdot %b}\hat\delta(q\cdot v_1)\hat\delta(q\cdot v_2)\int_{\ell_{1},\ell_{2}}\frac{\hat\delta(\ell_{1}\cdot v_2)\hat\delta(\ell_{2}\cdot v_1)\,\mathcal{N}}{(\ell_{1}\cdot v_1)^2(\ell_{2}\cdot v_2)^2\,\ell_{1}^2\,\ell_{2}^2\,(\ell_{1}+\ell_{2}-q)^2}
%\end{align}
%where upto first order in spin,\quad $\mathcal{N}=\mathcal{N}\Big|_{\mathcal{S}^0}+\mathcal{N}\Big|_{\mathcal{S}^1}\,,$
%\begin{align}
  %  \begin{split}
 %   &  \mathcal{N}\Big|_{\mathcal{S}^0}=  , \\&
%\mathcal{N}\Big|_{\mathcal{S}^1}=-\frac{i}{4}(q-\ell_1-\ell_2)\cdot\mathcal{S}\cdot v_2\Bigg[2(\ell_1\cdot v_1)(\ell_2\cdot v_2)\left(q^2-(\ell_1-q)^2-(\ell_2-q)^2\right)+\gamma(\ell_2-q)^2\left((\ell_1-q)^2-4(\ell_2\cdot v_2)^2\right)\Bigg]\\ &
%\hspace{0.35cm}+\frac{i}{2(d-2)}(q-\ell_1-\ell_2)\cdot \mathcal{S}_1\cdot \ell_1\,(\ell_1\cdot v_1)(\ell_1-q)^2-\frac{i\gamma}{2} (q-\ell_1-\ell_2)\cdot \mathcal{S}_1\cdot \ell_2\,(\ell_2\cdot v_2)(\ell_2-q)^2\\ &
%\hspace{0.35cm}+\frac{i}{4}(q-\ell_1-\ell_2)\cdot \mathcal{S}_2\cdot v_1\,\Bigg[2(\ell_1\cdot v_1)(\ell_2\cdot v_2)\left(q^2-(\ell_1-q)^2-(\ell_2-q)^2\right)-\gamma(\ell_1-q)^2\left(-(\ell_2-q)^2+4(\ell_1\cdot v_1)^2\right)\Bigg]\\ &
%\hspace{0.35cm}+\frac{i\gamma}{2}(q-\ell_1-\ell_2)\cdot \mathcal{S}_2\cdot \ell_1 (\ell_1\cdot v_1)(\ell_1-q)^2+\frac{i}{2(2-d)}(q-\ell_1-\ell_2)\cdot \mathcal{S}_2\cdot \ell_2\,(\ell_2\cdot v_2)(\ell_2-q)^2 
 %   \end{split}
%\end{align}
      %\\ &
%\mathcal{N}[q,l_i,v_i]\Big|_{\mathcal{S}^1}=
%\frac{-1}{2}i\Bigg[8\gamma (-q\cdot \ell_{1}+\ell_{1}^2+\ell_{1}\cdot \ell_{2})\Big((-\ell_{1}-\ell_{2}+q)\cdot\mathcal{S}_1\cdot v_{2}\Big)(\ell_{2}\cdot v_2)^2+4\gamma (q\cdot \ell_{2}-\ell_{2}^2-\ell_{1}\cdot \ell_{2})\Big((-\ell_{1}-\ell_{2}\\&+q)\cdot\mathcal{S}_2\cdot v_{1}\Big)(\ell_{1}\cdot v_1)^2+4\gamma (q\cdot \ell_{1}-\ell_{1}^2-\ell_{1}\cdot \ell_{2})\Big((-\ell_{1}+q)\cdot\mathcal{S}_1\cdot \ell_{2}\Big)\,(\ell_{2}\cdot v_2)-4\gamma (q\cdot \ell_{1}-\ell_{1}^2-\ell_{1}\cdot \ell_{2})\\&(-q\cdot \ell_{2}+\ell_{2}^2+\ell_{1}\cdot \ell_{2})\Big((q-\ell_{2}-\ell_{1})\cdot\mathcal{S}_1\cdot v_{2}\Big)+2 (q\cdot \ell_{2}-\ell_{2}^2-\ell_{1}\cdot \ell_{2})\Big(l_{1}\cdot\mathcal{S}_1\cdot(-\ell_{2}+q)\Big)(\ell_{1}\cdot v_1)\\&+4(\ell_{1}\cdot \ell_{2})(\ell_{2}\cdot v_2)(\ell_{1}\cdot v_1)\Big((-\ell_{1}-\ell_{2}+q)\cdot\mathcal{S}_1\cdot {v}_{2}\Big)+2\gamma (q\cdot \ell_{2}-\ell_{2}^2-\ell_{1}\cdot \ell_{2})(\ell_{1}\cdot v_1)\Big(l_{1}\cdot\mathcal{S}_2\cdot(-\ell_{2}+q)\Big)\\&+2\gamma (q\cdot \ell_{1}-\ell_{1}^2-\ell_{1}\cdot \ell_{2})(-q\cdot \ell_{2}+\ell_{2}^2+\ell_{1}\cdot \ell_{2})\Big((-\ell_{1}-\ell_{2}+q)\cdot\mathcal{S}_2\cdot v_{1}\Big)\,-\,2(\ell_{1} \cdot \ell_{2})(\ell_{2} \cdot v_2)(\ell_{1}\cdot v_1)\Big((-\ell_{1}-\ell_{2}+q)\cdot\mathcal{S}_2\cdot v_{1}\Big)\\&+(q\cdot \ell_{1}-\ell_{1}^2-\ell_{1}\cdot \ell_{2})(\ell_{2}\cdot v_2)\Big((q-\ell_{1})\cdot  \mathcal{S}_2\cdot \ell_{2}\Big)\Bigg]
 %    \\ &
%\mathcal{N}[q,l_i,v_i]\Big|_{\mathcal{S}^2}= \textcolor{black}{\Bigg[-\frac{1}{4} \Bigg(2 \left(-q\cdot \ell_{2}+\ell_{2}^2+\ell_{1} \cdot \ell_{2}\right) \left(\ell_{1} \cdot v_1\right)^2-\left(-q\cdot \ell_{1} +\ell_{1}^2+ \ell_{1}\cdot \ell_{2}\right) \left(-q.\ell_{2}+\ell_{2}^2+\ell_{1} \cdot \ell_{2}\right)\Bigg)\mathcal{S}_{1\,\mu\nu}\mathcal{S}_{1,\alpha}{}^{\nu}}\\&\textcolor{black}{+\frac{1}{4}\Bigg( -2 \left(-q\cdot \ell_{1}+\ell_{1}^2+\ell_{1}.\cdot \ell_{2}\right) \left(\ell_{2}\cdot v_2\right)^2+\left(- q \cdot \ell_{1}+\ell_{1}^2+\ell_{1}\cdot \ell_{2}\right) \left(-q\cdot \ell_{2}+\ell_{2}^2+\ell_{1}\cdot \ell_{2}\right)\Bigg)\mathcal{S}_{2\,\mu\nu}\mathcal{S}_{2,\alpha}{}^{\nu}}\\&\textcolor{black}{+\frac{1}{4}\Bigg( 4 \left(-q\cdot \ell_{2} +\ell_{2}^2+\ell_{1} \cdot  \ell_{2}\right)\left( \ell_{1} \cdot v_1\right)^2-4 \left(-q\cdot  \ell_{1}+\ell_{1}^2+\ell_{1}\cdot \ell_{2}\right) \left(\ell_{2}\cdot v_2\right)^2+8 \left(\ell_{2} \cdot v_2\right)^2\left(\ell_{1} \cdot v_1\right)^2}\\&\textcolor{black}{-2 \left(-q.\ell_{2}+\ell_{2}^2+\ell_{1}\cdot \ell_{2}\right) \left(-q\cdot \ell_{1} +\ell_{1}^2 + \ell_{1} \cdot \ell_{2}\right)\Bigg)\mathcal{S}_{1\,\mu\nu}\mathcal{S}_{1,\alpha \beta}v_2{}^{\nu}v_2{}^{\beta}}\\&\textcolor{black}{+\frac{1}{4}\Bigg(- 4 \left(-q\cdot \ell_{2} +\ell_{2}^2+\ell_{1} \cdot  \ell_{2}\right)\left( \ell_{1} \cdot v_1\right)^2+4 \left(-q\cdot  \ell_{1}+\ell_{1}^2+\ell_{1}\cdot \ell_{2}\right) \left(\ell_{2}\cdot v_2\right)^2+8 \left(\ell_{2} \cdot v_2\right)^2\left(\ell_{1} \cdot v_1\right)^2}\\&\textcolor{black}{-2 \left(-q.\ell_{2}+\ell_{2}^2+\ell_{1}\cdot \ell_{2}\right) \left(-q\cdot \ell_{1} +\ell_{1}^2 + \ell_{1} \cdot \ell_{2}\right)\Bigg)\mathcal{S}_{2\,\mu\nu}\mathcal{S}_{2,\alpha \beta}v_1{}^{\nu}v_1{}^{\beta}}\\&\textcolor{black}{+ \Bigg(\left(  \ell_{1} \cdot \ell_{2}\right) \left(  \ell_{1} \cdot v_1\right) \left(  \ell_{2} \cdot v_2\right)-2\gamma  \left(  \ell_{1}\cdot q-\ell_{1}^2-\ell_{2}\cdot  \ell_{1}\right) \left(-  \ell_{2} \cdot q+\ell_{2}^2+ \ell_{1} \cdot \ell_{2}\right)\Bigg)\mathcal{S}_{1\,\mu\nu}\mathcal{S}_{2,\alpha}{}^{\nu}}\\&\textcolor{black}{+\Bigg(\left(-q\cdot \ell_{1} +\ell_{1}^2+\ell_{2}\cdot \ell_{1}\right) \left(-q\cdot \ell_{2}+\ell_{2}^2+\ell_{1}\cdot \ell_{2}\right)\Bigg)\mathcal{S}_{1\,\mu\nu}\mathcal{S}_{2,\alpha\beta}v_2{}^{\nu}v_1{}^{\beta}\Bigg]}(q-\ell_{1}-\ell_{2})^{\mu}(q-\ell_{1}-\ell_{2})^{\alpha}\\&\textcolor{red}{-\Bigg[ (\ell_{1}\cdot v_1)(\ell_{2}\cdot v_2)\mathcal{S}_{1\,\alpha\beta}\mathcal{S}_{2\,\mu\nu}  - (\ell_{1}\cdot v_1)(\ell_{2}\cdot v_2)\mathcal{S}_{1\,\mu\nu} \mathcal{S}_{2\,\alpha\beta} \Big)\Bigg]}(q-\ell_{1})^{\mu} \ell_{2}{}^{\nu}(q-\ell_{1}-\ell_{2})^{\alpha}\ell_{1}{}^{\beta}\\&-\Bigg(\left(- q\cdot \ell_{1} +\ell_{1}^2+\ell_{1}\cdot \ell_{2}\right) \left( \ell_{2}\cdot v_2\right)\Bigg)\Big( (q-\ell_{1}) \cdot \mathcal{S}_1 \cdot \ell_{2}\Big)\Big((q-\ell_{1}-\ell_{2})\cdot \mathcal{S}_2\cdot v_1\Big)\\&-\Bigg(\left(- q\cdot \ell_{2} +\ell_{2}^2+\ell_{1}\cdot \ell_{2}\right) \left( \ell_{1}\cdot v_1\right)\Bigg)\Big( (q-\ell_{2}) \cdot \mathcal{S}_2 \cdot \ell_{1}\Big)\Big((q-\ell_{1}-\ell_{2})\cdot \mathcal{S}_1\cdot v_2\Big)\\&-\Bigg(-\left(-q\cdot \ell_{2}+\ell_{2}^2+\ell_{1}\cdot \ell_{2}\right) \left(\ell_{1}\cdot v_1\right)+2\left(\ell_{1}\cdot v_1\right) \left(\ell_{2} \cdot v_2\right)^2\Bigg)\Big( (q-\ell_{2}) \cdot \mathcal{S}_2 \cdot \ell_{1}\Big)\Big((q-\ell_{1}-\ell_{2})\cdot \mathcal{S}_2\cdot v_1\Big)+\\& \Bigg(\left(-q\cdot \ell_{1}+\ell_{1}^2+\ell_{1}\cdot \ell_{2}\right) \left(\ell_{2}\cdot v_2\right)-2\left(\ell_{2}\cdot v_2\right) \left(\ell_{1} \cdot v_1\right)^2\Bigg)\Big( (q-\ell_{1}) \cdot \mathcal{S}_1 \cdot \ell_{2}\Big)\Big((q-\ell_{1}-\ell_{2})\cdot \mathcal{S}_1\cdot v_2\Big)
%    \end{split}
%\end{align}
%\noindent
%First we write down the results for the first diagram.
 At $\mathcal{O}(\mathcal{S}^0)$ we get the total contribution from these two diagrams after doing the IBP reduction, 
\begin{align}
    \begin{split}
\chi\Big|_{\mathcal{S}^0}=-\frac{ m_1^2 m_2^2 s_1^2 s_2^2}{256 m_p^6}\,\int e^{-i q \cdot b}\, \hat\delta(q\cdot v_1)\hat\delta(q\cdot v_2)&\Big(\boldsymbol{\tilde h}_{55}\,\,\mathbfcal{M}_{0,0;0,0,1,1,1}+\boldsymbol{\tilde h}_{56}\,\,q^2\mathbfcal{M}_{0,0;0,0,2,1,1}\\&+\boldsymbol{\tilde h}_{57}\,\,q^2\mathbfcal{M}_{0,0;0,0,1,2,1}+\boldsymbol{\tilde h}_{58}\,\,q^2\mathbfcal{M}^{(++)}_{1,1;0,0,1,1,1}\Big)\,,
   \end{split}
\end{align}
and proceeding as before we get at $\mathcal{O}(\mathcal{S}^1)$, 
\begin{align}
    \begin{split}
\chi\Big|_{\mathcal{S}^1}=-\frac{ m_1^2 m_2^2 s_1^2 s_2^2}{64m_p^6}\,\int e^{-i q \cdot b}\, \hat\delta(q\cdot v_1)\hat\delta(q\cdot v_2)&\Bigg[i\,\Big(\boldsymbol{\tilde h}_{59}\,\,\mathbfcal{M}_{0,0;0,0,1,1,1}+\boldsymbol{\tilde h}_{60}\,\,q^2\mathbfcal{M}_{0,0;0,0,2,1,1}\\&+\boldsymbol{\tilde h}_{61}\,\,q^2\mathbfcal{M}_{0,0;0,0,1,2,1}-\frac{1}{2}\boldsymbol{\tilde h}_{62}\,\,q^2\mathbfcal{M}^{(++)}_{1,1;0,0,1,1,1}\\&-\frac{1}{4}\boldsymbol{\tilde h}_{63}\,\,q^2\mathbfcal{M}^{(+-)}_{1,1;0,0,1,1,1}\Big)\big(q\cdot\mathcal{S}_1\cdot v_2\big)\\&+i\,\Big(\boldsymbol{\tilde h}_{64}\,\,\mathbfcal{M}_{0,0;0,0,1,1,1}+\boldsymbol{\tilde h}_{65}\,\,q^2\mathbfcal{M}_{0,0;0,0,2,1,1}\\&+\boldsymbol{\tilde h}_{66}\,\,q^2\mathbfcal{M}_{0,0;0,0,1,2,1}-\frac{1}{2}\boldsymbol{\tilde h}_{67}\,\,q^2\mathbfcal{M}_{1,1;0,0,1,1,1}\\&-\frac{1}{4}\boldsymbol{\tilde h}_{68}\,\,q^2\mathbfcal{M}^{(+-)}_{1,1;0,0,1,1,1}\Big)\big(q\cdot\mathcal{S}_2\cdot v_1\big)\Bigg]\,.
 \end{split}
\end{align}
%where, $\{\boldsymbol{h_{9}}\to \boldsymbol{h_{12}}\}$ are defined in \eqref{E.1}.
%\newpage
Finally, using \eqref{MasterM} we get,

\begin{align}
    \begin{split} \label{eqt7b}
        \chi\Big|_{\mathcal{S}^0}=\frac{m_1^2m_2^2}{m_p^6}&\Bigg[\Big(\Delta_{\mathfrak{12}}^{\textrm{poly}}+\Delta_{\mathfrak{05}}^{\textrm{log}}\log (x)\Bigg] 
        \end{split}
\end{align}
where, 
\begin{align}
    \begin{split}
   & \Delta_{\mathfrak{12}}^{\textrm{poly}}=-\frac{s_1^2s_2^2 (1+x^2)}{8192 \pi ^3 |\boldsymbol{b}|^2 \left(x^2-1\right)^5}\Bigg[48 \left(3 x^4-2 x^2+3\right) x^2 \log (|\boldsymbol{b}|)+x^8+24 x^6 \left(5 \gamma_E +1+\log \left(\frac{\pi }{16}\right)\right)\\&\hspace{0.8cm}-2 x^4 (40 \gamma_E +17-32 \log (2)+8 \log (\pi ))+24 x^2 \left(5 \gamma_E +1+\log \left(\frac{\pi }{16}\right)\right)+1\Bigg]\,,\\&
   \Delta_{\mathfrak{05}}^{\textrm{log}}=-\frac{ s_1^2 s_2^2 \left(3 x^4+22 x^2+3\right)}{6144 \pi ^3 |\boldsymbol{b}|^2 \left(x^2-1\right)^2 }\,.
      \end{split}
\end{align}
The total contribution at $\mathcal{O}(\mathcal{S}^1)$ is, 
\begin{align}
\begin{split} \label{eqt8}
\chi\Big|_{\mathcal{S}^1} =\frac{m_1^2m_2^2}{m_p^6}&\Bigg[\Big(\Delta_{\mathfrak{13}}^{\textrm{poly}}+\Delta_{\mathfrak{06}}^{\textrm{log}}\log (x)\Big)\Big(\hat{b}\cdot \mathcal{S}_1\cdot v_2\Big)+\Big(\Delta_{\mathfrak{14}}^{\textrm{poly}}+\Delta_{\mathfrak{07}}^{\textrm{log}}\log (x)\Big)\Big(\hat{b}\cdot \mathcal{S}_2\cdot v_1\Big)\Bigg]
\end{split}
\end{align}
where,
\begin{align}
\begin{split}
&\Delta_{\mathfrak{13}}^{\textrm{poly}}=\frac{s_1^2s_2^2\,x}{3072 \pi ^3 |\boldsymbol{b}|^3 \left(x^2-1\right)^5}\Bigg[132 \left(x^3+x\right)^2 \log (|\boldsymbol{b}|)+3 x^8-6 x^7+2 x^6 (55 \gamma_E -38-44 \log (2)+11 \log (\pi ))\\&\hspace{0.8cm}-18 x^5+22 x^4 (10 \gamma_E -5-8 \log (2)+2 \log (\pi ))-18 x^3\\&\hspace{0.8cm}+2 x^2 (55 \gamma_E -38-44 \log (2)+11 \log (\pi ))-6 x+3\Bigg]\,,\\&
\Delta_{\mathfrak{14}}^{\textrm{poly}}=-\frac{s_1^2s_1^2 x}{6144 \pi ^3 |\boldsymbol{b}|^3 \left(x^2-1\right)^5 }\Bigg[48 \left(48 x^4+55 x^2+48\right) x^2 \log (|\boldsymbol{b}|)+15 x^8+2 x^6 (960 \gamma_E -359-768 \log (2)\\&\hspace{0.8cm}+192 \log (\pi ))-24 x^5+10 x^4 (220 \gamma_ E -49-176 \log (2)+44 \log (\pi ))-24 x^3\\&\hspace{0.8cm}+2 x^2 (960 \gamma_E -359-768 \log (2)+192 \log (\pi ))+15\Bigg]\,,\\&
\Delta_{\mathfrak{6}}^{\textrm{log}}=\frac{3s_1^2 s_2^2 x^3 \left(x^4-1\right)}{256 \pi ^3 |\boldsymbol{b}|^3 \left(x^2-1\right)^5 }\,,\Delta_{\mathfrak{7}}^{\textrm{log}}=-\frac{s_1^2 s_2^2 x \left(x^8-10 x^6+x^4\right)}{128 \pi ^3 |\boldsymbol{b}|^3 \left(x^2-1\right)^6 \left(x^2+1\right) }\,.
\end{split}
\end{align}
%\hfsetfillcolor{gray!10}
%\hfsetbordercolor{black!150}
%\begin{align}
 %   \begin{split}
 %       \tikzmarkin[disable rounded corners=false]{e67}(4.5,-1)(-0.1,1)
%i\chi\Big|_{\mathcal{S}^0}& =-\frac{i m_1^2 m_2^2 s_1^2 s_2^2}{1536\,\pi^3\,|\boldsymbol{b}|^2\,(1-x^2)^5\, m_p^6}\,\Bigg(4 \left(x^8+x^2\right) \log \left(\frac{\pi  |\boldsymbol{b}|^6}{16}\right)+4 \left((5 \gamma_E -1) x^6+x^4+x^2+5 \gamma_E -1\right) x^2\\&\hspace{0.8cm}+(x^2-1) \left(x^4+8 x^2+1\right) \left(3 x^4-4 x^2+3\right) \log (x)\Bigg)\,.
%\tikzmarkend{e67}
%  \end{split}
%\end{align}
%Then we are left with the following set of terms, 
%\begin{align}
   % \begin{split}
% \mathcal{N}[q,l_i,v_i]\Big|_{\mathcal{S}^1}=
%-i\Bigg[& 2\Bigg(2\gamma (-q\cdot \ell_{1}+\ell_{1}^2+\ell_{1}\cdot \ell_{2})(\ell_{2}\cdot v_2)^2-\gamma (q\cdot \ell_{1}-\ell_{1}^2-\ell_{1}\cdot \ell_{2})(-q\cdot \ell_{2}+\ell_{2}^2+\ell_{1}\cdot \ell_{2})\\&+(\ell_{1}\cdot \ell_{2})(\ell_{2}\cdot v_2)(\ell_{1}\cdot v_1)\Bigg)\Big((-\ell_{1}-\ell_{2}+q)\cdot\mathcal{S}_1\cdot v_{2}\Big)\\&+\Bigg(2\gamma (q\cdot \ell_{2}-\ell_{2}^2-\ell_{1}\cdot \ell_{2})(\ell_{1}\cdot v_1)^2+\gamma (q\cdot \ell_{1}-\ell_{1}^2-\ell_{1}\cdot \ell_{2})(-q\cdot \ell_{2}+\ell_{2}^2+\ell_{1}\cdot \ell_{2})\\&-(\ell_{1} \cdot \ell_{2})(\ell_{2} \cdot v_2)(\ell_{1}\cdot v_1)\Bigg)\Big((-\ell_{1}-\ell_{2}+q)\cdot\mathcal{S}_2\cdot v_{1}\Big)\Bigg]
  %  \end{split}
%\end{align}
%Next we proceed to compute the eikonal phase at $\mathcal{O}(\mathcal{S}^1).$ To do that we will need to evaluate the following tensor integrals using PV reduction. We list them below.
%Ignoring the tadpoles we have the following numerator,
%\begin{align}
 %   \begin{split}
%\mathcal{N}[q,l_i,v_i]\Big|_{\mathcal{S}^1}=
%-i\Bigg[& 2\Bigg(-\gamma (\ell_{2}-q)^2(\ell_{2}\cdot v_2)^2+\frac{\gamma}{4} (\ell_{2}-q)^2 (\ell_{1}-q)^2\\&+\frac{1}{2}\Big(q^2 - (q-\ell_{1})^2 - (q-\ell_{2})^2\Big)(\ell_{2}\cdot v_2)(\ell_{1}\cdot v_1)\Bigg)\Big((-\ell_{1}-\ell_{2}+q)\cdot\mathcal{S}_1\cdot v_{2}\Big)\\&+\Bigg(\gamma (\ell_{1}-q)^2(\ell_{1}\cdot v_1)^2-\frac{\gamma}{4} (\ell_{2}-q)^2(\ell_{1}-q)^2)\\&-\frac{1}{2}\Big(q^2 - (q-\ell_{1})^2 - (q-\ell_{2})^2\Big)(\ell_{2}\cdot v_2)(\ell_{1}\cdot v_1)\Bigg)\Big((-\ell_{1}-\ell_{2}+q)\cdot\mathcal{S}_2\cdot v_{1}\Big)\Bigg]\,.
 %  \end{split}
%\end{align}
%We will now face  the following tensor integrals,
%\begin{align}
 %   \begin{split}
        %\mathcal{K}_1^{\delta}&\,\equiv\,\int_{\ell_{1},\ell_{2}}\frac{\hat\delta(\ell_{1}\cdot v_2)\hat\delta(\ell_{2}\cdot v_1)\,(\ell_1+\ell_2-q)^\delta}{(\ell_{1}\cdot v_1)^2(\ell_{2}\cdot v_2)^2\ell_{1}^2 \ell_{2}^2 (\ell_{1}+\ell_{2}-q)^2}\,\Bigg[2(\ell_1\cdot v_1)(\ell_2\cdot v_2)\left(q^2-(\ell_1-q)^2-(\ell_2-q)^2\right)\\ &
 %       \hspace{0.8cm}+\gamma(\ell_2-q)^2\left((\ell_1-q)^2-4(\ell_2\cdot v_2)^2\right)\Bigg]\\ &
 %       =\Bigg(\boldsymbol{h_{13}}\,  \mathbfcal{M}_{0,0;0,0,1,1,1}
  %      + \boldsymbol{h_{14}}q^2 \mathbfcal{M}_{0,0;0,0,2,1,1}+\,\,\boldsymbol{h_{15}} q^2\mathbfcal{M}_{0,0;0,0,1,2,1}+\boldsymbol{h_{16}}\,q^2\mathbfcal{M}_{1,1;0,0,1,1,1}\Bigg)q^\delta\\ &
 %       \hspace{0.35cm}+ \boldsymbol{h_{17}}\,q^2 \mathbfcal{M}^{(++)}_{0,1;0,0,1,1,1}\,v_1^\delta+\boldsymbol{h_{18}}\,q^2 \mathbfcal{M}^{(++)}_{0,1;0,0,1,1,1}\,v_2^\delta
%  \end{split}
%\end{align}
%and,
%\begin{align}
%    \begin{split}
%{\mathcal{K}}_2^{\delta}=&\int_{\ell_{1},\ell_{2}}\frac{\hat\delta(\ell_{1}\cdot v_2)\hat\delta(\ell_{2}\cdot v_1)\,(\ell_{1}+\ell_{2}-q)^{\delta}}{(\ell_{1}\cdot v_1)^2(\ell_{2}\cdot v_2)^2\,\ell_{1}^2\,\ell_{2}^2\,(\ell_{1}+\ell_{2}-q)^2}\Bigg[2(\ell_1\cdot v_1)(\ell_2\cdot v_2)\left(q^2-(\ell_1-q)^2-(\ell_2-q)^2\right)\\ &
%\hspace{0.8cm}-\gamma(\ell_1-q)^2\left(-(\ell_2-q)^2+4(\ell_1\cdot v_1)^2\right)\Bigg]\,,\\ &
% =\Bigg(\boldsymbol{h_{19}}
%      \mathbfcal{M}_{0,0;0,0,1,1,1}
 %     \hspace{0 cm}+ \boldsymbol{h_{20}}\,q^2 \mathbfcal{M}_{0,0;0,0,2,1,1}
 %   \hspace{0 cm}  +\boldsymbol{h_{21}}
%      q^2\,\mathbfcal{M}_{0,0;0,0,1,2,1}
%      \hspace{0 cm}+\boldsymbol{h_{22}}\,q^2\mathbfcal{M}^{(++)}_{1,1;0,0,1,1,1}\Bigg) q^\delta \\ &+\boldsymbol{h_{23}}\,q^2 \mathbfcal{M}^{(++)}_{0,1;0,0,1,1,1}v_1^\delta+\boldsymbol{h_{24}}\,q^2 \mathbfcal{M}^{(++)}_{0,1;0,0,1,1,1}v_1^\delta
%       \end{split}
%\end{align}
%Furthermore we will have two two-index integrals. The details are given below,
%\begin{align}
%\begin{split}
%       \Delta_{(1)}^{\eta\delta}&=\int_{\ell_{1},\ell_{2}}\frac{\hat\delta(\ell_{1}\cdot v_2)\hat\delta(\ell_{2}\cdot v_1)(\ell_{1}+\ell_2-q)^{[\eta}\ell_{1}^{\delta]}}{(\ell_1\cdot v_1)\,(\ell_2\cdot v_2)^2\,\ell_{1}^2 \ell_{2}^2 (\ell_{1}+\ell_{2}-q)^2}(\ell_1-q)^2\,,\\ &
%       =\Bigg(\boldsymbol{h_{25}}\,\mathbfcal{M}_{0,0;0,0,1,1,1}+ \boldsymbol{h_{26}} q^2\mathbfcal{M}_{0,0;0,0,2,1,1}+\boldsymbol{h_{27}}\,q^2\mathbfcal{M}_{0,0;0,0,1,2,1} +\boldsymbol{h_{28}}q^2\mathbfcal{M}^{(++)}_{1,1;0,0,1,1,1}\Bigg) q^{[\eta}v_1^{\delta]}\\ &+\Bigg(\boldsymbol{h_{29}}\mathbfcal{M}_{0,0;0,0,1,1,1}+\boldsymbol{h_{30}}\,q^2\mathbfcal{M}_{0,0;0,0,2,1,1}+\boldsymbol{h_{31}}\,q^2\mathbfcal{M}_{0,0;0,0,2,1,1}+\boldsymbol{h_{32}}\,q^2\mathbfcal{M}^{)(++)}_{1,1;0,0,1,1,1}\Bigg)q^{[\eta}v_2^{\delta]}\,.
%    \end{split}
%\end{align}
%and,
%\begin{align}
 %   \begin{split}
        %\Delta_{(2)}^{\eta\delta}&=\int_{\ell_{1},\ell_{2}}\frac{\hat\delta(\ell_{1}\cdot v_2)\hat\delta(\ell_{2}\cdot v_1)(\ell_{1}+\ell_2-q)^{[\eta}\ell_2^{\delta]}}{(\ell_1\cdot v_1)^2 (\ell_2\cdot v_2)\,\ell_{1}^2 \ell_{2}^2 (\ell_{1}+\ell_{2}-q)^2}(\ell_2-q)^2\,,\\ &
%=\Bigg(\boldsymbol{h_{33}}\,\mathbfcal{M}_{0,0;0,0,1,1,1}+\boldsymbol{{h_{34}}}\,q^2\mathbfcal{M}_{0,0;0,0,2,1,1}+\boldsymbol{h_{35}}\,q^2\mathbfcal{M}_{0,0;0,0,1,2,1}+\boldsymbol{h_{36}}\,q^2\mathbfcal{M}^{(++)}_{1,1;0,0,1,1,1}\Bigg)q^{[\eta}v_1^{\delta]}\\ &+\Bigg(\boldsymbol{h_{37}}\,\mathbfcal{M}_{0,0;0,0,1,1,1}+\boldsymbol{h_{38}}\,q^2\mathbfcal{M}_{0,0;0,0,2,1,1}\,+\boldsymbol{h_{39}},q^2\mathbfcal{M}_{0,0;0,0,1,2,1}\\ &+\boldsymbol{\boldsymbol{h_{40}}}\,q^2\mathbfcal{M}^{(++)}_{1,1;0,0,1,1,1}\Bigg)q^{[\eta}v_2^{\delta]}
%    \end{split}
%\end{align}
%Finally putting everything together we get the eikonal phase at $\mathcal{O}(\mathcal{S}^1)$, \\
%\hfsetfillcolor{gray!10}
%\hfsetbordercolor{black!150}
%\begin{align}
%    \begin{split}
%    \tikzmarkin[disable rounded corners=false]{e7}(3.3,-1)(-0.1,1)
%i\chi\Big|_{\mathcal{S}^1}&=-\frac{i m_1^2 m_2^2 s_1^2 s_2^2\,x}{1536\pi ^3 |\boldsymbol{b}|^3 \left(1-x^2\right)^5 \left(x^2+1\right)\, m_p^6}\Bigg(4 x^2 \big(\left(3 x^6+5 x^4+5 x^2+3\right) \log \left(\pi  |\boldsymbol{b}|^6\right)\\&\hspace{0.8cm}+2 \left(-3 x^6-7 x^4+7 x^2+3\right) \log (x)\big)\\&\hspace{0.8cm}+\left(x^2+1\right) \big(21 x^8+6 x^6 (10 \gamma_E +15-8 \log (2))\\&\hspace{0.8cm}+x^4 (40 \gamma_E -30-32 \log (2))+6 x^2 (10 \gamma_E +15-8 \log (2))\\&\hspace{0.8cm}+21\big)\Bigg)\Big(\hat{b}\cdot\mathcal{S}_1\cdot v_2\Big)\\&-\frac{i m_1^2 m_2^2 s_1^2 s_2^2\,x}{3072  \pi ^3 |\boldsymbol{b}|^3 \left(1-x^2\right)^5\,\,m_p^6}\Bigg(-4 x^2\left(1+x^2\right)^2 \log \left(\pi | \boldsymbol{b}|^6\right)+x^2 \big(-20 \gamma_E  \left(x^2+1\right)^2\\&\hspace{0.8cm}+x^2 \left(x^2 \left(3 x^2+34+4 \log (16)\right)+134+32 \log (2)\right)+34\\&\hspace{0.8cm}+4 \log (16)\big)+40 \left(x^4-1\right) x^2 \log (x)+3\Bigg)\Big(\hat{b}\cdot \mathcal{S}_2\cdot v_1\Big)\,.
%\tikzmarkend{e7}
%\end{split}
%\end{align}
%The Spinless part of the \textcolor{red}{BHAI} diagram:
%\begin{align}
 %   \begin{split}
 %   &  \textcolor{red}{\mathcal{N}\Big|_{\mathcal{S}^0}}=  \frac{\left(2 \ell_1^2+q^2\right) \left(2 \ell_2^2 \left(2 \gamma ^2 (\epsilon -1)+1\right)+8 \gamma  (\epsilon -1) \left( \ell_2 \cdot v_2\right) \left(\gamma  \left( \ell_2 \cdot v_2\right)- \ell_1\cdot v_1\right)+q^2 \left(2 \gamma ^2 (\epsilon -1)+1\right)\right)}{8 (\epsilon -1)}\\ &
 %   \end{split}
%\end{align}
%\begin{align}
%    \begin{split}
%i\chi\Big|_{\mathcal{S}^0}=-\frac{\textcolor{red}{i m_1^2 m_2^2 s_1^2 s_2^2}}{256 m_p^6}\,\int e^{-i q \cdot b}\, \hat\delta(q\cdot v_1)\hat\delta(q\cdot v_2)&\Big(\boldsymbol{h_{42}}\,\,\mathbfcal{M}_{0,0;0,0,1,1,1}+\\&\boldsymbol{h_{43}}\,\,q^2\mathbfcal{M}_{0,0;0,0,2,1,1}+\boldsymbol{h_{44}}\,\,q^2\mathbfcal{M}_{0,0;0,0,1,2,1}+\boldsymbol{h_{45}}\,\,q^2\mathbfcal{M}^{(+-)}_{1,1;0,0,1,1,1}\Big)\,,
%   \end{split}
%\end{align}
%Similarly the spinning part of the diagram has the following contribution to the eikonal;
%\begin{align}
%\begin{split}
%\textcolor{red}{\mathcal{N}\Big|_{\mathcal{S}^1}}&=(\ell_2-q)\cdot \mathcal{S}_2\cdot (q-\ell_1-\ell_2)\Bigg[\frac{i \left(2 \ell_1^2+q^2\right) \left(\ell_2\cdot v_2\right)}{4 (\epsilon -1)}\Bigg]+\\&(\ell_2-q)\cdot \mathcal{S}_2\cdot v_1 \Bigg[\frac{1}{4} i \left(2 \ell_1^2+q^2\right) \Bigg((-2 \ell_1\cdot v_1)\left(\ell_2\cdot v_2\right)+2 \gamma  \ell_2^2+\gamma  q^2\Bigg)\Bigg]+(q-\ell_1-\ell_2)\cdot \mathcal{S}_1\cdot v_2
%\Bigg[\frac{1}{2} i \gamma  \ell_2^2 \left(2 \ell_1^2+q^2\right)\Bigg]
%\end{split}   
%\end{align}
%Correspondingly, in the spinning eikonal phase we encounter  the following vector integrals,
%\begin{align}
%    \begin{split}
 %     \boldsymbol{  \mathscr{H}}_{q.S_2.v_1}&=\frac{-i}%{4}\int\frac{\hat\delta(\ell_1\cdot v_2)\hat\delta(\ell_2\cdot v_1)}{(\ell_1\cdot v_1+i\varepsilon)^2\,(\ell_2\cdot v_2-i\varepsilon)^2\,(\ell_1+\ell_2-q)^2\,(\ell_1-q)^2\,(\ell_2-q)^2}\Bigg[\Bigg]\\&
%      =\frac{-i}{4}\mathfrak{B}_1
%    \end{split}
%\end{align}
%where 
%\begin{align}
%    \begin{split}
%&\mathfrak{B}_1=\boldsymbol{h_{146}}\mathbfcal{M}_{0,0;0,0,1,1,1}+\boldsymbol{h_{147}}\,q^2\mathbfcal{M}_{0,0;0,0,2,1,1}+\boldsymbol{h_{148}}\,q^2\mathbfcal{M}_{0,0;0,0,1,2,1}+\boldsymbol{h_{149}}\,q^2\mathbfcal{M}^{(+-)}_{1,1;0,0,1,1,1}
%  \\ &
 % \hspace{0cm}
%    \end{split}
%\end{align}
%\begin{align}
%\begin{split}
%&h_{146}=\frac{152 \gamma ^2 \epsilon ^3+156 \gamma ^2 \epsilon ^2+4 \left(5 \gamma ^2-1\right) \epsilon +4 \left(\gamma ^2-1\right)}{3 \gamma  \left(\gamma ^2-1\right)^2 \epsilon ^2 (\epsilon +1)}\\&
%h_{147}=\frac{2 \left(114 \gamma ^2-114 \gamma ^4\right) \epsilon ^4+2 \left(155 \gamma ^2-155 \gamma ^4\right) \epsilon ^3+2 \left(-43 \gamma ^4+42 \gamma ^2+1\right) \epsilon ^2+2 \left(\gamma ^4-2 \gamma ^2+1\right) \epsilon}{3 \gamma  \left(\gamma ^2-1\right)^2 \epsilon ^2 (\epsilon +1)}\\&
%h_{148}=\frac{456 \gamma ^2 \epsilon ^4+544 \gamma ^2 \epsilon ^3+2 \left(65 \gamma ^2-2\right) \epsilon ^2+2 \left(7 \gamma ^2-3\right) \epsilon +2 \left(\gamma ^2-1\right)}{3 \gamma  \left(\gamma ^2-1\right)^2 \epsilon ^2 (\epsilon +1)}\\&
%h_{149}=\frac{2 \left(38 \gamma ^3-38 \gamma ^5\right) \epsilon ^4+2 \left(-58 \gamma ^5+63 \gamma ^3-5 \gamma \right) \epsilon ^3+2 \left(-26 \gamma ^5+31 \gamma ^3-5 \gamma \right) \epsilon ^2}{3 \gamma  \left(\gamma ^2-1\right)^2 \epsilon ^2 (\epsilon +1)}
%\end{split}
%\end{align}
%\begin{align}
 %   \begin{split}
 %     \boldsymbol{  \mathscr{H}}_{q.S_1.v_2}&=\frac{i\gamma}%{2}\int\frac{\hat\delta(\ell_1\cdot v_2)\hat\delta(\ell_2\cdot v_1)}{(\ell_1\cdot v_1+i\varepsilon)^2\,(\ell_2\cdot v_2-i\varepsilon)^2\,(\ell_1+\ell_2-q)^2\,(\ell_1-q)^2\,(\ell_2-q)^2}\Bigg[\Bigg]\\&
 %     =\frac{i\gamma}{2}\mathfrak{B}_2
%    \end{split}
%\end{align}
%where 
%\begin{align}
 %   \begin{split}
%&\mathfrak{B}_2=\boldsymbol{h_{150}}\mathbfcal{M}_{0,0;0,0,1,1,1}+\boldsymbol{h_{151}}\,q^2\mathbfcal{M}_{0,0;0,0,2,1,1}+\boldsymbol{h_{152}}\,q^2\mathbfcal{M}_{0,0;0,0,1,2,1}+\boldsymbol{h_{153}}\,q^2\mathbfcal{M}^{(+-)}_{1,1;0,0,1,1,1}
 % \\ &
%  \hspace{0cm}
%    \end{split}
%\end{align}
%\begin{align}
%\begin{split}
%&h_{150}=\frac{8 \epsilon  (5 \epsilon +1)}{\gamma ^2-1},\,
%h_{151}=\frac{60 \epsilon +32}{3-3 \gamma ^2},\,\,
%h_{152}=\frac{4 (5 \epsilon +1) (6 \epsilon +1)}{3 \left(\gamma ^2-1\right)^2 \epsilon },\,\,
%h_{153}=\frac{2 \gamma  (10 \epsilon +9)}{3 \left(\gamma ^2-1\right)}
%\end{split}
%\end{align}
%\begin{align}
 %   \begin{split}
 %     \boldsymbol{  \mathscr{H}}_{q.S_2.l_1}^{\gamma}&=\frac{i}{4(\epsilon-1)}\int\frac{\hat\delta(\ell_1\cdot v_2)\hat\delta(\ell_2\cdot v_1)(\ell_1)^{\gamma}(2\ell_1^2+q^2)(\ell_2\cdot v_2)}{(\ell_1\cdot v_1+i\varepsilon)^2\,(\ell_2\cdot v_2-i\varepsilon)^2\,(\ell_1+\ell_2-q)^2\,(\ell_1-q)^2\,(\ell_2-q)^2}\\&
 %     =a_1 (v_1^{\gamma}-\gamma v_2^\gamma)
%    \end{split}
%\end{align}
%where 
%\begin{align}
 %   \begin{split}
%a_1&=\boldsymbol{h_{46}} \mathbfcal{M}_{0,0;0,0,1,1,1}+\boldsymbol{h_{47}}\,q^2\mathbfcal{M}_{0,0;0,0,2,1,1}+\boldsymbol{h_{48}}\,q^2\mathbfcal{M}_{0,0;0,0,1,2,1}+\boldsymbol{h_{49}}\,q^2\mathbfcal{M}^{(+-)}_{1,1;0,0,1,1,1}
%  \\ &
%  \hspace{0cm}
%    \end{split}
%\end{align}
%\begin{align}
%    \begin{split}
 %     \boldsymbol{  \mathscr{H}}_{\ell_2.S_2.v_1}^{\delta}&=\frac{i}{4}\int\frac{\hat\delta(\ell_1\cdot v_2)\hat\delta(\ell_2\cdot v_1)(\ell_2)^{\delta}\Bigg(-2(\ell_1\cdot v_1)(\ell_2\cdot v_2)+\gamma\left(2 \ell_2^2+q^2\right) \Bigg)(2\ell_1^2+q^2)\left(\ell_2\cdot v_2\right)}{(\ell_1\cdot v_1+i\varepsilon)^2\,(\ell_2\cdot v_2-i\varepsilon)^2\,(\ell_1+\ell_2-q)^2\,(\ell_1-q)^2\,(\ell_2-q)^2}\\&
%      =\left(\hat{a}_3 q^\delta\right)
%    \end{split}
%\end{align}
%where,
%\begin{align}
%    \begin{split}
%&%\hat{a}_2=\Big(\boldsymbol{K_{42}}q^2\,\,\mathbfcal{M}_{0,0;0,0,1,1,1}+\\&\boldsymbol{K_{43}}\,\,q^4\mathbfcal{M}_{0,0;0,0,2,1,1}+\boldsymbol{K_{44}}\,\,q^4\mathbfcal{M}_{0,0;0,0,1,2,1}+\boldsymbol{K_{45}}\,\,q^4\mathbfcal{M}^{(+-)}_{1,1;0,0,1,1,1}\Big),\,\,\\&
%\hat{a}_3=\Big(\boldsymbol{K_{46}}\,\,\mathbfcal{M}_{0,0;0,0,1,1,1}+\\&\boldsymbol{K_{47}}\,\,q^2\mathbfcal{M}_{0,0;0,0,2,1,1}+\boldsymbol{K_{48}}\,\,q^2\mathbfcal{M}_{0,0;0,0,1,2,1}+\boldsymbol{K_{49}}\,\,q^2\mathbfcal{M}^{(+-)}_{1,1;0,0,1,1,1}\Big)
%    \end{split}
%\end{align}
%\begin{align}
%    \begin{split}
%      \boldsymbol{  \mathscr{H}}_{\ell_1.S_1.v_2}^{\gamma}&=\frac{-i\gamma}{2}\int\frac{\hat\delta(\ell_1\cdot v_2)\hat\delta(\ell_2\cdot v_1)\ell_1^{\gamma} \Bigg(2\ell_1^2+q^2)(\ell_2\cdot v_2)\Bigg)\ell_2^2}{(\ell_1\cdot v_1+i\varepsilon)^2\,(\ell_2\cdot v_2-i\varepsilon)^2\,(\ell_1+\ell_2-q)^2\,(\ell_1-q)^2\,(\ell_2-q)^2}\\&
%      =-\left(\hat{a}_3 q^\gamma\right)
%    \end{split}
%\end{align}
%where 
%\begin{align}
%    \begin{split}
%&\hat{a}_1=\frac{i \left(15 \gamma ^2+1\right) \epsilon +i \left(1-5 \gamma ^2\right)}{12 (\gamma -1)^2 (\gamma +1)^2 \epsilon ^2 (3 \epsilon -1)} q^2 \mathbfcal{M}_{ 0, 0; 0, 0, 1, 1, 1}\\&+\frac{-i \left(\gamma ^2-1\right) \epsilon ^2-i \left(\gamma ^2-1\right) \epsilon}{12 (\gamma -1)^2 (\gamma +1)^2 \epsilon ^2 (3 \epsilon -1)}q^4\mathbfcal{M}_{ 0, 0; 0, 0, 2, 1, 1}+\frac{-2 i \epsilon ^2-3 i \epsilon -i}{12 (\gamma -1)^2 (\gamma +1)^2 \epsilon ^2 (3 \epsilon -1)}q^4\mathbfcal{M}_{ 0, 0; 0, 0, 1, 2, 1}\\&+\frac{2 i \left(\gamma ^3-\gamma \right) \epsilon ^2-10 i \left(\gamma ^3-\gamma \right) \epsilon ^3}{12 (\gamma -1)^2 (\gamma +1)^2 \epsilon ^2 (3 \epsilon -1)}q^4\mathbfcal{M}_{ 1, 1; 0, 0, 1, 1, 1}\\&
%\hat{a}_3 = \frac{-2 i \gamma  \left(40 (\gamma -2) \epsilon ^2+2 (14 \gamma +5) \epsilon -9 \gamma \right)}{6 (\gamma -1) (\gamma +1)^2 (6 \epsilon -1)}q^2\mathbfcal{M}_{ 0, 1; 0, 0, 1, 1, 1} 
 %   \end{split}
%\end{align}
%\begin{align}
%    \begin{split}
%      \boldsymbol{  \mathscr{H}}_{-\ell_2.S_1.v_2}^{\delta}&=\frac{i\gamma}{2}\int\frac{\hat\delta(\ell_1\cdot v_2)\hat\delta(\ell_2\cdot v_1)\ell_2^{\gamma} \Bigg((2\ell_1^2+q^2)(\ell_2\cdot v_2)\Bigg)\ell_2^2}{(\ell_1\cdot v_1+i\varepsilon)^2\,(\ell_2\cdot v_2-i\varepsilon)^2\,(\ell_1+\ell_2-q)^2\,(\ell_1-q)^2\,(\ell_2-q)^2}\\&
%      =\frac{1}{2}\left(\hat{A}_3q^\gamma\right)
%    \end{split}
%\end{align}
%\begin{align}
%    \begin{split}
%\hat{A}_3 = \frac{9 i \gamma ^2-80 i (\gamma -2) \gamma  \epsilon ^2-6 i (9 \gamma +4) \gamma  \epsilon }{6 (\gamma -1) (\gamma +1)^2 (2 \epsilon -1) (6 \epsilon -1)}q^2\mathbfcal{M}_{ 0, 1; 0, 0, 1, 1, 1} 
%    \end{split}
%\end{align}
%\textcolor{red}{Two index object}:
%\begin{align}
 %   \begin{split}
 %     \boldsymbol{  \mathscr{H}}^{\eta\delta}_{-\ell_2\cdot \mathcal{S}_2\cdot\ell_1}&=\frac{i}{4(\epsilon-1)}\int\frac{\hat\delta(\ell_1\cdot v_2)\hat\delta(\ell_2\cdot v_1)\ell_2^{\delta}\ell_1^{\eta}\left(\ell_2\cdot v_2\right)(2\ell_1^2+q^2)}{(\ell_1\cdot v_1+i\varepsilon)^2\,(\ell_2\cdot v_2-i\varepsilon)^2\,(\ell_1+\ell_2-q)^2\,(\ell_1-q)^2\,(\ell_2-q)^2}\\&
%      =\Bigg(\tilde{a}_1 q^{[\eta}v_1^{\delta]}+\tilde{a}_2 q^{[\eta}v_2^{\delta]}\Bigg)
%    \end{split}
%\end{align}
%where 
%\begin{align}
%    \begin{split}
%&\tilde{a}_1=\tilde h_{61}\,\mathbfcal{M}_{0,0;0,0,1,1,1}+\tilde h_{62}\,q^2\mathbfcal{M}_{0, 0; 0, 0, 1, 2, 1, }+\tilde h_{63}\,q^2\mathbfcal{M}_{0, 0; 0, 0, 2, 1, 1, }+\tilde h_{64}\,q^2\mathbfcal{M}^{(+-)}_{1,1;0,0,1,1,1}\\&
%\tilde{a}_2=\tilde h_{65}\,\mathbfcal{M}_{0, 0, 0, 0, 1, 1, 1 }+\tilde h_{66}\,q^2\mathbfcal{M}_{0, 0; 0, 0, 1, 2, 1}+\tilde h_{67}\,q^2\mathbfcal{M}_{0, 0; 0, 0, 2, 1, 1}+\tilde h_{68}\,q^2\mathbfcal{M}^{(+-)}_{1, 1; 0, 0, 1, 1, 1}
%    \end{split}
%\end{align}
%\newpage
\textbullet $\,\,$ Now we consider the following diagram. It comes from the 3-point interaction vertex between the graviton and the kinetic term of the scalar field. 
\begin{align}
    \begin{split}
         \begin{minipage}[h]{0.08\linewidth}
	\vspace{4pt}
	\scalebox{1.65}{\includegraphics[width=\linewidth]{3PM_3.png}}
 \end{minipage}&
    \end{split}
\end{align}
Proceeding as before, we get the following,
%\begin{align}
 %   \begin{split}
 %      i \chi=i\frac{s_1 s_2^3m_1^2 m_2^2}{32 m_p^6} \int_{k_i}e^{-i (k_4-k_1)\cdot b}\frac{\hat\delta(k_1\cdot v_1)\hat\delta(k_2\cdot v_2)\hat\delta(k_4\cdot v_1)\hat\delta(k_4\cdot v_2-k_1\cdot v_2)}{(k_4\cdot v_2+i\varepsilon)^2\,k_1^2\, k_2^2 \,(k_1+k_2)^2\, k_4^2}\mathcal{N}
%    \end{split}
%\end{align}
%where the numerator takes the form,
%\begin{align}
%    \mathcal{N}=\frac{1}{2}(k_1+k_2)\cdot k_4\,\Bigg(-2(k_2\cdot v_1)^2-2i\, k_1\cdot \mathcal{S}_1\cdot k_2 (k_2\cdot v_1)+k_2^2\Bigg)\,.\label{3.34}
%\end{align}
%Now redefining the loop momenta as $k_4-k_1\to q,\,k_2 \to \ell_{1},\, k_1\to \ell_{2}$ we get,
%\begin{align}
 %   \begin{split}
 %  i     \chi=i\frac{s_1 s_2^3m_1^2 m_2^2}{32 m_p^6}\int e^{-iq\cdot b}\hat\delta(q\cdot v_1)\hat\delta(q\cdot v_2)\int_{\ell_{1},\ell_{2}}\frac{\hat\delta(\ell_{1}\cdot v_2)\hat\delta(\ell_{2}\cdot v_1)}{(\ell_{2}\cdot v_2+i\varepsilon)^2 \ell_{1}^2\,\ell_{2}^2\,(\ell_1+\ell_2-q)^2\,(\ell_{2}-q)^2 }\mathcal{N}\,.\label{4.35}
%    \end{split}
%\end{align}
%Separating the numerator in terms of spin and spinless part, we get
%\begin{align}
%    \begin{split}
%&\mathcal{N}\Big|_{\mathcal{S}^0}=\frac{1}{2}(\ell_1\cdot v_1)^2\,\Big[(\ell_1-q)^2-(\ell_1+\ell_2-q)^2-\ell_2^2\Big]-\frac{1}{2(d-2)}\ell_1^2\,(\ell_1-q)^2\,,\\ &
%\mathcal{N}\Big|_{\mathcal{S}^1}=\frac{i}{2}(\ell_2-q)\cdot \mathcal{S}_1\cdot \ell_1\,(\ell_1\cdot v_1)\Big[(\ell_1-q)^2-(\ell_1+\ell_2-q)^2-\ell_2^2\Big] \,.
%    \end{split}
%\end{align}
%Now to write the denominator of \eqref{4.35} in terms of the basis to carry out the IBP reduction,  we again make a field redefinition $\ell_{2}\to \ell_{2}-q\,.$ After that, ignoring the tadpoles, we get the following for the spinless part of the eikonal phase,\\
%\hfsetfillcolor{gray!10}
%\hfsetbordercolor{black!150}
\begin{align}
    \begin{split}
i\chi\Big|_{\mathcal{S}^{0}}&=i\frac{s_1 s_2^3m_1^2 m_2^2}{128 m_p^6}\int e^{-iq\cdot b} \,\hat\delta(q\cdot v_1)\hat\delta(q\cdot v_2)\int_{\ell_{1},\ell_{2}}\left(\frac{1}{(\ell_2\cdot v_2+i\varepsilon)^2}+\frac{1}{(\ell_2\cdot v_2-i\varepsilon)^2}\right)\frac{\hat\delta(\ell_{1}\cdot v_2)\hat\delta(\ell_{2}\cdot v_1)\,}{\ell_{1}^2 \ell_{2}^2 (\ell_{1}+\ell_{2}-q)^2(\ell_{2}-q)^2}\\ &
         \hspace{0.8cm}\Bigg((\ell_{1}\cdot v_1)^2\Big[(q-\ell_{1})^2-\ell_2^2-(\ell_1+\ell_2-q)^2\Big]-\frac{1}{d-2}\ell_1^2(\ell_1-q)^2\Bigg)\,,\\ &
        =i\frac{s_1 s_2^3m_1^2 m_2^2}{64 m_p^6}\int e^{-iq\cdot b} \,\hat\delta(q\cdot v_1)\hat\delta(q\cdot v_2)\Bigg[\boldsymbol{\tilde h}_{69}\,\mathbfcal{M}_{0,0;0,0,1,1,1}+\boldsymbol{\tilde h}_{70}\,\,q^2
      \mathbfcal{M}_{0,0;0,0,2,1,1} +\boldsymbol{\tilde h}_{71}\,\,q^2\mathbfcal{M}_{0,0;0,0,1,2,1}\Bigg]\,.
         \end{split}
\end{align}
    %  -\frac{-4 \gamma ^4 (\epsilon -2) (\epsilon -1) \epsilon +2 \gamma ^2 (\epsilon -1) ((2 \epsilon -9) \epsilon +1)+2 (\epsilon -3) \epsilon +1}{2 (\gamma^2 -1) (\epsilon -1) \epsilon  (2 \epsilon 
      %-3)
    %  \mathbfcal{M}_{0,0;0,1,1,0,1}\Bigg],
    %  \\&
  % \,.
    %  \tikzmarkend{e8}
Then using \eqref{MasterM}. we get, 
\begin{align}
    \begin{split}  \label{eqt9}      \chi\Big|_{\mathcal{S}^{0}}&=\frac{m_1^2m_2^2}{m_p^6}\Delta_{\mathfrak{15}}^{\textrm{poly}}\,,\,\, \Delta_{\mathfrak{15}}^{\textrm{poly}}=\Bigg(\frac{s_{1}s_{2}^3 \left(x^2+1\right)}{4096 \pi ^3 |\boldsymbol{b}|^2 \left(x^2-1\right)}\Bigg)\,.
         \end{split}
\end{align}
%Surprisingly, in this diagram, at order $\mathcal{O}(\mathcal{S}^0)$, terms consist of MPls cancels out and we left with a very simple dependency on kinematic variable $x$.
%\textcolor{red}{Final $q$ integral.}\\
Another term that appears in the spinless part from a term proportional to $\ell_{1}^2$ in the numerator. In that case, the integral will be proportional $\mathbfcal{M}_{0,0;0,1,0,1,1}$ which vanishes identically.\par 
Now, the spinning part of the eikonal phase is given by,
\begin{align}
    \begin{split}
       \chi\Big|_{\mathcal{S}^1}=&-\frac{s_1 s_2^3m_1^2 m_2^2}{128 m_p^6}\, S_{1}^{\eta\delta}\int e^{-iq\cdot b} \,\hat\delta(q\cdot v_1)\hat\delta(q\cdot v_2)\int_{\ell_{1},\ell_{2}}\sum_{(\pm)}\frac{\hat\delta(\ell_{1}\cdot v_2)\hat\delta(\ell_{2}\cdot v_1)\,(\ell_{1}\cdot v_1) \,(\ell_2-q)_{[\eta} \, \ell_{1\delta]}}{(\ell_2\cdot v_2\pm i\varepsilon)^2\ell_{1}^2 \,\ell_{2}^2\,(\ell_{1}+\ell_{2}-q)^2}\,\Big[(\ell_{1}-q)^2\\ &
        \hspace{0.8cm}-(\ell_1+\ell_2-q)^2-\ell_2^2\Big]\,,\\&
        =-\frac{s_1 s_2^3m_1^2 m_2^2}{64 m_p^6}\int e^{-iq\cdot b} \,\hat\delta(q\cdot v_1)\hat\delta(q\cdot v_2)\Bigg[i \Bigg(\boldsymbol{\tilde h}_{72}\,\mathbfcal{M}_{0,0;0,0,1,1,1}+\boldsymbol{\tilde h}_{73}\,\,q^2
      \mathbfcal{M}_{0,0;0,0,2,1,1}\\& \hspace{0.8cm}+\boldsymbol{\tilde h}_{74}\,\,q^2\mathbfcal{M}_{0,0;0,0,1,2,1}+\boldsymbol{\tilde h}_{75}\,\,q^2\mathbfcal{M}^{(++)}_{1,1;0,0,1,1,1}
      \Bigg)\big(q\cdot\mathcal{S}_1\cdot v_2\big)\Bigg]\,,\\&=\frac{m_1^2m_2^2}{m_p^6}\Bigg[\Delta_{\mathfrak{16}}^{\textrm{poly}}+\Delta_{\mathfrak{08}}^{\textrm{log}}\log(x)\Bigg],\label{eqt10}      
    \end{split}
\end{align}
where,
\begin{align}
    \begin{split}
    &\Delta_{\mathfrak{16}}^{\textrm{poly}}=\Bigg(\frac{s_1 s_2^3}{131072 \pi ^3 |\boldsymbol{b}|^3 \left(x^2-1\right)^6 \left(x^2+1\right)}\Bigg)\Bigg[x\bigg\{6 \left(65 x^{12}-387 x^8+387 x^4-65\right) \log (|\boldsymbol{b}|)\\ & \hspace{0.8cm}+\left(x^4-1\right) \big(x^8 (86-260 \log (2)+65 \log (\pi ))\\& \hspace{0.8cm}+1712 x^6+2 x^4 (1466+644 \log (2)-161 \log (\pi ))\\&\hspace{0.8cm}+1712 x^2+5 \gamma_E  \left(65 x^8-322 x^4+65\right)+86-260 \log (2)\\& \hspace{0.8cm}+65 \log (\pi )\bigg\}\Bigg], \\&\,\Delta_{\mathfrak{08}}^{\textrm{log}}=\Bigg(\frac{s_1 s_2^3}{131072 \pi ^3 |\boldsymbol{b}|^3 \left(x^2-1\right)^6 \left(x^2+1\right)}\Bigg)\Bigg[8 x\Bigg\{20 x^{12}+239 x^{10}+700 x^8+1154 x^6+700 x^4+239 x^2+20\Bigg\}\Bigg]\,.
    \end{split}
\end{align}
%The integral in \eqref{4.38} can be evaluated by using PV reduction. Noting that the tensor integral should be antisymmetric in $\eta\leftrightarrow \delta$ we first evaluate the following integral,
%\begin{align}
%    \begin{split}
%\theta_{[\eta\delta]}&:=\int_{\ell_{1},\ell_{2}}\frac{\hat\delta(\ell_{1}\cdot v_2)\hat\delta(\ell_{2}\cdot v_1)\,(\ell_{1}\cdot v_1) \,(\ell_2-q)_{[\eta} \, \ell_{1\delta]}}{(\ell_{2}\cdot v_2)^2\,\ell_{1}^2 \,\ell_{2}^2\,(\ell_{1}+\ell_{2}-q)^2(\ell_{2}-q)^2}\,\Big[(\ell_{1}-q)^2-(\ell_1+\ell_2-q)^2-\ell_2^2\Big]\,,\\ &
 %      =j_1\,q_{\eta}v_{1\delta}+j_2\,q_{\eta}v_{2\delta}
 %   \end{split}
%\end{align}
%\newpage
%where the scalar integrals $j_i$ s takes the form,
%\begin{align}
%    \begin{split}
 %   & j_1=\boldsymbol{h_{54}}\,\mathbfcal{M}_{0,0;0,0,1,1,1}+\boldsymbol{h_{55}}\,\mathbfcal{M}_{0,0;0,0,2,1,1}\,q^2\,
%   + \boldsymbol{h_{56}}\,q^2\mathbfcal{M}_{0,0;0,0,1,2,1}\,,\\ &
%  j_2=\boldsymbol{h_{57}}\mathbfcal{M}_{0,0;0,0,1,1,1}
%  +\boldsymbol{h_{58}}\,q^2\mathbfcal{M}_{0,0;0,0,2,1,1} +\boldsymbol{h_{59}}\,\,q^2\mathbfcal{M}_{0,0;0,0,1,2,1}
  %+\frac{1}{4 (\gamma^2 -1) (\epsilon -1) \epsilon  (2 \epsilon -3) (4 \epsilon -3)}\Big[\gamma  \Big(4 \left(\gamma ^4+2 \gamma ^2-3\right) \epsilon ^4-4 \left(4 \gamma ^4+\gamma ^2-13\right) \epsilon ^3+2 \left(10 \gamma ^4-7 \gamma ^2-27\right) \epsilon ^2\\ &\hspace{0.5 cm}+\left(-8 \gamma ^4+10 \gamma ^2+20\right) \epsilon -3\Big)\Big]\,\mathbfcal{M}_{0,0;0,1,1,0,1}\,.
 %   \end{split}
%\end{align}
%Then we get, 
%\hfsetfillcolor{gray!10}
%\hfsetbordercolor{black!150}
%\begin{align}
%    \begin{split}
%         \tikzmarkin[disable rounded corners=false]{e65}(0.2,-1)(-0.1,1)i\chi\Big|_{\mathcal{S}^1}
%        \tikzmarkend{e65}
%         \end{split}
%\end{align} 
\par
%\textcolor{red}{Final $q$ integral.}\\
\textbullet $\,\,$ The next diagrams mentioned below come from the two scalar-graviton interaction vertices (and graviton three vertex) and have $\mathbb{H}$-type topological structure. This diagram will only contribute to the spinless part of the 3PM eikonal phase. For $\mathbb{H}$-type diagram, the worldline propagator will not appear in the amplitude, and we get five independent propagators: $\ell_{1}^2,\,\ell_{2}^2,\,(\ell_{1}+\ell_{2}-q)^2,\,(\ell_{1}-q)^2,\,(q-\ell_{2})^2$. The diagrams are in the following form,
\begin{center}
\includegraphics[width=0.38\linewidth]{H_f.png}
\end{center}
From the first diagram, the contribution to the eikonal phase is given by,
%\begin{align}
%    \begin{split}
 %     i  \chi =i \frac{s_1^2 s_2^2 m_1^2 m_2^2}{64 m_p^6} \int e^{-iq\cdot b}\hat\delta(q\cdot v_1)\hat\delta(q\cdot v_2)\int_{\ell_{1},\ell_{2}}\frac{\hat\delta(\ell_{1}\cdot v_2)\hat\delta(\ell_{2}\cdot v_1)}{\ell_{1}^2\,\ell_{2}^2\,(\ell_{1}+\ell_{2}-q)^2\,(\ell_{1}-q)^2\,(q-\ell_{2})^2}\mathcal{N}
%    \end{split}
%\end{align}
%where the numerator (ignoring the tadpoles) takes the form,
%\begin{align}
 %   \begin{split}
%        \mathcal{N}=\,.
%    \end{split}
%\end{align}
%Therefore, using IBP reduction we get,\\
%\hfsetfillcolor{gray!10}
%\hfsetbordercolor{black!150}
\begin{align}
    \begin{split}
    % \tikzmarkin[disable rounded corners=false]{e10}(7,-1)(-0.1,0.7)  
     i   \chi\Big|_{\mathcal{S}^0}
     %= i \frac{s_1^2 s_2^2 m_1^2 m_2^2}{256 m_p^6} \int e^{-iq\cdot b}\hat\delta(q\cdot v_1)\hat\delta(q\cdot v_2)\Bigg(\int_{\ell_{1},\ell_{2}}\frac{\hat\delta(\ell_{1}\cdot v_2)\hat\delta(\ell_{2}\cdot v_1)}{\ell_{1}^2 \ell_{2}^2 (\ell_{1}+\ell_{2}-q)^2}-q^2\int_{\ell_{1},\ell_{2}}\frac{\hat\delta(\ell_{1}\cdot v_2)\hat\delta(\ell_{2}\cdot v_1)}{\ell_{1}^2\,\ell_{2}^2\,(\ell_{1}-q)^2(\ell_{2}-q)^2}\Bigg)\,,\\ &
        =i \frac{s_1^2 s_2^2 m_1^2 m_2^2}{64 m_p^6} \int e^{-iq\cdot b}\hat\delta(q\cdot v_1)\hat\delta(q\cdot v_2)\,\Bigg[&\frac{\epsilon }{4 (\epsilon -1)}\mathbfcal{M}_{0,0;0,0,1,1,1}-\frac{3 q^2}{8 \epsilon }\mathbfcal{M}_{0,0;0,0,2,1,1}\\&+\frac{q^2 (3 \epsilon -4)}{16 (\epsilon -1)}\mathbfcal{M}_{0,0;1,1,0,1,1}+\frac{q^4}{4}\mathbfcal{M}_{0,0;1,1,1,1,1}\Bigg]\,,
       % \tikzmarkend{e10}
    \end{split}
\end{align}
Then using \eqref{MasterM} we get,

\begin{align}
    \begin{split} \label{eqt11}   
     \chi\Big|_{\mathcal{S}^0}=-\Big(\frac{ m_1^2 m_2^2}{64 m_p^6}\Big)\Delta_{\mathfrak{17}}^{\textrm{poly}}\,, \Delta_{\mathfrak{17}}^{\textrm{poly}}=-\frac{5 x^2s_1^2 s_2^2 \log (x)}{32 \pi ^3 |\boldsymbol{b}|^2 \left(x^2-1\right)^2}\,.
        \end{split}
\end{align}
%\newpage
%\textcolor{red}{Final $q$ integral.}\\
%Now we focus on the computation of eikonal phase at order of $\mathcal{S}^2$. Start with the term:
%\begin{align}
   % \begin{split}
%\mathcal{N}\Big|_{\mathcal{S}^2}^{(1)}&= -(\ell_{1}\cdot v_1)(\ell_{2}\cdot v_2)\Bigg[ \mathcal{S}_{1\,\alpha\beta}\mathcal{S}_{2\,\mu\nu}  - \mathcal{S}_{1\,\mu\nu} \mathcal{S}_{2\,\alpha\beta} \Bigg](q-\ell_{1})^{\mu} \ell_{2}^{\nu}(q-\ell_{1}-\ell_{2})^{\alpha}\ell_{1}^{\beta}\\ &
%=-(\ell_{1}\cdot v_1)(\ell_{2}\cdot v_2)\Bigg[ \underbrace{\mathcal{S}_{1\,\alpha\beta}\mathcal{S}_{2\,\mu\nu}  - \mathcal{S}_{1\,\mu\nu} \mathcal{S}_{2\,\alpha\beta}}_{T_{\alpha\beta;\mu\nu}} \Bigg]\,\Big(q^\mu q^\alpha \ell_{2}^\nu \ell_{1}^\beta-q^\mu \ell_{2}^\nu \ell_{2}^\alpha \ell_{1}^\beta-q^\alpha \ell_{1}^\mu \ell_{1}^\beta \ell_{2}^\nu\Big)
  %  \end{split}
%\end{align}
%We start with the three index object,
%\begin{align}
   % \begin{split}
%\chi\Big|_{\mathcal{S}^2}^{1,3-         \textrm{ind}}&\sim T_{\alpha\beta;\mu\nu} \int e^{i q\cdot b}\,\hat\delta(q\cdot v_1)\hat\delta(q\cdot v_2) \,q^\mu\int_{\ell_{1}}\frac{\hat\delta(\ell_{1}\cdot v_2)\,\ell_{1}^\beta}{(\ell_{1}\cdot v_1)\,\ell_{1}^2} \int_{\ell_{2}}\frac{\hat\delta(\ell_{2}\cdot v_1)\ell_{2}^{\nu}\ell_{2}^\alpha}{(\ell_{2}\cdot v_2)\,\ell_{2}^2 (\ell_{2}+\ell_{1}-q)^2}\label{3.71}
  %  \end{split}
%\end{align}
%In \eqref{3.71} we first do the tensor integral reduction of the $\ell_{2}$ integral.
%\begin{align}
  %  \begin{split}
        %\Lambda^{\nu\alpha}&:=\int_{\ell_{2}}\frac{\hat\delta(\ell_{2}\cdot v_1)\ell_{2}^{\nu}\ell_{2}^\alpha}{(\ell_{2}\cdot v_2)\,\ell_{2}^2 (\ell_{2}+\ell_{1}-q)^2}\\ &
        %=C_{00}\eta^{\nu\alpha}+\sum_{i,j=1}^4 C_{ij}\, p_{i}^{(\nu} p_{i}^{\alpha)},\,\, p_i\in(q,\ell_{1},v_1,v_2)
   % \end{split}
%\end{align}
%One should note that after integrating the result has divergent at the near static limit $x_0\to 1$, so on needs to properly regularize the integral to get rid of the divergent piece. The integral in 
%where, $A$ is some number that need to be fixed to cancel the divergences. For example the  at $\mathcal{O}(\epsilon)$,
%\begin{align}
  %  \begin{split}
       % \mathbb{B}(x,x_0)|_{\epsilon}&=\epsilon\left(
%\begin{array}{cccc}
 %\frac{9}{4} \left(3 \log \left(1-x^2\right)-2 \log (x)\right) & \frac{1}{8} \left(5 \log \left(1-x^2\right)-6 \log (x)\right) & -3 \log \left(x^2-1\right) & 0 \\
 %15 \log (x)-\frac{21}{2} \log \left(1-x^2\right) & \frac{5 \log (x)}{2}-\frac{11}{4} \log \left(1-x^2\right) & 6 \log \left(x^2-1\right) & 0 \\
 %6 \log (x) & -\frac{\log (x)}{3} & -2 \log (x) & 0 \\
% 12 \log (x) & 10 \log (x) & -12 \log (x) & 0 \\
%\end{array}
%\right)
%    \end{split}
%\end{align}
%Note that after taking all the suitable limits $\tilde{\mathbb{B}}(x,x_0)$ should be independent of $\epsilon$. Now first expanding $\tilde{\mathbb{B}}(x,x_0)$ around $\epsilon=0$ (upto $\mathcal{O}(\epsilon)$) we get,
 The second one with $\mathbb{H}$-type topology may appear at 3PM where we have one scalar-scalar-graviton and one graviton 3-point vertex. This diagram contributes to the spinless as well as the spinning part of the eikonal phase because of the.
The graviton 3-point vertex takes the form,
\begin{align}
    S_{EH}\Big|_{h^3}=\int d^4x \,\,\,\mathcal{U}^{\mu\nu\,\,\alpha\beta\rho\,\,\gamma\delta\sigma}\,h_{\mu\nu}\,\partial_{\rho}h_{\alpha\beta}\partial_{\sigma}h_{\gamma\delta}\,.
\end{align}
In the Fourier space, the vertex factor takes the following form,
\begin{align}
    \begin{split}
        \mathcal{V}(k_1,k_2,k_3)=-\hat\delta^{(4)}(k_1+k_2+k_3)\,\mathcal{U}^{\mu\nu\,\,\alpha\beta\rho\,\,\gamma\delta\sigma}\,k_{2\rho}k_{3\sigma}h_{\mu\nu}(-k_1)\,h_{\alpha\beta}(-k_2)\,h_{\gamma\delta}(-k_3)\,.
    \end{split}
\end{align}
Correspondingly, the contribution to the eikonal phase is given by,
\begin{align}
    \begin{split}
        i\chi=i\frac{m_1^2 m_2^2 s_1 s_2}{64 m_p^6}\int e^{-i q\cdot b}\hat\delta(q\cdot v_1)\hat\delta(q\cdot v_2)\int_{\ell_{1},\ell_{2}}\frac{\hat\delta(\ell_{1}\cdot v_2)\hat\delta(\ell_{2}\cdot v_1)}{\ell_{1}^2 \ell_{2}^2 (\ell_{1}+\ell_{2}-q)^2 (\ell_{1}-q)^2 (\ell_{2}-q)^2}\mathcal{N}\label{4.68j}
    \end{split}
\end{align}
where the numerator $\mathcal{N}$ has the following form,
\begin{align}
    \begin{split}
        \mathcal{N}= \Gamma_{(3)}^{\mu\nu\,\alpha\beta\,\gamma\delta}(k_2,k_4,k)\,k_1^a k_3^b,P_{a b;\mu\nu}P_{\alpha\beta,c d}P_{\gamma\delta;e f}(v_1^{c}v_1^d+i(k_2\cdot \mathcal{S}_1))^{(c}v_1^{d)}\,\,(v_2^{e}v_2^f+i(k_4\cdot \mathcal{S}_2)^{(e}v_2^{f)}),
    \end{split}
\end{align}
where the 3-point vertex function $\Gamma_{(3)}$ takes the form,
\begin{align}
    \begin{split}
    \Gamma^{\mu\nu\,\alpha\beta\,\gamma\delta}=&\textrm{sym}\Big[\frac{1}{2}k_{2}^\mu k_{4}^\nu\eta^{\alpha\beta}\eta^{\gamma\delta}-\frac{1}{4}k_2\cdot k_4\,\eta^{\mu\nu}\eta^{\alpha\beta}\eta^{\gamma\delta}+k_2\cdot k_4 \eta^{\alpha\nu}\eta^{\beta\mu}\eta^{\gamma\delta}-k_2^\delta k_4^\gamma \eta^{\alpha\mu}\eta^{\beta\nu}\\ &
        +\frac{1}{4}(k_2\cdot k_4)\eta^{\mu\nu}\eta^{\alpha\gamma}\eta^{\beta\delta}-k_2^\nu\,k_4^\beta \eta^{\gamma\delta}\eta^{\mu\alpha}-k_{2}^\mu k_4^\nu \eta^{\alpha\delta}\eta^{\beta\gamma}+\frac{1}{2}\,k_2^\alpha\,k_4^\beta\,\eta^{\mu\nu}\eta^{\gamma\delta}\\ &
        -k_2\cdot k_4\,\eta^{\alpha\mu}\eta^{\beta\delta}\eta^{\nu\gamma}-\frac{1}{2}k_2^\gamma\,k_4^\beta\,\eta^{\mu\nu}\eta^{\alpha\delta}+k_2^\gamma k_4^\alpha\,\eta^{\mu\beta}\eta^{\nu\delta}-\frac{1}{2}\,k_2^\mu k_4^\nu \eta^{\alpha\delta}\eta^{\beta \gamma}+2 k_{2}^\nu k_4^\alpha\eta^{\beta\gamma}\eta^{\mu\delta}+\textrm{permutations}\Big]\,.
    \end{split}
\end{align}
%Now with the following relabelling we get the integral form given in \eqref{4.68j}: $k_1\to -\ell_2,\,k_4\to \ell_1$ implying, $k_3\to q-\ell_1,\,k_2\to \ell_2-q$. With this we can identify the numerator with spinless and spinning part as,
%\begin{align}
  %  \begin{split}
 %       \mathcal{N}\Big|_{\mathcal{S}^0}=& (\ell_1-q)^2\Bigg[\frac{20 \gamma ^2+3 \gamma ^2 d^2-\left(16 \gamma ^2+3\right) d-2}{8 (d-2)^2}(\ell_1+\ell_2-q)^2+\frac{1}{2 (d-2)}(\ell_1\cdot v_1)^2+\frac{5}{2 (d-2)}(\ell_2\cdot v_2)^2\\ &-\frac{(3 d-14) \left(\gamma ^2 (d-2)-1\right)}{8 (d-2)^2}(\ell_2-q)^2-\frac{\gamma }{2}\gamma (\ell_1\cdot v_1)(\ell_2\cdot v_2)-\frac{(3 d-14) \left(\gamma ^2 (d-2)-1\right)}{8 (d-2)^2}\ell_1^2\\ &+\frac{1}{4} \left(\frac{d-6}{(d-2)^2}-\gamma ^2\right)\ell_2^2+\frac{3}{4} \left(\gamma ^2+\frac{1}{2-d}\right)q^2\Bigg]\\&+
%        (\ell_2-q)^2\Bigg[\frac{(3 d-14) \left(\gamma ^2 (d-2)-1\right)}{8 (d-2)^2}(\ell_1+\ell_2-q)^2+\frac{3}{4} \left(\frac{1}{d-2}-\gamma ^2\right)\ell_1^2+\frac{(14-3 d) \left(\gamma ^2 (d-2)-1\right)}{8 (d-2)^2}\ell_2^2\\ &+\frac{3}{4} \left(\gamma ^2+\frac{1}{2-d}\right)q^2\Bigg]+(\ell_1+\ell_2-q)^2\Bigg[\frac{1}{4-2 d}(\ell_1\cdot v_1)^2+\frac{1}{4-2 d}(\ell_2\cdot v_2)^2+\frac{\gamma ^2 (3 d-10) (d-2)-3 d-2}{8 (d-2)^2}\ell_2^2\\ &
 %       +\frac{(3 d-14) \left(\gamma ^2 (d-2)-1\right)}{8 (d-2)^2}\ell_1^2
 %       +\frac{\gamma }{2}(\ell_1\cdot v_1)(\ell_2\cdot v_2)+\frac{3}{4} \left(\frac{1}{d-2}-\gamma ^2\right)q^2\Bigg]
%        +\frac{1}{4} \left(\frac{d}{(d-2)^2}-\gamma ^2\right)(\ell_1-q)^4\\ &
%      +\frac{-\left(\gamma ^2 (d-6) (d-2)\right)+d+2}{8 (d-2)^2} (\ell_1+\ell_2-q)^4+\gamma(\ell_1\cdot v_1)(\ell_2\cdot v_2)\Bigg[q^2-\frac{\ell_2^2}{2}\Bigg]
%      +(\ell_1\cdot v_1)^2\Bigg[\frac{5\ell_2^2}{2(d-2)}-(\ell_2\cdot v_2)^2\Bigg]\\ &+\frac{1}{2(d-2)}(\ell_2\cdot v_2)^2\,\ell_2^2+\frac{3 \left(\gamma ^2 (d-2)-1\right)}{4 (d-2)}q^2(\ell_1^2+\ell_2^2)+\frac{3}{4} \left(\frac{1}{d-2}-\gamma ^2\right)q^4+\frac{1}{4} \left(\frac{d}{(d-2)^2}-\gamma ^2\right)\ell_2^4\\ &
%      -\frac{(3 d-14) \left(\gamma ^2 (d-2)-1\right)}{8 (d-2)^2}\ell_2^2\ell_1^2
%    \end{split}
%\end{align}
%Therefore the spinless part of the eikonal phase is given by,
Then, proceeding as before, we get, 
\begin{align}
    \begin{split}
        \chi\Big|_{\mathcal{S}^0}= \frac{m_1^2 m_2^2 s_1 s_2}{64 m_p^6}\int e^{-i q\cdot b}\hat\delta(q\cdot v_1)\hat\delta(q\cdot v_2)\Bigg(&\boldsymbol{\tilde h}_{76}\mathbfcal{M}_{0,0;0,0,1,1,1}+\boldsymbol{\tilde h}_{77}\,q^2\mathbfcal{M}_{0,0;0,0,2,1,1}+\boldsymbol{\tilde h}_{78}\,q^2\mathbfcal{M}_{0,0;0,0,1,2,1}\\ &+\boldsymbol{\tilde h}_{79}\,q^2\mathbfcal{M}_{0,0;1,1,0,1,1}+\boldsymbol{\tilde h}_{80}\,q^4\mathbfcal{M}_{0,0;1,1,1,1,1}+\boldsymbol{\tilde h}_{81}\,q^6\mathbfcal{M}_{0,0;1,1,2,1,1}\Bigg)
    \end{split}
\end{align}
%After evaluating the integral we get, 
%\begin{align}
%    \begin{split}
 %       \chi\Big|_{\mathcal{S}^0}\sim \frac{1}{32768 \pi ^3 b^2 x \left(x^2-1\right)}\Big(\frac{2 x}{1-x^2}\Big)&\Bigg[6 \left(55 x^4+36 x^3-36 x-55\right)+\Big(1764 x^2 \log \left(|\boldsymbol{b}|^2\right)+665 x^4+216 x^3\\&+6 x^2 (490 \gamma_E +475-392 \log (2)+98 \log (\pi ))+216 x+665\Big)\log(x)\\&+294x(1-x)\log(x)^2-\big(294 (x-1) x\big)\textbf{Li}_2\left(1-x^2\right)\Bigg]\,.
  %     \end{split}
%\end{align}  
%\newpage
%Now we go on to the computation of the spinning part of the eikonal phase. The spinning numerator takes the following form.
%\begin{align}
%    \begin{split}
%\mathcal{N}\Big|_{\mathcal{S}^1}=&\frac{i}{8(d-2)}(\ell_2-q)\cdot \mathcal{S}_1\cdot v_2\Bigg[\gamma(6  -  d)(\ell_1+\ell_2-q)^4+\gamma  (14-3 d)(\ell_1-q)^2(\ell_2-q)^2
%+\gamma  (14-3 d)\ell_1^2\ell_2^2\\ &+6 \gamma  (d-2)q^2 \ell_1^2+\gamma  (14-3 d)\ell_1^2(\ell_1-q)^2-6 \gamma  (d-2)\ell_1^2(\ell_2-q)^2+\gamma  (3 d-14)\ell_1^2 (\ell_1+\ell_2-q)^2-2 \gamma  (d-2)\ell_2^4\\ &+\gamma  (3 d-10)\ell_2^2 (\ell_1+\ell_2-q)^2-2 \gamma  (d-2)\ell_2^2 (\ell_1-q)^2+\gamma  (14-3 d)\ell_2^2(\ell_2-q)^2+(4-2 d)\ell_2^2 (\ell_1\cdot v_1)(\ell_2\cdot v_2)\\ &+6 \gamma  (d-2)\ell_2^2\,q^2+\gamma  (3 d-14)(\ell_2-q)^2(\ell_1+\ell_2-q)^2+6 \gamma  (d-2)(\ell_2-q)^2\,q^2+(2 d-4)(\ell_1+\ell_2-q)^2(\ell_1\cdot v_1)(\ell_2\cdot v_2)\\ &-6 \gamma  (d-2)(\ell_1+\ell_2-q)^2\,q^2+6 \gamma  (d-2)(\ell_1-q)^2\,q^2+(4-2 d)(\ell_1-q)^2(\ell_1\cdot v_1)(\ell_2\cdot v_2)-6 \gamma  (d-2)q^4\\ &
%+(4 d-8)q^2 (\ell_1\cdot v_1)(\ell_2\cdot v_2)-2 \gamma  (d-2)(\ell_1-q)^4+\gamma  (3 d-10)(\ell_1+\ell_2-q)^2(\ell_1-q)^2
%\Bigg]\\ &
%+\frac{i}{4(d-2)}\ell_2\cdot \mathcal{S}_1\cdot q\Bigg[(\ell_1\cdot v_1)\Big[4 (d-2)(\ell_2\cdot v_2)^2+(\ell_1+\ell_2-q)^2-(\ell_1-q)^2-7\ell_2^2+2 q^2\Big]\\ &+\gamma(\ell_2\cdot v_2)\Big[-4(\ell_1+\ell_2-q)^2+d(\ell_1-q)^2+4\ell_2^2-d \,q^2\Big]\Bigg]\\ &
%+\frac{i}{8(d-2)}(\ell_2-q)\cdot \mathcal{S}_1\cdot \ell_1 \Bigg[4(\ell_1\cdot v_1)\Big[-2(d-2)(\ell_2\cdot v_2)^2-(\ell_1+\ell_2-q)^2+(\ell_1-q)^2+5\ell_2^2\Big]\\ &-2\gamma (d-2)(\ell_2\cdot v_2)\Big[-(\ell_1+\ell_2-q)^2+(\ell_1-q)^2+\ell_2^2-2q^2\Big]\Bigg]\\ &
%+\frac{i}{4(d-2)}(\ell_2-q)\cdot \mathcal{S}_1\cdot \ell_2\Bigg[(\ell_1\cdot v_1)\Big[4 (d-2)(\ell_2\cdot v_2)^2-5(\ell_1+\ell_2-q)^2+5(\ell_1-q)^2\\ &-2(\ell_2-q)^2+2\ell_1^2+3 \ell_2^2-2q^2\Big]+\gamma\,(\ell_2\cdot v_2)\Big[-2(d-4)(\ell_1+\ell_2-q)^2-( d+2)(\ell_1-q)^2+\gamma(d-2)\ell_1^2\\ &+2(d-4)\ell_2^2-(d-2)q^2\Big]\Bigg]\\ &
%+\frac{i}{8(d-2)}\ell_1\cdot \mathcal{S}_2\cdot v_1\Bigg[\gamma  (14-3 d)\ell_1^2\ell_2^2-6 \gamma  (d-2)\ell_1^2(\ell_2-q)^2\\ &
%+\gamma  (3 d-14)\ell_1^2 (\ell_1+\ell_2-q)^2-\gamma  (3 d-14)\ell_1^2(\ell_1-q)^2+6 \gamma  (d-2)\ell_1^2\,q^2\\ &-2 \gamma  (d-2)\ell_2^4+\gamma  (14-3 d)\ell_2^2(\ell_2-q)^2+\gamma  (3 d-10)\ell_2^2(\ell_1+\ell_2-q)^2\\ &-2 \gamma  (d-2)\ell_2^2(\ell_1-q)^2+6 \gamma  (d-2)\ell_2^2\,q^2+(4-2 d))\ell_2^2(\ell_1\cdot v_1)(\ell_2\cdot v_2)\\ &+\gamma  (3 d-10)(\ell_2-q)^2(\ell_1+\ell_2-q)^2+\gamma  (14-3 d)(\ell_2-q)^2(\ell_1-q)^2+6 \gamma  (d-2)(\ell_2-q)^2\,q^2\\ &
%-(\gamma  (d-6))(\ell_1+\ell_2-q)^4+\gamma  (3 d-10)(\ell_1+\ell_2-q)^2(\ell_1-q)^2-6 \gamma  (d-2)(\ell_1+\ell_2-q)^2\,q^2\\ &+2 (d-2)(\ell_1+\ell_2-q)^2(\ell_1\cdot v_1)(\ell_2\cdot v_2)-2 \gamma  (d-2)(\ell_1-q)^4+6 \gamma  (d-2)(\ell_1-q)^2\,q^2\\ &+(4-2 d)(\ell_1-q)^2(\ell_1\cdot v_1)(\ell_2\cdot v_2)+4 (d-2)q^2(\ell_1\cdot v_1)(\ell_2\cdot v_2)-6 \gamma  (d-2)q^4\Bigg]
%    \end{split}
%\end{align}
%\begin{align}
%    \begin{split}
 % &  +    \frac{i}{8(d-2)}\ell_1\cdot \mathcal{S}_2\cdot q\Bigg[2\gamma(d-2)(\ell_1\cdot v_1)\Big[(\ell_1+\ell_2-q)^2-(\ell_1-q)^2+\ell_2^2\Big]+4(\ell_2\cdot v_2)\Big[-2(d-2)(\ell_1\cdot v_1)^2\\ &+3(\ell_1+\ell_2-q)^2-3(\ell_1-q)^2-(\ell_2-q)^2+2\ell_1^2-3\ell_2^2\Big]\Bigg]\\ &
%    +\frac{i}{8(d-2)}\ell_1\cdot \mathcal{S}_2\cdot \ell_2\Bigg[-2\gamma(d-2)(\ell_1\cdot v_1)\Big[-\gamma(\ell_1+\ell_2-q)^2+(\ell_1-q)^2+\ell_2^2-2q^2\Big]\\ &
 %   +4(\ell_2\cdot v_2)\Big[2(d-2)(\ell_1\cdot v_1)^2-(\ell_1+\ell_2-q)^2+5(\ell_1-q)^2+\ell_2^2\Big]\Bigg]
%    \end{split}
%\end{align}
%In order to do the integrals we will face the following vector and tensor integrals,
%\begin{align}
%    \begin{split}
 %      \bf {\Gamma_2}^{\eta}&= \int \frac{\hat\delta(\ell_1\cdot v_2)\hat\delta(\ell_2\cdot v_1)\,\ell_2^\eta}{\ell_1^2\,\ell_2^2\,(\ell_1+\ell_2-q)^2\,(\ell_1-q)^2\,(\ell_2-q)^2}\Bigg[(\ell_1\cdot v_1)\Big[4 (d-2)(\ell_2\cdot v_2)^2+(\ell_1+\ell_2-q)^2-(\ell_1-q)^2-7\ell_2^2+2 q^2\Big]\\ &+\gamma(\ell_2\cdot v_2)\Big[-4(\ell_1+\ell_2-q)^2+d(\ell_1-q)^2+4\ell_2^2-d \,q^2\Big]\Bigg]\\ &
 %       =\frac{\gamma v_1^\eta-v_2^\eta}{\gamma^2-1}\Bigg[\boldsymbol{h_{68}}\,\mathbfcal{M}_{0,0;0,0,1,1,1}    +\boldsymbol{h_{69}}\,\,q^2\mathbfcal{M}_{0,0;0,0,2,1,1}+\boldsymbol{h_{70}}\,\,q^2\mathbfcal{M}_{0,0;0,0,1,2,1} +\boldsymbol{\mathfrak{h}_1} q^2 \mathbfcal{M}_{0,0;1,1,0,1,1}\\&+ \boldsymbol{\mathfrak{h}_2}q^4 \mathbfcal{M}_{0,0;1,1,1,1,1}
 %       +\boldsymbol{\mathfrak{h}_3}q^6 \mathbfcal{M}_{0,0;1,1,2,1,1}
 %       \Bigg]\\ &
%\hspace{0cm}\boldsymbol{\Gamma_1}^{\eta}=\int \frac{\hat\delta(\ell_1\cdot v_2)\hat\delta(\ell_2\cdot v_1)(\ell_2-q)^{\eta}}{\ell_1^2\,\ell_2^2\,(\ell_1+\ell_2-q)^2\,(\ell_1-q)^2\,(\ell_2-q)^2}\Bigg[\gamma(6  -  d)(\ell_1+\ell_2-q)^4+\gamma  (14-3 d)(\ell_1-q)^2(\ell_2-q)^2
%+\gamma  (14-3 d)\ell_1^2\ell_2^2\\ &+6 \gamma  (d-2)q^2 \ell_1^2+\gamma  (14-3 d)\ell_1^2(\ell_1-q)^2-6 \gamma  (d-2)\ell_1^2(\ell_2-q)^2+\gamma  (3 d-14)\ell_1^2 (\ell_1+\ell_2-q)^2-2 \gamma  (d-2)\ell_2^4\\ &+\gamma  (3 d-10)\ell_2^2 (\ell_1+\ell_2-q)^2-2 \gamma  (d-2)\ell_2^2 (\ell_1-q)^2+\gamma  (14-3 d)\ell_2^2(\ell_2-q)^2+(4-2 d)\ell_2^2 (\ell_1\cdot v_1)(\ell_2\cdot v_2)\\ &+6 \gamma  (d-2)\ell_2^2\,q^2+\gamma  (3 d-14)(\ell_2-q)^2(\ell_1+\ell_2-q)^2+6 \gamma  (d-2)(\ell_2-q)^2\,q^2+(2 d-4)(\ell_1+\ell_2-q)^2(\ell_1\cdot v_1)(\ell_2\cdot v_2)\\ &-6 \gamma  (d-2)(\ell_1+\ell_2-q)^2\,q^2+6 \gamma  (d-2)(\ell_1-q)^2\,q^2+(4-2 d)(\ell_1-q)^2(\ell_1\cdot v_1)(\ell_2\cdot v_2)-6 \gamma  (d-2)q^4\\ &
%+(4 d-8)q^2 (\ell_1\cdot v_1)(\ell_2\cdot v_2)-2 \gamma  (d-2)(\ell_1-q)^4+\gamma  (3 d-10)(\ell_1+\ell_2-q)^2(\ell_1-q)^2
%\Bigg]\\  &
%=\Bigg[\boldsymbol{h_{71}}\,\,\mathbfcal{M}_{0,0;0,0,1,1,1}+\boldsymbol{h_{72}}\,q^2\mathbfcal{M}_{0,0;0,0,2,1,1}+\boldsymbol{h_{73}}\,q^2 \mathbfcal{M}_{0,0;0,0,1,2,1}+\boldsymbol{h_{74}}\,q^2\mathbfcal{M}_{0,0;1,1,0,1,1}+\boldsymbol{h_{75}}\,q^4\mathbfcal{M}_{0,0;1,1,1,1,1}\\ &+\boldsymbol{h_{76}}\,\,q^6\mathbfcal{M}_{0,0;1,1,2,1,1}\Bigg]q^{\eta}\\ & 
%     \hspace{0cm}\boldsymbol{\Gamma}_{4}^{\eta}=\int \frac{\hat\delta(\ell_1\cdot v_2)\hat\delta(\ell_2\cdot v_1)\ell_2^{\eta}}{\ell_1^2\,\ell_2^2\,(\ell_1+\ell_2-q)^2\,(\ell_1-q)^2\,(\ell_2-q)^2}\Bigg[(\ell_1\cdot v_1)\Big[4 (d-2)(\ell_2\cdot v_2)^2-5(\ell_1+\ell_2-q)^2+5(\ell_1-q)^2\\ &-2(\ell_2-q)^2+2\ell_1^2+3 \ell_2^2-2q^2\Big]+\gamma\,(\ell_2\cdot v_2)\Big[-2(d-4)(\ell_1+\ell_2-q)^2-( d+2)(\ell_1-q)^2+\gamma(d-2)\ell_1^2\\ &+2(d-4)\ell_2^2-(d-2)q^2\Big]\Bigg]\\ &
%     =\frac{\gamma v_1^\eta-v_2^\eta}{\gamma^2-1}\Bigg[\boldsymbol{h_{77}}\,\mathbfcal{M}_{0,0;0,0,1,1,1}    +\boldsymbol{h_{78}}\,q^2\mathbfcal{M}_{0,0;0,0,2,1,1}
%     +\boldsymbol{h_{79}}\,q^2 \mathbfcal{M}_{0,0;0,0,1,2,1}+\boldsymbol{h_{80}}\,q^2\mathbfcal{M}_{0,0;1,1,0,1,1}\\&+\boldsymbol{h_{81}}\,q^4\mathbfcal{M}_{0,0;1,1,1,1,1}
%+\boldsymbol{h_{82}}\,q^6\mathbfcal{M}_{0,0;1,1,2,1,1}\Bigg]\\ & 
%\hspace{0cm}\boldsymbol{\Gamma}_{6}^{\eta}=\int \frac{\hat\delta(\ell_1\cdot v_2)\hat\delta(\ell_2\cdot v_1)\ell_1^\mu}{\ell_1^2\,\ell_2^2\,(\ell_1+\ell_2-q)^2\,(\ell_1-q)^2\,(\ell_2-q)^2}\Bigg[2\gamma(d-2)(\ell_1\cdot v_1)\Big[(\ell_1+\ell_2-q)^2-(\ell_1-q)^2+\ell_2^2\Big]+4(\ell_2\cdot v_2)\Big[-2(d-2)(\ell_1\cdot v_1)^2\\ &+3(\ell_1+\ell_2-q)^2-3(\ell_1-q)^2-(\ell_2-q)^2+2\ell_1^2-3\ell_2^2\Big]\Bigg]
%    \end{split}
%\end{align}
%\begin{align}
%    \begin{split}
%        &=\Bigg(\boldsymbol{h_{83}}\,\mathbfcal{M}_{0,0;0,0,1,1,1}
%+\boldsymbol{h_{84}}\,q^2\,\mathbfcal{M}_{0,0;0,0,2,1,1}+\boldsymbol{h_{85}}\,q^2 \mathbfcal{M}_{0,0;0,0,1,2,1}+\boldsymbol{h_{86}}\,q^2 \mathbfcal{M}_{0,0;1,1,0,1,1}+\boldsymbol{h_{87}}\,q^4 \, \mathbfcal{M}_{0,0;1,1,1,1,1}\\ &+\boldsymbol{h_{88}}\,q^6\,\mathbfcal{M}_{0,0;1,1,2,1,1}\Bigg)\frac{\gamma v_2^\eta-v_1^\eta}{\gamma^2-1}\\ & \hspace{0cm}\boldsymbol{\Gamma}_5^\eta=\int\frac{\hat\delta(\ell_1\cdot v_2)\hat\delta(\ell_2\cdot v_1)\ell_1^\eta}{\ell_1^2\,\ell_2^2\,(\ell_1+\ell_2-q)^2\,(\ell_1-q)^2\,(\ell_2-q)^2}\Bigg[\gamma  (14-3 d)\ell_1^2\ell_2^2-6 \gamma  (d-2)\ell_1^2(\ell_2-q)^2\\ &
%+\gamma  (3 d-14)\ell_1^2 (\ell_1+\ell_2-q)^2-\gamma  (3 d-14)\ell_1^2(\ell_1-q)^2+6 \gamma  (d-2)\ell_1^2\,q^2\\ &-2 \gamma  (d-2)\ell_2^4+\gamma  (14-3 d)\ell_2^2(\ell_2-q)^2+\gamma  (3 d-10)\ell_2^2(\ell_1+\ell_2-q)^2\\ &-2 \gamma  (d-2)\ell_2^2(\ell_1-q)^2+6 \gamma  (d-2)\ell_2^2\,q^2+(4-2 d))\ell_2^2(\ell_1\cdot v_1)(\ell_2\cdot v_2)\\ &+\gamma  (3 d-10)(\ell_2-q)^2(\ell_1+\ell_2-q)^2+\gamma  (14-3 d)(\ell_2-q)^2(\ell_1-q)^2+6 \gamma  (d-2)(\ell_2-q)^2\,q^2\\ &
%-(\gamma  (d-6))(\ell_1+\ell_2-q)^4+\gamma  (3 d-10)(\ell_1+\ell_2-q)^2(\ell_1-q)^2-6 \gamma  (d-2)(\ell_1+\ell_2-q)^2\,q^2\\ &+2 (d-2)(\ell_1+\ell_2-q)^2(\ell_1\cdot v_1)(\ell_2\cdot v_2)-2 \gamma  (d-2)(\ell_1-q)^4+6 \gamma  (d-2)(\ell_1-q)^2\,q^2\\ &+(4-2 d)(\ell_1-q)^2(\ell_1\cdot v_1)(\ell_2\cdot v_2)+4 (d-2)q^2(\ell_1\cdot v_1)(\ell_2\cdot v_2)-6 \gamma  (d-2)q^4\Bigg]\\ &
%=\Bigg[\boldsymbol{h_{89}}\,\mathbfcal{M}_{0,0;0,0,1,1,1}+\boldsymbol{h_{90}}\,q^2\mathbfcal{M}_{0,0;0,0,2,1,1}+ \boldsymbol{h_{91}}\,q^2\,\mathbfcal{M}_{0,0;0,0,1,2,1}\\ &+\boldsymbol{h_{92}}\,q^2\,\mathbfcal{M}_{0,0;1,1,0,1,1}+\boldsymbol{h_{93}}\,q^4\,\mathbfcal{M}_{0,0;1,1,1,1,1}+\boldsymbol{h_{94}}\,q^6\,\mathbfcal{M}_{0,0;1,1,2,1,1}\Bigg]q^\eta
 %   \end{split}
%\end{align}
%Apart from the vector integrals we will face the following tensor integrals,
%\begin{align}
%    \begin{split}
        %\boldsymbol{\Gamma}_3^{\eta\delta}&=\int\frac{\hat\delta(\ell_1\cdot v_2)\hat\delta(\ell_2\cdot v_1)(\ell_2-q)^{\eta}\ell_1^{\delta}}{\ell_1^2\,\ell_2^2\,(\ell_1+\ell_2-q)^2\,(\ell_1-q)^2\,(\ell_2-q)^2}\Bigg[4(\ell_1\cdot v_1)\Big[-2(d-2)(\ell_2\cdot v_2)^2-(\ell_1+\ell_2-q)^2+(\ell_1-q)^2+5\ell_2^2\Big]\\ &-2\gamma (d-2)(\ell_2\cdot v_2)\Big[-(\ell_1+\ell_2-q)^2+(\ell_1-q)^2+\ell_2^2-2q^2\Big]\Bigg]\\ &
        %=\left(\Delta_1(q,\epsilon,\gamma)\,q^{[\eta}v_1^{\delta]}+\Delta_2(q,\epsilon,\gamma)\,q^{[\eta}v_2^{\delta]}\right)
  %  \end{split}
%\end{align}
%where,
%\begin{align}
%    \begin{split}
%        \hspace{0cm}\Delta_1(q,\epsilon,\gamma)=&\boldsymbol{h_{95}}\,
%\mathbfcal{M}_{0,0;0,0,1,1,1}+\boldsymbol{h_{96}}\,q^2\mathbfcal{M}_{0,0;0,0,2,1,1}
%+\boldsymbol{h_{97}}\,q^2\mathbfcal{M}_{0,0;0,0,1,2,1}+\boldsymbol{h_{98}}\,q^2\,\mathbfcal{M}_{0,0;1,1,0,1,1}
%+\boldsymbol{h_{99}}\,q^4\,\mathbfcal{M}_{0,0;1,1,1,1,1}\\ &+\boldsymbol{h_{100}}\,q^6
%\mathbfcal{M}_{0,0;1,1,2,1,1}
 %   \end{split}
%\end{align}
%and,
%\begin{align}
%    \begin{split}
 %       \hspace{0cm}\Delta_2(q,\epsilon,\gamma)=&\boldsymbol{h_{101}}        \mathbfcal{M}_{0,0;0,0,1,1,1}+\boldsymbol{h_{102}}\,q^2\mathbfcal{M}_{0,0;0,0,2,1,1}
        %+\boldsymbol{h_{103}}\,q^2\mathbfcal{M}_{0,0;0,0,1,2,1}+\boldsymbol{h_{104}}\,q^2 \mathbfcal{M}_{0,0;1,1,0,1,1}
  %      +\boldsymbol{h_{105}}\,q^4 \mathbfcal{M}_{0,0;1,1,1,1,1}\\ &+\boldsymbol{h_{106}}\,q^6\,\mathbfcal{M}_{0,0;1,1,2,1,1}
%    \end{split}
%\end{align}
%\newpage
%\begin{align}
%    \begin{split}
%        \boldsymbol{\Gamma}_7^{\eta\delta}&=\int\frac{\hat\delta(\ell_1\cdot v_2)\hat\delta(\ell_2\cdot v_1)\ell_1^\eta \ell_2^\delta}{\ell_1^2\,\ell_2^2\,(\ell_1+\ell_2-q)^2\,(\ell_1-q)^2\,(\ell_2-q)^2}\Bigg[-2\gamma(d-2)(\ell_1\cdot v_1)\Big[-\gamma(\ell_1+\ell_2-q)^2+(\ell_1-q)^2+\ell_2^2-2q^2\Big]\\ &
 %   +4(\ell_2\cdot v_2)\Big[2(d-2)(\ell_1\cdot v_1)^2-(\ell_1+\ell_2-q)^2+5(\ell_1-q)^2+\ell_2^2\Big]\Bigg]\\ &
        %=\Bigg(\Lambda_1(q,\epsilon,\gamma)\,q^{[\eta}v_1^{\delta]}+\Lambda_2(q,\epsilon,\gamma)\,q^{[\eta}v_2^{\delta]}\bigg)
 %   \end{split}
%\end{align}
%where,
%\begin{align}
%    \begin{split}
%    \Lambda_1(q,\epsilon,\gamma)=&\boldsymbol{h_{107}}\,\mathbfcal{M}_{0,0;0,0,1,1,1}+\boldsymbol{h_{108}}\,\,q^2 \mathbfcal{M}_{0,0;0,0,2,1,1}+\boldsymbol{h_{109}}\,q^2\mathbfcal{M}_{0,0;0,0,1,2,1}+\boldsymbol{h_{110}}\,q^2 \mathbfcal{M}_{0,0;1,1,0,1,1}\\ &
%+\boldsymbol{h_{111}}\,q^4 \mathbfcal{M}_{0,0;1,1,1,1,1}+\boldsymbol{h_{112}}\,q^6 \mathbfcal{M}_{0,0;1,1,2,1,1}
 %   \end{split}
%\end{align}
%\newpage  
and,
\begin{align}
    \begin{split}
        \chi\Big|_{\mathcal{S}^1}= \frac{m_1^2 m_2^2 s_1 s_2}{64 m_p^6}\int e^{-i q\cdot b}\hat\delta(q\cdot v_1)\hat\delta(q\cdot v_2)\Bigg[i\,\Bigg(&\boldsymbol{\tilde h}_{82}\mathbfcal{M}_{0,0;0,0,1,1,1}+\boldsymbol{\tilde h}_{83}\,q^2\mathbfcal{M}_{0,0;0,0,2,1,1}\\ &+\boldsymbol{\tilde h}_{84}\,q^2\mathbfcal{M}_{0,0;0,0,1,2,1}+\boldsymbol{\tilde h}_{85}\,q^2\mathbfcal{M}_{0,0;1,1,0,1,1}\\&+\boldsymbol{\tilde h}_{86}\,q^4\mathbfcal{M}_{0,0;1,1,1,1,1}+\boldsymbol{\tilde h}_{87}\,q^6\mathbfcal{M}_{0,0;1,1,2,1,1}\Bigg)\big(q\cdot\mathcal{S}_1\cdot v_2\big)+\\&i\,\Bigg(\boldsymbol{\tilde h}_{88}\mathbfcal{M}_{0,0;0,0,1,1,1}+\boldsymbol{\tilde h}_{89}\,q^2\mathbfcal{M}_{0,0;0,0,2,1,1}\\ &+\boldsymbol{\tilde h}_{90}\,q^2\mathbfcal{M}_{0,0;0,0,1,2,1}+\boldsymbol{\tilde h}_{91}\,q^2\mathbfcal{M}_{0,0;1,1,0,1,1}\\&+\boldsymbol{\tilde h}_{92}\,q^4\mathbfcal{M}_{0,0;1,1,1,1,1}+\boldsymbol{\tilde h}_{93}\,q^6\mathbfcal{M}_{0,0;1,1,2,1,1}\Bigg)\big(q\cdot\mathcal{S}_2\cdot v_1\big)\Bigg]\,.
    \end{split}
\end{align}
Then using \eqref{MasterM} we get,

\begin{align}
    \begin{split} \label{eqt12}   
        \chi\Big|_{\mathcal{S}^0}=\frac{m_1^2m_2^2}{m_p^6}\Bigg[\Delta_{\mathfrak{18}}^{\textrm{poly}}+\Delta_{\mathfrak{09}}^{\textrm{log}}\log(x)+\Delta_{\mathfrak{03}}^{\textrm{Polylog}}\Big(\log(x)^2+ \textbf{Li}_2\left(1-x^2\right)\Big)\Bigg]\,,
          \end{split}
\end{align}
\enlargethispage{2\baselineskip}
where, 
\begin{thesiscompactdisplay}
\begin{align}
    \begin{split} 
  &  \Delta_{\mathfrak{18}}^{\textrm{poly}}= \frac{s_1s_2 (3-3 x^4)}{16384 \pi ^3 |\boldsymbol{b}|^2 \left(x^2-1\right)^2}\,,   \Delta_{\mathfrak{09}}^{\textrm{log}}=\frac{s_1s_2\bigg\{4 \left(-2 x^2 \log \left(\pi  |\boldsymbol{b}|^6\right)+5 x^4+x^2 (-10 \gamma_E -4+\log (256))+5\right)\bigg\}}{16384 \pi ^3 |\boldsymbol{b}|^2 \left(x^2-1\right)^2}\,,   \\& \Delta_{\mathfrak{03}}^{\textrm{Polylog}}=\frac{4\,s_1s_2\,(x-1) x}{16384 \pi ^3 |\boldsymbol{b}|^2 \left(x^2-1\right)^2}
      \end{split}
\end{align}
\end{thesiscompactdisplay}
and 
\begin{align}
    \begin{split} \label{eqt13}  
        \hspace{0cm}\chi\Big|_{\mathcal{S}^1}=&\frac{m_1^2m_2^2}{m_p^6}\Bigg\{\Bigg[\Delta_{\mathfrak{19}}^{\textrm{poly}}+\Delta_{\mathfrak{10}}^{\textrm{log}}\log(x)+\Delta_{\mathfrak{04}}^{\textrm{Polylog}}\Big(\log(x)^2+ \textbf{Li}_2\left(1-x^2\right)\Big)\Bigg]\big(\hat{b}\cdot\mathcal{S}_1\cdot v_2\big)\\&\hspace{1cm}+\Bigg[\Delta_{\mathfrak{20}}^{\textrm{poly}}+\Delta_{\mathfrak{11}}^{\textrm{log}}\log(x)+\Delta_{\mathfrak{05}}^{\textrm{Polylog}}\Big(\log(x)^2+ \textbf{Li}_2\left(1-x^2\right)\Big)\Bigg]\big(\hat{b}\cdot\mathcal{S}_2\cdot v_1\big)\Bigg\}
          \end{split}
\end{align}
\enlargethispage{2\baselineskip}
where,
\begin{thesiscompactdisplay}
\begin{align}
\begin{split}
&\Delta_{\mathfrak{19}}^{\textrm{poly}}=\frac{s_{1} s_{2} \left(x^6-10 x^5+3 x^4-49 x^3+3 x^2-10 x+1\right)}{8192 \pi ^3 |\boldsymbol{b}|^3 \left(x^2-1\right)^3},\\ &
\Delta_{\mathfrak{4}}^{\textrm{Polylog}}=\frac{7 x^2\,s_1s_2\, (x-1) x^2 \left(x^4+4 x^2+1\right)}{8192 \pi ^3 |\boldsymbol{b}|^3 \left(x^2-1\right)^4 \left(x^2+1\right)},\\ &
\Delta_{\mathfrak{10}}^{\textrm{log}}=-\Bigg(\frac{s_1s_2}{{16384 \pi ^3 |\boldsymbol{b}|^3 \left(x^2-1\right)^4 \left(x^2+1\right)}}\Bigg)\Bigg[x \big(28 \left(x^6+4 x^4+x^2\right) \log \left(\pi  |\boldsymbol{b}|^6\right)-37 x^8+8 x^7+4 x^6 (35 \gamma_E -8\\&-28 \log (2))+24 x^5+x^4 (560 \gamma_E -278-448 \log (2))+24 x^3+4 x^2 (35 \gamma_E -8-28 \log (2))+8 x-37\big)\Bigg]\,,\\ &
\Delta_{\mathfrak{11}}^{\textrm{log}}=\Bigg(\frac{s_1s_2}{16384 \pi ^3 |\boldsymbol{b}|^3 \left(x^2-1\right)^6 \left(x^2+1\right)^2}\Bigg)\Bigg[2 \big(48 \left(x^2-1\right)^2 \left(x^6+5 x^4+5 x^2+1\right) x^3 \log ( |\boldsymbol{b}|)-18 x^{16}-13 x^{15}+96 x^{14}\\
&+x^{13} (40 \gamma_E -33-32 \log (2)+8 \log (\pi ))+280 x^{12}+x^{11} (120 \gamma_E -47-96 \log (2)+24 \log (\pi ))\\
&-992 x^{10}+x^9 (-160 \gamma_E +93+128 \log (2)-32 \log (\pi ))-2604 x^8\\
&+x^7 (-160 \gamma_E +93+128 \log (2)-32 \log (\pi ))-992 x^6+x^5 (120 \gamma_E -47-96 \log (2)+24 \log (\pi ))\\
&+280 x^4+x^3 (40 \gamma_E -33-32 \log (2)+8 \log (\pi ))+96 x^2-13 x-18\big)\Bigg]\,,\\ &
\Delta_{\mathfrak{5}}^{\textrm{Polylog}}=-\frac{8 (x-1) x^2 \left(x^6+5 x^4+5 x^2+1\right)\, s_1s_2}{16384 \pi ^3 |\boldsymbol{b}|^3 \left(x^2-1\right)^4 \left(x^2+1\right)^2}\,,\\ &
\Delta_{\mathfrak{20}}^{\textrm{poly}}=\frac{x  \left(19 x^4-218 x^2+19\right)\,s_1s_2}{16384 \pi ^3 |\boldsymbol{b}|^3 \left(x^2-1\right)^3 }\,.
\end{split}
\end{align}
\end{thesiscompactdisplay}

%\begin{align}
%    \begin{split}
        %\Lambda_2(q,\gamma,\epsilon)=&\boldsymbol{h_{113}}\,\mathbfcal{M}_{0,0;0,0,1,1,1}+\boldsymbol{h_{114}}\,\,q^2 \mathbfcal{M}_{0,0;0,0,2,1,1}
  %      +\boldsymbol{h_{115}}\,q^2 %\mathbfcal{M}_{0,0;0,0,1,2,1}+\boldsymbol{h_{116}}\,q^2\mathbfcal{M}_{0,0;1,1,0,1,1}\\ &
        %+\boldsymbol{h_{117}}\,q^4\,\mathbfcal{M}_{0,0;1,1,1,1,1}+\boldsym%bol{h_{118}}\,\,q^6\,\mathbfcal{M}_{0,0;1,1,2,1,1}
%    \end{split}
%\end{align}
%Then we get following two terms proportional two $\mathcal{S}_1$ and $\mathcal{S}_2$, 
%\begin{align}
%    \begin{split}
 %       \chi^{(1)}\Big|_{\mathcal{S}^1}\sim 
%       \end{split}
%\end{align}  
%and 
%\begin{align}
%    \begin{split}
 %       \chi^{(2)}\Big|_{\mathcal{S}^1}\sim \,\frac{\hat{b}\cdot %\mathcal{S}_{2}\cdot v_{1}}{1024 \pi ^3 b^3 \left(x^2-1\right)^4(1+x^2)}\Big(\frac{2x}{1-x^2}\Big)&\Bigg[\Bigg]\,,
  %     \end{split}
%\end{align}  
%\\ 
\textbullet $\,\,$ At 3PM we face a $\mathbb{N}$-type topology involving scalar and graviton lines. It gives the following three diagrams. 
\begin{center}
\includegraphics[width=0.46\linewidth]{N_topo.png}
\end{center}
The spinless contributions come from the first and the second diagrams only. We get, 
%The contribution to the eikonal phase from the first diagram is given by (relabelling $k_1\to \ell_1,\,k_3\to \ell_2,\,k_2\to q-\ell_1-\ell_2$),
\begin{align}
\begin{split} \label{eqt14}  
    \chi\Big|_{\mathcal{S}^0}&=\frac{m_1^2 m_2^2}{16\,m_p^6}\int e^{-i q\cdot b}\hat\delta(q\cdot v_1)\hat\delta(q\cdot v_2)\Bigg(s_1s_2g_1g_2+s_1^2s_2^2\,\boldsymbol{\tilde h}_{94}\Bigg)\mathbfcal{M}_{0,0;0,0,1,1,1}\,,\\&
    =\frac{m_1^2 m_2^2}{m_p^6}\Delta_{\mathfrak{12}}^{\textrm{log}}\,,\,\,\Delta_{\mathfrak{12}}^{\textrm{log}}=\Bigg(\frac{s_{1}s_{2} \log (x) \left(8 g_{1} g_{2} x^2+s_{1} s_{2} \left(x^4+1\right)\right)}{1024 \pi ^3 |\boldsymbol{b}|^2 \left(x^2-1\right)^2}\Bigg)\,.
    \end{split}
\end{align}
%\hat\delta(q\cdot v_2)\int \frac{\hat\delta(\ell_1\cdot v_2)\hat\delta(\ell_2\cdot v_1)}{\ell_1^2 \ell_2^2 (\ell_1+\ell_2-q)^2}\\ &
%    =i\frac{m_1^2 m_2^2 s_1 s_2 g_1 g_2}{16 m_p^6}\int e^{-i q\cdot b}\,

%The contribution from the second diagram takes the following form,
%\begin{align}
 %   \begin{split}
%\textcolor{red}
%{\mathcal{N}|_{\mathcal{S}^0}}=
%\gamma ^2+\frac{1}{2 (1-\frac{d}{2})}
 %  \end{split}
%\end{align}

%\begin{align}
%   \begin{split}
%\mathcal{N}|_{\mathcal{S}^1}=(\ell_2-q)\cdot \mathcal{S}_1\cdot v_2(-i\gamma)-(\ell_2-q)\cdot \mathcal{S}_2\cdot v_1(-i\gamma)
 %   \end{split}
%\end{align}

%\textcolor{black}{AB: Spinning part needs to be checked. Then only we can proceed with this diagram. }
%Hence the contribution to the eikonal phase from the second diagram is given by,

%\begin{align}
 %   \begin{split}
 %     \textcolor{blue} {i \chi_2^{(0)}=i\Bigg(\gamma ^2+\frac{1}{2 (\epsilon -1)}\Bigg)\frac{m_1^2 m_2^2 s_1^2  s_2^2}{32 m_p^4}\int e^{-i q\cdot b}\hat\delta(q\cdot v_1)\hat\delta(q\cdot v_2)\,\mathbfcal{M}_{0,0;0,0,1,1,1}},
%    \end{split}
%\end{align}
The spinning contributions come from the second and third diagrams. We get,
%=-i\gamma\frac{m_1^2 m_2^2 s_1^2  s_2^2}{32 m_p^6}\int e^{-i q\cdot b}\hat\delta(q\cdot v_1)\hat\delta(q\cdot v_2)\,\int \frac{\hat\delta(\ell_1\cdot v_2)\hat\delta(\ell_2\cdot v_1)(\ell_2-q)^\delta}{\ell_1^2 \ell_2^2 (\ell_1+\ell_2-q)^2}\\&
\begin{align}
    \begin{split}
       & \chi\Big|_{\mathcal{S}^{1}}=-\frac{m_1^2 m_2^2 s_1^2  s_2^2}{32 m_p^6}\int e^{-i q\cdot b}\hat\delta(q\cdot v_1)\hat\delta(q\cdot v_2)\Bigg[i\,\Bigg(\boldsymbol{\tilde h}_{95}\mathbfcal{M}_{0,0;0,0,1,1,1}+\boldsymbol{\tilde h}_{96}\,q^2\mathbfcal{M}_{0,0;0,0,1,2,1}\\ &\hspace{0.8cm}+\boldsymbol{\tilde h}_{97}\,q^2\mathbfcal{M}_{0,0;0,0,2,1,1}\Bigg)\big(q\cdot\mathcal{S}_1\cdot v_2\big)+\\&\hspace{0.8cm}i\,\Bigg\{\Big(\frac{1}{2 s_1 s_2}\boldsymbol{\tilde h}_{98}-\boldsymbol{\tilde h}_{95}\Big)\mathbfcal{M}_{0,0;0,0,1,1,1}+\Big(\frac{1}{2 s_1 s_2}\boldsymbol{\tilde h}_{99}-\boldsymbol{\tilde h}_{97}\Big)\,q^2\mathbfcal{M}_{0,0;0,0,2,1,1}\\ &\hspace{0.8cm}+\Big(\frac{1}{2 s_1 s_2}\boldsymbol{\tilde h}_{100}-\boldsymbol{\tilde h}_{96}\Big)\,q^2\mathbfcal{M}_{0,0;0,0,1,2,1}\Bigg)\big(q\cdot\mathcal{S}_2\cdot v_1\big)\Bigg\}\Bigg]\,.
    \end{split}
\end{align}
Now using (\ref{MasterM}) we get, 
\begin{align}
    \begin{split} \label{eqt15}     \chi\Big|_{\mathcal{S}^{1}}=\frac{m_1^2m_2^2}{m_p^6}\Bigg\{&\Bigg[\Delta_{\mathfrak{13}}^{\textrm{log}}\log(x)+\Delta_{\mathfrak{06}}^{\textrm{Polylog}}\Big(\log(x)^2+ \textbf{Li}_2\left(1-x^2\right)\Big)\Bigg]\big(\hat{b}\cdot\mathcal{S}_1\cdot v_2\big)+\\&\Bigg[\Delta_{\mathfrak{21}}^{\textrm{poly}}+\Delta_{\mathfrak{14}}^{\textrm{log}}\log(x)+\Delta_{\mathfrak{07}}^{\textrm{Polylog}}\Big(\log(x)^2+ \textbf{Li}_2\left(1-x^2\right)\Big)\Bigg]\big(\hat{b}\cdot\mathcal{S}_2\cdot v_1\big)\Bigg\}\,,
         \end{split}
\end{align}
where, 
\begin{align}
    \begin{split}
 &   \Delta_{\mathfrak{13}}^{\textrm{log}}=-\Bigg(\frac{s_{1}^2 s_{2}^2}{2048 \pi ^3 |\boldsymbol{b}|^3 \left(x^2-1\right)^4}\Bigg)\Bigg[x \left(x^2+1\right) \big(-12 \left(x^4+1\right) \log (|\boldsymbol{b}|)+x^4-10 \gamma_E  \left(x^4+1\right)+8 x^4 \log (2)\\&\hspace{0.8cm}-2 \left(x^4+1\right) \log (\pi )+18 x^2+1+\log (256)\big)\Bigg]\,,\\&
    \Delta_{\mathfrak{06}}^{\textrm{Polylog}}=-\frac{3 s_{1}^2 s_{2}^2 x \left(x^2+1\right) \left(x^4+1\right)}{1024 \pi ^3 |\boldsymbol{b}|^3 \left(x^2-1\right)^4}\,,
     \Delta_{\mathfrak{21}}^{\textrm{poly}}=-\frac{41\,s_{1} s_{2} x \left(x^2+1\right)^2}{2048 \pi ^3 |\boldsymbol{b}|^3 \left(x^2-1\right)^3}\,,\\&
       \Delta_{\mathfrak{14}}^{\textrm{log}}=\Bigg(\frac{s_{1}s_{2}}{4096 \pi ^3 |\boldsymbol{b}|^3 \left(x^2-1\right)^4}\Bigg)\Bigg[-x \left(x^2+1\right) \big(12 \left(x^4+1\right) \log (|\boldsymbol{b}|) (2\,s_{1}s_{2}+1)+10 \gamma_E  \left(x^4+1\right) (2\,s_{1} s_{2}+1)\\&\hspace{0.8cm}+2 \log (\pi ) \left(2\,s_{1} s_{2} \left(x^4+1\right)+x^4\right)-2\,s_{1}s_{2} \left(x^4+8 \left(x^4+1\right) \log (2)+18 x^2+1\right)\\&\hspace{0.8cm}+x^4 (5-8 \log (2))-18 x^2+5-8 \log (2)+2 \log (\pi )\big)\Bigg]\,,\\&
       \Delta_{\mathfrak{07}}^{\textrm{Polylog}}=\frac{3 \,s_{1} s_{2} x \left(x^2+1\right) \left(x^4+1\right) (2 s_{1}s_{2}+1)}{2048 \pi ^3 |\boldsymbol{b}|^3 \left(x^2-1\right)^4 }\,.
      \end{split}
\end{align}
%where
%\begin{align}
%\begin{split}
%a_2=0\end{split}\end{align}
%    \begin{align}
%    \begin{split}
%    a_3&=\frac{-28 \gamma ^3 \epsilon ^3+8 \gamma ^3 \epsilon ^2+7 \gamma ^3 \epsilon -2 \gamma ^3-20 \gamma  \epsilon ^3-12 \gamma  \epsilon ^2+\gamma  \epsilon +\gamma}{4 (\gamma^2-1)  (\epsilon -1)   (2 \epsilon -1) (3 \epsilon -1)}\mathbfcal{M}_{0, 0, 0, 0, 1, 1, 1}\\&+\frac{16 \gamma  q^2 \epsilon ^3+12 \gamma  q^2 \epsilon ^2+(-\gamma ) q^2}{8 (\gamma^2-1)  (\epsilon -1) \epsilon  (2 \epsilon -1) (3 \epsilon -1)}\mathbfcal{M}_{0, 0, 0, 0, 1, 2, 1}+\frac{4 \gamma ^3 q^2 \epsilon ^3-\gamma ^3 q^2 \epsilon -4 \gamma  q^2 \epsilon ^3+\gamma  q^2 \epsilon}{8 (\gamma^2-1)  (\epsilon -1) \epsilon  (2 \epsilon -1) (3 \epsilon -1)}\mathbfcal{M}_{0, 0, 0, 0, 2, 1, 1}
%\end{split}
%\end{align}
%From the third diagram we have the following contribution,
%\begin{align}
 %   \begin{split}
%\mathcal{N}|_{\mathcal{S}^1}=\textcolor{red}{i}(\ell_2-q)\cdot %\mathcal{S}_2\cdot v_1\Bigg[\frac{\gamma(5 - d)}{2(1-\frac{d}{2})}\Bigg]
%    \end{split}
%\end{align}
%-\frac{\gamma(5-d) \,m_1^2 m_2^2 s_1 s_2}{64(1-d/2) m_p^6}\int e^{-i q\cdot b}\hat\delta(q\cdot v_1)\hat\delta(q\cdot v_2)\int \frac{\hat\delta(\ell_1\cdot v_2)\hat\delta(\ell_2\cdot v_1)\,(v_1\cdot \mathcal{S}_2\cdot (q-\ell_2))}{\ell_1^2\,\ell_2^2\,(\ell_1+\ell_2-q)^2}\\ &
%\begin{align}
 %   \begin{split}
% \chi_3&
%        =-\frac{m_1^2 m_2^2 s_1^2 s_2^2}{32 m_p^6}\int e^{-i q\cdot b}\hat\delta(q\cdot v_1)\hat\delta(q\cdot v_2)\,\Bigg[i\,\Bigg((\frac{1}{2 s_1 s_2}\boldsymbol{\tilde h}_{98}-\boldsymbol{\tilde h}_{95})\mathbfcal{M}_{0,0;0,0,1,1,1}+(\frac{1}{2 s_1 s_2}\boldsymbol{\tilde h}_{99}-\boldsymbol{\tilde h}_{97})\,q^2\mathbfcal{M}_{0,0;0,0,2,1,1}\\ &\hspace{0.8cm}+(\frac{1}{2 s_1 s_2}\boldsymbol{\tilde h}_{100}-\boldsymbol{\tilde h}_{96})\,q^2\mathbfcal{M}_{0,0;0,0,1,2,1}\Bigg)\big(q\cdot\mathcal{S}_2\cdot v_1\big)\Bigg]
 %   \end{split}
%\end{align}
%\textcolor{red}{diagram?}Fourth N-topology diagram contribution at $\mathcal{O}(\mathcal{S}_0)$;
%\begin{align}
%\mathcal{N}|_{\mathcal{S}_0}=\Big(\gamma^2-\frac{1}{2}\Big)
%\end{align}
%The spinning eikonal at $\mathcal{O}(S_1)$ is given by,
%\begin{align}
%\mathcal{N}|_{\mathcal{S}_1}=i\gamma\Big(v_2\cdot \mathcal{S}_1\cdot (l_1)-\,v_1\cdot \mathcal{S}_2\cdot (l_1)\Big)
%\end{align}
%As the spinning eikonal contribution vanishes, hence the eikonal phase is given by,

%\begin{align}
   % \begin{split}
   %     i \chi_4=i(\gamma^2-1/2)\textcolor{red}{\frac{m_1^2 m_2^2 s_1^2  s_2^2}{32 m_p^4}\int e^{-i q\cdot b}\hat\delta(q\cdot v_1)\hat\delta(q\cdot v_2)\,\mathbfcal{M}_{0,0;0,0,1,1,1}},
   % \end{split}
%\end{align}
\textbullet $\,\,$ We face the following topology with single 3-point kind of interaction vertex,
\begin{center}
\includegraphics[width=0.46\linewidth]{new_top.png}
\end{center}
Before quoting the result for the contribution to the eikonal phase, we note that because of the presence of graviton two-point worldline vertex in the second diagram, this diagram only contributes to the spinning eikonal, and also, there is no term linear in the spin of the first black hole coming from this diagram. Now, we first write down the total spinless contribution from these three diagrams, 
\begin{align}
    \begin{split}  \label{eqt16} 
   \chi\Big|_{\mathcal{S}^0}&=\frac{m_1^2 m_2^2}{m_p^6}\int e^{-i q\cdot b}\hat\delta(q\cdot v_1)\hat\delta(q\cdot v_2)\Bigg[\Bigg(\frac{s_1^2 s_2^2}{128 (1-\epsilon)}-\frac{s_1 s_2}{32} \boldsymbol{\tilde h}_{101}\Bigg)\mathbfcal{M}_{0,0;0,0,1,1,1}-\frac{s_1 s_2}{32}\boldsymbol{\tilde h}_{102}q^2\mathbfcal{M}_{0,0;0,0,1,2,1}\Bigg)\Bigg]\,,\\&
   =\frac{m_1^2 m_2^2}{m_p^6}\Bigg(\Delta_{\mathfrak{22}}^{\textrm{poly}}+\Delta_{\mathfrak{15}}^{\textrm{log}}\log(x)\Bigg)
  \end{split}
\end{align}    
where, 
\begin{align}
    \begin{split}
      \Delta_{\mathfrak{22}}^{\textrm{poly}}= \frac{s_{1} s_{2} \left(-x^8-x^6+x^2+1\right)}{4096 \pi ^3 |\boldsymbol{b}|^2 x^2 \left(x^2-1\right)^2}\,,\Delta_{\mathfrak{15}}^{\textrm{log}}=-\frac{s_{1} s_{2} \left(x^2 (5-s_{1}s_{2})+2 x^4+2\right)}{1024 \pi ^3 |\boldsymbol{b}|^2 \left(x^2-1\right)^2}\,.
    \end{split}
\end{align}
The total contribution to the spinning part is
%\textcolor{black}{AB: Fill the blanks after fixing the vertex factors of the second diagram}
\begin{align}
\begin{split}
\chi\Big|_{\mathcal{S}^1}
=&\frac{m_1^2 m_2^2}{m_p^6}\Big(-i\frac{s_1s_2}{32}\Big)\int e^{-i q\cdot b}\hat\delta(q\cdot v_1)\hat\delta(q\cdot v_2)\Bigg[\Bigg(\boldsymbol{\tilde h}_{103}\mathbfcal{M}_{0,0;0,0,1,1,1}
+\boldsymbol{\tilde h}_{104}q^2\mathbfcal{M}_{0,0;0,0,2,1,1}\\
&+\boldsymbol{\tilde h}_{105}q^2\mathbfcal{M}_{0,0;0,0,1,2,1}\Bigg)\big(q\cdot\mathcal{S}_1\cdot v_2\big)\\
&+\Bigg((\boldsymbol{\tilde h}_{106}+\boldsymbol{\tilde h}_{109})\mathbfcal{M}_{0,0;0,0,1,1,1}
+(\boldsymbol{\tilde h}_{107}+\boldsymbol{\tilde h}_{110})q^2\mathbfcal{M}_{0,0;0,0,2,1,1}\\
&+(\boldsymbol{\tilde h}_{108}+\boldsymbol{\tilde h}_{111})q^2\mathbfcal{M}_{0,0;0,0,1,2,1}\Bigg)\big(q\cdot\mathcal{S}_2\cdot v_1\big)\Bigg]\,,\\
=&\frac{m_1^2 m_2^2}{m_p^6}\Bigg\{\Bigg(\Delta_{\mathfrak{23}}^{\textrm{poly}}+\Delta_{\mathfrak{16}}^{\textrm{log}}\log(x)\Bigg)(\hat b\cdot \mathcal{S}_1\cdot v_2)\\
&\qquad\qquad\qquad\;
+\Bigg(\Delta_{\mathfrak{24}}^{\textrm{Poly}}+\Delta_{\mathfrak{17}}^{{\textrm{log}}}\log(x)\Bigg)(\hat b\cdot \mathcal{S}_2\cdot v_1)\Bigg\}\,.
\end{split}
\end{align}
where, 
\begin{align}
    \begin{split}
   & \Delta_{\mathfrak{23}}^{\textrm{poly}}=-\frac{s_{1}s_{2} \left(x^{10}+9 x^8+32 x^6-32 x^4-9 x^2-1\right)}{2048 \pi ^3 |\boldsymbol{b}|^3 x \left(x^2-1\right)^4 }\,,\Delta_{\mathfrak{16}}^{\textrm{log}}=-\frac{s_{1} s_{2} x \left(x^2+1\right)^3}{128 \pi ^3 |\boldsymbol{b}|^3 \left(x^2-1\right)^4 }\,,
   \\& \Delta_{\mathfrak{24}}^{\textrm{poly}}=\frac{s_{1} s_{2} \left(113 x^6-32 x^5+555 x^4+555 x^2-32 x+113\right)}{32768 \pi ^3 |\boldsymbol{b}|^3 \left(x^2-1\right)^3 }\,,\\&\Delta_{\mathfrak{17}}^{{\textrm{log}}}=\frac{s_{1} s_{2} \left(17 x^8-8 x^7+153 x^6-8 x^5+328 x^4-8 x^3+153 x^2-8 x+17\right)}{8192 \pi ^3 |\boldsymbol{b}|^3 \left(x^2-1\right)^4}\,.  \label{eqt17} 
     \end{split}
\end{align} 
%The amplitude corresponding to the above diagrams has the following form (after relabelling $k_1\to \ell_2,\,k_3\to \ell_1,\,k_2\to \ell_2-q$),
%\begin{align}
%    \begin{split}
%        i\chi=i\frac{m_1^2 s_1^2 m_2^2 s_2^2}{32 m_p^6} \int e^{-i q\cdot b}\hat\delta(q\cdot v_1)\hat\delta(q\cdot v_2)\int \frac{\hat\delta(\ell_1\cdot v_2)\hat\delta(\ell_2\cdot v_1)}{\ell_1^2\,\ell_2^2\,(\ell_1+\ell_2-q)^2\,(\ell_2-q)^2}\mathcal{N}
%    \end{split}
%\end{align}
%where the numerator of the first diagram takes the form (note that we use all the delta function constraint while evaluating the numerator),
%\begin{align}
%    \begin{split}
 %       \mathcal{N}\Big|_{(1)}=-(v_2^\mu v_2^\nu+i (q-\ell_1-\ell_2)\cdot %\mathcal{S}_2^{(\mu}v_2^{\nu)})\,P_{\mu\nu,\alpha\beta}(\ell_2-q)^\alpha \ell_1^\beta&\to \frac{-(\ell_1+\ell_2-q)^2+(\ell_2-q)^2+\ell_1^2}{4-2d}\\ &+\frac{i}{2}(\ell_2\cdot v_2)(\ell_1+\ell_2-q)\cdot \mathcal{S}_2\cdot \ell_1
%    \end{split}
%\end{align}
%Therefore, the spinless eikonal phase coming from the first diagram takes the following form,
%\begin{align}
 %   \begin{split}
  %     i\chi\Big|_{\mathcal{S}^{0}}&=i\frac{m_1^2 s_1^2 m_2^2 s_2^2}{32 %m_p^6} \int e^{-i q\cdot b}\hat\delta(q\cdot v_1)\hat\delta(q\cdot v_2)\int_{\ell_1,\ell_2} \frac{\hat\delta(\ell_1\cdot v_2)\hat\delta(\ell_2\cdot v_1)}{\ell_1^2\,\ell_2^2\,(\ell_1+\ell_2-q)^2\,(\ell_2-q)^2}\Bigg[\frac{-(\ell_1+\ell_2-q)^2+(\ell_2-q)^2+\ell_1^2}{4-2d}\Bigg],\\ &
%       =i\frac{m_1^2 s_1^2 m_2^2 s_2^2}{32 m_p^6}\frac{1}{4(1-\epsilon)}\int e^{-i q\cdot b}\hat\delta(q\cdot v_1)\hat\delta(q\cdot v_2)\mathbfcal{M}_{0,0;0,0,1,1,1}
%    \end{split}
%\end{align}
%Furthermore the spinning counterpart has the following contribution,
%\begin{align}
%    \begin{split}
%        i\chi\Big|_{\mathcal{S}^{1}}&=i\frac{m_1^2 s_1^2 m_2^2 s_2^2}{32 m_p^6}\times\frac{i}{2} \int e^{-iq\cdot b}\hat\delta(q\cdot v_1)\hat\delta(q\cdot v_2)\int \frac{\hat\delta(\ell_1\cdot v_2)\hat\delta(\ell_2\cdot v_1)\,
 %       \ell_2\cdot v_2}{\ell_1^2\, \ell_2^2\, (\ell_1+\ell_2-q)^2\,(\ell_2-q)^2}(\ell_2-q)\cdot \mathcal{S}_2\cdot \ell_1
%    \end{split}
%\end{align}
%where we will encounter the following tensor integral,
%\begin{align}
%    \begin{split}
%        \boldsymbol{\Upsilon}^{\eta\delta}=\int %\frac{\hat\delta(\ell_1\cdot v_2)\hat\delta(\ell_2\cdot v_1)(\ell_2-q)^{\eta}\ell_1^\delta}{\ell_1^2\,\ell_2^2\,(\ell_1+\ell_2-q)^2\,(\ell_2-q)^2}\ell_2\cdot v_2:=\kappa_1\,q^{[\eta]}v_1^{\delta]}+\kappa_2\,q^{[\eta]}v_2^{\delta]}
%    \end{split}
%\end{align}
%where,
%\begin{align}
%&\kappa _1=0,\\ &
%\kappa _2=\boldsymbol{h_{122}}\,\mathbfcal{M}_{0,0;0,0,1,1,1}+\boldsymbol{h_{123}}\,q^2\,\mathbfcal{M}_{0,0;0,0,2,1,1}+\boldsymbol{h_{124}}q^2\,\mathbfcal{M}_{0,0;0,0,1,2,1}
%\end{align}
%From the second diagram we will have the following contribution to the eikonal phase,
%\begin{align}
%    \begin{split}
%        \chi\sim \frac{m_1^2 m_2^2}{m_p^6}\Big(-i\frac{s_1s_2}{32})\int e^{-i q\cdot b}\hat\delta(q\cdot v_1)\hat\delta(q\cdot v_2)\int_{\ell_1,\ell_2}\frac{\hat\delta(\ell_1\cdot v_2)\hat\delta(\ell_2\cdot v_2)}{\ell_1^2\,\ell_2^2\,(\ell_1+\ell_2-q)^2\,(\ell_2-q)^2}\mathcal{N},\label{4.84}
%   \end{split}
%\end{align}
%where the numerator $\mathcal{N}$ takes the following form,
%\begin{align}
%    \begin{split}
  %      \mathcal{N}&=\frac{(q-\ell_1)\cdot \mathcal{S}_2\cdot v_1}{2(d-%2)}\Bigg[\gamma (\ell_1+\ell_2-q)^2-\gamma (\ell_2-q)^2-\gamma \ell_1^2\Bigg]+\frac{\gamma}{2}(q-\ell_1)\cdot \mathcal{S}_2\cdot q\,(\ell_1\cdot v_1)+\frac{(q-\ell_1)\cdot \mathcal{S}_2\cdot \ell_1}{2(d-2)}\,\ell_2\cdot v_2\\ &
 %       \hspace{0.8 cm}-\frac{\gamma}{2}(q-\ell_1)\cdot \mathcal{S}_2\cdot \ell_2\,\ell_1\cdot v_1\label{4.87q}
 %   \end{split}
%\end{align}
%Note that, because of the presence of graviton two point worldline vertex, the diagram only contributes to the spinning eikonal and also the spin of the first black hole does not come in to the play at first order in spin. We compute the phase by calculating the the two loop integral given in \eqref{4.84}.
%\begin{align}
 %   \begin{split}
 %       \mathcal{P}_1^\eta&=\gamma\int \frac{\hat\delta(\ell_1\cdot v_2)\hat\delta(\ell_2\cdot v_1)\,(q-\ell_1)^\eta}{\ell_1^2\,\ell_2^2\,(\ell_1+\ell_2-q)^2\,(\ell_2-q)^2}\Bigg[(\ell_1+\ell_2-q)^2-(\ell_2-q)^2-\ell_1^2\Bigg]\\ &
 %       =\gamma\Bigg[\boldsymbol{h_{125}}\,\mathbfcal{M}_{0,0;0,0,1,1,1}+\boldsymbol{h_{126}}\,q^2 \,\mathbfcal{M}_{0,0;0,0,2,1,1}
%       + \boldsymbol{h_{127}}\,q^2\mathbfcal{M}_{0,0;0,0,1,2,1}\Bigg]\,q^\eta\\ &
%   \hspace{0cm}     \mathcal{P}_2^\eta=\int \frac{\hat\delta(\ell_1\cdot v_2)\hat\delta(\ell_2\cdot v_1)\,\ell_1^\eta}{\ell_1^2\,\ell_2^2\,(\ell_1+\ell_2-q)^2\,(\ell_2-q)^2}\Bigg[\frac{\gamma}{2}\ell_1\cdot v_1+\frac{1}{2(d-2)}\ell_2\cdot v_2\Bigg]\\ &
 %  =\Bigg[\boldsymbol{h_{128}}\mathbfcal{M}_{0,0;0,0,1,1,1}+\boldsymbol{h_{129}}\,q^2\mathbfcal{M}_{0,0;0,0,2,1,1}+\boldsymbol{h_{130}}\,q^2\mathbfcal{M}_{0,0;0,0,1,2,1}\,\Bigg]\frac{\gamma v_2^\eta-v_1^\eta}{\gamma^2-1}\\ &
%  \hspace{0cm} %\mathcal{P}_3^\eta=\int \frac{\hat\delta(\ell_1\cdot v_2)\hat\delta(\ell_2\cdot v_1)\,\ell_2^\eta}{\ell_1^2\,\ell_2^2\,(\ell_1+\ell_2-q)^2\,(\ell_2-q)^2}\Bigg[-\frac{\gamma}{2}(\ell_2-q)^2+\frac{\gamma}{4}(\ell_1+\ell_2-q)^2-\ell_2^2-\ell_1^2+q^2\Bigg]\\ &
  %=\Bigg[\boldsymbol{h_{130}}\,\mathbfcal{M}_{0,0;0,0,1,1,1}+\boldsymbol{h_{131}}\,q^2\mathbfcal{M}_{0,0;0,0,2,1,1}
  %+\boldsymbol{h_{132}}\,q^2\mathbfcal{M}_{0,0;0,0,1,2,1}\Bigg]q^\eta\\ &
  %\hspace{0cm}\mathcal{P}_4^\eta=\frac{\gamma}{2}\int \frac{\hat\delta(\ell_1\cdot v_2)\hat\delta(\ell_2\cdot v_1)\,\ell_2^\eta}{\ell_1^2\,\ell_2^2\,(\ell_1+\ell_2-q)^2\,(\ell_2-q)^2}\ell_1\cdot v_1\\ &
  %=\frac{\gamma v_1^\eta-v_2^\eta}{\gamma^2-1}\Bigg[\boldsymbol{h_{133}}\,\mathbfcal{M}_{0,0;0,0,1,1,1}+\boldsymbol{h_{134}}\,\,q^2\mathbfcal{M}_{0,0;0,0,2,1,1}+\boldsymbol{h_{135}}\,q^2\mathbfcal{M}_{0,0;0,0,1,2,1}\Bigg]\label{4.87}
 %   \end{split}
%\end{align}
%Apart from the vector integrals in \eqref{4.87q} we will face the following two index tensor integral,
%\begin{align}
 %   \begin{split}
  %      \mathcal{P}_3^{\eta\delta}&=\int \frac{\hat\delta(\ell_1\cdot %v_2)\hat\delta(\ell_2\cdot v_1)\,(q-\ell_1)^\eta\,\ell_2^\delta}{\ell_1^2\,\ell_2^2\,(\ell_1+\ell_2-q)^2\,(\ell_2-q)^2}\ell_1\cdot v_1\\ &
 %       =\Bigg[h_{136}\,\mathbfcal{M}_{0,0;0,0,1,1,1}+\boldsymbol{h_{137}}\,\,q^2\mathbfcal{M}_{0,0;0,0,2,1,1}+\boldsymbol{h_{138}}\,q^2\mathbfcal{M}_{0,0;0,0,1,2,1}\Bigg]q^{[\eta}v_2^{\delta]}+\Bigg[\boldsymbol{h_{139}}\,\mathbfcal{M}_{0,0;0,0,1,1,1}\\ &+\boldsymbol{h_{140}}\,q^2\mathbfcal{M}_{0,0;0,0,2,1,1}+\boldsymbol{h_{141}}\,q^2\mathbfcal{M}_{0,0;0,0,1,2,1}\Bigg]q^{[\eta}v_1^{\delta]}
 %   \end{split}
%\end{align}
%We now move on to the calculation of eikonal phase coming from the three graviton vertex. Contribution to the eikonal phase has the following form,
%\begin{align}
%\begin{split}
%    \chi =-i\frac{m_1^2 s_1  m_2^2 s_2}{32 m_p^6}  \int e^{-iq\cdot b}\hat\delta(q\cdot v_1)\hat\delta(q\cdot v_2)\underbrace{\int_{\ell_1,\ell_2}\frac{\hat\delta(\ell_1\cdot v_2)\hat\delta(\ell_2\cdot v_1)}{\ell_1^2\,\ell_2^2\,(\ell_1+\ell_2-q)^2(\ell_2-q)^2}\mathcal{N}}_{\mathcal{M}=\mathcal{M}[\mathcal{S}^0]+\mathcal{M}[\mathcal{S}^1 ]},
%\end{split}
%\end{align}
%where the numerator $\mathcal{N}=\mathcal{N}^{\mathcal{S}^0}+\mathcal{N}^{\mathcal{S}^1}$ can be decomposed into spinless and spinning contribution,
%\begin{align}
%    \begin{split}
%       &\mathcal{N}^{\mathcal{S}^0}=\frac{1}{(d-2)^3}\Bigg[2 (d-2) \left(\gamma ^2 (d-2)^2-1\right)(\ell_2\cdot v_2)^2+3 (d-3) (d-2)^2(\ell_1\cdot v_1)^2-2 \gamma  (d-4) (d-2)^2(\ell_1\cdot v_1)(\ell_2\cdot v_2)\Bigg]\\ &\hspace{1.1 cm}+(\gamma ^2 (3 d-10) (d-2)^2+d)(\ell_2-q)^2\\ &
 %      \mathcal{N}^{\mathcal{S}^1}=-\frac{i}{d-2}(\ell_2-q)\cdot \mathcal{S}_1\cdot v_2\,\Bigg[( \gamma  (10-3 d)(\ell_2-q)^2-2 \gamma  (d-2)(\ell_2\cdot v_2)^2+(d-4)(\ell_1\cdot v_1)(\ell_2\cdot v_2)\Bigg]\\ &
%    \hspace{1.1 cm}  -\frac{i}{2(d-2)}(\ell_2-q)\cdot \mathcal{S}_1\cdot %q\,\Bigg[3(d-3)(\ell_1\cdot v_1)-4\gamma(d-4)(\ell_2\cdot v_2)\Bigg]\\ &
%    \hspace{1.1 cm}
%    +\frac{i}{d-2}(\ell_2-q)\cdot \mathcal{S}_1\cdot \ell_1\Bigg[3(d-3)%(\ell_1\cdot v_1)-\gamma(d-4)(\ell_2\cdot v_2)\Bigg]\\ &
%    \\ &
%    \hspace{1.1 cm}
%    +\frac{i}{d-2}(\ell_2-q)\cdot \mathcal{S}_1\cdot \ell_2\Bigg[3(d-3)(\ell_1\cdot v_1)-4\gamma(d-4)(\ell_2\cdot v_2)\Bigg] \\ &
 %   \hspace{1.1 cm}
%    +\frac{i}{2(d-2)}(\ell_1+\ell_2-q)\cdot \mathcal{S}_2\cdot v_1\Bigg[2(\ell_1\cdot v_1)^2-6 (d-3)(\ell_2\cdot v_2)^2+2 \gamma  (d-2)(\ell_1\cdot v_1)(\ell_2\cdot v_2)\\ &\hspace{1.1 cm}+\frac{8 \left(5 \gamma ^2+2\right)-3 \gamma ^2 d^3+\left(22 \gamma ^2+4\right) d^2-2 \left(26 \gamma ^2+9\right) d}{(d-2)^2}(\ell_2-q)^2\Bigg]\\ &
%    \hspace{1.1 cm}
%    -\frac{i}{2(d-2)^2}(\ell_1+\ell_2-q)\cdot \mathcal{S}_2\cdot q\Bigg[(4 \gamma ^2+\gamma ^2 d^2-\left(4 \gamma ^2+5\right) d+14)(\ell_1\cdot v_1)+2 \gamma  \left(d^2-6 d+8\right))(\ell_2\cdot v_2)\Bigg]\\ &
%    \hspace{1.1 cm}
 %   -\frac{i}{(d-2)^2}(\ell_1+\ell_2-q)\cdot\mathcal{S}_2\cdot %\ell_1\Bigg[(d-4)(\ell_1\cdot v_1)+2 \gamma  (d-2)(\ell_2\cdot v_2)\Bigg]\\ &
%  \hspace{1.1 cm}  +\frac{i}{2(d-2)^2}(\ell_1+\ell_2-q)\cdot\mathcal{S}_2\cdot \ell_2\Bigg[(4 \gamma ^2+\gamma ^2 d^2-\left(4 \gamma ^2+5\right) d+14)(\ell_1\cdot v_1)+(2 \gamma  \left(d^2-6 d+8\right))(\ell_2\cdot v_2)\Bigg]\\ &
%  \hspace{1.1 cm}+\frac{i}{2(d-2)}\ell_1\cdot \mathcal{S}_2\cdot v_1\,\Bigg[\gamma(3d-10)(\ell_2-q)^2\Bigg]+i \ell_1\cdot \mathcal{S}_2\cdot (q-\ell_2) \Bigg[\left(\frac{1}{d-2}-2 \gamma ^2\right)(\ell_2\cdot v_2)+\gamma(\ell_1\cdot v_1)\Bigg]
%       \end{split}
%\end{align}
%Therefore,
%\begin{align}
%    \begin{split}
%        \mathcal{M}[\mathcal{S}^0]=&\boldsymbol{h_{142}}\,\mathbfcal{M}_{0,0;0,0,1,1,1}+\boldsymbol{h_{143}}\,q^2\mathbfcal{M}_{0,0;0,0,2,1,1}+\boldsymbol{h_{144}}\,q^2\mathbfcal{M}_{0,0;0,0,1,2,1}
%    \end{split}
%\end{align}
%Now, in the spinning part of the eikonal phase we will encounter the following tensor integrals,
%\begin{align}
%    \begin{split}
%  \boldsymbol{  \zeta}^{\eta}_1&=\int \frac{\hat\delta(\ell_1\cdot v_2)\hat\delta(\ell_2\cdot v_1)\,(\ell_2-q)^{\eta}}{\ell_1^2\,\ell_2^2\,(\ell_1+\ell_2-q)^2\,(\ell_2-q)^2}\Bigg[( \gamma  (10-3 d)(\ell_2-q)^2-2 \gamma  (d-2)(\ell_2\cdot v_2)^2+(d-4)(\ell_1\cdot v_1)(\ell_2\cdot v_2)\Bigg]\\ &
%  = %\Big(\boldsymbol{h_{145}}\mathbfcal{M}_{0,0;0,0,1,1,1}+\boldsymbol{h_{146}}\,q^2\mathbfcal{M}_{0,0;0,0,2,1,1}+\boldsymbol{h_{147}}\,q^2\mathbfcal%{M}_{0,0;0,0,1,2,1}\Big)q^\eta\\ &
%  \hspace{0cm}\boldsymbol{\zeta}_2^{\eta}=\int \frac{\hat\delta(\ell_1\cdot v_2)\hat\delta(\ell_2\cdot v_1)(\ell_2-q)^\eta}{\ell_1^2\,\ell_2^2\,(\ell_1+\ell_2-q)^2\,(\ell_2-q)^2}\Bigg[3(d-3)(\ell_1\cdot v_1)-4\gamma(d-4)(\ell_2\cdot v_2)\Bigg]
%  \\ & 
%  =\frac{(\gamma v_1^\mu-v_2^\mu)}{\gamma^2-1}\Bigg[\boldsymbol{h_{148}}\mathbfcal{M}_{0,0;0,0,1,1,1}+\boldsymbol{h_{149}}\,q^2\mathbfcal{M}_{0,0;0,0,2,1,1}+\boldsymbol{h_{150}}\,q^2\,\mathbfcal{M}_{0,0;0,0,1,2,1}\Bigg]\\ &
%  \hspace{0cm}\boldsymbol{\zeta}^{\eta}_4=\int\frac{\hat\delta(\ell_1\cdot v_2)\hat\delta(\ell_2\cdot v_1)\ell_2^\eta}{\ell_1^2\,\ell_2^2\,(\ell_1+\ell_2-q)^2\,(\ell_2-q)^2}\Bigg[3(d-3)(\ell_1\cdot v_1)-4\gamma(d-4)(\ell_2\cdot v_2)\Bigg]\equiv\boldsymbol{\zeta}_2^\eta
%  \\ &
%  \hspace{0cm}\boldsymbol{\zeta}^{\eta}_5= \int\frac{\hat\delta(\ell_1\cdot v_2)\hat\delta(\ell_2\cdot v_1)(\ell_1+\ell_2-q)^\eta}{\ell_1^2\,\ell_2^2\,(\ell_1+\ell_2-q)^2\,(\ell_2-q)^2}\Bigg[2(\ell_1\cdot v_1)^2-6 (d-3)(\ell_2\cdot v_2)^2+2 \gamma  (d-2)(\ell_1\cdot v_1)(\ell_2\cdot v_2)\\ &\hspace{1.1 cm}+\frac{8 \left(5 \gamma ^2+2\right)-3 \gamma ^2 d^3+\left(22 \gamma ^2+4\right) d^2-2 \left(26 \gamma ^2+9\right) d}{(d-2)^2}(\ell_2-q)^2\Bigg]\\ &
       %=\Bigg(\boldsymbol{h_{151}}\,\mathbfcal{M}_{0,0;0,0,1,1,1}+\boldsymbol{h_{152}}\,q^2\,\mathbfcal{M}_{0,0;0,0,2,1,1}+\boldsymbol{h_{153}}\,q^2\,\mathbfcal{M}_{0,0;0,0,1,2,1}\Bigg)\,q^\eta
%       \\ &
%       \hspace{0cm}
%\boldsymbol{\zeta}_{6}^{\eta}=\int \frac{\hat\delta(\ell_1\cdot v_2)\hat\delta(\ell_2\cdot v_1)(\ell_1+\ell_2-q)^\eta}{\ell_1^2\,\ell_2^2\,(\ell_1+\ell_2-q)^2\,(\ell_2-q)^2}\Bigg[(4 \gamma ^2+\gamma ^2 d^2-\left(4 \gamma ^2+5\right) d+14)(\ell_1\cdot v_1)+2 \gamma  \left(d^2-6 d+8\right))(\ell_2\cdot v_2)\Bigg]\\ &
%=\Bigg[\left(\boldsymbol{h_{154}}\,\mathbfcal{M}_{0,0;0,0,1,1,1}+\boldsymbol{h_{155}}\,q^2\,\mathbfcal{M}_{0,0;0,0,2,1,1}+\boldsymbol{h_{156}}\,q^2\,\mathbfcal{M}_{0,0;0,0,1,2,1}\right)v_1^\eta+\Big(\boldsymbol{h_{157}}\,\mathbfcal{M}_{0,0;0,0,1,1,1}\\ &+\boldsymbol{h_{158}}\,q^2\,\mathbfcal{M}_{0,0;0,0,2,1,1}+\boldsymbol{h_{159}}\,q^2\,\mathbfcal{M}_{0,0;0,0,1,2,1}\Big)v_2^\eta \Bigg]\\ &
%  \hspace{0cm}  \boldsymbol{\zeta}_{9}^{\eta}= \int \frac{\hat\delta(\ell_1\cdot v_2)\hat\delta(\ell_2\cdot v_1)\ell_1^\eta}{\ell_1^2\,\ell_2^2\,(\ell_1+\ell_2-q)^2\,(\ell_2-q)^2}\Bigg[\gamma (3d-10)(\ell_2-q)^2\Bigg]\\ &=\Bigg(\boldsymbol{h_{160}}\,\mathbfcal{M}_{0,0;0,0,1,1,1}+\boldsymbol{h_{161}}\,q^2\mathbfcal{M}_{0,0;0,0,2,1,1}+\boldsymbol{\mathfrak{h}}_{4}q^2\mathbfcal{M}_{0,0,0,0,1,2,1}\Bigg)q^\eta
 %   \end{split}
%\end{align}
%Apart from the vector integrals we have the following tensor integrals,
%\begin{align}
%    \begin{split}
%\boldsymbol{\zeta}_3^{\eta\delta}&=\int \frac{\hat\delta(\ell_1\cdot v_2)\hat\delta(\ell_2\cdot  v_1)(\ell_2-q)^{\mu}\ell_1^\nu}{\ell_1^2\,\ell_2^2\,(\ell_1+\ell_2-q)^2\,(\ell_2-q)^2} \Bigg[3(d-3)(\ell_1\cdot v_1)-\gamma (d-4)(\ell_2\cdot v_2)\Bigg] \\ &
%=\Bigg(\boldsymbol{h_{168}}\,\mathbfcal{M}_{0,0;0,0,1,1,1}+\boldsymbol{h_{169}}\,q^2\,\mathbfcal{M}_{0,0;0,0,2,1,1}
%+\boldsymbol{h_{170}}\,q^2\mathbfcal{M}_{0,0;0,0,1,2,1}\Bigg)q^{[\mu}v_1^{\%nu]}\\ &
%+\Bigg(\boldsymbol{h_{171}}\,\mathbfcal{M}_{0,0;0,0,1,1,1}+\boldsymbol{h_{1%72}}\,q^2\,\mathbfcal{M}_{0,0;0,0,2,1,1}+\boldsymbol{h_{173}}\,q^2\mathbfcal{M}_{0,0;0,0,1,2,1}\Bigg)q^{[\mu}v_2^{\nu]}
%        \\ &\hspace{0cm}\boldsymbol{\zeta}_7^{\eta\delta}=\int \frac{\hat\delta(\ell_1\cdot v_2)\hat\delta(\ell_2\cdot v_1)(\ell_1+\ell_2-q)^\mu \ell_1^\nu}{\ell_1^2\,\ell_2^2\,(\ell_1+\ell_2-q)^2\,(\ell_2-q)^2}\,\Bigg[(d-4)(\ell_1\cdot v_1)+2 \gamma  (d-2)(\ell_2\cdot v_2)\Bigg]\\ &
%=\Bigg(\boldsymbol{h_{174}}\,\mathbfcal{M}_{0,0;0,0,1,1,1}+\boldsymbol{h_{175}}\,q^2\,\mathbfcal{M}_{0,0;0,0,2,1,1}
%+\boldsymbol{h_{176}}\,q^2\mathbfcal{M}_{0,0;0,0,1,2,1}\Bigg)q^{[\eta}v_1^{\delta]}\\ &+\Bigg(\boldsymbol{h_{177}}\,\mathbfcal{M}_{0,0;0,0,1,1,1}+\boldsymbol{h_{178}}\,q^2\,\mathbfcal{M}_{0,0;0,0,2,1,1}
%+\boldsymbol{h_{179}}\,q^2\mathbfcal{M}_{0,0;0,0,1,2,1}\Bigg)q^{[\eta}v_2^{\delta]}
%        \\ & \hspace{0cm}\boldsymbol{\zeta}_{8}^{\eta\delta}=\int \frac{\hat\delta(\ell_1\cdot v_2)\hat\delta(\ell_2\cdot v_1)(\ell_1+\ell_2-q)^
 %      \eta\,\ell_2^\delta}{\ell_1^2\,\ell_2^2\,(\ell_1+\ell_2-q)^2\,(\ell_2-q)^2}\Bigg[(4 \gamma ^2+\gamma ^2 d^2-\left(4 \gamma ^2+5\right) d+14)(\ell_1\cdot v_1)+(2 \gamma  \left(d^2-6 d+8\right))(\ell_2\cdot v_2)\Bigg]\\ &
       %=\Bigg(\boldsymbol{h_{180}}\,\mathbfcal{M}_{0,0;0,0,1,1,1}+\boldsymbol{h_{181}}\,q^2\mathbfcal{M}_{0,0;0,0,2,1,1}+\boldsymbol{h_{182}}\,q^2\mathbfcal{M}_{0,0;0,0,2,1,1}\Bigg)q^{[\mu}v_1^{\nu]}\\ &+\Bigg(\boldsymbol{h_{183}}\,\mathbfcal{M}_{0,0;0,0,1,1,1}+\boldsymbol{h_{184}}\,q^2\mathbfcal{M}_{0,0;0,0,2,1,1}+\boldsymbol{h_{185}}\,q^2\mathbfcal{M}_{0,0;0,0,2,1,1}\Bigg)q^{[\mu}v_2^{\nu]} \\ & \hspace{0cm}\boldsymbol{\zeta}_{10}^{\eta\delta}=\int \frac{\hat\delta(\ell_1\cdot v_2)\hat\delta(\ell_2\cdot v_1)\ell_1^\eta\,(\ell_2-q)^\delta}{\ell_1^2\,\ell_2^2\,(\ell_1+\ell_2-q)^2\,(\ell_2-q)^2}\Bigg[\left(\frac{1}{d-2}-2 \gamma ^2\right)(\ell_2\cdot v_2)+\gamma(\ell_1\cdot v_1)\Bigg]\\ &
       %=\Bigg(\boldsymbol{h_{162}}\,\mathbfcal{M}_{0,0;0,0,1,1,1}+\boldsymbol{h_{163}}\,q^2\mathbfcal{M}_{0,0;0,0,2,1,1}+\boldsymbol{h_{164}}\,q^2\mathbfcal{M}_{0,0;0,0,2,1,1}\Bigg)q^{[\mu}v_1^{\nu]}\\ &+\Bigg(\boldsymbol{h_{165}}\,\mathbfcal{M}_{0,0;0,0,1,1,1}+\boldsymbol{h_{166}}\,q^2\mathbfcal{M}_{0,0;0,0,2,1,1}+\boldsymbol{h_{167}}\,q^2\mathbfcal{M}_{0,0;0,0,2,1,1}\Bigg)q^{[\mu}v_2^{\nu]}
 %   \end{split}
%\end{align}
\textbullet $\,\,$ We have the following topologies both involving scalar-graviton and graviton 3-point vertices.
\begin{center}
\includegraphics[width=0.38\linewidth]{Y_type.png}
\end{center}
We can easily check that, after evaluating the corresponding integrals for both the spinless and spin-dependent parts, they vanish identically up to $\mathcal{O}(\epsilon^0)\,.$ Hence, this class of diagrams does not contribute to the eikonal phase in our case. 
%The contribution to the eikonal phase from the first diagram has the following form,
%\begin{align}
%    \begin{split}
  %      \chi&=-i\frac{m_1^2 s_1^2 m_2^2 s_2^2}{64 m_p^6}\int e^{-iq\cdot b}\frac{\hat\delta(q\cdot v_1)\hat\delta(q\cdot v_2)}{q^2}\int \frac{\hat\delta(\ell_1\cdot v_2)\hat\delta(\ell_2\cdot v_1)}{\ell_1^2\,\ell_2^2\,(\ell_1-q)^2\,(\ell_2-q)^2}\Big[(\ell_1+\ell_2-q)^4+q^2(\ell_1+\ell_2-q)^2+\frac{d-4}{2(d-2)}q^4\Big]\\ &
 %       =-i\frac{m_1^2 s_1^2 m_2^2 s_2^2}{64 m_p^6}\int e^{-iq\cdot b}{\hat\delta(q\cdot v_1)\hat\delta(q\cdot v_2)}\,q^2\,\boldsymbol{h_{186}}\,\mathbfcal{M}_{0,0;1,1,0,1,1}
 %   \end{split}
%\end{align}
%The second diagram involving one graviton and one scalar-graviton three point vertices contributes to the spinless eikonal phase as,
%\begin{align}
%    \begin{split}
  %      \chi=&\frac{m_1^2 s_1^2 m_2^2}{64 m_p^6}\frac{1}{8(d-2)^2} \int e^{-iq\cdot b}\,\frac{\hat\delta(q\cdot v_1)\hat\delta(q\cdot v_2)}{q^2}\int \frac{\hat\delta(\ell_1\cdot v_2)\hat\delta(\ell_2\cdot v_1)}{\ell_1^2\,\ell_2^2\,(\ell_1-q)\,(\ell_2-q)^2}\Bigg[ 6(d^2-5d+6)(\ell_1+\ell_2-q)^4\\ &+6(d^2-5d+6)q^2(\ell_1+\ell_2-q)^2+8(d^2-8d+12)q^2(\ell_2\cdot v_2)^2+(d^2-9d+10)q^4\Bigg]\\ &
 %      =\frac{m_1^2 s_1^2 m_2^2}{64 m_p^6}\int e^{-i q\cdot b}\hat\delta(q\cdot v_1)\hat\delta(q\cdot v_2)\,\boldsymbol{h_{187}}\,q^2\,\mathbfcal{M}_{0,0;1,1,0,1,1}
%    \end{split}
%\end{align}
%The spinning contribution has the following form,
%\begin{align}
 %   \begin{split}
  %      \chi^{\mathcal{S}^1}&=\frac{m_1^2 s_1^2 m_2^2}{64 m_p^6} \int e^{-%i q\cdot b}\frac{\hat\delta(q\cdot v_2)\hat\delta(q\cdot v_1)}{q^2}\int_{\ell_1,\ell_2}\frac{\hat\delta(\ell_1\cdot v_2)\hat\delta(\ell_2\cdot v_1)}{\ell_1^2\,\ell_2^2\,(\ell_1-q)^2\,(\ell_2-q)^2}\mathcal{N}
 %   \end{split}
%\end{align}
%where the numerator has the following form,
%\begin{align}
 %   \begin{split}
 %       \mathcal{N}&=\frac{-i}{4(d-2)}\Bigg[(\ell_1\cdot \mathcal{S}_2\cdot q )d(\ell_2\cdot v_2)\Big[2(\ell_1+\ell_2-q)^2+q^2\Big]+q\cdot \mathcal{S}_2\cdot \ell_2\Big[(2d-12)q^2(\ell_2\cdot v_2)\Big]\Bigg]
%    \end{split}
%\end{align}
%We have the following vector integrals to do,
%\begin{align}
%    \begin{split}
  %      \mathbf{M}_1^\eta&=d\int \frac{\hat\delta(\ell_1\cdot %v_2)\hat\delta(\ell_2\cdot v_1)(\ell_2\cdot v_2)\ell_1^\eta}{\ell_1^2\,\ell_2^2\,(\ell_1-q)^2\,(\ell_2-q)^2}\Bigg[\left(2(\ell_1+
  %      \ell_2-q)^2+q^2\right)\Bigg]\\ &
  %      =\frac{d}{\gamma^2-1}(\gamma v_2^\eta-v_1^\eta) \Bigg[\frac{\gamma  \left(\gamma ^2-1\right)}{16 (\epsilon -1)^2}\mathbfcal{M}_{0,0;1,1,0,1,1}\Bigg](q^2)^2
  %      \\ &
%        \hspace{0cm}\mathbf{M}_2^\eta=2(d-6)q^2\int %\frac{\hat\delta(\ell_1\cdot v_2)\hat\delta(\ell_2\cdot v_1)(\ell_2\cdot v_2)\ell_2^\eta}{\ell_1^2\,\ell_2^2\,(\ell_1-q)^2\,(\ell_2-q)^2}\\ &
 %       =2(d-6)\frac{\gamma v_2^\eta-v_1^\eta}{\gamma^2-1}\Bigg[-\frac{\gamma ^2-1}{4 (\epsilon -1)}\mathbfcal{M}_{0,0;1,1,0,1,1}\Bigg](q^2)^2
%    \end{split}
%\end{align}
%Therefore,
%\begin{align}
%    \begin{split}
 %       \chi^{\mathcal{S}^1}\sim 
 %   \end{split}
%\end{align}
\par
\textbullet $\,\,$
We now discuss the scalar-graviton diagrams where we have spin contribution coming from the scalar-fermion vertex. The diagrams look like the following. 
\begin{center}
\includegraphics[width=0.38\linewidth]{3PM_5.png}
\end{center}
One may notice that, unlike other diagrams, this one consists of $\bar\Psi^\mu\Psi^\nu$. To recover the spin, one needs to take the spinor flow in the opposite direction. So, for this particular diagram, we need to take the antisymmetric combination of $\bar\Psi^\mu\Psi^\nu$ to recover the classical spin.  Moreover, we can easily check that the contribution coming from the first diagram to the eikonal phase gets cancelled by the second diagram. Hence, together, they don't contribute. %\textcolor{black}{AB: Crosscheck once!}
%The contribution to the eikonal phase from the first diagram has the following form,
%\begin{align}
%    \begin{split}
 %       i\chi=-i\frac{ m_1 m_2^3 s_1^2 s_2^2}{32 m_p^6} \int_{k_i}&e^{-i(k_1+k_2+k_3)\cdot b}{\hat\delta(k_1\cdot v_2)\hat\delta(k_2\cdot v_2)\hat\delta(k_3\cdot v_2)\hat\delta(\omega_1-k_1\cdot v_1)\hat\delta(-k_2\cdot v_1-\omega_1+\omega_2)\hat\delta(\omega_2+k_3\cdot v_1)}\\ &\hspace{0 cm}
%        \times\frac{\mathcal{N}}{\omega_1\,\omega_2\,k_1^2\, k_2^2\, k_3^2}\,,
%    \end{split}
%\end{align}
%where the numerator is given by,
%\begin{align}
 %   \begin{split}
 %       \mathcal{N}= \bar\Psi_{\eta}\omega_1\omega_2\,\Big(-k_{3[\rho}\delta_{\sigma]}^{(\mu}v_1^{\nu)}\Big)\Psi^{\sigma}\eta^{\eta\delta}\eta_\delta^{\rho}P_{\mu\nu;\alpha \beta} v_2^\alpha v_2^\beta\to{i\gamma}(k_3\cdot \mathcal{S}_1\cdot v_2)\omega_1\omega_2\,.
%    \end{split}
%\end{align}
%Now, relabelling $k_1\to \ell_{1},\,k_3\to \ell_{2}$ and $k_1+k_2+k_2\to q\,,$ we have,\\
%\hfsetfillcolor{gray!10}
%\hfsetbordercolor{black!150}
%\begin{align}
%    \begin{split}
%     \tikzmarkin[disable rounded corners=false]{e11}(7.2,-0.7)(-0.2,1)  
 %     i  \chi\Big|_{\mathcal{S}^1}&= \frac{\gamma m_1 m_2^3 s_1^2 s_2^2}{8 m_p^6}\mathcal{S}_{1\eta\delta}v_2^\delta\int_{q}e^{-iq\cdot b}\hat\delta(q\cdot v_1)\hat\delta(q\cdot v_2)\int_{\ell_{1},\ell_{2}}\frac{\hat\delta(\ell_{1}\cdot v_2)\hat\delta(\ell_{2}\cdot v_2)\ell_{2}^\eta}{\ell_{1}^2 \ell_{2}^2 (\ell_{1}+\ell_{2}-q)^2}\,,\\ &
 %       =\frac{\gamma m_1 m_2^3 s_1^2 s_2^2}{96 m_p^6}\int_{q}e^{-iq\cdot b}\hat\delta(q\cdot v_1)\hat\delta(q\cdot v_2)\,{(q\cdot \mathcal{S}_1\cdot v_2)}\,\mathbfcal{L}_{0,0;0,0,1,1,1},\\&
 %       =i\Bigg(\frac{1+x^2}{1-x^2}\Bigg)\,\frac{ m_1 m_2^3 s_1^2 s_2^2}{1536  \pi^3 |\boldsymbol{b}|^3 m_p^6}\,\Bigg(\hat{b}\cdot\mathcal{S}_1\cdot v_2\Bigg)\,.
 %       \tikzmarkend{e11}
 %   \end{split}
%\end{align}
%The contribution to the Eikonal phase from the second diagram has the following form,
%\begin{align}
%    \begin{split}
%        \chi=-i\frac{ m_1 m_2^3 s_1^2 s_2^2}{16 m_p^6} \int e^{-iq\cdot b}\hat\delta(q\cdot v_1)\hat\delta(q\cdot v_2)\int \frac{\hat\delta(\ell_1\cdot v_2)\hat\delta(\ell_2\cdot v_2)\mathcal{N}}{\ell_1^2\,\ell_2^2\,(\ell_1+\ell_2-q)^2}\,,
%    \end{split}
%\end{align}
%where, the numerator $\mathcal{N}$ has the following form,
%\begin{align}
 %   \begin{split}
 %       \mathcal{N}=-\frac{i\gamma}{2}(q-\ell_1-\ell_2)\cdot\mathcal{S}_1\cdot v_2
 %   \end{split}
%\end{align}
%Therefore the Eikonal phase is given by,
%\begin{align}
  %  \begin{split}
  %      \chi&\sim  -\frac{i\gamma}{2}\mathcal{S}_{1\eta\delta}\,v_2^\delta\int e^{-iq\cdot b}\hat\delta(q\cdot v_1)\hat\delta(q\cdot v_2)\int \frac{\hat\delta(\ell_1\cdot v_2)\hat\delta(\ell_2\cdot v_2)(q-\ell_1-\ell_2)^{\eta}}{\ell_1^2\,\ell_2^2\,(\ell_1+\ell_2-q)^2}\\ &
 %       = -\frac{i\gamma}{6}\int e^{-iq\cdot b}\hat\delta(q\cdot v_1)\hat\delta(q\cdot v_2)\,(q\cdot \mathcal{S}_1\cdot v_2)\,\mathbfcal{L}_{0,0;0,0,1,1,1}
%    \end{split}
%\end{align}
\par
\textbullet $\,\,$ Another 3PM spinor diagram that will contribute has the topology of the following form,
\begin{center}
\includegraphics[width=0.38\linewidth]{3PM_6.png}
\end{center}
Let's start our computations from the first diagram.   The 3-point spin-spin-graviton interaction vertex which contributes to this term comes from the following part of the action,
\begin{align}
    \begin{split}
        S\subset& \,i\frac{m_i}{m_p}\int  d\tau\, \dot x^\mu \omega^{ab}_{\mu}\bar\psi_{a} \psi_{b}\,,\\ &
        = \frac{i m_i}{m_p}v^\mu\int d\tau\, \partial^{[a}h^{b]}_\mu \bar\psi_{[a}\psi_{b]} \xrightarrow[]{F.S.}-\frac{m_a}{m_p}\int_{k,\omega_i}e^{ik\cdot b}\hat\delta(k\cdot v-\omega_1+\omega_2) v^\mu\,k^{[a}h^{b]}_\mu (-k)\bar\psi_{[a}(-\omega_1)\,\psi_{b]}(-\omega_2)  \,.
        \end{split}
\end{align}
Therefore, the vertex factor is given by
\begin{align}
    V: e^{ik\cdot b}\hat\delta(k\cdot v-\omega_1+\omega_2)\,v_i^\mu \,k^a+\textrm{other 3 combinations}\,.
\end{align}
Schematically, the amplitude looks like
\begin{align}
    \begin{split}
       \mathcal{A}\sim\int_{k_i,\omega_i}\Big(\prod\boldsymbol{\delta}^{(6)}\Big)\bar\Psi^\eta\Psi^\sigma\omega_1\langle \psi_\eta\,\bar\psi_{[[a}\rangle_{\omega_1}k_{2[a}\langle  h_{b]\mu}h_{\alpha\beta}\rangle_{k_2} \langle \psi_{b]]}\,\bar\psi_\rho\rangle k_{3[\rho}\delta^{(\chi}_{\sigma]}v_1^{\pi)} \langle  h_{\chi\pi}h_{\kappa\theta}\rangle_{k_3}v_2^\alpha v_2^\beta v_2^\kappa v_2^\theta\,.\nonumber
    \end{split}
\end{align}
Here, $\prod \boldsymbol\delta^6$ stands for the product of all relevant delta functions along with $e^{ik\cdot b}$ coming from the worldline vertices. \textcolor{black}{Note that we take only the antisymmetric combination of $\eta,\sigma$ to relate the combination of background superfield to classical spin, which is eventually important to define the Pauli-Lubanski vector.} 
%Finally, the contribution to the eikonal phase has the following form, %\textcolor{black}{AB: add the contribution from the second term},
%\begin{align}
%    \begin{split}
 %      i \chi\Big|_{\mathcal{S}^1}&= -\frac{ m_1 m_2^3 s_1 s_2}{1024 m_p^6 (d-2)^2} \mathcal{S}_{1\eta \delta}\int e^{-iq\cdot b} \hat\delta(q\cdot v_1)\hat\delta(q\cdot v_2)\int_{\ell_{i}}\sum_{\mp}\frac{\hat\delta(\ell_{1}\cdot v_2)\hat\delta(\ell_{2}\cdot v_2) }{(\mp\ell_{2}\cdot v_1+i\varepsilon)\ell_{1}^2 \ell_{2}^2 (\ell_{1}+\ell_{2}-q)^2}\\ &
 %      \hspace{2 cm}\Bigg((\gamma ^2 \left(d^2-6 d+8\right)+1)(\ell_1-q)^\eta\ell_2^{\delta}-\gamma  (d-2)(\ell_2\cdot v_1)(q-\ell_1)^\eta\,v_2^\delta+\gamma  (d-2)(\ell_1\cdot v_1)\ell_2^\eta v_2^\delta\Bigg)\,.
%    \end{split}
%\end{align}
Then proceeding as before we get, 
\begin{align}
    \begin{split}
i\,\chi\Big|_{\mathcal{S}^1}&=\frac{ m_1 m_2^3 s_1 s_2}{16 m_p^6}\int e^{-iq\cdot b} \hat\delta(q\cdot v_1)\hat\delta(q\cdot v_2)\boldsymbol{\tilde h}_{112} \mathbfcal{L}_{0,0;0,0,1,1,1}\big(q\cdot\mathcal{S}_1\cdot v_2\big)\,.
       \end{split}
\end{align}
%We need to do the following two tensor integrals,
%\begin{align}
%    \begin{split}
%        \Sigma^{[\eta\delta]}_{(1)}&=\int_{\ell_{1},\ell_{2}}\frac{\hat\delta(\ell_{1}\cdot v_2)\hat\delta(\ell_{2}\cdot v_2)}{(\ell_{2}\cdot v_1)\,\ell_{1}^2\,\ell_{2}^2\,(\ell_{1}+\ell_{2}-q)^2}(\ell_1-q)^\eta\,\ell_2^\delta\,,\\ &
 %       =\frac{1}{\gamma^2-1}\mathbfcal{L}_{0,0;0,0,1,1,1}(q^{[\eta}v_1^{\delta]}-\gamma q^{[\eta}v_2^{\delta]}),\,\,
 %   \\ &    \hspace{0cm}\Sigma_{(2)}^{\delta}\equiv \int_{\ell_{1},\ell_{2}}\frac{\hat\delta(\ell_{1}\cdot v_2)\hat\delta(\ell_{2}\cdot v_2)}{\ell_{1}^2\,\ell_{2}^2\,(\ell_{1}+\ell_{2}-q)^2}(q-\ell_{1})^\delta\,,\\ &
 %       =\frac{2}{3}\mathbfcal{L}_{0,0;0,0,1,1,1}q^{\delta}\,,
 %         \\ &    \hspace{0cm}\Sigma_{(3)}^{\delta}\equiv \int_{\ell_{1},\ell_{2}}\frac{\hat\delta(\ell_{1}\cdot v_2)\hat\delta(\ell_{2}\cdot v_2)(\ell_1\cdot v_1)}{(\ell_2\cdot v_1-i\varepsilon)\ell_{1}^2\,\ell_{2}^2\,(\ell_{1}+\ell_{2}-q)^2}\ell_2^\delta\\ &
  %        =-\frac{1}{6}\mathbfcal{L}_{0,0;0,0,1,1,1}q^{\delta}.
 %       \label{4.72w}
%    \end{split}
%\end{align}
%Therefore it is quite easy to perform the $q$  integral using (\ref{4.72w}) and the result can be casted as,
%\hfsetfillcolor{gray!10}
%\hfsetbordercolor{black!150}
%\begin{align}
%\begin{split}
%\centering
%     \tikzmarkin[disable rounded corners=false]{e100}(1.4,-0.7)(-0.6,0.8) 
%i\chi\Big|_{\mathcal{S}^1}= -i\frac{m_1m_2^3s_1s_2^2}{12288 \,\pi^3 \,|\boldsymbol{b}|^3\,m_p^6}\Big(\frac{2 x}{1-x^2}\Big)\Big(\hat{b}\cdot \mathcal{S}_1\cdot v_2\Big)\,.
%\tikzmarkend{e100}
%\end{split}
%\end{align}
Correspondingly, the contribution from the second diagram has the following form,
\begin{align}
    \begin{split}
i\,\chi\Big|_{\mathcal{S}^1}&=\frac{ m_1 m_2^3 s_1 s_2}{8 m_p^6}\int e^{-iq\cdot b} \hat\delta(q\cdot v_1)\hat\delta(q\cdot v_2)\boldsymbol{\tilde h}_{113} \mathbfcal{L}_{0,0;0,0,1,1,1}\big(q\cdot\mathcal{S}_1\cdot v_2\big)\,.
       \end{split}
\end{align}
%where,
%\begin{align}
 %   \begin{split}
 %       \mathcal{N}=&\frac{i}{2}\ell_1\cdot \mathcal{S}_1\cdot \ell_2\,\Bigg[\frac{\gamma^2(d-4)(d-2)+1}{4(d-2)^2}\Bigg]+\frac{i}{2}\ell_2\cdot \mathcal{S}_1\cdot q\,\Bigg[\frac{\gamma^2(d^2-6d+8)+1}{4(d-2)^2}\Bigg]+\frac{i\gamma}{8(8-4d)}\ell_2\cdot \mathcal{S}_1\cdot v_2 (\ell_1\cdot v_1)\\ &+\frac{i\gamma}{8(d-2)}\ell_1\cdot \mathcal{S}_1\cdot v_2(\ell_2\cdot v_1)+\frac{i\gamma}{2(8-4d)}q\cdot \mathcal{S}_1\cdot v_2\,(\ell_2\cdot v_1)
%    \end{split}
%\end{align}
%where,
%\begin{align}
 %   \begin{split}
  %      &\mathbfcal{F}_1^{\eta\delta}=\int \frac{\hat\delta(\ell_1\cdot v_2)\hat\delta(\ell_2\cdot v_2)\,\ell_1^\eta\,\ell_2^\delta}{(\ell_2\cdot v_1+i\varepsilon)\,\ell_1^2\,\ell_2^2\,(\ell_1+\ell_2-q)^2}\\ &
 %       \hspace{0.8 cm}=\frac{1}{2(\gamma^2-1)}\left(\gamma q^{[\eta]}v_2^{\delta]}-q^{[\eta]}v_1^{\delta]}\right)\,\mathbfcal{L}_{0,0,0,0,1,1,1}\\ &
  %      \mathbfcal{F}_2^{\eta}=\int \frac{\hat\delta(\ell_1\cdot %v_2)\hat\delta(\ell_2\cdot v_2)\,\ell_2^\eta}{(\ell_2\cdot v_1+i\varepsilon)\,\ell_1^2\,\ell_2^2\,(\ell_1+\ell_2-q)^2}=\frac{\gamma v_2^\eta-v_1^\eta}{\gamma^2-1}\,\mathbfcal{L}_{0,0,0,0,1,1,1}\\ &
 %       \mathbfcal{F}_3^{\eta}=\int \frac{\hat\delta(\ell_1\cdot v_2)\hat\delta(\ell_2\cdot v_2)\,\ell_1^\eta}{(\ell_2\cdot v_1+i\varepsilon)\,\ell_1^2\,\ell_2^2\,(\ell_1+\ell_2-q)^2}=\frac{3-2 \epsilon }{12 (\epsilon -1)}\,\mathbfcal{L}_{0,0,0,0,1,1,1} q^\eta
%        \\ &\mathbfcal{F}_4^{\eta}=\int \frac{\hat\delta(\ell_1\cdot v_2)\hat\delta(\ell_2\cdot v_2)\,\ell_1^\eta}{(\ell_2\cdot v_1+i\varepsilon)\,\ell_1^2\,\ell_2^2\,(\ell_1+\ell_2-q)^2}=\frac{1}{3}\,\mathbfcal{L}_{0,0,0,0,1,1,1}\,q^\eta
%        \\ &
%        \mathbfcal{F}_5=\int \frac{\hat\delta(\ell_1\cdot v_2)\hat\delta(\ell_2\cdot v_2)}{\ell_1^2\,\ell_2^2\,(\ell_1+\ell_2-q)^2}=\mathbfcal{L}_{0,0,0,0,1,1,1}
%    \end{split}
%\end{align}
%\newpage
Then, the total contribution to the spinning part is,
\begin{align}
    \begin{split} \label{eqt18} 
 \chi\Big|_{\mathcal{S}^1}&=\frac{m_1m_2^3}{m_p^6}\Delta_{\mathfrak{25}}^{\textrm{poly}}\,,\quad\,\,\Delta_{\mathfrak{25}}^{\textrm{poly}}=\frac{ s_{1} s_{2}\,9 \left(x^2+1\right) \left(7 x^4+2 x^2+7\right)}{32768 \pi ^3 |\boldsymbol{b}|^3 \left(x^2-1\right)^3}\,.
      \end{split}
\end{align}
Here we have used (\ref{MasterL}).

\textbullet $\,\,$ Another 3PM diagram having one worldline and one super-field propagator that will contribute has the following topology %(\textcolor{red}{Need to do independently!!}),
\begin{center}
\includegraphics[width=0.38\linewidth]{3PM_7.png}
\end{center}
It is clear from the structure of this diagram that due to the presence of a spinning worldline propagator, we need to take only the spinless terms from the other vertices. Like the previous diagram, this one also has $\bar\Psi^\mu \Psi^\nu$, but for this case, we can add another counterpart, which will have opposite spinor flow. So, in this case, we can also take the antisymmetric combination of $\bar\Psi^\mu \Psi^\nu$ to obtain the classical spin.  Taking all of those things into account, the total contribution to the eikonal phase coming from these two diagrams is given by,
\begin{align}
    \begin{split}
     \chi\Big|_{\mathcal{S}^1}=\frac{m_1^2 m_2^2}{m_p^6}\Big(-i\frac{s_1s_2}{16}\Big)\int e^{-i q\cdot b}\hat\delta(q\cdot v_1)\hat\delta(q\cdot v_2)&\Bigg[\Bigg(\boldsymbol{\tilde h}_{114}\mathbfcal{M}_{0,0;0,0,1,1,1}+\boldsymbol{\tilde h}_{115}q^2\mathbfcal{M}_{0,0;0,0,2,1,1}\\&\hspace{0.25cm}+\boldsymbol{\tilde h}_{116}q^2\mathbfcal{M}_{0,0;0,0,1,2,1}-\frac{1}{2}\boldsymbol{\tilde h}_{117}q^2\mathbfcal{M}^{(++)}_{1,1;0,0,1,1,1}\Bigg)\big(q\cdot\mathcal{S}_1\cdot v_2\big)\,.
     \end{split}
\end{align}
Finally, we get, 
\begin{align}
    \begin{split}  \label{eqt19} 
     \chi\Big|_{\mathcal{S}^1}=\frac{m_1^2 m_2^2}{m_p^6}\Bigg(\Delta_{\mathfrak{26}}^{\textrm{poly}}+\Delta_{\mathfrak{18}}^{\textrm{log}}\log(x)\Bigg)\big(\hat b\cdot\mathcal{S}_1\cdot v_2\big)\,,
      \end{split}
\end{align} 
where, 
\begin{align}
    \begin{split}
  &\Delta_{\mathfrak{26}}^{\textrm{poly}}=\frac{s_1s_2}{32768 \pi ^3 |\boldsymbol{b}|^3 \left(x^2-1\right)^5}\Bigg[-192 \left(x^4-1\right)^2 x \log (|\boldsymbol{b}|)-48 x^{10}+x^9 (-160 \gamma_E +407+128 \log (2)\\&\hspace{0.8cm}-32 \log (\pi ))-144 x^8+1520 x^7+192 x^6+x^5 (320 \gamma_E +2226-256 \log (2)\\&\hspace{0.8cm}+64 \log (\pi ))+192 x^4+1520 x^3-144 x^2+x (-160 \gamma_E +407+128 \log (2)\\&\hspace{0.8cm}-32 \log (\pi ))-48\Bigg]\,,\\&
\Delta_{\mathfrak{18}}^{\textrm{log}}=\frac{s_{1} s_{2} x}{4096 \pi ^3 |\boldsymbol{b}|^3 \left(x^2-1\right)^6 \left(x^2+1\right)}\Bigg(20 x^{12}-32 x^{11}+239 x^{10}-32 x^9+700 x^8+64 x^7+1154 x^6\\&\hspace{0.8cm}+64 x^5+700 x^4-32 x^3+239 x^2-32 x+20\Bigg)\,.
      \end{split}
\end{align}
%\begin{align}
%    \begin{split}
%      i   \chi= i\frac{m_1^2 s_1 m_2^2 s_2}{16 m_p^6}\int e^{-iq \cdot b}\prod \hat\delta(q\cdot v_i)\int_{\ell_{1},\ell_{2}}\frac{\hat\delta(\ell_{1}\cdot v_2)\hat\delta(\ell_{2}\cdot v_1)}{(\ell_{1}\cdot v_1)\,(\ell_{2}\cdot v_2)^2\,\ell_{1}^2\,\ell_{2}^2\,(\ell_{1}+\ell_{2}-q)^2}\mathcal{N}
 %   \end{split}
%\end{align}
%where the numerator $\mathcal{N}$ takes the form \textcolor{black}{AB: %check the sign of the "i" in $\mathcal{N}$},
%\begin{align}
 %   \begin{split}
 %       -i\,\mathcal{N}=&\frac{\ell_{1}\cdot \mathcal{S}_1\cdot \ell_{2}}{8(d-2)^2}\Big[-4 \gamma ^2 \left(d^2-5 d+6\right)(\ell_{2}\cdot v_2)^2+\left(\gamma ^2 \left(d^2-6 d+8\right)+1\right)(\ell_{1}-q)^2+2 \gamma  (d-2)(\ell_{1}\cdot v_1)(\ell_{2}\cdot v_2)\Big]\\ &
 %       +\frac{\ell_{1}\cdot \mathcal{S}_1\cdot v_2}{8(d-2)}\Big[4 \gamma ^2 (d-2)(\ell_{2}\cdot v_2)^3-4\gamma(\ell_{1}\cdot v_1)(\ell_{2}\cdot v_2)^2+\frac{\gamma^2}{8}(\ell_{2}\cdot v_2)\Big(8(d-2)q^2+(15-8d)(\ell_{2}-q)^2\\ &-16(d-2)(\ell_{1}-q)^2\Big)\Big]\\ &
 %    +  \frac{q\cdot\mathcal{S}_1\cdot \ell_{1}}{4(d-2)^2}\,\Big[\left(\gamma ^2 \left(d^2-6 d+8\right)+1\right)(\ell_{1}-q)^2-2 \gamma ^2 \left(d^2-6 d+8\right)(\ell_2\cdot v_2)^2\Big]\\ &+
%        (\ell_{2}\cdot \mathcal{S}_1\cdot v_2)\Big[\frac{\gamma^2}{8}\Big(q^2-(\ell_{1}-q)^2\Big)(\ell_{2}\cdot v_2)-\frac{\gamma}{8(d-2)}(\ell_{1}\cdot v_1)(\ell_{1}-q)^2\Big]\\ &
 %     +  (q\cdot \mathcal{S}_1\cdot v_2)\Big[\frac{\gamma}{4(d-2)}(\ell_{1}\cdot v_1)\,(q-\ell_{1})^2-\frac{\gamma^2}{8}(\ell_{2}\cdot v_2)\Big(2q^2-(\ell_{1}-q)^2-(\ell_{2}-q)^2\Big)\Big]\,.
 %   \end{split}
%\end{align}
%Therefore we need to do the following tensor integral ,
%\begin{align}
%    \begin{split}
 %       \Theta^{(1)}_{[\eta\delta]}&\equiv \int_{\ell_{1},\ell_{2}}\frac{\hat\delta(\ell_{1}\cdot v_2)\hat\delta(\ell_{2}\cdot v_1)\,\ell_{1[\eta}\ell_{2\delta]}}{(\ell_1\cdot v_1)(\ell_2\cdot v_2)^2\ell_{1}^2\,\ell_{2}^2\,(\ell_{1}+\ell_{2}-q)^2}\Big[-4 \gamma ^2 \left(d^2-5 d+6\right)(\ell_{2}\cdot v_2)^2\\& \hspace{0.7 cm}+\left(\gamma ^2 \left(d^2-6 d+8\right)+1\right)(\ell_{1}-q)^2+2 \gamma  (d-2)(\ell_{1}\cdot v_1)(\ell_{2}\cdot v_2)\Big]\,,\\ &
 %        =\Bigg(\boldsymbol{h_{188}}\,\mathbfcal{M}_{0,0;0,0,1,1,1}
%+\boldsymbol{h_{189}}\,q^2\mathbfcal{M}_{0,0;0,0,2,1,1}+\boldsymbol{h_{190}}\,q^2\mathbfcal{M}_{0,0;0,0,1,2,1}+\boldsymbol{h_{191}}\,q^2\mathbfcal{M}_{1,1;0,0,1,1,1}\Bigg)\,q^{[\eta]}v_1^{\delta]}\\ &+\Bigg(\boldsymbol{h_{188}}\,\mathbfcal{M}_{0,0;0,0,1,1,1}
%+\boldsymbol{h_{189}}\,q^2\mathbfcal{M}_{0,0;0,0,2,1,1}+\boldsymbol{h_{190}}\,q^2\mathbfcal{M}_{0,0;0,0,1,2,1}+\boldsymbol{h_{191}}\,q^2\mathbfcal{M}_{1,1;0,0,1,1,1}\Bigg)\,q^{[\eta]}v_2^{\delta]}
%    \end{split}
%\end{align}
%and the following vector integrals,
%\newpage
%\begin{align}
%    \begin{split}
  %      \Theta^{(2)}_{\eta}&\equiv %\int_{\ell_{1},\ell_{2}}\frac{\hat\delta(\ell_{1}\cdot v_2)\hat\delta(\ell_{2}\cdot v_1) \ell_{1}^\eta}{(\ell_1\cdot v_1)(\ell_2\cdot v_2)^2\,\ell_{1}^2\, \ell_{2}^2\, (\ell_{1}+\ell_{2}-q)^2}\Big[4 \gamma ^2 (d-2)(\ell_{2}\cdot v_2)^3-4\gamma(\ell_{1}\cdot v_1)(\ell_{2}\cdot v_2)^2+\frac{\gamma^2}{8}(\ell_{2}\cdot v_2)\\ &\Big(8(d-2)q^2+(15-8d)(\ell_{2}-q)^2-16(d-2)(\ell_{1}-q)^2\Big)\Big]\,,\\&
 %       =\Bigg(\boldsymbol{h}_{196}\mathbfcal{M}_{0,0;0,0,1,1,1}+\boldsymbol{h}_{197}\,q^2\,\mathbfcal{M}_{0,0;0,0,2,1,1}
       %+\boldsymbol{h}_{198}\,q^2\,\mathbfcal{M}_{0,0;0,0,1,2,1}+\boldsymbol{h}_{199}\,q^2\,\mathbfcal{M}_{1,1;0,0,1,1,1}\Bigg)\,q^\eta\\ &\hspace{0.8 cm}+\boldsymbol{h}_{200}\,q^2\,\mathbfcal{M}_{0,1;0,0,1,1,1}\,\frac{\gamma v_2^\eta-v_1^\eta}{\gamma^2-1}\,,\\ &
   %     \hspace{0cm} \Theta_{\eta}^{(3)}\equiv \int_{\ell_{1},\ell_{2}}\frac{\hat\delta(\ell_{1}\cdot v_2)\hat\delta(\ell_{2}\cdot v_1)\,\ell_{1}^\eta}{(\ell_{1}\cdot v_1)(\ell_{2}\cdot v_2)^2\ell_{1}^2 \ell_{2}^2 (\ell_{1}+\ell_{2}-q)^2}\Big[\left(\gamma ^2 \left(d^2-6 d+8\right)+1\right)(\ell_{1}-q)^2-2 \gamma ^2 \left(d^2-6 d+8\right)(\ell_2\cdot v_2)^2\Big]\\ &
  %      =\Bigg(\boldsymbol{h}_{201}\,\mathbfcal{M}_{0,0;0,0,1,1,1}+\boldsymbol{h}_{202}\,q^2\,\mathbfcal{M}_{0,0;0,0,2,1,1}+\boldsymbol{h}_{203}\,q^2\,\mathbfcal{M}_{0,0;0,0,1,2,1}\Bigg)\frac{\gamma v_2^\delta-v_1^\delta}{\gamma^2-1}
%       +\boldsymbol{h}_{204}\,\mathbfcal{M}_{0,1;0,0,1,1,1}\, q^{\delta}\,,\\ &
%        \hspace{0cm}\Theta_{\eta}^{(4)}=\int_{\ell_{1},\ell_{2}}\frac{\hat\delta(\ell_{1}\cdot v_2)\hat\delta(\ell_{2}\cdot v_1)\ell_{2}^\eta}{(\ell_1\cdot v_1)(\ell_2\cdot v_2)^2\,\ell_{1}^2 \ell_{2}^2 (\ell_{1}+\ell_{2}-q)^2}\Big[\frac{\gamma^2}{8}\Big(q^2-(\ell_{1}-q)^2\Big)(\ell_{2}\cdot v_2)-\frac{\gamma}{8(d-2)}(\ell_{1}\cdot v_1)(\ell_{1}-q)^2\Big]\\ &
  %      %=\boldsymbol{h}_{205}\,q^2\,\mathbfcal{M}_{0,1;0,0,1,1,1}\frac{\gamma v_1^\eta-v_2^\eta}{\gamma^2-1}+\Bigg(\boldsymbol{h}_{206}\,\mathbfcal{M}_{0,0;0,0,1,1,1}+\boldsymbol{h}_{207}\,q^2\,\mathbfcal{M}_{0,0;0,0,2,1,1}+\boldsymbol{h}_{208}\,q^2\,\mathbfcal{M}_{0,0;0,0,1,2,1}\\ &+\boldsymbol{h}_{209}\,q^2\,\mathbfcal{M}_{1,1;0,0,1,1,1}\Bigg)q^{\eta}
%    \end{split}
%\end{align}
%and the scalar last scalar integral in the numerator takes the form,
%\begin{align}
 %   \begin{split}
 %   \centering
 %       \Theta^{(5)}&= \int_{\ell_{1},\ell_{2}}\frac{\hat\delta(\ell_{1}\cdot v_2)\hat\delta(\ell_{2}\cdot v_1)}{(\ell_1\cdot v_1)(\ell_2\cdot v_2)^2\ell_{1}^2 \ell_{2}^2 (\ell_{1}+\ell_{2}-q)^2}\Big[\frac{\gamma}{4(d-2)}(\ell_{1}\cdot v_1)\,(q-\ell_{1})^2-\frac{\gamma^2}{8}(\ell_{2}\cdot v_2)\Big(2q^2-(\ell_{1}-q)^2-(\ell_{2}-q)^2\Big)\Big]\,\\ &
 %    =  \boldsymbol{h}_{210}\,\mathbfcal{M}_{0,0;0,0,1,1,1}+\boldsymbol{h}_{211}\,q^2\,\mathbfcal{M}_{0,0;0,0,2,1,1}+\boldsymbol{h}_{212}\,q^2\,\mathbfcal{M}_{0,0;0,0,1,2,1}
%+   \boldsymbol{h}_{213}\, q^2\,\mathbfcal{M}_{1,1;0,0,1,1,1}\,.
%    \end{split}
%\end{align}
%Using all these and doing the $q$ integral we get, 
%\newpage
%\hfsetfillcolor{gray!10}
%\hfsetbordercolor{black!150}
%\begin{align}
 %   \begin{split}
 %    \centering
%\tikzmarkin[disable rounded corners=false]{e12}(3.8,-1)(-0.6,0.8)  
 %     \hspace{0cm}i   \chi\Big|_{\mathcal{S}^1}= \frac{-i\,m_1^3 m_2 s_1 s_2\,x}{24576\pi^3\,|\boldsymbol{b}|^3\,\left(1-x^2\right)^6\,m_p^6}&\Bigg[-2 \left(x^4-1\right)^2 \big(6 \left(x^2-1\right) \log(|\boldsymbol{b}|(1- x))+x \log (x) \big(2 (6 \log (|\boldsymbol{b}| (x+1))+5 \gamma_E -7)\\&+\log \left(\frac{\pi ^2}{256 x^6}\right)\big)+6 x \,\textbf{Li}_2\left(x^2\right)+5 \gamma_E  \left(x^2-1\right)+x \left(x \left(\log \left(\frac{\pi }{16}\right)-4\right)-\pi ^2\right)\\&+4+\log \left(\frac{16}{\pi }\right)\big)+3 \left(x^2+1\right)^2 \left(x^2-1\right)^3+192 \left(x^3+x\right)^2 \left(x^2-1\right)\\&-6 \left(x^2+1\right) \left(x^2-1\right)^2 \left(1-x^4+\left(x^4+6 x^2+1\right) \log (x)\right)\\&-12 \big(x^{10}+11 x^8+6 x^6-6 x^4-11 x^2\\&-2 \left(x^2+1\right) \left(5 x^4+14 x^2+5\right) x^2 \log (x)-1\big)\Bigg]\Bigg(\hat{b}\cdot\mathcal{S}_1\cdot v_2\Bigg)\,.
 %     \tikzmarkend{e12}
 %     \end{split}
 %     \end{align}
   %   The contribution to the eikonal phase form the second diagram has the following form,
 %     \begin{align}
   %       \begin{split}
    %          i\chi=-i \frac{m_1^2 s_1m_2^2 s_2}{16 m_p^6} \int e^{iq\cdot %b}\hat\delta(q\cdot v_1)\hat\delta(q\cdot v_2)\int \frac{\hat\delta(\ell_1\cdot v_2)\hat\delta(\ell_2\cdot v_1)\mathcal{N}}{(\ell_1\cdot v_1+i\varepsilon)(\ell_2\cdot v_2-i\varepsilon)^2\,(\ell_1-q)^2\,(\ell_2-q)^2\,(\ell_1+\ell_2-q)^2}
  %        \end{split}
  %    \end{align}
 %     where, the numerator $\mathcal{N}$ has the following form,
 %     \begin{align}
%          \begin{split}
 %             \mathcal{N}=\frac{i}{4(d-2)}(\ell_1-q)\cdot \mathcal{S}_1\cdot v_2\,(\ell_1\cdot v_1)\,\Big[(\gamma ^2 (d-2)-1)\ell_1^2+4\gamma(d-2)(\ell_2\cdot v_2)(\ell_1\cdot v_1-\gamma \ell_2\cdot v_2)\Big]
 %         \end{split}
 %     \end{align}
 %     We need the following vector integral,
%      \begin{align}
  %        \begin{split}
  %            \mathbf{Q}^{\eta}&=\int \frac{\hat\delta(\ell_1\cdot v_2)\hat\delta(\ell_2\cdot v_1)\,(\ell_1-q)^{\eta}}{(\ell_2\cdot v_2-i\varepsilon)^2\,(\ell_1-q)^2\,(\ell_2-q)^2\,(\ell_1+\ell_2-q)^2}\Big[(\gamma ^2 (d-2)-1)\ell_1^2+4\gamma(d-2)(\ell_2\cdot v_2)\\&\hspace{0.8cm}(\ell_1\cdot v_1-\gamma \ell_2\cdot v_2)\Big]\\ &
  %            =\Bigg(\boldsymbol{h}_{214}\,\mathbfcal{M}_{0,0;0,0,1,1,1}+\boldsymbol{h}_{215}\,q^2\,\mathbfcal{M}_{0,0;0,0,2,1,1}+\boldsymbol{h}_{216}\,q^2\,\mathbfcal{M}_{0,0;0,0,1,2,1}\Bigg)\,q^\eta-\frac{\gamma v_2^\eta-v_1^\eta}{\gamma^2-1}\boldsymbol{h}_{217}\,q^2\,\mathbfcal{M}^{(+-)}_{0,1;0,0,1,1,1}
 %         \end{split}
%      \end{align}
\textbullet $\,\,$ Two 3PM diagrams coming from graviton-scalar and scalar-scalar worldline interaction vertex have the following topology,
\begin{center}
\includegraphics[width=0.38\linewidth]{3PM_89.jpg}
\end{center}
% The corresponding amplitude is given by,
 %\begin{align}
 %    \begin{split}
%         i\chi= -i\frac{s_1^2 s_2^2 m_1^2 m_2^2}{32 m_p^6}\int_{k_i,\omega} e^{-i(k_1+k_2+k_3)\cdot b}\hat\delta(k_1\cdot v_1+k_2\cdot v_1)\hat\delta(k_3\cdot v_1)\hat \delta(k_1\cdot v_2)\hat\delta(k_2\cdot v_2+\omega)\hat\delta(k_3\cdot v_2-\omega)\frac{\mathcal{N}}{\omega ^2\,k_1^2\, k_2^2\, k_3^2 }\label{4.68}
%     \end{split}
% \end{align}
% where the numerator has the form,
% \begin{align}
  %   \begin{split}
 %        \mathcal{N}=\Big(v_1^\mu v_1^\nu-i(k_1\cdot \mathcal{S}_1)^{(\mu}v_1^{\nu)}\Big)P_{\mu\nu;\alpha\beta}\Big(v_2^\alpha v_2^\beta+i(k_1\cdot \mathcal{S}_2)^{(\alpha}v_2^{\beta)}\Big)(2\omega v_2^\rho+k_2^\rho)(-2\omega v_{2\rho}+k_3^\rho)\,.
 %    \end{split}
% \end{align}
% Now relabelling, $k_1\to \ell_{1}, k_3\to \ell_{2}$ and $k_2\to q-\ell_{1}-\ell_{2}\,,$ we will have (recovering the vertex factors),
% \begin{align}
%     \begin{split}
 %       i \chi=-i\frac{s_1^2 s_2^2 m_1^2 m_2^2}{32 m_p^6}\int e^{-iq\cdot b}\hat\delta(q\cdot v_1)\hat\delta(q\cdot v_2)\int_{\ell_{1},\ell_{2}}\frac{\hat\delta(\ell_{1}\cdot v_2)\hat\delta(\ell_{2}\cdot v_1)}{(\ell_{2}\cdot v_2)^2\ell_{1}^2 \ell_{2}^2 (\ell_{1}+\ell_{2}-q)^2}\mathcal{N}
  %   \end{split}
% \end{align}
% where the numerator can be separated into spinning and non-spinning part (eliminating the tadpoles),
% \begin{align}
%     \begin{split}
%         &\mathcal{N}\Big|_{\mathcal{S}^0}=\frac{\gamma^2(d-2)-1}{2(d-2)}(\ell_{1}-q)^2\,,\\ &
%         \mathcal{N}\Big|_{\mathcal{S}^1}=\Bigg(-\frac{i\gamma}%{2}\,\ell_1\cdot \mathcal{S}_1\cdot v_2-i\frac{3-d}{2(d-2)}\ell_1\cdot \mathcal{S}_2\cdot v_1\Bigg)(\ell_1-q)^2
%     \end{split}
% \end{align}
For the first diagram, we get the following contribution for the spinless part of the eikonal phase,

% \hfsetfillcolor{gray!10}
%\hfsetbordercolor{black!150}
 \begin{align}
     \begin{split}
         i\chi\Big|_{\mathcal{S}^{0}}=& -i\frac{s_1^2 s_2^2 m_1^2 m_2^2}{32 m_p^6} \frac{\gamma^2(d-2)-1}{2(d-2)}\int e^{-iq\cdot b}\hat\delta(q\cdot v_1)\hat\delta(q\cdot v_2)\Bigg(\frac{\epsilon  (2 \epsilon +1)}{2 \epsilon -1}\mathbfcal{M}_{0,0;0,0,1,1,1}+\frac{q^2 (2 \epsilon +1)}{2-4 \epsilon }\mathbfcal{M}_{0,0;0,0,2,1,1}\\ &\hspace{0.8cm}+\frac{q^2 (2 \epsilon +1)}{2 \left(\gamma ^2-1\right) \epsilon }\mathbfcal{M}_{0,0;0,0,1,2,1}\Bigg)\,.
     \end{split}
 \end{align}
 Using (\ref{MasterM}) we get, 
 \begin{align}
     \begin{split}  \label{eqt20} 
     \chi\Big|_{\mathcal{S}^{0}}=-\frac{m_1^2m_2^2}{m_p^6}\Delta_{\mathfrak{27}}^{\textrm{poly}},\,\,\,\Delta_{\mathfrak{27}}^{\textrm{poly}}=\Bigg(\frac{s_{1}^2 s_{2}^2 \left(x^2+1\right) \left(x^4+1\right)}{2048 \pi ^3 |\boldsymbol{b}|^2 \left(x^2-1\right)^3}\Bigg).&
      \end{split}
 \end{align}
The contribution to the spinning part is given by,
 \begin{align}
     \begin{split}
         i\chi\Big|_{\mathcal{S}^{1}}= -\frac{\gamma s_1^2 s_2^2 m_1^2 m_2^2}{64 m_p^6}\int e^{-i q\cdot b}\hat\delta(q\cdot v_1)\hat\delta(q\cdot v_2)\,&\Bigg(\boldsymbol{\tilde h}_{118}\mathbfcal{M}_{0,0;0,0,1,1,1}+q^2 \boldsymbol{\tilde h}_{119}\mathbfcal{M}_{0,0;0,0,2,1,1} +q^2 \boldsymbol{\tilde h}_{120}\mathbfcal{M}_{0,0;0,0,1,2,1}\Bigg)\\&\Bigg(q\cdot\mathcal{S}_1\cdot v_2- \frac{1-2\epsilon}{4\,\gamma\,(1-\epsilon)}q\cdot\mathcal{S}\cdot v_1\Bigg)\,.
     \end{split}
 \end{align}
 %where,
% \begin{align}
  %   \begin{split}
 %        \Theta(q,\gamma)=&\int_{\ell_{1},\ell_{2}}\frac{\hat\delta(\ell_{1}\cdot v_2)\hat\delta(\ell_{2}\cdot v_1)(\ell_{1}-q)^2 \ell_{1}^\eta}{(\ell_{2}\cdot v_2)^2 \ell_{1}^2\,\ell_{2}^2\,(\ell_{1}+\ell_{2}-q)^2}\\ &
%         =\Bigg(\boldsymbol{\tilde h}_{114}\mathbfcal{M}_{0,0;0,0,1,1,1}+q^2 \boldsymbol{\tilde h}_{115}\mathbfcal{M}_{0,0;0,0,2,1,1}\\ &
 %       \hspace{0.6 cm} +q^2 \boldsymbol{\tilde h}_{116}\mathbfcal{M}_{0,0;0,0,1,2,1}\Bigg)q^{\eta}\,.
%     \end{split}
% \end{align}
 Then using \eqref{MasterM} we get,

 \hfsetfillcolor{gray!10}
\hfsetbordercolor{black!150}
\begin{align}
     \begin{split}  \label{eqt21} 
 \chi\Big|_{\mathcal{S}^{1}}= \frac{m_1^2 m_2^2}{m_p^6}\Bigg(\Delta_{\mathfrak{28}}^{\textrm{poly}}(\hat{b}\cdot\mathcal{S}_1\cdot v_2)+\Delta_{\mathfrak{29}}^{\textrm{poly}} (\hat{b}\cdot\mathcal{S}_2\cdot v_1)\Bigg),\,
     \end{split}
 \end{align}
 where,
 \begin{align}
     \Delta_{\mathfrak{28}}^{\textrm{poly}}=-\Bigg(\frac{s_1^2s_2^2\, x(x^2+1)^2}{1024 \pi ^3 |\boldsymbol{b}|^3 \left(x^2-1\right)^3}\Bigg), \quad \Delta_{\mathfrak{29}}^{\textrm{poly}}= -\frac{x}{2(x^2+1)}\Delta_{\mathfrak{28}}^{\textrm{poly}}.
 \end{align}
 Proceeding as before, for the second diagram, we have only  $\mathcal{O}(\mathcal{S}^0)\,.$  contribution.
 \begin{align}
     \begin{split}  \label{eqt22} 
    \chi \Big|_{\mathcal{S}^0}=\frac{m_1m_2^3}{m_p^6}\Delta_{\mathfrak{30}}^{\textrm{poly}},\quad \Delta_{\mathfrak{30}}^{\textrm{poly}}=-\Bigg(\frac{g_{1} s_{1} s_{2}^3 x^2}{256 \pi ^3 |\boldsymbol{b}|^2 \left(x^2-1\right)^2}\Bigg)\,.
    \end{split}
 \end{align}
 %   i\frac{s_1 s_2^3 g_1 m_1 m_2^3}{32 m_p^6} \int e^{-i(k_1+k_2+k_3)\cdot b}  \,\hat\delta(\omega-k_1\cdot v_1-k_2\cdot v_1) \,\hat\delta(\omega+k_3\cdot v_1)\hat\delta(k_1\cdot v_2)\hat\delta(k_2\cdot v_2)\hat\delta(k_3\cdot v_2) \frac{\mathcal{N}}{\omega^2 k_1^2 k_2^2 k_3^2}
%where the vertex factor takes the form,
%\begin{align}
%    \begin{split}
%        \mathcal{N}=\Big(-(k_{1\rho}+k_{2\rho})+2\omega v_{1\rho}\Big)(-2\omega v_1^\rho-k_{3\rho})\,.
%    \end{split}
%\end{align}
%Now again relabelling $k_1\to \ell_{1},\,k_3\to \ell_{2}$ and $k_2\to q-\ell_{1}-\ell_{2}$ we have,\\
%\hfsetfillcolor{gray!10}
%\hfsetbordercolor{black!150}
%\begin{align}
%    \begin{split}
 %    \tikzmarkin[disable rounded corners=false]{e14}(10,-1)(-0.1,1)  
 %       i\chi\Big|_{\mathcal{S}^0}&=i\frac{s_1 s_2^3 g_1 m_1 m_2^3}{32 %m_p^6}\int e^{-i q\cdot b}\hat\delta(q\cdot v_1)\hat\delta(q\cdot v_2)\int_{\ell_{1},\ell_{2}}\frac{\hat\delta(\ell_{1}\cdot v_2)\hat\delta(\ell_{2}\cdot v_2)}{(\ell_{2}\cdot v_1)^2 \ell_{1}^2 \ell_{2}^2 (q-\ell_{1}-\ell_{2})^2}\Bigg(q^2-(q-\ell_{1})^2-(q-\ell_{2})^2\Bigg)\,,\\ &
%        =i\frac{s_1 s_2^3 g_1 m_1 m_2^3}{32 m_p^6}\int e^{-i q\cdot b}\hat\delta(q\cdot v_1)\hat\delta(q\cdot v_2)\Bigg(\frac{6 \epsilon -1}{\gamma ^2-1}\Bigg)\mathbfcal{L}_{0,0;0,0,1,1,1}\,,\\&
%        =-i\frac{s_1 s_2^3 g_1 m_1 m_2^3}{32 m_p^6}\Bigg(\frac{x^3}{4 \pi ^3 |\boldsymbol{b}|^2 \left(1-x^2\right)^3}\Bigg)\,.
 %       \tikzmarkend{e14}
 %   \end{split}
%\end{align}
 Again, one also needs to add the contributions from 3PM diagrams obtained by interchanging the worldlines one and two in all the diagrams mentioned above. We will not show them explicitly as the contributions from those diagrams will be similar to those we have already computed and can be obtained easily by exchanging labels one and two. 
 %\\\\
 %\newpage
% \textcolor{red}{Final $q$ integral.}
\par
\textbf{3PM eikonal phase from Chern-Simons interaction vertices:}
Now let us focus on the dCS contribution to the eikonal vertex. The corresponding vertex term is given by,
 \begin{align}
    \begin{split}
        \mathcal{V}_{\textrm{dCS}}(k_1,k_2,q)
        =&\,\epsilon^{\chi\varepsilon\mu\nu}\Bigg(
        k_{1\mu}k_{1\beta}k_{2\chi}k_{2\sigma}h^{\sigma}_{\nu}(-k_1)h_{\delta}^{\beta}(-k_2)\\
        &\qquad\qquad
        -k_{1\mu}(k_1\cdot k_2)\,k_{2\chi}h^{\sigma}_{\nu}(-k_1)h_{\delta\sigma}(-k_2)
        \Bigg)\varphi(-q)\,\hat\delta^{(4)}(k_1+k_2+q)\,.
    \end{split}
\end{align}
With this vertex, we will now proceed with the computation of the contribution of this dCS term to the eikonal phase. Below, we list all possible diagrams for the same. 
 %\newpage
\par
\textbullet $\,\,$ The following two diagrams will contribute to the eikonal phase at 3 PM involving Chern-Simons interaction  :
\begin{center}
\includegraphics[width=0.38\linewidth]{CS_f.jpg}
\end{center}
Contribution to the eikonal phase from the first diagram is given by,
\begin{align}
    \begin{split}
      i  \chi=-l_{dCS}^2\frac{s_1 m_1 m_2^2}{8 m_p^6} \int e^{iq\cdot b}\frac{\hat\delta(q\cdot v_1)\hat\delta(q\cdot v_2)}{q^2}\int_{k_2}\hat\delta(k_2\cdot v_2)\frac{\mathcal{N}}{k_2^2 (k_2+q)^2}
    \end{split}
\end{align}
where the numerator takes the form,
\begin{align}
\begin{split}
    \mathcal{N}&=-\frac{i}{2}\epsilon^{\gamma\delta\tau\sigma}(q+k_2)\cdot k_2(q+k_2)_{\tau}k_{2\gamma}\,\Big(v_{2\sigma}(k_2\cdot \mathcal{S}_2)_{\delta}-[(q+k_2)\cdot \mathcal{S}_2 ]_{\sigma}\,\,v_{2\delta}\Big)\,,\\ &
  \to  \frac{i}{2}\epsilon^{\gamma\delta\tau\sigma}q^2\, q_\tau \,k_{2\gamma} v_{2\sigma}(k_2\cdot \mathcal{S}_2)_{\delta}\,\,(\textrm{ignoring the massless tadpoles})\,.
  \end{split}
\end{align}
Now, doing the $k_2$  and $q$ integral and using the anti-symmetric property of epsilon, it is easy to show that the contribution to the eikonal phase from this diagram is given by,

\begin{align}
\begin{split} \label{eqt23} 
 \chi\Big|_{\mathcal{S}^1}&=l_{dCS}^2\frac{s_1 m_1 m_2^2}{8 m_p^6}\frac{ 3 x\,}{64\pi (1-x^2)\,|\boldsymbol{b}|^4}\,\epsilon^{\gamma\delta\tau\sigma}v_{2\sigma}\mathcal{S}_{2\gamma\delta}\,\hat{b}_\tau\,,\\ &=\Delta_{\mathfrak{1}}^{\textrm{dCS}}\left(\mathcal{S}_2\wedge \hat b\right).
\end{split}
\end{align}
where,
\begin{align}
\Delta_{\mathfrak{1}}^{\textrm{dCS}}=l_{dCS}^2\frac{s_1 }{8 }\frac{ 3 x\,}{64\pi (1-x^2)\,|\boldsymbol{b}|^4}
\end{align}
%\newpage
 %We now compute$\,\,$ $\mathcal{O}(\mathcal{S}_1\mathcal{S}_2)$ diagram from Chern-Simons term. 
Correspondingly, the second diagram gives,
%\begin{align}
  %  \begin{split}
    %     \begin{minipage}[h]{0.08\linewidth}
	%\vspace{4pt}
	%\scalebox{1.5}{\includegraphics[width=\linewidth]{CS_2.png}}
% \end{minipage}&\nonumber
 %\end{split}
 %\end{align}
%The corresponding amplitude is given by,
\begin{align}
    \begin{split}
      i  \chi=-i\,l_{dCS} ^2\frac{s_2m_1 m_2^2}{m_p^6}\int_{k_1}e^{ik_1\cdot b}\frac{\hat\delta(k_1\cdot v_1)\hat\delta(k_1\cdot v_2)}{k_1^2}\int_{k_2}\hat\delta(k_2\cdot v_2)\frac{\mathcal{N}}{k_2^2(k_1+k_2)^2}
    \end{split}
\end{align}
where the numerator is given by,
\begin{align}
    \begin{split}
        \mathcal{N}=&\epsilon^{\gamma\delta\tau\sigma}\Bigg(v_1^\mu v_1^\nu+i(k_1\cdot \mathcal{S}_1^{(\mu})\,v_1^{\nu)}-\frac{1}{2}(k_1\cdot \mathcal{S}_1)^\mu (k_1\cdot \mathcal{S}_1)^{\nu}\Bigg)\,P_{\mu\nu,\rho\sigma}\Bigg(k_{1\tau}k_{2\gamma}k_1^{\chi}k_2^\rho-\eta^{\rho\chi}(k_1\cdot k_2)\Bigg)\\ &
        \times P_{\chi\varepsilon,\alpha\beta}\Bigg(v_2^\alpha v_2^\beta+i(k_2\cdot \mathcal{S}_2^{(\alpha})\,v_2^{\beta)}-\frac{1}{2}(k_2\cdot \mathcal{S}_2)^\alpha (k_2\cdot \mathcal{S}_2)^{\beta}\Bigg)\,.
    \end{split}
\end{align}
The numerator can be decomposed in different orders of spin and are given by,
\begin{align}
    \begin{split}
&\mathcal{N}\Big|_{\mathcal{S}^0}=\gamma  \left(  k_1 \cdot k_2\right)  \epsilon_{\alpha \beta\rho\sigma} k_1^\alpha k_2^\beta v_1^\rho v_2^\sigma\,, \\ &
\mathcal{N}\Big|_{\mathcal{S}^1}=\frac{1}{4}i\Big(2\gamma(k_1\cdot k_2)(k_1\cdot \mathcal{S}_1)^\mu\epsilon_{\mu \beta\rho\sigma}k_1^\beta k_2^\rho v_2^\sigma+2(k_1\cdot k_2)\epsilon_{\alpha \beta\rho\sigma} k_1^\alpha k_2^\beta v_1^\rho v_2^\sigma (k_1\cdot \mathcal{S}_1\cdot v_2)\,\\&
\hspace{1.2 cm}-2\gamma(k_1\cdot k_2) (k_2\cdot{\mathcal{S}_2})^\alpha\epsilon_{\alpha\beta\rho\sigma}{k_1}^\beta {k_2}^\rho {v_1}^\sigma+2 (k_1\cdot \mathcal{S}_2\cdot k_2)(k_2\cdot v_1)\epsilon_{\alpha\beta\rho\sigma} k_1^\alpha k_2^\beta v_1^\rho v_2^\sigma\\&\hspace{1.0 cm}+2(k_1\cdot k_2)(k_2\cdot \mathcal{S}_2\cdot v_1) \epsilon_{\alpha \beta\rho\sigma} k_1^\alpha k_2^\beta v_1^\rho v_2^\sigma \Big)\,.
\end{split}
\end{align}
Therefore, the contribution to the eikonal phase takes the form, %\vspace{-2cm}

\begin{align}
\begin{split}
i\chi\Big|_{\mathcal{S}^0}&
%=\epsilon_{\alpha\beta\rho\sigma}\,v_1^\rho v_2^\sigma\int_{k_1}e^{ik_1\cdot b}\frac{\hat\delta(k_1\cdot v_1)\hat\delta(k_1\cdot v_2)}{k_1^2}\int_{k_2}\hat\delta(k_2\cdot v_2)\frac{\gamma(k_1\cdot k_2)k_1^\alpha k_2^\beta}{k_2^2(k_1+k_2)^2}\,,\\ &
    \to\frac{\gamma}{2}\epsilon_{i j\rho\sigma}\delta_{mn}\,v_1^\rho v_2^\sigma\int_{k_1}e^{ik_1\cdot b}{\hat\delta(k_1\cdot v_1)\hat\delta(k_1\cdot v_2)k_1^i\,k_1^{m}}\underbrace{\int_{\vec k_2}\frac{k_2^j \,k_2^n}{\vec k_2^2(\vec k_2+\vec k_1)^2}}_{\propto ()\delta^{jn}+()k_1^j k_1^n}=0\,.\label{3.79}
    \end{split}
\end{align}
%\vspace{-1cm}
%Here we have used the fact that the epsilon tensor is antisymmetric. 
%Hence, it is evident that for nonspinning black holes the Chern-Simons interaction does not contribute to the scattering angle. 
%
%\begin{align}
%\begin{split}
%i\chi\Big|_{\mathcal{S}^0}&=0\,.
 %\end{split}
%\end{align}
%\newpage
Now, let us concentrate on the spin-orbit contribution to the eikonal phase coming from the Chern-Simons interaction vertices. %The first and second order spin term are given by,
%\vspace{-1cm}
\begin{align}
    \begin{split}
        \chi\Big|_{\mathcal{S}^1}=\sum_{i=1}^{3}\chi_{\mathcal{S}^1}^{(i)}
    \end{split}
\end{align}
%\vspace{-2 cm}
%\newpage
where,
\begin{align}
    \begin{split}
      %  \bullet \,\,\, \chi_{\mathcal{S}^1}^{(1)}&\sim -i\frac{\gamma}{2}\epsilon_{\mu\beta\rho\sigma}v_2^\sigma\, \mathcal{S}_{1}^{\eta\mu}\int_{k_1}e^{ik_1\cdot b}{\hat\delta(k_1\cdot v_1)\hat\delta(k_1\cdot v_2)k_{1\eta}\,k_1^\beta}\int_{k_2}\frac{\hat\delta(k_2\cdot v_2)k_2^\rho}{k_2^2 \,(k_2+k_1)^2}\\ &  =i\frac{\gamma}{4}\epsilon_{\mu\beta\rho\sigma}v_2^\sigma \mathcal{S}_1^{\mu\eta}\int_{k_1}e^{ik_1\cdot b}\hat\delta(k_1\cdot v_1)\hat\delta(k_1\cdot v_2)k_{1\eta}k_1^\beta k_1^\rho \int_{\vec k_2}\frac{1}{\vec k_2^2(\vec k_2+\vec k_1)^2}\to 0,\\ &
  %    =  i\frac{\Gamma(1/2+\epsilon)\Gamma(1/2-\epsilon)^2}{\Gamma(1-2\epsilon)(4\pi)^{3/2-\epsilon}}\frac{\gamma}{4}\epsilon_{\mu\beta\rho\sigma}v_2^\sigma (\mathcal{S}_1)^{\mu}_{\,\,\,\,\eta}\int_{k_1}e^{ik_1\cdot b}\frac{\hat\delta(k_1\cdot v_1)\hat\delta(k_1\cdot v_2)}{(-k^2)^{\epsilon+\frac{1}{2}}}k_1^{\eta}k_1^\beta k_1^\rho
  \\
      \bullet \,\,\, \chi_{\mathcal{S}^1}^{(1)}&\sim -\frac{i}{4}\epsilon_{\alpha\beta\rho\sigma}(\mathcal{S}_2)^{\eta\tau}\,v_{1\tau}v_1^\rho v_2^\sigma \int_{k_1}e^{ik_1\cdot b}\hat\delta(k_1\cdot v_1)\hat\delta(k_1\cdot v_2)k_1^\alpha \int_{k_2}\frac{\hat\delta(k_2\cdot v_2)k_2^\beta k_{2\eta}}{k_2^2(k_2+k_1)^2}\,,\\ &
     % =\frac{i}{4}\epsilon_{\alpha i \rho\sigma}(\mathcal{S}_2)^{j\tau}\,v_{1\tau}v_1^\rho v_2^\sigma\int_{k_1}e^{ik_1\cdot b}\hat\delta(k_1\cdot v_1)\hat\delta(k_1\cdot v_2)k_1^\alpha \int_{\vec k_2}\frac{k_2^i \,k_{2j}}{\vec k_2^2(\vec k_2+\vec k_1)^2}\,,\\ &
   %   =\frac{i\,2^{2 \epsilon -3} \pi ^{\epsilon -\frac{3}{2}} \Gamma \left(\frac{1}{2}-\epsilon \right) \Gamma \left(\frac{3}{2}-\epsilon \right) \Gamma \left(\epsilon -\frac{1}{2}\right)}{4\Gamma (3-2 \epsilon )}\epsilon_{m i \rho\sigma}(\mathcal{S}_2)^{j\tau}\,v_{1\tau}v_1^\rho v_2^\sigma\int \frac{d^2k_1}{(2\pi)^2}\frac{e^{-i\vec k_1\cdot \boldsymbol{b}}\,k_1^m}{(\vec k_1)^{2(\epsilon-1/2)}}\delta^i_j\,,\\ &
      \xrightarrow[]{\epsilon\to 0}\frac{3}{128\pi |\boldsymbol{b}|^4} \Bigg(\frac{2x}{1-x^2}\Bigg)(\overrightarrow{\mathcal{S}_2\cdot v_1})\cdot (\vec v_1\times \hat b)\,.
      \\ &
       %    =  i\frac{\Gamma(1/2+\epsilon)\Gamma(1/2-\epsilon)^2}{\Gamma(1-2\epsilon)(4\pi)^{3/2-\epsilon}}\frac{\gamma}{4}\epsilon_{\mu\beta\rho\sigma}v_2^\sigma (\mathcal{S}_1)^{\mu}_{\,\,\,\,\eta}\int_{k_1}e^{ik_1\cdot b}\frac{\hat\delta(k_1\cdot v_1)\hat\delta(k_1\cdot v_2)}{(-k^2)^{\epsilon+\frac{1}{2}}}k_1^{\eta}k_1^\beta k_1^\rho
  \\
 \bullet \,\,\, \chi_{\mathcal{S}^1}^{(2)}&\sim \frac{i}{2}\epsilon_{\alpha\beta \rho\sigma}v_1^\rho v_2^\sigma v_{1\chi}(\mathcal{S}_2)_{\eta\delta}\int e^{i k_1\cdot b}\frac{\hat\delta(k_1\cdot v_1)\hat\delta(k_1\cdot v_2)\,k_1^\eta k_1^\alpha}{k_1^2}\int \frac{\hat\delta(k_2\cdot v_2)k_2^\delta\, k_2^\beta k_2^\chi}{k_2^2\,(k_2+k_1)^2}\,,\\ &
% \to \frac{i}{64}\epsilon_{\alpha j \rho\sigma}v_1^\rho v_2^\sigma v_{1l} (\mathcal{S}_2)_{\eta i}\int e^{i k_1\cdot b}\frac{\hat\delta(k_1\cdot v_1)\hat\delta(k_1\cdot v_2)k_1^\eta k_1^\alpha}{k_1^2}\Bigg(\frac{1}{4}|\vec k_1|(\delta^{ij}k_1^l+\delta^{il}k_1^j+\delta^{jl}k_1^i)-\frac{5}{4|\vec k_1|}k_1^i k_1^j k_1^l\Bigg)\,,\\ &
 %=-\frac{i}{256}\epsilon_{m j \rho\sigma}v_1^\rho v_2^\sigma v_{1l} (\mathcal{S}_2)_{n i}\int \frac{d^2k_1}{(2\pi)^2}e^{-i \vec k_1\cdot \boldsymbol{b}}\frac{k_1^m k_1^n k_1^l}{|\vec k_1|}\delta^{ij}\\ & 
 %=-\frac{3}{128\pi b^4}\epsilon_{mjp0}v_1^p v_{1l}(\mathcal{S}_2)_n^{\,\,\,\,j}\Bigg(\frac{1}{2}(\delta^{mn}\hat b^l+\delta^{ml}\hat b^n+\delta^{nl}\hat b^m)-3\hat b^m \hat b^n \hat b^l\Bigg)\\ &
 =-\frac{3}{128\pi |\boldsymbol{b}|^4}\Bigg(\frac{2x}{1-x^2}\Bigg)\Bigg(\frac{1}{2}\vec v_1\cdot \hat b\times (\overrightarrow {v_1\cdot \mathcal{S}_2})\Bigg)=\frac{3}{256\pi |\boldsymbol{b}|^4}(\overrightarrow{\mathcal{S}_2\cdot v_1})\cdot (\vec v_1\times \hat b  )\,.\nonumber
\end{split}
\end{align}
\begin{align}
\begin{split}
\hspace{0.7 cm} \bullet \,\,\, \chi_{\mathcal{S}^1}^{(3)}&\sim -\frac{i \gamma}{2}\epsilon_{\alpha\beta \rho\sigma}v_1^\sigma \mathcal{S}_2^{\eta\alpha}\int e^{ik_1\cdot b}\frac{\hat\delta(k_1\cdot v_1)\hat\delta(k_1\cdot v_2)k_1^\beta}{k_1^2}\int \frac{\hat\delta(k_2\cdot v_2)\,k_1\cdot k_2}{k_2^2(k_2+k_1)^2}k_2^\eta k_{2\rho}\,,\\ &
 %=-\frac{i \gamma}{2}\epsilon_{\alpha\beta \rho\sigma}v_1^\sigma \mathcal{S}_2^{\eta\alpha}\int e^{ik_1\cdot b}\frac{\hat\delta(k_1\cdot v_1)\hat\delta(k_1\cdot v_2)k_1^\beta}{k_1^2}\\&
 %=\frac{i\gamma}{256}\epsilon_{\alpha j i\sigma}v_1^\sigma \mathcal{S}_2^{i\alpha}\int e^{i k_1\cdot b}\frac{\hat\delta(k_1\cdot v_1)\hat\delta(k_1\cdot v_2) k_1^j}{(-k_1^2)^{-\frac{1}{2}}}\,,\\ &
% =\frac{1}{512 \pi|\boldsymbol{b}|^2\sqrt{\gamma^2-1}}\epsilon_{\alpha j i\sigma}v_1^\sigma  \mathcal{S}_2^{i\alpha}\,\hat b^j
 \\ &=-\frac{3}{512 \pi|\boldsymbol{b}|^4\sqrt{\gamma^2-1}}\Big(\hat b\cdot (\vec{\mathcal{S}}_2\times \vec v_1)+\gamma \hat{\mathcal{S}}_2\wedge \hat b\Big),\,\,\textrm{with},\,\,\hat{\mathcal{S}}_2\wedge \hat b=\hat b^j\,\epsilon_{jik}\,\mathcal{S}_2^{ik}
% =\frac{3}{512\pi |\boldsymbol{b}|^4
 %}\Bigg(\frac{2x}{1-x^2}\Bigg)\Big(\big(\frac{1+x^2}{2\,x}\big)^2\, \vec{\mathcal {S}}_2\cdot \epsilon\cdot \hat b \,+\big(\frac{1-x^4}{4\,x^2}\big)\, \hat b\cdot \vec {\mathcal{S}}_2\times \vec v_1\Big)\,.
      \label{3.81}
    \end{split}
\end{align}
  Therefore, the total contribution coming from the dCS term at $\mathcal{O}(\mathcal{S}^1)$ after inserting all the vertex factors is given by,
\begin{align}
    \begin{split}
 \hspace{0cm}i\chi\Big|_{\mathcal{S}^1}&=-il_{dCS} ^2\frac{s_2m_1 m_2^2}{m_p^6}\Bigg(\frac{2x}{1-x^2}\Bigg)\Bigg[\frac{9}{256 \pi |\boldsymbol{b}|^4} (\overrightarrow{\mathcal{S}_2\cdot v_1})\cdot (\vec v_1\times \hat b)-\frac{3}{512\,\pi |\boldsymbol{b}|^4}\Big(\hat b\cdot (\vec{\mathcal{S}}_2\times \vec v_1)+\frac{1+x^2}{2x} \hat{\mathcal{S}}_2\wedge \hat b\Big)\Bigg]\\ &
 =
 \frac{m_1 m_2^2}{m_p^6}\left[\Delta_{\mathfrak{2}}^{\textrm{dCS}}(\overrightarrow{\mathcal{S}_2\cdot v_1})\cdot (\vec v_1\times \hat b)+\Delta_{\mathfrak{3}}^{\textrm{dCS}}\left(\hat b\cdot \left(\vec{\mathcal{S}}_2\times \vec v_1\right)\right)+\Delta_{\mathfrak{4}}^{\textrm{dCS}}\left(\hat{\mathcal{S}}_2\wedge \hat b\right)\right]
    \end{split} \label{eqt24} 
\end{align}
where, $(\overrightarrow{\mathcal{S}_2\cdot v_1})\equiv (\mathcal{S}_2\cdot v_1)^{i},\,\,\vec {\mathcal{S}}_{1,2}\equiv \mathcal{S}_{1,2}^{0i},\,\textrm{and,} \,\,\mathcal{S}_2\wedge \hat b=\epsilon^{\tau\gamma\delta\sigma}v_{2\sigma}\mathcal{S}_{\gamma\delta}\hat b_{\tau}$ 
and,
\begin{align}
    \begin{split}
\Delta_{\mathfrak{2}}^{\textrm{dCS}}=-l_{\textrm{dCS}}^2 \frac{18 xs_2}{256\pi |\boldsymbol{b}|^4(1-x^2)},\,\Delta_{\mathfrak{3}}^{\textrm{dCS}}=l_{\textrm{dCS}}^2 \frac{6 s_2 x}{512 \pi |\boldsymbol{b}|^4(1-x^2)},\,\textrm{and,}\,\Delta_{\mathfrak{4}}^{\textrm{dCS}}=l_{\textrm{dCS}}^2\frac{3s_2(1+x^2)}{512 \pi|\boldsymbol{b}|^4 (1-x^2)}.
    \end{split}
\end{align}
% \textcolor{red}{Finish rest of the term+ insert vertex factor+Final $q$ integral.}\\
%Finish 
%$\epsilon_{ijl}S_2^{ij}\hat b^l=\vec {\mathcal{S}}_2\wedge \hat b$\\
%\newpage
Let us make some comments on the computation of \eqref{3.81}:
\begin{itemize}
\item Due to the presence of the antisymmetric $\epsilon$ tensor, there is no contribution from the dCS term to the eikonal phase at $\mathcal{O}(\mathcal{S}^0)\,.$ It contributes only in the presence of spin.
    \item We ignore the contribution of the two-loop integral proportional to $k_1^i k_1^j$ because of the antisymmetric properties of the epsilon tensor.
    \item Also it is evident from \eqref{eqt24}, that if we take (anti-)aligned spin, the contribution will again give zero due to the presence of binormal like term. In the case of aligned spins, we define
    \begin{align}
\mathcal{S}_1^{\mu\nu}=\frac{2\,\tilde{s}_1}{|\boldsymbol{b}|\sqrt{\gamma^2-1} }b^{[\mu}(\gamma v_1-v_2)^{\nu]} \,,\hspace{0.4 cm }\mathcal{S}_2^{\mu\nu}=\frac{2\,\tilde{s}_2}{|\boldsymbol{b}|\sqrt{\gamma^2-1} }b^{[\mu}( v_1-\gamma v_2)^{\nu]}, \,\,\, \textrm{with},\,\, {\textrm{tr\,}\,(\mathcal{S}_i\cdot\mathcal{S}_i)=-2 \,\tilde{s}_i^2}.
    \end{align}
      \end{itemize}
      \vspace{-0.89 cm}
   After using this parametrization, one finds that the \textit{3PM contribution from the Chern-Simons-type term (arising at the bulk vertices) vanishes.} Therefore, if the two black holes being scattered have aligned spins in dCS theory, there is no contribution to the eikonal phase and hence no correction to the scattering angle.
Now, adding all spinless contributions at 3PM, we have,

\begin{tcolorbox}[thesisresultbox, title=\textit{Spinless Eikonal at 3PM}]
\restorethesisbodyformat
\begin{align}  \label{dCSne3waaa}
 \chi\Big|_{\mathcal{S}^0}=&\frac{m_1m_2^3}{m_p^6}\left[\Delta_{\mathfrak{01}}^{\textrm{poly}}+\Delta_{\mathfrak{02}}^{\textrm{poly}}+\Delta_{\mathfrak{05}}^{\textrm{poly}}\right]
+\\&\frac{m_1^2 m_2^2}{m_p^6}\Big[\Delta_{\mathfrak{08}}^{\textrm{poly}}+\Delta_{\mathfrak{09}}^{\textrm{poly}}+\Delta_{\mathfrak{12}}^{\textrm{poly}}+\Delta_{\mathfrak{15}}^{\textrm{poly}}+\Delta_{\mathfrak{17}}^{\textrm{poly}}+\Delta_{\mathfrak{17}}^{\textrm{poly}}+\Delta_{\mathfrak{18}}^{\textrm{poly}}+\Delta_{\mathfrak{22}}^{\textrm{poly}}+\Delta_{\mathfrak{27}}^{\textrm{poly}}+\Delta_{\mathfrak{30}}^{\textrm{poly}}\nonumber+\\ &\hspace{1.2 cm}\log(x)\Big(\Delta_{\mathfrak{01}}^{\textrm{log}}+\Delta_{\mathfrak{02}}^{\textrm{log}}+\Delta_{\mathfrak{05}}^{\textrm{log}}+\Delta_{\mathfrak{09}}^{\textrm{log}}+\Delta_{\mathfrak{12}}^{\textrm{log}}+\Delta_{\mathfrak{15}}^{\textrm{log}}\Big)+(\log^2(x)+\textbf{Li}_2(1-x^2))\left(\Delta_{\mathfrak{03}}^{\textrm{Polylog}}\right)\Big]\,.\nonumber
    \end{align}
\end{tcolorbox}
\restorethesisbodyformat
Finally, adding all the spinning contributions to the eikonal phase at 3PM we will get,
\begin{tcolorbox}[thesisresultbox, title=\textit{Spinning Eikonal at 3PM}]
\restorethesisbodyformat
\begin{align}\label{dC7Snew}
 \chi\Big|_{\mathcal{S}^1}=&\frac{m_1 m_2^3}{m_p^6}\Big[(\hat b\cdot\mathcal{S}_1\cdot v_2)\Big\{ \Delta_{\mathfrak{03}}^{\textrm{poly}}+\Delta_{\mathfrak{06}}^{\textrm{poly}} \Big\}+(\hat b\cdot\mathcal{S}_2\cdot v_1)\Big\{ \Delta_{\mathfrak{04}}^{\textrm{poly}}+\Delta_{\mathfrak{07}}^{\textrm{poly}} \Big\}+\Delta_{\mathfrak{25}}^{\textrm{poly}}\Big]+ \\ &
\frac{m_1^2 m_2^2}{m_p^6}\Big[(\hat b\cdot \mathcal{S}_1\cdot v_2)\Big\{\Delta_{\mathfrak{10}}^{\textrm{poly}}+\Delta_{\mathfrak{13}}^{\textrm{poly}}+\Delta_{\mathfrak{16}}^{\textrm{poly}}+\Delta_{\mathfrak{19}}^{\textrm{poly}}+\Delta_{\mathfrak{23}}^{\textrm{poly}}+\Delta_{\mathfrak{26}}^{\textrm{poly}}+\Delta_{\mathfrak{28}}^{\textrm{poly}}\nonumber+\\ &\hspace{0.8cm}\log(x)\Big(\Delta_{\mathfrak{3}}^{\textrm{log}}+\Delta_{\mathfrak{06}}^{\textrm{log}}+\Delta_{\mathfrak{08}}^{\textrm{log}}+\Delta_{\mathfrak{10}}^{\textrm{log}}+\Delta_{\mathfrak{13}}^{\textrm{log}}+\Delta_{\mathfrak{16}}^{\textrm{log}}+\Delta_{\mathfrak{18}}^{\textrm{log}}\Big)\nonumber+\\ &\hspace{0.8cm}
(\log^2 x+\textbf{Li}_2(1-x^2))\Big(\Delta_{\mathfrak{1}}^{\textrm{Polylog}}+\Delta_{\mathfrak{4}}^{\textrm{Polylog}}+\Delta_{\mathfrak{6}}^{\textrm{Polylog}}\Big)+\boldsymbol{\mathfrak{U}_1}^{\textrm{LI}_3}\,\boldsymbol{\mathfrak{(UT)}}_3\Big\}\nonumber+\\ & \hspace{1.2cm}(\hat b\cdot \mathcal{S}_2\cdot v_1)\Big\{\Delta_{\mathfrak{11}}^{\textrm{poly}}+\Delta_{\mathfrak{14}}^{\textrm{poly}}+\Delta_{\mathfrak{20}}^{\textrm{poly}}+\Delta_{\mathfrak{21}}^{\textrm{poly}}+\Delta_{\mathfrak{24}}^{\textrm{poly}}+\Delta_{\mathfrak{29}}^{\textrm{poly}}\nonumber+\\&\hspace{0.8cm}\log(x)\Big(\Delta_{\mathfrak{4}}^{\textrm{log}}+\Delta_{\mathfrak{7}}^{\textrm{log}}+\Delta_{\mathfrak{11}}^{\textrm{log}}+\Delta_{\mathfrak{14}}^{\textrm{log}}+\Delta_{\mathfrak{17}}^{\textrm{log}}\Big)\nonumber+\\ &\hspace{0.8cm}(\log^2 x+\textbf{Li}_2(1-x^2))\Big(\Delta_{\mathfrak{2}}^{\textrm{Polylog}}+\Delta_{\mathfrak{5}}^{\textrm{Polylog}}+\Delta_{\mathfrak{7}}^{\textrm{polylog}}\Big)+\boldsymbol{\mathfrak{U}_2}^{\textrm{LI}_3}\,\boldsymbol{\mathfrak{(UT)}}_3\Big\}\Big]\nonumber+\\ &
 \frac{m_1 m_2^2 }{m_p^6}\left[\Delta_{\mathfrak{1}}^{\textrm{dCS}}(\mathcal{S}_2\wedge \hat b )+\Delta_{\mathfrak{2}}^{\textrm{dCS}}(\overrightarrow{\mathcal{S}_2\cdot v_1})\cdot (\vec v_1\times \hat b)+\Delta_{\mathfrak{3}}^{\textrm{dCS}}\left(\hat b\cdot \left(\vec{\mathcal{S}}_2\times \vec v_1\right)\right)+\Delta_{\mathfrak{4}}^{\textrm{dCS}}\left(\hat{\mathcal{S}}_2\wedge \hat b\right)\right]
\,\nonumber
    \end{align} 
\end{tcolorbox}
\restorethesisbodyformat

\noindent

where, $\boldsymbol{\mathfrak{(UT)}}_3$ a polynomial consisting of a particular combination of functions with transcendental weight three (UT-3) and defined as,
\begin{align}
    \begin{split}
\boldsymbol{\mathfrak{(UT)}}_3&=54\,\textbf{Li}_3\left(x^2\right)-216\, \textbf{Li}_3\left(\frac{x+1}{1-x}\right)
\\ &+432\, \textbf{Li}_3(1-x)+216\, \textbf{Li}_3\left(\frac{x+1}{x-1}\right)+432\, \textbf{Li}_3(x+1)+216 \Big(-\textbf{Li}_2\left(\frac{x+1}{1-x}\right)-\textbf{Li}_2(1-x)\\ &+\textbf{Li}_2\left(\frac{x+1}{x-1}\right)+\textbf{Li}_2(x+1)\Big) \log \left(\frac{2}{x+1}-1\right)-54 \left(\pi ^2-2 \log ^2(x)\right) \log \left(1-x^2\right)-20 \log ^3(x)\\ &+108 i \pi  \log ^2(1-x)+108 i \pi  \log ^2(x+1)-11 \pi ^2 \log (x)-432 \zeta (3)\label{4.133a}\,,
    \end{split}
\end{align}
and, also note that one needs to add the terms ($1\leftrightarrow 2$) with \eqref{dCSne3waaa} and \eqref{dC7Snew} to get the complete answer.
On the first note,  it may seem that the polynomial $\boldsymbol{\mathfrak{(UT)}_3}$ in \eqref{4.133a} is complex, in general. However, it can be easily checked (numerically) that it is actually real in the domain $x\in(0,1)$. \textit{To the best of our knowledge, this is the first example for which, unlike GR, we get UT-3 polynomials in the Eikonal phase at 3PM.} 
Again to remind the readers, $\gamma=\frac{1+x^2}{2\,x}\,.$ As we notice that in the expression of eikonal phase, we encounter $\log(b)$ type terms, and one should put an IR cutoff, $\mu_{IR}$, as $\log(\mu_{IR}b)$ to make the argument dimensionless. However, for computing observables like impulses, one needs to take a derivative with respect to $|b|$, and hence, it is independent of the IR cutoff. Also note that, the $\Delta's$ mentioned above are defined in \cref{eqt1,eqt2,eqt3,eqt4,eqt5,eqt6,eqt7,eqt7b,eqt8,eqt9,eqt10,eqt11,eqt12,eqt13,eqt14,eqt15,eqt16,eqt17,eqt18,eqt19,eqt20,eqt21,eqt22,eqt23}.
  
\section{Conclusion and outlook}\label{con}
In \cite{Jakobsen:2021zvh}, the authors introduced spin degrees of freedom in WQFT by implementing  $\mathcal{N}=2$ supersymmetry as an alternative approach to describe a compact spinning object, considering terms up to quadratic-in-spin. This work motivates us to study the scattering event of two spinning black holes up to $3$PM in a gravity theory modified by a dynamical Chern-Simons (dCS) term using the WQFT approach. %We begin by introducing the supersymmetric description of the spin-$1/2$ particles in WQFT in the modified gravity background. It turns out that the $\mathcal{N}=1$ supersymmetry perfectly fits our agenda.
The introduction of a parity-violating Chern-Simons term in Einstein gravity, which couples to gravity via a scalar field, opens up new avenues in the study of astrophysical events, especially gravitational waves. While the Chern-Simons term does not affect the non-rotating black hole solutions, it induces a correction to the rotating metric. Therefore, it becomes important for the study of spinning black holes. 

\textbullet$\,\,$ We start with the Einstein-Hilbert action modified by the dCS term and formulate the worldline theory for a spinning particle by introducing the anti-commuting worldline vectors $\psi^{a}$. Upon modification, $\mathcal{N}=1$  supersymmetry is broken (approximate), but without loss of generality, we define the classical spin of the body $(\mathcal{S}^{ab}=-2i\psi^{[a}\,\psi^{b]})$. Although, as we have mentioned, the SSC is also approximate due to the broken (approximate) supersymmetry, it holds for small coupling (sensitivity parameters) and leading order in a spin throughout the trajectory, but the asymptotic SUSY is still there. It provides the framework to study the scattering of two spinning black holes using the WQFT formalism. 

\textbullet $\,\,$ Our computation deals with two-loop Feynman integrals. We discuss the mathematical framework \cite{Dlapa:2023hsl,Henn:2014qga} used to solve multi-loop integrals via IBP reduction and differential equations, and then analyze the two-loop integrals required for our computation. \textit{We illustrate these techniques by solving \textit{four} master integrals to higher orders in $\epsilon$}.

\textbullet$\,\,$Next, we explicitly determined the eikonal phase for the spinning black-hole system up to 3PM order. We have shown the correction to the spinning eikonal due to the presence of a scalar field up to 3PM. \textit{To the best of our knowledge, this is the first result for the spinning eikonal in the presence of a massless scalar field at 3PM order. We also find from the eikonal phase that the Chern-Simons interaction does not contribute to non-spinning black holes, nor does it contribute when one considers (anti-)aligned spins.}

\textbf{\textit{$i\varepsilon$ prescription and the IR problem}:}

In the eikonal phase computation, one is expected to use the Feynman $i\varepsilon$ prescription to get sensible results. However,  a naive choice of Feynman $i\varepsilon$ prescription could result in unwanted IR divergences ($\propto \frac{1}{\epsilon}$) in the eikonal phase, which we have also encountered in our computations, specifically for those diagrams with the worldline propagator. More precisely, let's say, after IBP reduction, we face the following master integral (in some diagram): $\mathbfcal{M}_{1,1;0,0,1,1,1}$. Now, using the Feynman prescription amounts to the following substitution:
\begin{align}
    \mathbfcal{M}_{1,1;0,0,1,1,1}\to \frac{1}{2}\left( \mathbfcal{M}_{1,1;0,0,1,1,1}^{(++)}-\mathbfcal{M}_{1,1;0,0,1,1,1}^{(+-)}\right)\,. \label{6.1k}
\end{align}
Then, for some (or most) of the cases, this gives rise to IR divergences. As a possible resolution, the authors of \cite{Kim:2024svw} have proposed the notion of Magnus expansion which fixes the choice of the $i\varepsilon$ prescription for cancelling IR divergences. In general, one needs to consider the following prescription for precise IR cancellation,
\begin{align}
    \mathbfcal{M}_{1,1;0,0,1,1,1}\to \left( w_{(++)}\mathbfcal{M}_{1,1;0,0,1,1,1}^{(++)}+w_{(+-)}\mathbfcal{M}_{1,1;0,0,1,1,1}^{(+-)}\right)\,, \label{6.2k}
\end{align}
where the weight factors ($w$ 's) are determined by the Magnus expansion. However, these weight factors may depend on the underlying input, or it might happen that the prescription also may not cancel the IR divergences (i.e., the theory itself has the IR problem). So, it needs to be properly investigated, especially for the theory considered in this paper. For the moment, we have only considered the finite part of the regulated (dimensionally regularized) integral, giving rise to the eikonal phase and leaving aside a more systematic approach to cancel IR divergences for future studies.

\textbf{\textit{Comments on computation of observables from eikonal:}}

In this paper, we mainly computed the eikonal phase using WQFT. Now, the question is, how does one extract the observables, such as impulse (scattering angle), from the eikonal? A possible answer to this question has been discussed recently in \cite{Kim:2024grz,Kim:2024svw}. The eikonal phase can be understood as a generator of the canonical transformations that can transform the states from in-state to out-state in some scattering event. Any instantaneous change of observable in the classical limit can be written as,
    \begin{align}
        \begin{split}
            \Delta\mathcal{O}=\lim_{\hbar\to0}\left[\langle\psi|S^\dagger\mathcal{O}S|\psi\rangle-\langle\psi|O|\psi\rangle\right]\,.
        \end{split}
    \end{align}
 We know that in the classical limit, the S-matrix enjoys eikonal exponentiation, $S=e^{i\chi/\hbar}$. Therefore, 
 \begin{align}
     \begin{split}
         \mathcal{O}_{out}&=e^{i\chi/\hbar}\mathcal{O}_{in}e^{i\chi/\hbar}\\ &=\mathcal{O}_{in}+\frac{1}{i\hbar}[\chi,\mathcal{O}_{in}]+\frac{1}{2(i\hbar)^2}[\chi,[\chi,\mathcal{O}_{in}]]+\cdots\\ &
         \xrightarrow[]{\hbar\to 0}\mathcal{O}_{in}+\{\chi,\mathcal{O}_{in}\}_{P.B.}+\frac{1}{2}\{\chi,\{\chi,\mathcal{O}_{in}\}\}_{P.B.}+\cdots
     \end{split}
 \end{align}
 Now, in WQFT formalism, the eikonal phase can be identified as the free energy of the theory, $\chi=-i\log\,Z_{WQFT}$. Therefore, the observable, e.g impulse, can be computed as,
 \begin{align}
     \Delta p_i^\mu=\{\chi,p_i^\mu\}_{P.B.}+\frac{1}{2}\{\chi,\{\chi,p_i^\mu\}\}_{P.B.}+\cdots
 \end{align}
 In Post-Minkowskian formalism, one can expand the eikonal perturbatively in $G_N$,
 \begin{align}
     \begin{split}
         \chi=\chi_{1PM}+\chi_{2PM}+\chi_{3PM}+\cdots
     \end{split}
 \end{align}
 Therefore, the impulse can be computed as,
 \begin{align}
     \begin{split}
         &\Delta_{(1PM)}\,p_i^\mu=\{\chi_{(1PM)},p_i^\mu\},\\ &
         \Delta_{(2PM)}p_i^\mu=\{\chi_{(2PM)},p_i^\mu\}+\frac{1}{2}\{\chi_{(1PM)},\{\chi_{(1PM)},p_i^\mu\},\\ &
         \Delta_{(3PM)}p_i^\mu=\{\chi_{(3PM)},p_i^\mu\}+\frac{1}{2}\{\chi_{(2PM)},\{\chi_{(1PM)},p_i^\mu\}+\frac{1}{2}\{\chi_{(1PM)},\{\chi_{(2PM)},p_i^\mu\}+\frac{1}{3!}\{\chi_{1PM},\{\chi_{1PM},\{\chi_{1PM},p_i^\mu\}\}\},\\ &
         \vdots\label{5.5a}
     \end{split}
 \end{align}
 To compute the brackets in \eqref{5.5a}, one needs to use the fundamental Poisson bracket relations mentioned in \cite{Kim:2024grz}. We hope to use this methodology to compute the impulse for our case and report it in the near future.


Our work shows the path of several follow-up opportunities. The immediate one is to find the Eikonal phase in $\mathcal{N}=2$ supersymmetric worldline theory, which deals with the interaction of spin-1 particles or quadratic order in spin, including the scalar dipole term in the action. One can also apply the spinning WQFT formalism to find out observables for the dCS theory for higher PM order. Our work can also be naturally extended to find out the observables, e.g., Impulse $(\Delta P^{\mu}),$ the spin kick $(\Delta S^{\mu\nu})$ from the Eikonal Phase and the waveform up to 3PM. More importantly, one needs to be careful about the $i\varepsilon$ prescription and find a way to handle the IR divergences along the way of \cite{Kim:2024svw}.  We hope to report on this soon.


%\newpage
\appendix 
\section{Useful Feynman integrals} \label{ch3:app:A}
In the main text, we extensively use the following integrals. %especially the Fourier ones. 
 %Below, we list master integrals that will be useful in our context. 
 We take the integrals from \cite{Levi:2011eq}.
%\vspace{-1.3cm}
\begingroup
\small
\begin{eqnarray}
\label{ch3:tensorfourierindentity}
\hspace{0cm}I&\equiv& \int \frac{d^d \boldsymbol{k}}{(2\pi)^d}\frac{e^{i\bf{k}\cdot\bf{r}}}{({\bf{k}}^2)^\alpha}=\frac{1}{(4\pi)^{d/2}}\frac{\Gamma(d/2-\alpha)}{\Gamma(\alpha)}\Bigg(\frac{\bf{r}^2}{4}\Bigg)^{\alpha-d/2}, \label{int1}\\ 
I^i\equiv\rmint\frac{d^d\bf{k}}{(2\pi)^d}\frac{k^ie^{i\bf{k}\cdot\bf{r}}}{({\bf{k}}^2)^\alpha}&=&\frac{i}{(4\pi)^{d/2}}\frac{\Gamma(d/2-\alpha+1)}{\Gamma(\alpha)}\left(\frac{{\bf{r}}^2}{4}\right)^{\alpha-d/2-1/2}n^i, \label{int}\\
I^{ij}\equiv\rmint\frac{d^d\bf{k}}{(2\pi)^d}\frac{k^ik^je^{i\bf{k}\cdot\bf{r}}}{({\bf{k}}^2)^\alpha}&=&\frac{1}{(4\pi)^{d/2}}\frac{\Gamma(d/2-\alpha+1)}{\Gamma(\alpha)}\left(\frac{{\bf{r}}^2}{4}\right)^{\alpha-d/2-1}\left(\frac{1}{2}\delta^{ij}+(\alpha-1-d/2)n^in^j\right),\\
I^{ijl}\equiv\rmint\frac{d^d\bf{k}}{(2\pi)^d}\frac{k^ik^jk^le^{i\bf{k}\cdot\bf{r}}}{({\bf{k}}^2)^\alpha}&=&\frac{i}{(4\pi)^{d/2}}\frac{\Gamma(d/2-\alpha+2)}{\Gamma(\alpha)}\left(\frac{{\bf{r}}^2}{4}\right)^{\alpha-d/2-3/2}\nonumber\\
&&\quad
\times\left(\frac{1}{2}\left(\delta^{ij}n^l+\delta^{il}n^j+\delta^{jl}n^i\right)+(\alpha-d/2-2)n^in^jn^l\right),\\
I^{ijlm}\textstyle{\equiv\rmint\frac{d^d\bf{k}}{(2\pi)^d}\frac{k^ik^jk^lk^me^{i\bf{k}\cdot\bf{r}}}{({\bf{k}}^2)^\alpha}}&=&\textstyle{\frac{1}{(4\pi)^{d/2}}\frac{\Gamma(d/2-\alpha+2)}{\Gamma(\alpha)}\left(\frac{{\bf{r}}^2}{4}\right)^{\alpha-d/2-2}\left(\frac{1}{4}\left(\delta^{ij}\delta^{lm}+\delta^{il}\delta^{jm}+\delta^{im}\delta^{jl}\right)\right.} \\
&&
\textstyle{+\frac{\alpha-d/2-2}{2}\left(\delta^{ij}n^ln^m+\delta^{il}n^jn^m+\delta^{im}n^jn^l+\delta^{jl}n^in^m+\delta^{jm}n^in^l+\delta^{lm}n^in^j\right) }\nonumber\\
&&\left.\quad
+(\alpha-d/2-2)(\alpha-d/2-3)n^in^jn^ln^m\right).\nonumber
\end{eqnarray}
\endgroup
We use the d-dimensional master formula for one-loop scalar integrals given by ,
\begin{align}
\begin{split}
   & J\equiv\rmint \frac{d^d\bf{k}}{(2\pi)^d}\frac{1}{\left[{\bf{k}}^2\right]^\alpha\left[({\bf{k}-\bf{q}})^2\right]^\beta}=  \frac{1}{(4\pi)^{d/2}}\frac{\Gamma(\alpha+\beta-d/2)}{\Gamma(\alpha)\Gamma(\beta)}\frac{\Gamma(d/2-\alpha)\Gamma(d/2-\beta)}{\Gamma(d-\alpha-\beta)}\left(q^2\right)^{d/2-\alpha-\beta}.\label{ch3:eq:1loop}
    \end{split}
\end{align}
The d-dimensional master formula for one-loop tensor integrals is taken from \cite{Levi:2011eq}. 
Similarly, one can also derive the following d-dimensional formulas for the one-loop tensor integrals: 
\begin{align}
\begin{split}
J^i\equiv\rmint \frac{d^d\bf{k}}{(2\pi)^d}\frac{k^i}{\left[{\bf{k}}^2\right]^\alpha\left[({\bf{k}-\bf{q}})^2\right]^\beta}=\frac{1}{(4\pi)^{d/2}}\frac{\Gamma(\alpha+\beta-d/2)}{\Gamma(\alpha)\Gamma(\beta)}\frac{\Gamma(d/2-\alpha+1)\Gamma(d/2-\beta)}{\Gamma(d-\alpha-\beta+1)}\\ \left(q^2\right)^{d/2-\alpha-\beta}q^i\,,\end{split}
\end{align}
\begin{align}
\begin{split}
J^{ij}\equiv\rmint \frac{d^d\bf{k}}{(2\pi)^d}\frac{k^ik^j}{\left[{\bf{k}}^2\right]^\alpha\left[({\bf{k}-\bf{q}})^2\right]^\beta}=\frac{1}{(4\pi)^{d/2}}\frac{\Gamma(\alpha+\beta-d/2-1)}{\Gamma(\alpha)\Gamma(\beta)}\frac{\Gamma(d/2-\alpha+1)\Gamma(d/2-\beta)}{\Gamma(d-\alpha-\beta+2)}\left(q^2\right)^{d/2-\alpha-\beta}\\
\times\left(\frac{d/2-\beta}{2}q^2\delta^{ij}+(\alpha+\beta-d/2-1)(d/2-\alpha+1)q^iq^j\right),\end{split}
\end{align}
\begin{align}
\begin{split}
J^{ijl}\equiv\rmint \frac{d^d\bf{k}}{(2\pi)^d}\frac{k^ik^jk^l}{\left[{\bf{k}}^2\right]^\alpha\left[({\bf{k}-\bf{q}})^2\right]^\beta}=\frac{1}{(4\pi)^{d/2}}\frac{\Gamma(\alpha+\beta-d/2-1)}{\Gamma(\alpha)\Gamma(\beta)}\frac{\Gamma(d/2-\alpha+2)\Gamma(d/2-\beta)}{\Gamma(d-\alpha-\beta+3)}\left(q^2\right)^{d/2-\alpha-\beta} \\ 
\times\left(\frac{d/2-\beta}{2}q^2\left(\delta^{ij}q^l+\delta^{il}q^j+\delta^{jl}q^i\right)\right.  \left. 
+(\alpha+\beta-d/2-1)(d/2-\alpha+2)q^iq^jq^l\right),
\end{split}
\end{align}


\begin{align}
\begin{split}
&J^{ijlm}\equiv\rmint \frac{d^d\bf{k}}{(2\pi)^d}\frac{k^ik^jk^lk^m}{\left[{\bf{k}}^2\right]^\alpha\left[({\bf{k}-\bf{q}})^2\right]^\beta}=\frac{1}{(4\pi)^{d/2}}\frac{\Gamma(\alpha+\beta-d/2-2)}{\Gamma(\alpha)\Gamma(\beta)}\\& \frac{\Gamma(d/2-\alpha+2)\Gamma(d/2-\beta)}{\Gamma(d-\alpha-\beta+4)}\left(q^2\right)^{d/2-\alpha-\beta}\,\,\,\,
\times\left(\frac{(d/2-\beta)(d/2-\beta+1)}{4}q^4\left(\delta^{ij}\delta^{lm}+\delta^{il}\delta^{jm}+\delta^{jl}\delta^{im}\right)\right.\\
& \textstyle{
+(\alpha+\beta-d/2-2)(d/2-\alpha+2)\frac{d/2-\beta}{2}q^2 
\times\left(\delta^{ij}q^lq^m+\delta^{il}q^jq^m+\delta^{im}q^jq^l+\delta^{jl}q^iq^m+\delta^{jm}q^iq^l+\delta^{lm}q^iq^j\right)}\\
&\left.\,\,\,\,\,\,\,\,\,
+(\alpha+\beta-d/2-2)(\alpha+\beta-d/2-1)(d/2-\alpha+2)(d/2-\alpha+3)\right.
\left.
\times q^iq^jq^lq^m\right).
\end{split}
\end{align}
More details of  Feynman integrals can be found in \cite{smirnov}.
Another important two loop integral that we use to compute in the potential region is given by,
\begin{align}
    \begin{split}
        \mathbfcal{M}^{(++)}=\int_{\ell_{1},\ell_{2}}\frac{1}{(\ell_{1}\cdot n+i\varepsilon)(\ell_{2}\cdot n+i\varepsilon)}\frac{1}{\ell_{1}^2\, \ell_{2}^2\, (\ell_{1}+\ell_{2}-q)^2}\,.\label{intextra}
    \end{split}
\end{align}
Now, introducing an auxiliary delta functio,n we can write the above integral as \cite{Saotome:2012vy, Parra-Martinez:2020dzs},
\begin{align}
    \begin{split}
        \mathbfcal{M}^{(++)}=\frac{1}{3!}\int_{\ell_{1},\ell_{2},l_3}\Bigg(\frac{1}{(\ell_{1}\cdot n+i\varepsilon)(\ell_{2}\cdot n+i\varepsilon)}+\textrm{perms.}\Bigg)\frac{\delta^{(d)}(l_{123}-q)}{\ell_{1}^2\, \ell_{2}^2\, l_3^2}\,.
    \end{split}
\end{align}
Now using the following $\delta$-function identities,
\begin{align}
    \begin{split}
        &\delta(z_1+z_2+z_3)\Bigg(\frac{1}{z_1+i\varepsilon}\frac{1}{z_{12}+i\varepsilon}+\textrm{perms.}\Bigg)=(2\pi i)^2 \delta(z_1)\delta(z_2)\delta(z_3)\,,\\ &
\delta(z_1+z_2+z_3)\Bigg(\frac{1}{z_1+i\varepsilon}\frac{1}{z_{2}+i\varepsilon}+\textrm{perms.}\Bigg)=2(2\pi i)^2 \delta(z_1)\delta(z_2)\delta(z_3),
    \end{split}
\end{align}
one can rewrite the integral as,
\begin{align}
    \begin{split}
        \mathbfcal{M}^{(++)}=\frac{(2\pi i)^2}{6}\int_{\boldsymbol{\ell_1},\boldsymbol{\ell_2}}\frac{1}{\boldsymbol{\ell_1}^2\,\boldsymbol{\ell_2}^2\,(\boldsymbol{\ell_1}+\boldsymbol{\ell_2}-\boldsymbol{q})^2}\sim -\frac{\Gamma^3(-\epsilon)\Gamma(1+2\epsilon)}{\Gamma(-3\epsilon)}\,.
    \end{split}
\end{align}
In principle, the integral \eqref{intextra}, can be computed directly using Schwinger parametrization, but the above-mentioned approach is useful when we go beyond two loops.
\section{Boundary integrals in potential mode relevant for solving the differential equation}\label{ch3:app:B}
\textbf{Computation of boundary integrals in potential region:}\\ Now we compute the boundary integrals relevant to the four master integrals as mentioned previously for our subsequent computations.
We start with the master with two linear propagators:
$\,\,$\begin{align}
    \begin{split}
        \mathbfcal{M}_{1,1;0,0,1,1,1}^{(-+)}\Big|_{v_{\infty}\to 0}^{\textrm{pot.}}&=-v_{\infty}^{-2}\int_{\boldsymbol{\ell_1},\boldsymbol{\ell_2}}\frac{1}{(-\boldsymbol{\ell_1}\cdot \boldsymbol{n}+i\varepsilon)(\boldsymbol{\ell_2}\cdot \boldsymbol{n}+i\varepsilon)(\boldsymbol{\ell_1}+\boldsymbol{\ell_2}-\boldsymbol{q})^2(\boldsymbol{q}-\boldsymbol{\ell_1})^2(\boldsymbol{q}-\boldsymbol{\ell_2})^2}\,.\label{2.30}
    \end{split}
\end{align}
In \eqref{2.30}, the $i\varepsilon$  prescription is very important, and one needs to carefully take care of it. Our goal is to make relation with $\mathcal{M}^{(++)}$ with $\mathcal{M}^{+-}$. Now doing the following change of variable ($\boldsymbol{q}-\boldsymbol{\ell_i}\to \boldsymbol{\ell_i}$) we left with 
\begin{align}
    \begin{split}
        \mathbfcal{M}_{1,1;0,0,1,1,1}^{(-+)}\Big|_{v_{\infty}\to 0}^{\textrm{pot.}}&=-v_{\infty}^{-2}\int_{\boldsymbol{\ell_1},\boldsymbol{\ell_2}}\frac{1}{(\boldsymbol{\ell_1}\cdot \boldsymbol{n}+i\varepsilon)(-\boldsymbol{\ell_2}\cdot \boldsymbol{n}+i\varepsilon)(\boldsymbol{\ell_1}+\boldsymbol{\ell_2}-\boldsymbol{q})^2\,\boldsymbol{\ell_1}^2\,\boldsymbol{\ell_2}^2}\,.
    \end{split}
\end{align}
Now again we do the following relabelling: $\boldsymbol{\ell_1}\to \boldsymbol{\ell_1}+\boldsymbol{\ell_2}-\boldsymbol{q}$, $\boldsymbol{\ell_2}\to -\boldsymbol{\ell_2}$ and, ${\boldsymbol{q}\to -\boldsymbol{q}}\,.$ The we get,
\begin{align}
    \begin{split}
        \mathbfcal{M}_{1,1;0,0,1,1,1}^{(-+)}\Big|_{v_{\infty}\to 0}^{\textrm{pot.}}=-v_\infty^{-2}\int_{\boldsymbol{\ell_1},\boldsymbol{\ell_2}}\frac{1}{(\boldsymbol{\ell_1}\cdot \boldsymbol{n}+\boldsymbol{\ell_2}\cdot \boldsymbol{n}+i\varepsilon)(\boldsymbol{\ell_2}\cdot \boldsymbol{n}+i\varepsilon)(\boldsymbol{\ell_1}+\boldsymbol{\ell_2}-\boldsymbol{q})^2\,\boldsymbol{\ell_1}^2\,\boldsymbol{\ell_2}^2}\,.
    \end{split}
\end{align}
Now interchanging $\boldsymbol{\ell_{1}}\leftrightarrow \boldsymbol{\ell}_2$ and adding we will get,
\begin{align}
    \begin{split}
        & 2\mathbfcal{M}_{1,1;0,0,1,1,1}^{(-+)}\Big|_{v_{\infty}\to 0}^{\textrm{pot.}}=-v_\infty^{-2}\int_{\boldsymbol{\ell_{1},\ell_{2}}}\frac{1}{(\boldsymbol{\ell_1}\cdot \boldsymbol{n}+i\varepsilon)(\boldsymbol{\ell_2}\cdot \boldsymbol{n}+i\varepsilon)(\boldsymbol{\ell_1}+\boldsymbol{\ell_2}-\boldsymbol{q})^2\,\boldsymbol{\ell_1}^2\,\boldsymbol{\ell_2}^2}\equiv\mathbfcal{M}_{1,1;0,0,1,1,1}^{(++)}\Big|_{v_{\infty}\to 0}^{\textrm{pot.}}\,
    \end{split}
\end{align}
and the integral is given by,
\begin{align}
    \mathbfcal{M}^{(++)}_{1,1;0,0,1,1,1}=\frac{1}{v_\infty^2}\frac{1}{(4\pi)^{2-2\epsilon}}\frac{\Gamma^3(-\epsilon)\Gamma(1+2\epsilon)}{3\Gamma(-3\epsilon)}\,.
\end{align}
In the main text, we mainly face the following form of two-loop boundary integral.
\begin{align}
    &\mathbfcal{M}^{(\pm\pm)\textrm{pot.}}_{\alpha_1,\alpha_2;\beta_1,\cdots \beta_5}=\int_{\boldsymbol{\ell_2},\boldsymbol{\ell_2}}\frac{1}{(\pm \boldsymbol{\ell_2}\cdot \boldsymbol{n}+i\varepsilon)^{\alpha_1}(\pm \boldsymbol{\ell_2}\cdot \boldsymbol{n}+i\varepsilon)^{\alpha_2}\,\boldsymbol{D}_1^{\beta_1}\cdots \boldsymbol{D}_5^{\beta_5}}
\end{align}
To compute the integral, one can further use the IBP reduction procedure with the following 7 basis functions
\begin{align}
    \{D\}: \{\boldsymbol{\ell_1}\cdot \boldsymbol{n},\,\boldsymbol{\ell_2}\cdot \boldsymbol{n},\,\boldsymbol{\ell_2}^2,\,\boldsymbol{\ell_2}^2,\,(\boldsymbol{\ell_2}+\boldsymbol{\ell_2}-\boldsymbol{q})^2,\,(\boldsymbol{\ell_2}-\boldsymbol{q})^2,\,(\boldsymbol{\ell_2}-\boldsymbol{q})^2\}\,.
\end{align}
After doing the IBP reduction one will get 15 unique sectors with 22 master integrals.
Particularly, we need the following two loop boundary integrals,
\begin{align}
    \begin{split}
        &\mathbfcal{M}^{(\pm\pm)\textrm{pot.}}_{0,0;1,1,1,1,1}=-\frac{2 \left(1-36 \epsilon ^2\right)}{\left(\boldsymbol q^2\right)^2 (2 \epsilon +1)^2}\mathbfcal{M}^{(\pm\pm)\textrm{pot.}}_{0,0;0,0,1,1,1}-\frac{4 \epsilon }{\boldsymbol q^2 (2 \epsilon +1)}\mathbfcal{M}^{(\pm\pm)\textrm{pot.}}_{0,0;1,1,0,1,1}\,,\\ &
        \mathbfcal{M}^{(\pm\pm)\textrm{pot.}}_{0,0;1,1,2,1,1}=\frac{12 (\epsilon +1) (2 \epsilon -1) (36 \epsilon^2 -1) }{\left(\boldsymbol q^2\right)^3 (2 \epsilon +1) (2 \epsilon +3)^2}\mathbfcal{M}^{(\pm\pm)\textrm{pot.}}_{0,0;0,0,1,1,1}+\frac{8 \epsilon }{\left(\boldsymbol q^2\right)^2 (2 \epsilon +3)}\mathbfcal{M}^{(\pm\pm)\textrm{pot.}}_{0,0;1,1,0,1,1}\,,\\ &
  \mathbfcal{M}^{(\pm\pm)\textrm{pot.}}_{0,0;0,0,2,1,1}=      \frac{2 \epsilon  (6 \epsilon -1)}{\boldsymbol q^2 (2 \epsilon +1)}\mathbfcal{M}^{(\pm\pm)\textrm{pot.}}_{0,0;0,0,1,1,1}=\mathbfcal{M}^{(\pm\pm)\textrm{pot.}}_{0,0;0,0,1,2,1}
    \end{split}
\end{align}
where,
\begin{align}
    \begin{split}
        &\mathbfcal{M}^{(\pm\pm)\textrm{pot.}}_{0,0;0,0,1,1,1}= \frac{1}{(4\pi)^{3-2\epsilon}}\frac{\Gamma(1/2-\epsilon)^3\Gamma(2\epsilon)}{\Gamma(3/2-3\epsilon)}\frac{1}{\boldsymbol{q}^{4\epsilon}}\,,\\ &
        \mathbfcal{M}^{(\pm\pm)\textrm{pot.}}_{0,0;1,1,0,1,1}=\frac{1}{(4\pi)^{3-2\epsilon}}\frac{\Gamma^4(1/2-\epsilon)\Gamma^2(1/2+2\epsilon)}{\Gamma^2(1-2\epsilon)}\,.
    \end{split}
\end{align}
Apart from that we will face the following integrals,
\begin{align}
    \begin{split}
        &
        \mathbfcal{M}^{(\pm\pm)\textrm{pot.}}_{0,0;0,1,1,0,1}=0\,,\\ &
\mathbfcal{M}^{(\pm\pm)\textrm{pot.}}_{0,1;0,0,1,1,1}=-i \frac{\sqrt{\pi}}{(4\pi)^3}\frac{\Gamma(1/2-2\epsilon)\Gamma^2(1/2-\epsilon)\Gamma(-\epsilon)\Gamma(1/2+2\epsilon)}{\Gamma(1/2-3\epsilon)\Gamma(1-2\epsilon)}\,,\\ &    \mathbfcal{M}^{(\pm\pm)\textrm{pot.}}_{0,2;1,0,1,1,0}=-\frac{1}{(4\pi)^3}\frac{4\epsilon\Gamma(2\epsilon)\Gamma^2(-2\epsilon)\Gamma(1/2-\epsilon)\Gamma(1/2+\epsilon)}{\Gamma(-4\epsilon)}\,,\\ &
\mathbfcal{M}^{(\pm\pm)\textrm{pot.}}_{1,1;1,1,0,1,1}=-\frac{\pi}{(4\pi)^3}\frac{\Gamma^4(-\epsilon)\Gamma^2(\epsilon+1)}{\Gamma^2(-2\epsilon)}\,,\\ &
\mathbfcal{M}^{(\pm\pm)\textrm{pot.}}_{0,1;1,0,1,1,0}=-i\frac{\pi}{(4\pi)^3}\frac{\Gamma(\epsilon)\Gamma(1/2-2\epsilon)\Gamma(1/2+2\epsilon)}{\Gamma(1-\epsilon)}\,,\\ &
%\mathbfcal{M}^{(\pm\pm)\textrm{pot.}}_{0,1;1,1,0,1,1}=\\ &
\mathbfcal{M}^{(++)\textrm{pot.}}_{1,1;0,0,1,1,1}=2\mathbfcal{M}^{(+-)\textrm{pot.}}_{1,1;0,0,1,1,1}=\frac{1}{(4\pi)^{2-2\epsilon}}\frac{\Gamma^3(-\epsilon)\Gamma(1+2\epsilon)}{3\Gamma(-3\epsilon)}\,.
    \end{split}
\end{align}
\section{\texorpdfstring{$A(x,\epsilon)$ in $\epsilon$ factorized form}{A(x, epsilon) in epsilon-factorized form}} \label{ch3:app:C}
%\vspace{-0.5 cm}

As mentioned in the main text, the $A(x,\epsilon)$ matrix constructed from the coupled differential equations required for solving the master integrals reduces to an  $\epsilon$ factorized form ($\epsilon\,\mathbb{S}(x)$) in UT basis. Here, we explicitly  write down the expression for $\mathbb{S}(x)$, which we have used to solve the master integrals. 

%\begin{landscape}
\begin{equation}\resizebox{\linewidth}{!}{$
    \mathbb{S}(x)=\left(
\begin{array}
{cccccccccccccccc}
 -\frac{18 x^2+9}{2 x-2 x^3} & -\frac{2 x^2+3}{4 x-4 x^3} & -\frac{6 x}{x^2-1} & 0 & 0 & 0 & 0 & 0 & 0 & 0 & 0 & 0 & 0 & 0 & 0 & 0 \\
 \frac{6 x^2+15}{x-x^3} & \frac{6 x^2+5}{2 x-2 x^3} & \frac{12 x}{x^2-1} & 0 & 0 & 0 & 0 & 0 & 0 & 0 & 0 & 0 & 0 & 0 & 0 & 0 \\
 \frac{6}{x} & -\frac{1}{3 x} & -\frac{2}{x} & 0 & 0 & 0 & 0 & 0 & 0 & 0 & 0 & 0 & 0 & 0 & 0 & 0 \\
 0 & 0 & 0 & \frac{2 (x-4) x+2}{x \left(x^2-1\right)} & 0 & 0 & 0 & 0 & 0 & 0 & 0 & 0 & 0 & 0 & 0 & 0 \\
 0 & 0 & 0 & 0 & \frac{2 \left(x^2+1\right)}{x-x^3} & 0 & 0 & 0 & 0 & 0 & 0 & 0 & 0 & 0 & 0 & 0 \\
 0 & 0 & 0 & 0 & 0 & \frac{2 \left(x^2+1\right)}{x-x^3} & 0 & 0 & 0 & 0 & 0 & 0 & 0 & 0 & 0 & 0 \\
 0 & 0 & 0 & 0 & 0 & 0 & 0 & 0 & 0 & 0 & 0 & 0 & 0 & 0 & 0 & 0 \\
 0 & 0 & 0 & -\frac{1}{x} & 0 & 0 & 0 & 0 & 0 & 0 & 0 & 0 & 0 & 0 & 0 & 0 \\
 0 & 0 & 0 & 0 & 0 & -\frac{1}{x} & 0 & 0 & 0 & 0 & 0 & 0 & 0 & 0 & 0 & 0 \\
 0 & 0 & 0 & 0 & \frac{336}{x^3-x} & 0 & 0 & 0 & 0 & 0 & 0 & 0 & 0 & 0 & 0 & 0 \\
 \frac{211-286 x^2}{15 x-15 x^3} & \frac{117-242 x^2}{30 x-30 x^3} & -\frac{97-122 x^2}{5 x-5 x^3} & 0 & \frac{84}{x-x^3} & 0 & 0 & 0 & 0 & 0 & -\frac{1-5 x^2}{x-x^3} & -\frac{2}{3 x} & 0 & 0 & 0 & 0 \\
 -\frac{3983-3481 x^2}{20 x-20 x^3} & \frac{2687 x^2+231}{40 x-40 x^3} & \frac{1821-4371 x^2}{20 x-20 x^3} & 0 & \frac{126}{x-x^3} & 0 & \frac{24}{x^2-1} & 0 & 0 & 0 & \frac{3-75 x^2}{2 x-2 x^3} & \frac{1-5 x^2}{x-x^3} & 0 & 0 & 0 & 0 \\
 0 & 0 & 0 & -\frac{16 (x (7 x+6)+7)}{5 x \left(x^2-1\right)} & 0 & \frac{64}{x^2-1} & 0 & 0 & 0 & 0 & 0 & 0 & -\frac{4}{x^2-1} & 0 & 0 & 0 \\
 0 & 0 & 0 & 0 & 0 & 0 & 0 & 0 & 0 & 0 & 0 & 0 & 0 & 0 & 0 & 0 \\
 \frac{5}{x} & \frac{25}{6 x} & -\frac{5}{x} & 0 & 0 & 0 & 0 & 0 & 0 & 0 & 0 & 0 & 0 & 0 & 0 & 0 \\
 0 & 0 & 0 & 0 & 0 & 0 & 0 & 0 & 0 & 0 & 0 & 0 & 0 & 0 & 0 & 0 
\end{array}\right)\,.
$}\end{equation}
\section{Iterative UT integrals}\label{ch3:app:D}
\begingroup
\small
\begin{align}
    \begin{split}
     &   \boldsymbol{\mathscr{J}}(\{-1,-1,0\},x)\equiv  -\textbf{Li}_3\left(\frac{1}{x+1}\right)+\frac{1}{6} \log ^3(x+1)-\frac{1}{12} \pi ^2 \log (x+1)+\frac{7 \zeta (3)}{8},\\ &
     \boldsymbol{\mathscr{J}}(\{-1,0,0\},x)\equiv -\textbf{Li}_3(-x)+\textbf{Li}_2(-x) \log (x)+\frac{1}{2} \log (x+1) \log ^2(x)-\frac{3 \zeta (3)}{4},\\ &
    \boldsymbol{\mathscr{J}}(\{-1,1,0\},x)\equiv  \textbf{Li}_3\left(\frac{x+1}{2}\right)+\textbf{Li}_3\left(\frac{1+x}{x}\right)+\textbf{Li}_3(x)-\textbf{Li}_3\left(\frac{x+1}{2 x}\right)+\log (2) \textbf{Li}_2\left(\frac{x+1}{2}\right)\\ &-\log (2) \textbf{Li}_2(x)-\log (2) \textbf{Li}_2\left(\frac{x+1}{2 x}\right)-\textbf{Li}_2\left(\frac{x+1}{2}\right) \log (x+1)+\textbf{Li}_2\left(\frac{1+x}{x}\right) \log \left(\frac{2 x}{x+1}\right)-\textbf{Li}_2(1-x) \log (x+1)\\ &-\textbf{Li}_2(x) \log (x)-\textbf{Li}_2\left(\frac{x+1}{2 x}\right) \log (x)+\textbf{Li}_2\left(\frac{x+1}{2 x}\right) \log (x+1)-\log ^2(2) \log (x)-\frac{1}{2} \log (2) \log ^2(x)\\ &+\log (2) \log (x) \log (x+1)-\log (1-x) \log (x) \log (x+1)-\frac{15 \zeta (3)}{8}+\frac{1}{2} i \pi  \log ^2(2)-\frac{1}{12} \pi ^2 \log (2),\\ &
      \boldsymbol{\mathscr{J}}(\{0,-1,0\},x)\equiv 2 \textbf{Li}_3(-x)-\textbf{Li}_2(-x) \log (x)+\frac{1}{12} \pi ^2 \log (x)+\frac{3 \zeta (3)}{2},\\ &
       \boldsymbol{\mathscr{J}}(\{0,0,0\},x)\equiv \frac{\log ^3(x)}{6},\\ &
       \boldsymbol{\mathscr{J}}(\{0,1,0\},x)\equiv 2 \textbf{Li}_3(x)-\textbf{Li}_2(x) \log (x)-\frac{1}{6} \pi ^2 \log (x)-2 \zeta (3),\\ &
   \boldsymbol{\mathscr{J}}(\{1,0,0\},x)\equiv     -\textbf{Li}_3(x)+\textbf{Li}_2(x) \log (x)+\frac{1}{2} \log (1-x) \log ^2(x)+\zeta (3),\\ &
   \boldsymbol{\mathscr{J}}(\{1,1,0\},x)\equiv-\textbf{Li}_3(1-x),\\ &
   \boldsymbol{\mathscr{J}}(\{1,-1,0\},x)\equiv-i \pi  \textbf{Li}_2\left(\frac{x}{x-1}\right)+i \pi  \textbf{Li}_2\left(\frac{2 x}{x-1}\right)+i \pi  \textbf{Li}_2\left(\frac{1}{x+1}\right)+\textbf{Li}_3\left(\frac{1-x}{2}\right)+\textbf{Li}_3\left(\frac{x}{x-1}\right)-\textbf{Li}_3\left(\frac{2 x}{x-1}\right)\\ &+\textbf{Li}_3(-x)-\log (2) \textbf{Li}_2\left(\frac{x}{x-1}\right)+\log (2) \textbf{Li}_2\left(\frac{2 x}{x-1}\right)-\textbf{Li}_2\left(\frac{2 x}{x-1}\right) \log (1-x)-\textbf{Li}_2\left(\frac{x}{x-1}\right) \log (x)\\ &+\textbf{Li}_2\left(\frac{1-x}{2}\right) \left(\log \left(-\frac{2}{x-1}\right)+i \pi \right)+\textbf{Li}_2\left(\frac{x}{x-1}\right) \log (1-x)+\textbf{Li}_2\left(\frac{2 x}{x-1}\right) \log (x)+\textbf{Li}_2\left(\frac{1}{x+1}\right) \log (1-x)\\ &+\textbf{Li}_2(-x) (-\log (2 x)-2 i \pi )-\log ^2(2) \log (1-x)+\frac{1}{2} \log (2) \log ^2(1-x)+\frac{1}{2} i \pi  \log ^2(x+1)+\frac{1}{2} \log ^2(x+1) \log (1-x)\\ &-i \pi  \log (2) \log (1-x)-i \pi  \log (x) \log (x+1)-\frac{1}{12} \pi ^2 \log (1-x)-\log (x) \log (x+1) \log (1-x)+\frac{3 \zeta (3)}{4}-\frac{i \pi ^3}{4}\\ &+\frac{\log ^3(2)}{3}+\frac{1}{2} i \pi  \log ^2(2)-\frac{1}{4} \pi ^2 \log (2).
    \end{split}
\end{align}
\endgroup
%$\boldsymbol{\mathscr{J}}$
\section{\texorpdfstring{$\boldsymbol{\tilde{h}}$ coefficients}{tilde-h coefficients}} \label{ch3:app:E}
Note that, $\gamma=\frac{x^2+1}{2 x}\,.$
\begin{align}
    \begin{split}
   &  \boldsymbol{\tilde h_1}=4,\,\,
      \boldsymbol{\boldsymbol{\tilde h_2}}   =0,
     \boldsymbol{\tilde h}_3=\frac{72 \epsilon ^2-2}{12 \left(\gamma ^2-1\right)^2},\,\,
        \boldsymbol{\tilde h}_4=\frac{6 \gamma ^2+2 \epsilon -3}{12 \left(\gamma ^2-1\right)},\\& \boldsymbol{\tilde h}_5=\frac{1}{24 \left(\gamma ^2-1\right)^2 (\epsilon -1) \epsilon }\Bigg[6 \gamma ^4+72 \gamma ^2 \epsilon ^4-10 \gamma ^2-36 \left(\gamma ^4+\gamma ^2-1\right) \epsilon ^3\\&+\left(-36 \gamma ^6+150 \gamma ^4-116 \gamma ^2\right) \epsilon ^2+\left(36 \gamma ^6-120 \gamma ^4+110 \gamma ^2-25\right) \epsilon +4\Bigg],
       \boldsymbol{\tilde h}_6=-\frac{4 \gamma ^2 \epsilon ^2+2 \left(\gamma ^2+1\right) \epsilon -4 \gamma ^2+1}{24 \left(\gamma ^2-1\right)^2 (\epsilon -1)}\,,\\ &
       \boldsymbol{\tilde h}_7=\frac{-6 \gamma ^3 (\epsilon  (4 \epsilon -3)+1)-6 \gamma ^2 \epsilon +\gamma  \left(2 \epsilon  \left(\epsilon  \left(-24 \epsilon ^2+22 \epsilon +21\right)-14\right)+7\right)+3 \epsilon  (6 \epsilon +1)}{24 \left(\gamma ^2-1\right)^2 (\epsilon -1) \epsilon }\,,\\ &
   \boldsymbol{\tilde h}_8    =\frac{\gamma  \left(8 \epsilon ^3-13 \epsilon +4\right)}{48 \left(\gamma ^2-1\right) (\epsilon -1) (3 \epsilon -1)}\,,
    \boldsymbol{\tilde h}_9 =\frac{\gamma  \left(-9 \gamma ^2-2 (\epsilon -5) \epsilon +8\right)}{6 \left(\gamma ^2-1\right)^2}\,,
    \\&
    \boldsymbol{\tilde h}_{10}  =-\frac{\gamma}{24 \left(\gamma ^2-1\right)^2 (3 \epsilon -1)}\Bigg[8 \gamma ^4 \left(18 \epsilon ^2-21 \epsilon +5\right) \epsilon +3 \gamma ^2 \left(48 \epsilon ^4-208 \epsilon ^3+216 \epsilon ^2-63 \epsilon +4\right)\\&-288 \epsilon ^5+216 \epsilon ^4+464 \epsilon ^3-558 \epsilon ^2+173 \epsilon -14\Bigg]\,,
    \\ &\boldsymbol{\tilde h}_{11}=\frac{-2 \gamma ^4+72 \gamma ^2 \epsilon ^4+2 \gamma ^2-36 \left(\gamma ^4+\gamma ^2-1\right) \epsilon ^3+\left(12 \gamma ^6+6 \gamma ^4-20 \gamma ^2\right) \epsilon ^2+\left(-12 \gamma ^6+32 \gamma ^4-18 \gamma ^2-1\right) \epsilon }{24 \left(\gamma ^2-1\right)^2 (\epsilon -1) \epsilon}\,,\\ &
   \boldsymbol{\tilde h}_{12}=\frac{-8 \gamma ^4+4 \gamma ^2 \epsilon ^2+8 \gamma ^2+\left(8 \gamma ^4-6 \gamma ^2+2\right) \epsilon -3}{24 \left(\gamma ^2-1\right)^2 (\epsilon -1)}\,,\\&
  \boldsymbol{\tilde h}_{13}  =\frac{\gamma  \left(4 \gamma ^3 (\epsilon -1) (4 \epsilon -1)-12 \gamma ^2 \epsilon +2 \gamma  (\epsilon -1) (4 \epsilon -1) (4 \epsilon  (3 \epsilon +1)-3)+\epsilon  (42 \epsilon +5)\right)}{48 \left(\gamma ^2-1\right)^2 (\epsilon -1) \epsilon }\,,\\ &
 \boldsymbol{\tilde h}_{14}=\frac{\gamma  \left(8 \gamma  (\epsilon -1) \left(\gamma ^2 (8 \epsilon -2)-2 \epsilon  (\epsilon +5)+3\right)-\epsilon \right)}{96 \left(\gamma ^2-1\right) (\epsilon -1) (3 \epsilon -1)}\,,
  \\ &
  \boldsymbol{\tilde h}_{15} =\frac{\gamma  \left(4 \gamma ^4 \epsilon ^2+\gamma ^2 (2 \epsilon  (\epsilon  (12 \epsilon -11)-3)+2)+\epsilon  (16 \epsilon  (\epsilon  (2-3 \epsilon )+1)-1)-1\right)}{12 \left(\gamma ^2-1\right)^2 (\epsilon -1) \epsilon }\,,\\ &
  \boldsymbol{\tilde h}_{16}=\frac{\gamma  \left(\gamma ^4 (4-20 \epsilon )+4 \gamma ^2 (2 \epsilon  (\epsilon +6)-3)-\epsilon  (8 \epsilon +27)+8\right)}{24 \left(\gamma ^2-1\right)^2 (3 \epsilon -1)}\,,\\&\boldsymbol{\tilde h}_{17}=\frac{1}{8 \left(\gamma ^2-1\right)^2 (\epsilon -1)^2 (\epsilon +1) (2 \epsilon +1)}\Bigg[-4 \gamma ^6 \left(108 \epsilon ^5-132 \epsilon ^4-109 \epsilon ^3+142 \epsilon ^2+\epsilon -10\right)\\&+2 \gamma ^4 \left(432 \epsilon ^6-432 \epsilon ^5-792 \epsilon ^4+308 \epsilon ^3+539 \epsilon ^2-16 \epsilon -39\right)\\&+2 \gamma ^2 \left(576 \epsilon ^5+156 \epsilon ^4-596 \epsilon ^3-317 \epsilon ^2+20 \epsilon +21\right)+216 \epsilon ^4+324 \epsilon ^3+138 \epsilon ^2-9 \epsilon -4\Bigg]\,,\nonumber
       \end{split}
\end{align}
\begin{align}
    \begin{split}
&\boldsymbol{\tilde h}_{18}=-\frac{\left(2 \gamma ^2 (\epsilon -1)+1\right) \left(4 \gamma ^4 \left(\epsilon ^2-1\right)-2 \gamma ^2 \left(\epsilon  \left(6 \epsilon ^2+\epsilon -3\right)-2\right)-3 \epsilon  (2 \epsilon +1)+1\right)}{16 \left(\gamma ^2-1\right) (\epsilon -1)^2 (\epsilon +1)}\,,\\&\boldsymbol{\tilde h}_{19}=-\frac{1}{96 (\gamma -1)^2 (\gamma +1)^2 (\epsilon -1)^2 \epsilon  (\epsilon +1) (2 \epsilon +1)}\Bigg[2 \gamma ^5 \epsilon  \big(1152 \epsilon ^7+4312 \epsilon ^6-4644 \epsilon ^5-7706 \epsilon ^4+3361 \epsilon ^3\\&+3347 \epsilon ^2+107 \epsilon +23\big)+2 \gamma ^3 \big(1728 \epsilon ^9-16416 \epsilon ^8+1912 \epsilon ^7+26492 \epsilon ^6+3490 \epsilon ^5-12697 \epsilon ^4-5305 \epsilon ^3\\&+282 \epsilon ^2+35 \epsilon -1\big)+192 \gamma ^2 (\epsilon -1)^2 \epsilon  \left(72 \epsilon ^5+108 \epsilon ^4-70 \epsilon ^3-99 \epsilon ^2+13 \epsilon +6\right)+\gamma  \big(-10368 \epsilon ^9+11520 \epsilon ^8\\&-1888 \epsilon ^7-20272 \epsilon ^6+8000 \epsilon ^5+14004 \epsilon ^4+3092 \epsilon ^3-630 \epsilon ^2-93 \epsilon +1\big)+96 (\epsilon -1)^2 \epsilon  \big(72 \epsilon ^4+172 \epsilon ^3+82 \epsilon ^2\\&-27 \epsilon -9\big)\Bigg]\,,\\&\boldsymbol{\tilde h}_{20}=\frac{1}{96 \left(\gamma ^2-1\right)^2 (\epsilon -1)^2 \epsilon  (\epsilon +1)}\Bigg[\gamma  \big(12 \gamma ^8 \epsilon ^3 \left(\epsilon ^2-1\right)-36 \gamma ^6 \epsilon ^2 \left(2 \epsilon ^4-4 \epsilon ^3-9 \epsilon ^2+4 \epsilon +7\right)\\&+4 \gamma ^4 \epsilon  \left(90 \epsilon ^5-126 \epsilon ^4-276 \epsilon ^3+114 \epsilon ^2+187 \epsilon +11\right)-2 \gamma ^2 \left(252 \epsilon ^6-192 \epsilon ^5-726 \epsilon ^4+204 \epsilon ^3+439 \epsilon ^2+22 \epsilon +1\right)\\&+216 \epsilon ^6-180 \epsilon ^5-576 \epsilon ^4+192 \epsilon ^3+360 \epsilon ^2-11 \epsilon +1\big)\Bigg]\,,\\&\boldsymbol{\tilde h}_{21}=\frac{1}{48 (\gamma -1)^2 (\gamma +1)^2 (\epsilon -1)^2 \epsilon  (\epsilon +1) (2 \epsilon +1)}\Bigg[4 \gamma ^5 \big(1296 \epsilon ^7-360 \epsilon ^6-2348 \epsilon ^5+304 \epsilon ^4+1049 \epsilon ^3+48 \epsilon ^2-6 \epsilon\\& -1\big)+360 \gamma ^4 (\epsilon -1)^2 \epsilon  \left(2 \epsilon ^2+3 \epsilon +1\right)-4 \gamma ^3 \big(2592 \epsilon ^8+864 \epsilon ^7-7644 \epsilon ^6-1168 \epsilon ^5+5281 \epsilon ^4-24 \epsilon ^3-36 \epsilon ^2\\&+100 \epsilon -1\big)-2 \gamma ^2 \left(360 \epsilon ^6-60 \epsilon ^5-694 \epsilon ^4-193 \epsilon ^3+333 \epsilon ^2+253 \epsilon +1\right)-2 \gamma  \epsilon  \big(3456 \epsilon ^6+5136 \epsilon ^5-4484 \epsilon ^4\\&-5268 \epsilon ^3+1339 \epsilon ^2+45 \epsilon -188\big)-408 \epsilon ^5-324 \epsilon ^4+250 \epsilon ^3+323 \epsilon ^2+158 \epsilon +1\Bigg]\,,\\&\boldsymbol{\tilde h}_{22}=\frac{1}{48 \left(\gamma ^2-1\right)^2 (\epsilon -1)^2 \epsilon  (\epsilon +1) (2 \epsilon +1)}\Bigg[8 \gamma ^5 \left(6 \epsilon ^5-3 \epsilon ^4-9 \epsilon ^3+\epsilon ^2+3 \epsilon +2\right)-4 \gamma ^3 \big(36 \epsilon ^6-12 \epsilon ^5-45 \epsilon ^4\\&-21 \epsilon ^3+17 \epsilon ^2+18 \epsilon +7\big)+2 \gamma ^2 (\epsilon -1)^2 \left(2 \epsilon ^2+3 \epsilon +1\right)-6 \gamma  (2 \epsilon +1)^2 \left(3 \epsilon ^3-2 \epsilon ^2-\epsilon -1\right)+2 \epsilon ^3+\epsilon ^2-2 \epsilon -1\Bigg]\,,\\&
        \boldsymbol{\tilde h}_{23}=\frac{4 \gamma ^4 \left(2 \epsilon ^3-3 \epsilon ^2+5 \epsilon -4\right)+2 \gamma ^2 \left(4 \epsilon ^2-8 \epsilon +7\right)+8 \gamma ^6 (\epsilon -1)^2+2 \epsilon -3}{48 \left(\gamma ^2-1\right) (\epsilon -1)^2}\,,\\&\boldsymbol{\tilde h}_{24}-\frac{\left(2 \gamma ^2 (\epsilon -1)+1\right) \left(4 \gamma ^4 \left(\epsilon ^2-1\right)-2 \gamma ^2 \left(\epsilon  \left(6 \epsilon ^2+\epsilon -3\right)-2\right)-3 \epsilon  (2 \epsilon +1)+1\right)}{16 \left(\gamma ^2-1\right) (\epsilon -1)^2 (\epsilon +1)}\,,\\&  \boldsymbol{\tilde h}_{25}= \frac{1}{96 \left(\gamma ^2-1\right)^2 (\epsilon -1)^2 \epsilon  (3 \epsilon -1)}\Bigg[\gamma  \big(-2 \gamma ^2 \left(288 \epsilon ^5-552 \epsilon ^4+534 \epsilon ^3-515 \epsilon ^2+256 \epsilon -27\right) \epsilon\\& +\gamma ^4 \left(372 \epsilon ^4-658 \epsilon ^3+308 \epsilon ^2-28 \epsilon +2\right)-288 \epsilon ^5+288 \epsilon ^4-288 \epsilon ^3+208 \epsilon ^2-33 \epsilon -1\big)\Bigg]\,,\\&  \boldsymbol{\tilde h}_{26}=\frac{\gamma  \left(-4 \gamma ^2 \left(\epsilon  \left(4 \epsilon ^4-3 \epsilon ^3+\epsilon -4\right)+2\right)-2 \gamma ^4 (\epsilon -1)^2 (\epsilon  (6 \epsilon +5)-4)-\epsilon  \left(8 (\epsilon -1) \epsilon ^2+\epsilon +7\right)+4\right)}{96 \left(\gamma ^2-1\right) (\epsilon -1)^2 \epsilon  (3 \epsilon -1)}\,,\\& \boldsymbol{\tilde h}_{27}=\frac{\gamma}{24 \left(\gamma ^2-1\right)^2 (\epsilon -1)^2 \epsilon}\Bigg[-4 \gamma ^4 \left(24 \epsilon ^4-57 \epsilon ^3+41 \epsilon ^2-5 \epsilon -2\right)+\gamma ^2 \left(96 \epsilon ^5-88 \epsilon ^4-224 \epsilon ^3+290 \epsilon ^2-64 \epsilon -8\right)\\&+48 \epsilon ^4+12 \epsilon ^3-90 \epsilon ^2+31 \epsilon +1\Bigg]\,,\nonumber\end{split}\end{align}
        \begin{align}\begin{split}
    & \boldsymbol{\tilde h}_{28}= \frac{\gamma ^5 \left(4 \epsilon ^2-6 \epsilon +2\right)+\gamma ^3 \left(4 \epsilon  \left(4 \epsilon ^2-6 \epsilon +3\right)+1\right)+4 \gamma  \epsilon  (2 \epsilon -3)+\gamma }{48 \left(\gamma ^2-1\right) (\epsilon -1) (3 \epsilon -1)}\,,\\&
   \boldsymbol{\tilde h}_{29}=\frac{4 \epsilon  (2 \epsilon +1)}{\gamma ^2-1}\,, \boldsymbol{\tilde h}_{30}=-\frac{2 (6 \epsilon +5)}{3 \left(\gamma ^2-1\right)}\,, \boldsymbol{\tilde h}_{31}=\frac{2 \left(12 \epsilon ^2+8 \epsilon +1\right)}{3 \left(\gamma ^2-1\right)^2 \epsilon }\,, \boldsymbol{\tilde h}_{32}=-\frac{2 \gamma  (2 \epsilon +3)}{3 \left(\gamma ^2-1\right)}\,,\\& \boldsymbol{\tilde h}_{33}=\frac{18 \epsilon ^3+19 \epsilon ^2+3 \epsilon -1}{4 \left(\gamma ^2-1\right) (\epsilon +1)}\,, \boldsymbol{\tilde h}_{34}=-\frac{54 \epsilon ^3+75 \epsilon ^2+25 \epsilon +1}{24 \left(\gamma ^2-1\right) \epsilon  (\epsilon +1)}\,, \boldsymbol{\tilde h}_{35}=\frac{108 \epsilon ^4+132 \epsilon ^3+37 \epsilon ^2-3 \epsilon -1}{24 \left(\gamma ^2-1\right)^2 \epsilon ^2 (\epsilon +1)}\,,\\& \boldsymbol{\tilde h}_{36}=-\frac{\gamma  \left(18 \epsilon ^2+19 \epsilon +7\right)}{24 \left(\gamma ^2-1\right) (\epsilon +1)}\,,\boldsymbol{\tilde h}_{37}=\frac{\epsilon  (2 \epsilon +1)}{2 \epsilon -1},\,\,\boldsymbol{\tilde h}_{38}=\frac{ (2 \epsilon +1)}{2(1-2 \epsilon) },\,\,\boldsymbol{\tilde h}_{39}=\frac{ (2 \epsilon +1)}{2 \left(\gamma ^2-1\right) \epsilon }\,,\\&\boldsymbol{\tilde h}_{40}=\frac{2 \epsilon +1}{2 (1-3 \epsilon )}\,,\boldsymbol{\tilde h}_{41}=\frac{\epsilon  (2 \epsilon +1) \left(\gamma ^2 (\epsilon -1) (2 \epsilon -1)-\epsilon  (2 \epsilon +5)+1\right)}{\left(\gamma ^2-1\right) (\epsilon -1) (2 \epsilon -1) (3 \epsilon -1)}\,,\boldsymbol{\tilde h}_{42}=\frac{(2 \epsilon +1) (\epsilon  (4 (\epsilon -1) \epsilon +5)-1)}{2 \left(\gamma ^2-1\right) \epsilon  (\epsilon -1) (2 \epsilon -1) (3 \epsilon -1)}\,,\\&\boldsymbol{\tilde h}_{43}=\frac{\gamma ^4 \left(7 \epsilon ^3+4 \epsilon ^2-7 \epsilon -4\right)+\gamma ^2 \left(50 \epsilon ^4-3 \epsilon ^3-49 \epsilon ^2+2\right)+31 \epsilon ^3+31 \epsilon ^2+5 \epsilon +2}{\left(\gamma ^2-1\right) \left(\epsilon ^2-1\right)}\,,\\& \boldsymbol{\tilde h}_{44}=\frac{-21 \gamma ^4 \epsilon  \left(\epsilon ^2-1\right)+\gamma ^2 \left(-150 \epsilon ^4-41 \epsilon ^3+149 \epsilon ^2+44 \epsilon -2\right)-93 \epsilon ^3-135 \epsilon ^2-43 \epsilon +2}{6 \left(\gamma ^2-1\right) \epsilon  \left(\epsilon ^2-1\right)}\,,\\&\boldsymbol{\tilde h}_{45}=\frac{1}{6 \left(\gamma ^2-1\right)^2 \epsilon ^2 \left(\epsilon ^2-1\right)}\Bigg[3 \gamma ^4 \epsilon  \left(14 \epsilon ^3-3 \epsilon ^2-14 \epsilon +3\right)+\gamma ^2 \left(300 \epsilon ^5+32 \epsilon ^4-265 \epsilon ^3-45 \epsilon ^2-20 \epsilon -2\right)\\&+186 \epsilon ^4+225 \epsilon ^3+61 \epsilon ^2+9 \epsilon +2\Bigg]\,,\\& \boldsymbol{\tilde h}_{46}=-\frac{\gamma  \left(\gamma ^2 \left(50 \epsilon ^3+19 \epsilon ^2-44 \epsilon -25\right)+36 \epsilon ^2+43 \epsilon +13\right)}{6 \left(\gamma ^2-1\right) \left(\epsilon ^2-1\right)}\\&\boldsymbol{\tilde h}_{47}=\frac{1}{48 \gamma  \left(\gamma ^2-1\right)^2 (\epsilon -1) \epsilon ^2 (2 \epsilon -1) (3 \epsilon -1)}\Bigg[32 \gamma ^6 \epsilon ^2 \left(5 \epsilon ^2-7 \epsilon +2\right)+2 \gamma ^5 \left(60 \epsilon ^4-116 \epsilon ^3+41 \epsilon ^2+9 \epsilon -4\right)\\&+\gamma ^4 \left(2248 \epsilon ^5-4548 \epsilon ^4+3722 \epsilon ^3-2080 \epsilon ^2+700 \epsilon -96\right)+\gamma ^3 \left(60 \epsilon ^4+42 \epsilon ^3-40 \epsilon ^2-6 \epsilon +4\right)\\&+\gamma ^2 \left(-2248 \epsilon ^5+4388 \epsilon ^4-3434 \epsilon ^3+1884 \epsilon ^2-668 \epsilon +96\right)+\gamma  \left(-180 \epsilon ^4+150 \epsilon ^3+4 \epsilon ^2-31 \epsilon +7\right)\\&-16 \epsilon  \left(7 \epsilon ^2-9 \epsilon +2\right)\Bigg]\,,\\& \boldsymbol{\tilde h}_{48}=\frac{1}{48 (\gamma -1)^2 \gamma  (\gamma +1)^2 (\epsilon -1) \epsilon ^2 (2 \epsilon -1) (3 \epsilon -1)}\Bigg[24 \gamma ^6 \epsilon ^3 \left(30 \epsilon ^3-45 \epsilon ^2+16 \epsilon -1\right)\\&-2 \gamma ^5 \epsilon  \left(180 \epsilon ^4-504 \epsilon ^3+353 \epsilon ^2-81 \epsilon +4\right)+8 \gamma ^4 \epsilon  \left(120 \epsilon ^5-164 \epsilon ^4-9 \epsilon ^3+80 \epsilon ^2-30 \epsilon +3\right)\\&+\gamma ^3 \epsilon  \left(360 \epsilon ^4-648 \epsilon ^3+356 \epsilon ^2-63 \epsilon +1\right)-\gamma ^2 \left(1680 \epsilon ^6-2392 \epsilon ^5+440 \epsilon ^4+520 \epsilon ^3-252 \epsilon ^2+27 \epsilon +1\right)\\&+\gamma  \epsilon  \left(-360 \epsilon ^3+350 \epsilon ^2-99 \epsilon +7\right)+8 \epsilon ^2 \left(7 \epsilon ^2-9 \epsilon +2\right)\Bigg]\,,\\& \boldsymbol{\tilde h}_{49}=\frac{1}{48 (\gamma -1)^2 \gamma  (\gamma +1)^2 (\epsilon -1) \epsilon ^2 (2 \epsilon -1) (3 \epsilon -1)}\Bigg[-60 \gamma ^6 \epsilon ^2 \left(6 \epsilon ^3-11 \epsilon ^2+6 \epsilon -1\right)+30 \gamma ^5 \epsilon ^2 \left(6 \epsilon ^2-5 \epsilon +1\right)\\&+\gamma ^4 \epsilon  \left(240 \epsilon ^4+12 \epsilon ^3-306 \epsilon ^2+77 \epsilon -1\right)+\gamma ^3 \left(-720 \epsilon ^5+1056 \epsilon ^4-448 \epsilon ^3+94 \epsilon ^2-38 \epsilon +8\right)\\&+\gamma ^2 \left(-4800 \epsilon ^6+7864 \epsilon ^5-3904 \epsilon ^4+1058 \epsilon ^3-217 \epsilon ^2-47 \epsilon +24\right)+\gamma  \left(180 \epsilon ^4-82 \epsilon ^3+2 \epsilon ^2-17 \epsilon +7\right)\\&+8 \epsilon  \left(-14 \epsilon ^3+11 \epsilon ^2+5 \epsilon -2\right)\Bigg]\,,\nonumber
 \end{split}
\end{align}
\begin{align}
\begin{split}
&\boldsymbol{\tilde h}_{50}=\frac{-40 \gamma ^3 \left(3 \epsilon ^2-4 \epsilon +1\right)+\gamma ^2 \left(-400 \epsilon ^3+160 \epsilon ^2+372 \epsilon -125\right)+4 \gamma  (25 \epsilon -8)-8 (\epsilon -1) \epsilon }{48 (\gamma -1) (\gamma +1) (\epsilon -1) (3 \epsilon -1)}\,,\\&\boldsymbol{\tilde h}_{51}=\frac{1}{96 \gamma  \left(\gamma ^2-1\right)^2 (\epsilon -1)^2 \epsilon ^2 \left(6 \epsilon ^2-5 \epsilon +1\right)}\Bigg[-16 \gamma ^6 (\epsilon -1)^2 \epsilon ^2 \left(90 \epsilon ^3-75 \epsilon ^2-5 \epsilon +8\right)\\&+\gamma ^4 \left(2880 \epsilon ^7-4960 \epsilon ^6-1448 \epsilon ^5+8824 \epsilon ^4-9099 \epsilon ^3+5167 \epsilon ^2-1568 \epsilon +192\right)\\&-8 \gamma ^3 \epsilon  \left(12 \epsilon ^4-52 \epsilon ^3+67 \epsilon ^2-32 \epsilon +5\right)+\gamma ^2 \big(-1440 \epsilon ^7+880 \epsilon ^6+5080 \epsilon ^5-9562 \epsilon ^4+8769 \epsilon ^3-5245 \epsilon ^2\\&+1720 \epsilon -220\big)-8 \gamma  \left(6 \epsilon ^4-29 \epsilon ^3+39 \epsilon ^2-19 \epsilon +3\right)+2 \epsilon  \left(-103 \epsilon ^3+241 \epsilon ^2-173 \epsilon +32\right)\Bigg]\,,\\&\boldsymbol{\tilde h}_{52}=\frac{1}{96 (\gamma -1)^2 \gamma  (\gamma +1)^2 (\epsilon -1)^2 \epsilon ^2 \left(6 \epsilon ^2-5 \epsilon +1\right)}\Bigg[-4800 \gamma ^2 \left(\gamma ^2-1\right) \epsilon ^7-32 \gamma ^2 \big(15 \gamma ^4-9 \gamma ^3-364 \gamma ^2\\&+9 \gamma +349\big)\epsilon ^6+16 \left(72 \gamma ^6-78 \gamma ^5-384 \gamma ^4+63 \gamma ^3+319 \gamma ^2+15 \gamma -7\right) \epsilon ^5-8 \big(108 \gamma ^6-201 \gamma ^5+500 \gamma ^4\\&+128 \gamma ^3-567 \gamma ^2+73 \gamma -32\big) \epsilon ^4+2 \left(96 \gamma ^6-384 \gamma ^5+2172 \gamma ^4+132 \gamma ^3-2165 \gamma ^2+252 \gamma -91\right) \epsilon ^3\\&+\left(120 \gamma ^5-1120 \gamma ^4+64 \gamma ^3+1109 \gamma ^2-184 \gamma +39\right) \epsilon ^2+3 \left(24 \gamma ^4-8 \gamma ^3-26 \gamma ^2+8 \gamma +1\right) \epsilon -1\Bigg]\,,\\&\boldsymbol{\tilde h}_{53}=\frac{1}{96 (\gamma -1)^2 \gamma  (\gamma +1)^2 (\epsilon -1)^2 \epsilon ^2 \left(6 \epsilon ^2-5 \epsilon +1\right)}\Bigg[-60 \gamma ^8 \epsilon ^3 \left(6 \epsilon ^3-11 \epsilon ^2+6 \epsilon -1\right)\\&+60 \gamma ^6 \epsilon ^2 \left(30 \epsilon ^4-67 \epsilon ^3+52 \epsilon ^2-17 \epsilon +2\right)-3 \gamma ^4 \epsilon ^2 \left(520 \epsilon ^4-836 \epsilon ^3+172 \epsilon ^2+185 \epsilon -48\right)\\&-24 \gamma ^3 \epsilon  \left(24 \epsilon ^5-116 \epsilon ^4+162 \epsilon ^3-87 \epsilon ^2+18 \epsilon -1\right)+\gamma ^2 \big(9600 \epsilon ^7-26696 \epsilon ^6+26436 \epsilon ^5-12388 \epsilon ^4+3812 \epsilon ^3\\&-727 \epsilon ^2-117 \epsilon +60\big)+8 \gamma  \left(12 \epsilon ^5-16 \epsilon ^4+19 \epsilon ^3-29 \epsilon ^2+17 \epsilon -3\right)+\epsilon  \left(224 \epsilon ^4-400 \epsilon ^3+99 \epsilon ^2+107 \epsilon -31\right)\Bigg]\,,\\&\boldsymbol{\tilde h}_{54}=\frac{32 \gamma ^2 (\epsilon -1)^2 \left(25 \epsilon ^2+20 \epsilon -9\right)-40 \gamma  \left(6 \epsilon ^3-17 \epsilon ^2+14 \epsilon -3\right)+16 \epsilon ^3-32 \epsilon ^2+28 \epsilon -5}{96 (\gamma -1) (\gamma +1) (\epsilon -1)^2 (3 \epsilon -1)}\,,\\& \boldsymbol{\tilde h}_{55}=\frac{-2 \gamma ^4 \left(6 \epsilon ^4-7 \epsilon ^3-2 \epsilon ^2+7 \epsilon -4\right)+2 \gamma ^2 \epsilon  \left(104 \epsilon ^4-42 \epsilon ^3-99 \epsilon ^2+33 \epsilon +4\right)+104 \epsilon ^4+40 \epsilon ^3-36 \epsilon ^2+11 \epsilon -8}{2 \left(\gamma ^2-1\right) (\epsilon -1) (\epsilon +1) (2 \epsilon -1)}\,,\\&
       \boldsymbol{\tilde h}_{56}=\frac{1}{12 \left(\gamma ^2-1\right) (\epsilon -1) \epsilon  (\epsilon +1) (2 \epsilon -1)}\Bigg[6 \gamma ^4 \epsilon  \left(6 \epsilon ^3+\epsilon ^2-6 \epsilon -1\right)+\gamma ^2 \big(-624 \epsilon ^5+44 \epsilon ^4+738 \epsilon ^3-34 \epsilon ^2-132 \epsilon\\& +8\big)-312 \epsilon ^4-288 \epsilon ^3+100 \epsilon ^2+77 \epsilon -8\Bigg]\,,\\&\boldsymbol{\tilde h}_{57}=\frac{\gamma ^4 \left(-4 \epsilon ^3+2 \epsilon ^2+4 \epsilon -2\right)+\gamma ^2 \left(-192 \epsilon ^4-28 \epsilon ^3+178 \epsilon ^2+38 \epsilon +4\right)-96 \epsilon ^3-112 \epsilon ^2-22 \epsilon -1}{4 \left(\gamma ^2-1\right)^2 \epsilon  \left(\epsilon ^2-1\right)}\,,\\&\boldsymbol{\tilde h}_{58}=-\frac{\gamma  \left(2 \gamma ^2 \left(8 \epsilon ^3+5 \epsilon ^2-7 \epsilon -6\right)+4 \epsilon ^2+11 \epsilon +8\right)}{4 \left(\gamma ^2-1\right) (\epsilon -1) (\epsilon +1)}\,,\\& \boldsymbol{\tilde h}_{59}=\frac{1}{4 \gamma  \left(\gamma ^2-1\right) (\epsilon -1) (2 \epsilon -1) (3 \epsilon -1)}\Bigg[2 \gamma ^4 \epsilon  \left(4 \epsilon ^3-2 \epsilon ^2-3 \epsilon +1\right)-2 \gamma ^3 \epsilon  \left(4 \epsilon ^3-3 \epsilon ^2-2 \epsilon +1\right)\\&+\gamma ^2 \big(-16 \epsilon ^5+20 \epsilon ^4+48 \epsilon ^3-26 \epsilon ^2+7 \epsilon -1\big)-2 \gamma  \epsilon  \left(5 \epsilon ^2-6 \epsilon +1\right)+\epsilon  \left(-28 \epsilon ^2+31 \epsilon -7\right)\Bigg]\,,\nonumber
    \end{split}
\end{align}
\begin{align}
    \begin{split}
       & 
       \boldsymbol{\tilde h}_{60}=\frac{1}{24 \gamma  \left(\gamma ^2-1\right) (\epsilon -1) \epsilon  (2 \epsilon -1) (3 \epsilon -1)}\Bigg[-6 \gamma ^4 \epsilon  \left(4 \epsilon ^3-2 \epsilon ^2-3 \epsilon +1\right)+6 \gamma ^3 \epsilon  \left(4 \epsilon ^3-3 \epsilon ^2-2 \epsilon +1\right)\\&+\gamma ^2 \left(48 \epsilon ^5-44 \epsilon ^4+48 \epsilon ^3-62 \epsilon ^2+27 \epsilon -5\right)-2 \gamma  \epsilon  \left(5 \epsilon ^2-6 \epsilon +1\right)+\epsilon  \left(-28 \epsilon ^2+31 \epsilon -7\right)\Bigg]\,,
       \\&\boldsymbol{\tilde h}_{61}=-\frac{1}{\left.24 \gamma  \left(\gamma ^2-1\right)^2 (\epsilon -1) \epsilon ^2 (2 \epsilon -1) (3 \epsilon -1)\right)}\Bigg[(2 \epsilon +1) \big(-6 \gamma ^4 \epsilon  \left(4 \epsilon ^3-14 \epsilon ^2+7 \epsilon -1\right)\\&+6 \gamma ^3 (\epsilon -1)^2 \epsilon  (4 \epsilon -1)+\gamma ^2 \left(48 \epsilon ^5-76 \epsilon ^4-60 \epsilon ^3+98 \epsilon ^2-39 \epsilon +5\right)+2 \gamma  \epsilon  \left(5 \epsilon ^2-6 \epsilon +1\right)\\&+\epsilon  \left(28 \epsilon ^2-31 \epsilon +7\right)\big)\Bigg]\,,\\&\boldsymbol{\tilde h}_{62}=\frac{\gamma ^2 \left(4 \epsilon ^3-6 \epsilon +2\right)+2 \gamma  (\epsilon -1) \epsilon +\epsilon  (2 \epsilon -1)}{12 \left(\gamma ^2-1\right) (\epsilon -1) (3 \epsilon -1)}\,,\boldsymbol{\tilde h}_{63}=\frac{\gamma ^2 (4 \epsilon -1)}{3 \left(\gamma ^2-1\right) (3 \epsilon -1)}\,,\\&\boldsymbol{\tilde h}_{64}=\frac{1}{8 \gamma  \left(\gamma ^2-1\right)^2 (\epsilon -1) (\epsilon +1) (2 \epsilon -1) (3 \epsilon -1)}\Bigg[-4 \gamma ^6 \epsilon  \left(4 \epsilon ^4+2 \epsilon ^3-5 \epsilon ^2-2 \epsilon +1\right)\\&+\gamma ^4 \left(256 \epsilon ^6-184 \epsilon ^5-260 \epsilon ^4+27 \epsilon ^3+83 \epsilon ^2-45 \epsilon +7\right)-4 \gamma ^3 \epsilon ^2 \left(2 \epsilon ^2+\epsilon -1\right)\\&-2 \gamma ^2 \left(128 \epsilon ^6-100 \epsilon ^5-138 \epsilon ^4+27 \epsilon ^3+47 \epsilon ^2-29 \epsilon +5\right)+4 \gamma  \epsilon ^2 \left(2 \epsilon ^2+\epsilon -1\right)\\&-8 \epsilon ^4+7 \epsilon ^3+19 \epsilon ^2+3 \epsilon -1\Bigg]\,,\\&\boldsymbol{\tilde h}_{65}=\frac{1}{48 (\gamma -1) \gamma  (\gamma +1) (\epsilon -1) \epsilon  (\epsilon +1) (2 \epsilon -1) (3 \epsilon -1)}\Bigg[12 \gamma ^4 \epsilon  \left(4 \epsilon ^4+2 \epsilon ^3-5 \epsilon ^2-2 \epsilon +1\right)\\&+\gamma ^2 \left(-768 \epsilon ^6+344 \epsilon ^5+676 \epsilon ^4-281 \epsilon ^3-99 \epsilon ^2+39 \epsilon +5\right)+4 \gamma  \epsilon  \left(6 \epsilon ^3+5 \epsilon ^2-2 \epsilon -1\right)\\&+8 \epsilon ^4-7 \epsilon ^3-19 \epsilon ^2-3 \epsilon +1\Bigg]\,,\\& \boldsymbol{\tilde h}_{66}=\frac{1}{48 (\gamma -1)^2 \gamma  (\gamma +1)^2 (\epsilon -1) \epsilon ^2 (\epsilon +1) (2 \epsilon -1) (3 \epsilon -1)}\Bigg[2 \gamma ^4 \big(2400 \epsilon ^6-1180 \epsilon ^5-2616 \epsilon ^4\\&+1385 \epsilon ^3+345 \epsilon ^2-205 \epsilon +15\big) \epsilon +\gamma ^2 \big(6336 \epsilon ^7-4808 \epsilon ^6-4708 \epsilon ^5+3302 \epsilon ^4-401 \epsilon ^3+133 \epsilon ^2-21 \epsilon \\&-13\big)-4 \gamma  \left(12 \epsilon ^4+8 \epsilon ^3-7 \epsilon ^2-2 \epsilon +1\right) \epsilon -44 \epsilon ^5+40 \epsilon ^4+19 \epsilon ^3-77 \epsilon ^2-5 \epsilon +7\Bigg]\,,\\
    &\boldsymbol{\tilde h}_{67}=\frac{\gamma ^2 \left(-8 \epsilon ^3+13 \epsilon -5\right)+2 \gamma  \epsilon -\epsilon +1}{24 (\gamma -1) (\gamma +1) (\epsilon -1) (3 \epsilon -1)}\,,\\&\boldsymbol{\tilde h}_{68}=\frac{-8 \gamma ^2 \left(107 \epsilon ^4+27 \epsilon ^3-114 \epsilon ^2-45 \epsilon +25\right)+100 \epsilon ^3-74 \epsilon ^2-127 \epsilon +47}{24 (\gamma -1) (\gamma +1) (\epsilon -1) (\epsilon +1) (3 \epsilon -1)}\,,\\& \boldsymbol{\tilde h}_{69}=-\frac{\epsilon  \left(\gamma ^2 \left(2 \epsilon ^2-3 \epsilon +1\right)-10 \epsilon ^2-9 \epsilon -5\right)}{2 (\epsilon -1) (2 \epsilon -3) (2 \epsilon -1)}\,,\boldsymbol{\tilde h}_{70}=\frac{ \left(\gamma ^2 (\epsilon -1)-\epsilon -1\right)}{4 (\epsilon -1) (2 \epsilon -3)}\,,\boldsymbol{\tilde h}_{71}=-\frac{ \left(4 \epsilon ^3+5 \epsilon +3\right)}{4 (\epsilon -1) \epsilon  (2 \epsilon -3) (2 \epsilon -1)}\,,\\& \boldsymbol{\tilde h}_{72}=\frac{1}{256 \left(\gamma ^2-1\right)^2 (\epsilon -1)^2 (\gamma  \epsilon  (6 \epsilon -5)+\gamma )}\Bigg[-64 \gamma ^8 (\epsilon -1)^2 \epsilon ^2+\gamma ^6 \big(160 \epsilon ^5+1872 \epsilon ^4-5282 \epsilon ^3+4581 \epsilon ^2\\&-1490 \epsilon +159\big)-2 \gamma ^4 \left(80 \epsilon ^5-3816 \epsilon ^4+8202 \epsilon ^3-5819 \epsilon ^2+1453 \epsilon -120\right)+\gamma ^2 \big(800 \epsilon ^4-2506 \epsilon ^3+1597 \epsilon ^2\\&-184 \epsilon -27\big)-4 \left(22 \epsilon ^2+7 \epsilon -3\right)\Bigg]\,,\nonumber 
 \end{split}
\end{align}
\begin{align}
    \begin{split}
  &\boldsymbol{\tilde h}_{73}=\frac{1}{1536 \gamma  \left(\gamma ^2-1\right) (\epsilon -1)^2 \epsilon  (\epsilon  (6 \epsilon -5)+1)}\Bigg[192 \gamma ^6 (\epsilon -1)^2 \epsilon ^2+\gamma ^4 \big(-480 \epsilon ^5-944 \epsilon ^4+4078 \epsilon ^3-3551 \epsilon ^2\\&+896 \epsilon +1\big)+\gamma ^2 \left(-1440 \epsilon ^4+2762 \epsilon ^3-1733 \epsilon ^2+292 \epsilon -1\right)+4 \left(22 \epsilon ^2+7 \epsilon -3\right)\Bigg]\,,\\&\boldsymbol{\tilde h}_{74}=\frac{1}{1536 \left(\gamma ^2-1\right)^2 (\epsilon -1)^2 \epsilon ^2 (\gamma  \epsilon  (6 \epsilon -5)+\gamma )}\Bigg[(2 \epsilon +1) \big(-192 \gamma ^6 (\epsilon -1)^2 \epsilon ^2+\gamma ^4 \big(480 \epsilon ^5-5360 \epsilon ^4\\&+9646 \epsilon ^3-5663 \epsilon ^2+896 \epsilon +1\big)+\gamma ^2 \left(-480 \epsilon ^4+1850 \epsilon ^3-1493 \epsilon ^2+292 \epsilon -1\right)+4 \left(22 \epsilon ^2+7 \epsilon -3\right)\big)\Bigg]\,,\\&\boldsymbol{\tilde h}_{75}=\frac{4 (\epsilon -3)-\gamma ^2 (\epsilon -1) \left(\gamma ^2 \left(320 \epsilon ^2-418 \epsilon +97\right)-336 \epsilon ^2+466 \epsilon -97\right)}{768 \left(\gamma ^2-1\right) (\epsilon -1)^2 (3 \epsilon -1)}\,,\\&\boldsymbol{\tilde h}_{76}=\frac{\gamma ^2 \left(-200 \epsilon ^6+264 \epsilon ^5+75 \epsilon ^4-184 \epsilon ^3+34 \epsilon ^2+16 \epsilon -5\right)-72 \epsilon ^6-328 \epsilon ^5+117 \epsilon ^4+124 \epsilon ^3-57 \epsilon ^2+3 \epsilon +3}{8 (1-2 \epsilon )^2 (\epsilon -1)^2 (\epsilon +1) (2 \epsilon +1)}\,,\\&\boldsymbol{\tilde h}_{77}=\frac{\gamma ^2 (\epsilon -1)^2 \left(232 \epsilon ^4+90 \epsilon ^3-123 \epsilon ^2-10 \epsilon +19\right)-8 \epsilon ^6+110 \epsilon ^5-49 \epsilon ^4-132 \epsilon ^3+60 \epsilon ^2+16 \epsilon -9}{16 (1-2 \epsilon )^2 (\epsilon -1)^2 \epsilon  (\epsilon +1) (2 \epsilon +1)}\,,\\&\boldsymbol{\tilde h}_{78}=-\frac{64 \epsilon ^7-280 \epsilon ^6+22 \epsilon ^5+115 \epsilon ^4-60 \epsilon ^3+5 \epsilon ^2+7 \epsilon +1}{16 (1-2 \epsilon )^2 (\epsilon -1)^2 \epsilon ^2 (\epsilon +1) (2 \epsilon +1)}\,,\\&\boldsymbol{\tilde h}_{79}=\frac{-2 \gamma ^2 (1-2 \epsilon )^2 \left(10 \epsilon ^3-19 \epsilon ^2+4 \epsilon +8\right)+32 \epsilon ^5-168 \epsilon ^4+228 \epsilon ^3-106 \epsilon ^2+7 \epsilon +7}{32 (1-2 \epsilon )^2 (\epsilon -1)^2 (2 \epsilon +1)}\,,\\&\boldsymbol{\tilde h}_{80}=\frac{\left(4-36 \gamma ^2\right) \epsilon ^4+\left(74 \gamma ^2-26\right) \epsilon ^3+\left(21-48 \gamma ^2\right) \epsilon ^2+2 \left(5 \gamma ^2-2\right) \epsilon -1}{16 (1-2 \epsilon )^2 (\epsilon -1) \epsilon }\,,\boldsymbol{\tilde h}_{81}=-\frac{\left(\gamma ^2-1\right) \left(4 \epsilon ^3+\epsilon ^2-2 \epsilon +1\right)}{16 (1-2 \epsilon )^2 \epsilon  (2 \epsilon +1)}\,,\\&\boldsymbol{\tilde h}_{82}=\frac{1}{128 (\gamma^2 -1) \gamma \left(2 \epsilon ^2-3 \epsilon +1\right)^2 (\epsilon ^2 (6 \epsilon +7)-1)}\Bigg[2 \gamma ^4 \big(3624 \epsilon ^7-5990 \epsilon ^6+1394 \epsilon ^5\\&+2233 \epsilon ^4-1628 \epsilon ^3+362 \epsilon ^2+18 \epsilon -13\big)+16 \gamma ^3 (\epsilon -1)^2 \left(72 \epsilon ^5+36 \epsilon ^4-50 \epsilon ^3-5 \epsilon ^2+8 \epsilon -1\right)\\&+\gamma ^2 \left(-21216 \epsilon ^7+22288 \epsilon ^6+17516 \epsilon ^5-23897 \epsilon ^4+3172 \epsilon ^3+3674 \epsilon ^2-1392 \epsilon +143\right)\\&+3 \left(-3408 \epsilon ^7+4148 \epsilon ^6+1264 \epsilon ^5-2739 \epsilon ^4+652 \epsilon ^3+180 \epsilon ^2-80 \epsilon +7\right)\Bigg]\,,\\&\boldsymbol{\tilde h}_{83}=-\frac{1}{256 (\gamma^2 -1) \gamma \left(2 \epsilon ^2-3 \epsilon +1\right)^2 \epsilon  (\epsilon ^2 (6 \epsilon +7)-1)}\Bigg[2 \gamma ^4 \big(4380 \epsilon ^7-7886 \epsilon ^6+1546 \epsilon ^5\\&+4549 \epsilon ^4-2920 \epsilon ^3+10 \epsilon ^2+402 \epsilon -81\big)+\gamma ^2 \big(-6192 \epsilon ^7+5784 \epsilon ^6+8610 \epsilon ^5-12461 \epsilon ^4+3572 \epsilon ^3\\&+1740 \epsilon ^2-1226 \epsilon +197\big)+888 \epsilon ^7+388 \epsilon ^6-2142 \epsilon ^5+107 \epsilon ^4+1206 \epsilon ^3-514 \epsilon ^2+36 \epsilon +7\Bigg]\,,\\&\boldsymbol{\tilde h}_{84}=-\frac{1}{256 (\gamma^2 -1) \gamma \left(2 \epsilon ^2-3 \epsilon +1\right)^2 \epsilon ^2 (\epsilon ^2 (6 \epsilon +7)-1)}\Bigg[16 \gamma ^3 \left(-2 \epsilon ^2+\epsilon +1\right)^2 \left(6 \epsilon ^3+\epsilon ^2-4 \epsilon +1\right) \epsilon\\& +2 \gamma ^2 \left(-6096 \epsilon ^8+7564 \epsilon ^7+1670 \epsilon ^6-4134 \epsilon ^5+1011 \epsilon ^4+52 \epsilon ^3+36 \epsilon ^2-38 \epsilon +7\right)-9696 \epsilon ^8+10792 \epsilon ^7\\&+4068 \epsilon ^6-6210 \epsilon ^5+801 \epsilon ^4+272 \epsilon ^3+104 \epsilon ^2-66 \epsilon +7\Bigg]\,,\nonumber
    \end{split}
\end{align}
\begin{align}
    \begin{split}
    &\boldsymbol{\tilde h}_{85}=\frac{1}{256 \gamma  \left(\gamma ^2-1\right) (2 \epsilon +1) \left(2 \epsilon ^2-3 \epsilon +1\right)^2}\Bigg[\gamma ^4 (1-2 \epsilon )^2 \left(56 \epsilon ^3-144 \epsilon ^2+32 \epsilon +59\right)\\&+\gamma ^2 \left(-2048 \epsilon ^5+4672 \epsilon ^4-3376 \epsilon ^3+448 \epsilon ^2+444 \epsilon -143\right)-8 \epsilon  (1-2 \epsilon )^2 \left(15 \epsilon ^2-22 \epsilon +7\right)\Bigg]\,,\\&\boldsymbol{\tilde h}_{86}=\frac{1}{256 \gamma  \left(\gamma ^2-1\right) (1-2 \epsilon )^2 (\epsilon -1) \epsilon}\Bigg[4 \gamma ^4 \epsilon  \left(66 \epsilon ^3-133 \epsilon ^2+84 \epsilon -17\right)+\gamma ^2 \big(-720 \epsilon ^4+1104 \epsilon ^3-566 \epsilon ^2\\&+72 \epsilon +14\big)-(1-2 \epsilon )^2 \left(30 \epsilon ^2+\epsilon -7\right)\Bigg]\,,\\&\boldsymbol{\tilde h}_{87}=\frac{(\epsilon +1) \left(2 \gamma ^2 \left(42 \epsilon ^3-60 \epsilon ^2+33 \epsilon -7\right)+(15 \epsilon -7) (1-2 \epsilon )^2\right)}{256 \gamma  (1-2 \epsilon )^2 (\epsilon -1) \epsilon  (2 \epsilon +1)}\,,\\&\boldsymbol{\tilde h}_{88}=-\frac{1}{64 \gamma ^2 \left(\gamma ^2-1\right)^2 (1-2 \epsilon )^4 (\epsilon -1)^3 \left(\epsilon ^2 (6 \epsilon +7)-1\right)^2}\Bigg[\gamma ^8 \big(-156 \epsilon ^6+404 \epsilon ^5-332 \epsilon ^4+43 \epsilon ^3+92 \epsilon ^2\\&-63 \epsilon +12\big)^2+4 \gamma ^7 (\epsilon -1)^2 (2 \epsilon -1)^3 \big(1800 \epsilon ^8+396 \epsilon ^7-3464 \epsilon ^6-1260 \epsilon ^5+1299 \epsilon ^4+345 \epsilon ^3\\&-183 \epsilon ^2-21 \epsilon +8\big)+2 \gamma ^6 \big(112320 \epsilon ^{13}-537984 \epsilon ^{12}+889584 \epsilon ^{11}-486508 \epsilon ^{10}-237128 \epsilon ^9\\&+414152 \epsilon ^8-159781 \epsilon ^7-11177 \epsilon ^6+15554 \epsilon ^5+602 \epsilon ^4+4123 \epsilon ^3-5533 \epsilon ^2+2016 \epsilon -240\big)\\&-4 \gamma ^5 (\epsilon -1)^2 (2 \epsilon -1)^3 \left(7704 \epsilon ^8+8172 \epsilon ^7-5770 \epsilon ^6-5225 \epsilon ^5+2756 \epsilon ^4+1375 \epsilon ^3-511 \epsilon ^2-110 \epsilon +33\right)\\&+\gamma ^4 \big(518400 \epsilon ^{14}-2505600 \epsilon ^{13}+3769056 \epsilon ^{12}-1383360 \epsilon ^{11}-1057208 \epsilon ^{10}+1569336 \epsilon ^9-1608359 \epsilon ^8\\&+817324 \epsilon ^7+308191 \epsilon ^6-483302 \epsilon ^5+121399 \epsilon ^4+31648 \epsilon ^3-15071 \epsilon ^2+162 \epsilon +328\big)\\&+4 \gamma ^3 (\epsilon -1)^2 (2 \epsilon -1)^3 \left(4824 \epsilon ^8+7092 \epsilon ^7+508 \epsilon ^6-1106 \epsilon ^5+1265 \epsilon ^4+316 \epsilon ^3-358 \epsilon ^2-38 \epsilon +25\right)\\&-6 \gamma ^2 \big(172800 \epsilon ^{14}-826560 \epsilon ^{13}+1067520 \epsilon ^{12}+363056 \epsilon ^{11}-1442532 \epsilon ^{10}+407040 \epsilon ^9+477315 \epsilon ^8\\&-177793 \epsilon ^7-64099 \epsilon ^6-4775 \epsilon ^5+24081 \epsilon ^4-135 \epsilon ^3-3513 \epsilon ^2+703 \epsilon -20\big)+12 \gamma  \epsilon  (2 \epsilon -1)^3 (10 \epsilon -17)\\& \left(6 \epsilon ^4+\epsilon ^3-7 \epsilon ^2-\epsilon +1\right)^2+9 \left(-240 \epsilon ^7+620 \epsilon ^6-96 \epsilon ^5-565 \epsilon ^4+212 \epsilon ^3+76 \epsilon ^2-32 \epsilon +1\right)^2\Bigg]\,,\\& \boldsymbol{\tilde h}_{89}=\frac{1}{64 \gamma  \left(\gamma ^2-1\right) (1-2 \epsilon )^2 (\epsilon -1)^2 \epsilon  \left(\epsilon ^2 (6 \epsilon +7)-1\right)}\Bigg[\gamma ^4 \big(1416 \epsilon ^7-2900 \epsilon ^6+596 \epsilon ^5+1956 \epsilon ^4\\&-1131 \epsilon ^3-82 \epsilon ^2+175 \epsilon -30\big)+\gamma ^2 \left(-1344 \epsilon ^7+2432 \epsilon ^6+938 \epsilon ^5-3411 \epsilon ^4+1117 \epsilon ^3+500 \epsilon ^2-291 \epsilon +35\right)\\&+312 \epsilon ^7-716 \epsilon ^6-246 \epsilon ^5+975 \epsilon ^4-116 \epsilon ^3-248 \epsilon ^2+62 \epsilon +1\Bigg]\,,\\&\boldsymbol{\tilde h}_{90}=\frac{1}{64 \gamma  \left(\gamma ^2-1\right) (1-2 \epsilon )^2 (\epsilon -1)^2 \epsilon ^2 \left(\epsilon ^2 (6 \epsilon +7)-1\right)}\Bigg[\gamma ^2 \big(-1536 \epsilon ^8+1528 \epsilon ^7-244 \epsilon ^6-240 \epsilon ^5\\&+696 \epsilon ^4-323 \epsilon ^3-46 \epsilon ^2+19 \epsilon +2\big)-1152 \epsilon ^8+2232 \epsilon ^7+908 \epsilon ^6-2398 \epsilon ^5-303 \epsilon ^4+626 \epsilon ^3+78 \epsilon ^2-64 \epsilon +1\Bigg]\,,\\&\boldsymbol{\tilde h}_{91}=\frac{1}{64 \gamma  \left(\gamma ^2-1\right) (2 \epsilon +1) \left(2 \epsilon ^2-3 \epsilon +1\right)^2}\Bigg[\gamma ^4 \left(-(1-2 \epsilon )^2\right) \left(20 \epsilon ^3-48 \epsilon ^2+7 \epsilon +18\right)+\gamma ^2 \big(304 \epsilon ^5-816 \epsilon ^4\\&+648 \epsilon ^3-72 \epsilon ^2-97 \epsilon +30\big)+8 \epsilon  \left(2 \epsilon ^2-3 \epsilon +1\right)^2\Bigg]\,,\nonumber
    \end{split}
\end{align}
\begin{align}
    \begin{split}
  &\boldsymbol{\tilde h}_{92}=\frac{-4 \gamma ^4 \epsilon  \left(18 \epsilon ^2-19 \epsilon +5\right)+2 \gamma ^2 \left(64 \epsilon ^3-52 \epsilon ^2+11 \epsilon +1\right)+(2 \epsilon +1) (1-2 \epsilon )^2}{64 \gamma  \left(\gamma ^2-1\right) (1-2 \epsilon )^2 \epsilon }\,,\\
   & \boldsymbol{\tilde h}_{93}=-\frac{(\epsilon +1) \left(2 \gamma ^2 \left(6 \epsilon ^2-4 \epsilon +1\right)+(1-2 \epsilon )^2\right)}{64 \gamma  (1-2 \epsilon )^2 \epsilon  (2 \epsilon +1)}\,, \boldsymbol{\tilde h}_{94}=\gamma ^2+\frac{1}{2 (\epsilon -1)}\,,\\&\boldsymbol{\tilde h}_{95}=-\frac{\gamma  (2 \epsilon +1) \left(\gamma ^2 (2 \epsilon -1) (7 \epsilon -2)+10 \epsilon ^2+\epsilon -1\right)}{8 \left(\gamma ^2-1\right) (\epsilon -1) \epsilon  (2 \epsilon -1) (3 \epsilon -1)}\,,\\&\boldsymbol{\tilde h}_{96}=\frac{\gamma  (2 \epsilon +1)^2 (4 \epsilon -1)}{8 \left(\gamma ^2-1\right) (\epsilon -1) \epsilon  (2 \epsilon -1) (3 \epsilon -1)}\,,\boldsymbol{\tilde h}_{97}=\frac{2 \gamma  \epsilon +\gamma }{24 \epsilon ^2-32 \epsilon +8}\,,\\&\boldsymbol{\tilde h}_{98}=-\frac{\gamma^2  (2 \epsilon +1)^2 \left(\gamma ^2 (2 \epsilon -1) (7 \epsilon -2)+10 \epsilon ^2+\epsilon -1\right)}{8 \left(\gamma ^2-1\right) (\epsilon -1)^2 \epsilon  (\epsilon  (6 \epsilon -5)+1)}\,,\boldsymbol{\tilde h}_{99}=\frac{\gamma^2  (2 \epsilon +1)^2}{8 (\epsilon -1)^2 (3 \epsilon -1)}\,,\\&\boldsymbol{\tilde h}_{100}=\frac{\gamma^2  (2 \epsilon +1)^3 (4 \epsilon -1)}{8 \left(\gamma ^2-1\right) (\epsilon -1)^2 \epsilon  (\epsilon  (6 \epsilon -5)+1)}\,,\\&
    \boldsymbol{\tilde h}_{101}=\frac{8 \gamma ^2 (\epsilon -1)^2 \left(6 \epsilon ^2-6 \epsilon +1\right)+36 \epsilon ^4-60 \epsilon ^3+23 \epsilon ^2-\epsilon +1}{4 (\epsilon -1)^3 (2 \epsilon -1)}\,,\boldsymbol{\tilde h}_{102}=\frac{-4 \gamma ^2 (\epsilon -1)^2 (2 \epsilon +1)-12 \epsilon ^3+5 \epsilon +1}{8 (\epsilon -1)^2 \epsilon  (2 \epsilon -1)}\,,\\&\boldsymbol{\tilde h}_{103}=\frac{1}{8 \gamma  \left(\gamma ^2-1\right) (\epsilon -1)^2 (\epsilon  (6 \epsilon -5)+1)}\Bigg[-2 \gamma ^4 \left(228 \epsilon ^4-464 \epsilon ^3+313 \epsilon ^2-85 \epsilon +8\right)\\&+\gamma ^2 \epsilon  \left(-930 \epsilon ^3+1265 \epsilon ^2-454 \epsilon +47\right)-9 \epsilon  \left(18 \epsilon ^3-27 \epsilon ^2+13 \epsilon -2\right)\Bigg]\,,\\&\boldsymbol{\tilde h}_{104}=\frac{2 \gamma ^2 \left(18 \epsilon ^2-23 \epsilon +5\right)+3 \left(9 \epsilon ^2-9 \epsilon +2\right)}{16 \gamma  (\epsilon -1)^2 (3 \epsilon -1)}\,,\\&\boldsymbol{\tilde h}_{105}=\frac{(2 \epsilon +1) \left(4 \gamma ^4 \left(6 \epsilon ^3-11 \epsilon ^2+6 \epsilon -1\right)+2 \gamma ^2 \left(90 \epsilon ^3-119 \epsilon ^2+38 \epsilon -3\right)+54 \epsilon ^3-81 \epsilon ^2+39 \epsilon -6\right)}{16 (\gamma -1) \gamma  (\gamma +1) (\epsilon -1)^2 \epsilon  (2 \epsilon -1) (3 \epsilon -1)}\,,\\&\boldsymbol{\tilde h}_{106}=\frac{1}{32 \left(\gamma ^2-1\right) (\epsilon -1)^3 (\epsilon  (6 \epsilon -5)+1)}\Bigg[2 \gamma ^4 (\epsilon -1)^2 \left(228 \epsilon ^3-314 \epsilon ^2+131 \epsilon -16\right)\\&-4 \gamma ^3 (\epsilon -1)^2 \left(12 \epsilon ^2-8 \epsilon +1\right) \epsilon +\gamma ^2 \left(2640 \epsilon ^5-7108 \epsilon ^4+7052 \epsilon ^3-3179 \epsilon ^2+621 \epsilon -42\right)\\&-4 \gamma  (\epsilon -1)^2 \left(132 \epsilon ^3-76 \epsilon ^2+15 \epsilon -1\right)+2088 \epsilon ^5-4072 \epsilon ^4+2922 \epsilon ^3-1037 \epsilon ^2+162 \epsilon -7\Bigg]\,,\\&\boldsymbol{\tilde h}_{107}=\frac{-2 \gamma ^2 (\epsilon -1)^2 \left(84 \epsilon ^2-98 \epsilon +27\right)+4 \gamma  (\epsilon -1)^2 \left(12 \epsilon ^2-8 \epsilon +1\right)-264 \epsilon ^4+600 \epsilon ^3-482 \epsilon ^2+181 \epsilon -27}{64 (\epsilon -1)^3 \left(6 \epsilon ^2-5 \epsilon +1\right)}\,,\\&\boldsymbol{\tilde h}_{108}=-\frac{1}{\left.64 \left(\gamma ^2-1\right) (\epsilon -1)^3 \epsilon  (\epsilon  (6 \epsilon -5)+1)\right)}\Bigg[(2 \epsilon +1) \big(6 \gamma ^2 \left(64 \epsilon ^4-174 \epsilon ^3+169 \epsilon ^2-68 \epsilon +9\right)\\&-4 \gamma  (\epsilon -1)^2 \left(24 \epsilon ^2-10 \epsilon +1\right)+480 \epsilon ^4-932 \epsilon ^3+648 \epsilon ^2-215 \epsilon +27\big)\Bigg]\,,\\&\boldsymbol{\tilde h}_{109}=\frac{1}{64 \gamma  \left(\gamma ^2-1\right) (\epsilon -1) (2 \epsilon -1) (3 \epsilon -1)}\Bigg[\gamma ^5 \left(36 \epsilon ^3-62 \epsilon ^2+29 \epsilon -4\right)+8 \gamma ^4 \left(14 \epsilon ^2-11 \epsilon +2\right)\\&+\gamma ^3 \left(108 \epsilon ^3-106 \epsilon ^2+25 \epsilon -1\right)-8 \gamma ^2 \left(36 \epsilon ^3-19 \epsilon ^2-5 \epsilon +2\right)+24 \epsilon  (1-3 \epsilon )\Bigg]\,,\\&\boldsymbol{\tilde h}_{110}=\frac{\gamma ^3 \left(-12 \epsilon ^2+18 \epsilon -5\right)+\gamma ^2 (8-16 \epsilon )+24 \epsilon -8}{128 \gamma  (\epsilon -1) (2 \epsilon -1) (3 \epsilon -1)}\,,\\
&\boldsymbol{\tilde h}_{111}=-\frac{(2 \epsilon +1) \left(\gamma ^3 \left(24 \epsilon ^2-24 \epsilon +5\right)-8 \gamma ^2 \left(6 \epsilon ^2-6 \epsilon +1\right)-24 \epsilon +8\right)}{128 \gamma  \left(\gamma ^2-1\right) (\epsilon -1) \epsilon  (2 \epsilon -1) (3 \epsilon -1)}\,,\\
\nonumber
     \end{split}
\end{align}
\begin{align}
\begin{split}
&\boldsymbol{\tilde h}_{112}=\frac{\gamma  \left(-\left(\gamma ^2 (\epsilon -1) (12 \epsilon -5)\right)-5 \epsilon +2\right)}{12 \left(\gamma ^2-1\right) (\epsilon -1)^2}\,,\boldsymbol{\tilde h}_{113}=\frac{\gamma  \left(\gamma ^2 (6 (3-2 \epsilon ) \epsilon -5)-6 \epsilon +2\right)}{96 \left(\gamma ^2-1\right) (\epsilon -1)^2}\,,\\& \boldsymbol{\tilde h}_{114}=\frac{1}{256(\gamma -1)^2\gamma (\gamma +1)^2(\epsilon -1)^2\left(6 \epsilon ^2-5 \epsilon +1\right)}\Bigg[-64 \gamma ^8 (\epsilon -1)^2 \epsilon ^2-32 \gamma ^7 (\epsilon -1)^2 \epsilon  \left(4 \epsilon ^2+4 \epsilon -3\right)\\&+\gamma ^6 \left(160 \epsilon ^5+1872 \epsilon ^4-5282 \epsilon ^3+4581 \epsilon ^2-1490 \epsilon +159\right)+16 \gamma ^5 \left(16 \epsilon ^5-300 \epsilon ^4+556 \epsilon ^3-347 \epsilon ^2+83 \epsilon -8\right)\\&-2 \gamma ^4 \left(80 \epsilon ^5-3816 \epsilon ^4+8202 \epsilon ^3-5819 \epsilon ^2+1453 \epsilon -120\right)-32 \gamma ^3 \left(4 \epsilon ^5-112 \epsilon ^4+213 \epsilon ^3-149 \epsilon ^2+39 \epsilon -4\right)\\&+\gamma ^2 \left(800 \epsilon ^4-2506 \epsilon ^3+1597 \epsilon ^2-184 \epsilon -27\right)+16 \gamma  \epsilon  \left(68 \epsilon ^3-144 \epsilon ^2+69 \epsilon -11\right)-4 \left(22 \epsilon ^2+7 \epsilon -3\right)\Bigg]\,,\\& \boldsymbol{\tilde h}_{115}=\frac{1}{1536 (\gamma -1) \gamma  (\gamma +1) (\epsilon -1)^2 \epsilon  (2 \epsilon -1) (3 \epsilon -1)}\Bigg[192 \gamma ^6 (\epsilon -1)^2 \epsilon ^2+96 \gamma ^5 (\epsilon -1)^2 \epsilon  \left(4 \epsilon ^2+4 \epsilon -3\right)\\&+4 \left(22 \epsilon ^2+7 \epsilon -3\right)-48 \gamma  \epsilon  \left(28 \epsilon ^3-48 \epsilon ^2+23 \epsilon -3\right)+\gamma ^2 \left(-1440 \epsilon ^4+2762 \epsilon ^3-1733 \epsilon ^2+292 \epsilon -1\right)\\&-48 \gamma ^3 \epsilon  \left(8 \epsilon ^4-36 \epsilon ^3+34 \epsilon ^2-3 \epsilon -3\right)+\gamma ^4\left(-944 \epsilon ^4+4078 \epsilon ^3-3551 \epsilon ^2+896 \epsilon -480 \epsilon +1\right)\Bigg]\,,\\&\boldsymbol{\tilde h}_{116}=-\frac{1}{1536 (\gamma -1)^2 \gamma  (\gamma +1)^2 (\epsilon -1)^2 \epsilon ^2 \left(6 \epsilon ^2-5 \epsilon +1\right)}\Bigg[\left(2 \epsilon +1\right)\big(192 \gamma ^6 (\epsilon -1)^2 \epsilon ^2\\&+96 \gamma ^5 \epsilon  \left(4 \epsilon ^4-28 \epsilon ^3+43 \epsilon ^2-22 \epsilon +3\right)-\gamma ^4 \left(480 \epsilon ^5-5360 \epsilon ^4+9646 \epsilon ^3-5663 \epsilon ^2+896 \epsilon +1\right)\\&-48 \gamma ^3 \epsilon  \left(8 \epsilon ^4-36 \epsilon ^3+46 \epsilon ^2-25 \epsilon +3\right)+\gamma ^2 \left(480 \epsilon ^4-1850 \epsilon ^3+1493 \epsilon ^2-292 \epsilon +1\right)\\&+48 \gamma  \epsilon  \left(20 \epsilon ^3-40 \epsilon ^2+19 \epsilon -3\right)-4 \left(22 \epsilon ^2+7 \epsilon -3\right)\big)\Bigg]\,,\\&\boldsymbol{\tilde h}_{117}=\frac{4 (\epsilon -3)-\gamma ^2 (\epsilon -1) \left(\gamma ^2 \left(320 \epsilon ^2-418 \epsilon +97\right)-336 \epsilon ^2+466 \epsilon -97\right)}{768 \left(\gamma ^2-1\right) (\epsilon -1)^2 (3 \epsilon -1)}\,,\\&
\boldsymbol{\tilde h}_{118}=\frac{\epsilon  (2 \epsilon +1) \left(\gamma ^2 (\epsilon -1) (2 \epsilon -1)-\epsilon  (2 \epsilon +5)+1\right)}{2 (\gamma ^2-1) (\epsilon -1) (2 \epsilon -1) (3 \epsilon -1)}\,,\boldsymbol{\tilde h}_{119}=\frac{(2 \epsilon +1)}{4-12 \epsilon }\,,\\&\boldsymbol{\tilde h}_{120}=\frac{(2 \epsilon +1) (\epsilon  (4 (\epsilon -1) \epsilon +5)-1)}{4 (\gamma^2 -1)  (\epsilon -1) \epsilon  (2 \epsilon -1) (3 \epsilon -1)}\,.
\nonumber
   \end{split}
\end{align}
%\end{landscape}
%\restoregeometry
%\newpage
%The delta functions forces the integral to be reduced in two dimesnion and is given by,
%\begin{align}
    %\begin{split}
    %    \chi_{\mathcal{S}^1}^{(1)}&\to  i\frac{\Gamma(1/2+\epsilon)\Gamma(1/2-\epsilon)^2}{\Gamma(1-2\epsilon)(4\pi)^{3/2-\epsilon}}\frac{\gamma}{4}\epsilon_{\mu j l\sigma}v_2^\sigma (\mathcal{S}_1)^{\mu}_{\,\,\,\,i}\int \frac{d^2 k_1}{(2\pi)^2}e^{-i\vec k_1\cdot \boldsymbol{b}}\frac{k_1^i k_1^j k_1^l}{\vec k_1^{2(\epsilon+1/2)}}\\ &
      %  = -\frac{\gamma}{b^{2(2-2\epsilon)}}\left(\frac{4}{\pi }\right)^{1-\epsilon }\frac{\Gamma \left(\frac{5}{2}-2 \epsilon \right)\Gamma(1/2-\epsilon)^2}{4\Gamma(\epsilon+1/2)\Gamma(1-2\epsilon)(4\pi)^{3/2-\epsilon}}\epsilon_{\mu j l\sigma}v_2^\sigma (\mathcal{S}_1)^{\mu}_{\,\,\,\,i}\,\underbrace{\Big[\frac{1}{2}(\delta^{ij}\hat b^{l}+\delta^{il}\hat b^{j}+\delta^{jl}\hat b^i+(2\epsilon-3)\hat b^i \hat b^j \hat b^l\Big]}_{\mathcal{C}^{ijl}}\\ &
       % \xrightarrow[]{\epsilon\to 0}-\frac{3 \gamma }{8 \pi \,\boldsymbol{b}^4 }\, \epsilon_{\mu j l\sigma}\,v_2^\sigma \,(\mathcal{S}_1)^{\mu}_{\,\,\,\,i}\,\,\,\mathcal{C}^{ijl}\to \frac{3\gamma^2}{8\pi \boldsymbol{b}^4}\epsilon_{mjl0}\,(\mathcal{S}_1)^m_{\,\,\,i}\mathcal{C}^{ijl}
   % \end{split}
%\end{align}
%\newpage
%\appendix
%\section{Comments on the spin dependency of Feynman diagrams involving scalar fields}\label{A.1}
%Let us discuss some crucial points about the spin dependency of Feynman topology with only scalar lines (+ worldline propagators). From \eqref{2.18}, one may notice that scalar (without any graviton lines) can only couple with the spinor through $\varphi-\psi$ vertex, which means at 1PM, there will be no spin-dependent term that has only scalar propagator. Again, because of the presence of time derivative (in \eqref{2.13}), there would not be any vertex like $\varphi-\bar\psi$, which immediately implies the absence of spinning diagrams that have only scalars at the classical level. But we have seen that there can be vertex like $\bar\psi-\varphi-\psi.$ Hence, one can construct  spinning diagram with quantum (closed) loop that has only scalar lines as, 
%\begin{align}
 %   \begin{split}
   %      \begin{minipage}[h]{0.08\linewidth}
	%\vspace{4pt}
%	\scalebox{1.7}{\includegraphics[width=\linewidth]{quantum.png}}
% \end{minipage}&
%\hspace{0.35cm}:\,\,\mathcal{A} \sim \hbar \int_{q}e^{-iq\cdot b}\hat\delta(q\cdot v_1)\hat\delta(q\cdot v_2)\int_{\omega_1}\frac{\hat\delta(k_1\cdot v_2) \mathcal{N}(\omega_1,k_1)}{(\omega_1-k_1\cdot v_1)\omega_1k_1^2(q-k_1)^2}\,.
 %\end{split}
% \end{align}
%As  mentioned, when we have a scalar field in the theory we make the mass of the black holes as a function of the scalar field. Now if one concentrate on the the kinetic part of the superfield, we one notice that in the absence of $\varphi$ the term $\int d\tau\,\bar\psi_a\dot \psi^a$ is invariant under SUSY upto a total derivative term. As we are doing perturbation theory it is enough to show this at flat space limit. In flat space the susy transformations are,
%\begin{align}
 %   \begin{split}
  %      \delta\psi^a=-\epsilon \dot x^a,\,\,\delta\bar\psi^a=-\bar\epsilon \dot x^a,\,\,\delta x^\mu=i\bar\epsilon\psi^\mu+i\epsilon\bar\psi^\mu
   % \end{split}
%\end{align}
%Therefore,
%\begin{align}
 %   \begin{split}
   %     \int d\tau\, \delta(\bar\psi_a\,\dot \psi^a)=\int d\tau \Big(\bar\epsilon \,\,\dot x_a\,\dot \psi^a+\epsilon\bar\psi_a \ddot{x}^a\Big)
   % \end{split}
%\end{align}

%\begin{align}
 %   \begin{split}
   %     S\subset -i\int d\tau \varphi\Bigg(\frac{\bar\psi_a\dot \psi^a-\dot{\bar\psi}_a \psi^a}{2}\Bigg)\,.\label{4.77}
 %   \end{split}
%\end{align}
%\textcolor{red}{If we have only gravitons in our theory, then the mass $m$ can be taken out from the integral, and by doing integral by parts, we are left with only one term in the brackets of \label{4.77}, dropping out the total derivative.?? NOT CLEAR.}
